\documentclass[UTF8,12pt,oneside,fontset=none]{ctexbook}
\setcounter{tocdepth}{1}
%------------------------------------------------
% Packages
%------------------------------------------------
\usepackage{fontspec}
\setmainfont{Noto Serif}
\setCJKmainfont[AutoFakeSlant=0.15]{Noto Serif CJK SC}
\setCJKsansfont{Noto Sans CJK SC}
\setCJKmonofont{Noto Sans Mono CJK SC}
\usepackage{tikz}
\usetikzlibrary{arrows.meta}
\usepackage{amsmath,amssymb,amsthm,mathtools}
\usepackage{bm}
\usepackage[most]{tcolorbox}
\usepackage{graphicx}
\usepackage{float}
\usepackage{hyperref}
\usepackage[normalem]{ulem}
\usepackage{enumitem}
\usepackage{geometry}
\usepackage{fancyhdr}
\usepackage{textgreek}
\usepackage{xparse}
\usepackage{listings}
\usepackage{xcolor}
\usepackage[
    backend=biber,
    style=numeric,
    sorting=nyt
]{biblatex}
\addbibresource{references.bib}

%------------------------------------------------
% Page geometry
%------------------------------------------------
\geometry{margin=1in}
\setlength{\headheight}{15pt}
\emergencystretch=2em
% Allow three-digit page numbers in the table of contents.
\makeatletter
\renewcommand{\@pnumwidth}{2.5em}
\renewcommand{\@tocrmarg}{3.5em}
\makeatother
\pagestyle{fancy}
\fancyhf{}
\fancyhead[RO]{\thepage}
\fancyhead[LO]{\nouppercase{\rightmark}}

%------------------------------------------------
% Theorem-like environments
%------------------------------------------------
\newtheorem{definition}{定义}[chapter]
\newtheorem{example}{例}[chapter]
\newtheorem{theorem}{定理}[chapter]

% ---------- Global tcolorbox styles ----------
\tcbset{
  mybox/.style={
    enhanced,
    boxrule=0.8pt,
    arc=2mm,
    left=2mm, right=2mm, top=1mm, bottom=1mm,
    coltitle=black,
    fonttitle=\bfseries,
    title={#1},
    attach title to upper,
  },
  conceptbox/.style={
    mybox={#1},
    colback=blue!6,
    colframe=blue!55!black,
  },
  derivbox/.style={
    mybox={#1},
    colback=black!3,
    colframe=black!45,
  },
  insightbox/.style={
    mybox={#1},
    colback=green!6,
    colframe=green!50!black,
  },
}
\newtcolorbox{ConceptBox}[1]{conceptbox={#1}}
\newtcolorbox{DerivBox}[1]{derivbox={#1}}
\newtcolorbox{InsightBox}[1]{insightbox={#1}}

\NewDocumentEnvironment{faq}{+b}{
  \section*{常见问题}
  \begin{description}[leftmargin=2em,labelsep=1em]
  #1
  \end{description}
}{}
\newcommand{\question}[1]{\item[\textbf{问：}] #1}
\newcommand{\answer}[1]{\\ \textbf{答：} #1}

\newenvironment{exercises}{
  \section*{练习}
  \begin{enumerate}
}{
  \end{enumerate}
}
\newenvironment{summary}{
  \section*{要点总结}
  \begin{itemize}
}{
  \end{itemize}
}

\newcommand{\bookdate}{2026年9月8日}
\title{生成式人工智能基础小册：直观数学导论}
\author{
Tianhua Chen\\[0.3em]
计算与工程学院\\
哈德斯菲尔德大学\\[0.3em]
\texttt{t.chen@hud.ac.uk}\\[1.5em]
预印本中文翻译版\\
基于 arXiv:2605.29713v2\\
\bookdate
}
\date{}
\hypersetup{
  hidelinks,
  unicode=true,
  pdftitle={生成式人工智能基础小册：直观数学导论},
  pdfauthor={Tianhua Chen}
}

\begin{document}
\frontmatter
% The title page and preface both use i; only the preface needs a page anchor.
\hypersetup{pageanchor=false}
\maketitle
\hypersetup{pageanchor=true}
\chapter*{前言}
\addcontentsline{toc}{chapter}{前言}

近年来，生成式人工智能迅速发展，新模型、新系统和新应用以惊人的速度不断涌现。然而，在这一快速变化的表象之下，许多核心思想实际上建立在一组相对精炼的数学原理之上：潜变量、似然、变分界、可逆变换、随机加噪过程、得分场、对抗式比较以及能量景观。

本书《生成式人工智能基础小册：直观数学导论》（\emph{The Little Book of Generative AI Foundations: An Intuitive Mathematical Primer}）旨在为理解这些思想提供一条紧凑而严谨的学习路径。本书并不试图罗列每一种现代架构、实现细节或最新改进，而是聚焦于使生成模型各主要家族真正可理解的数学基础。我们的目标是在不削弱数学内涵的前提下，提高这些基础知识的可达性。本书面向希望从第一性原理理解生成建模、并对数学保持好奇的研究人员、实践者和学生，可作为一部用于夯实基础的入门读物。

因此，书名中的“小册”应理解为范围上的精炼，而非深度上的浅显。本书在内容选择上有意保持克制，但在所选主题内部则力求细致严谨。核心思想按步骤展开，目的在于显露其背后的数学结构，而不是把模型当作黑箱式配方。全书关注基础而非研究前沿：已经处于前沿研究中的读者可能会熟悉许多主题，但本书希望将这些主题背后的数学逻辑解释得更加显式、直观并且连贯。

参考文献也有意保持聚焦。其主要作用是支撑正文中的关键历史脉络和技术联系，而不是把本书扩展成一份篇幅庞大的综述型书目。

\section*{写作方法与内容组织}
%
%\addcontentsline{toc}{section}{写作方法与内容组织}

本书的核心目标，是用一条连贯的数学叙事把生成模型的主要家族联系起来。章节顺序的设计，是为了让数学概念尽可能自然、紧凑地展开，而不是严格遵循历史发展顺序。各类模型按照一种便于“后续思想建立在前述思想之上”的次序引入，使读者能够看到潜变量、变分目标、扩散过程、得分场、精确密度模型、对抗学习以及基于能量的建模之间的内在联系。

全书遵循一个简单原则：只有当某种数学工具对于理解一个建模思想变得必要时，才引入它。因此，所需的线性代数、概率论、微积分、高斯代数以及密度变换工具，都嵌入在具体的建模叙事之中，而不是被单独放在前面进行冗长的预备知识回顾。

与仅做高层概览相比，本书包含了更多的数学推导。这是有意为之：生成建模中的许多思想，只有把其数学目标逐层展开之后才会真正变得清晰。推导过程按步骤给出，并在需要时引入或回顾必要概念，使数学始终与建模直觉相连接，而不是成为彼此孤立的技术模块。

本书从主成分分析（PCA）和自编码器开始，借此引入线性变换、投影、重构以及潜在结构等概念。随后，概率主成分分析（Probabilistic PCA, PPCA）把这些思想转化为本书中的第一个潜变量生成模型。

接下来进入变分生成建模。我们通过概率潜变量引入证据下界（Evidence Lower Bound, ELBO）、期望最大化（Expectation--Maximisation, EM）算法以及变分推断，并进一步扩展到变分自编码器（Variational Autoencoder, VAE）。

随后，本书转向扩散模型。首先在离散时间框架下把去噪扩散概率模型（Denoising Diffusion Probabilistic Models, DDPM）作为序列潜变量模型引入；之后介绍连续时间生成建模所需的微积分基础，并进一步转向扩散的得分视角，在这一视角下，生成过程由学习得到的得分场和逆时间随机动力学共同驱动。

接下来的一章讨论精确密度模型，包括归一化流（normalising flows）和自回归分解（autoregressive factorisations）。这些模型展示了如何通过结构设计保持似然的可计算性：要么依赖可逆变换，要么利用概率链式法则。

最后一章讨论超越可计算似然的生成建模方法。我们将介绍生成对抗网络（GAN）、Wasserstein GAN 以及基于能量的模型（Energy-Based Models, EBM），说明生成学习如何不依赖显式似然求值，而通过比较、几何差异或标量能量景观来实现。

% 两个附录汇总了高斯代数和扩散逆后验相关的辅助推导。

\section*{起源与版本说明}
%
%\addcontentsline{toc}{section}{起源与版本说明}

本书是一项由作者独立完成的学术项目，其动机来自长期以来对“如何让现代生成式人工智能的数学基础更加透明、更加易于理解”的兴趣。本书并非某项受资助科研项目、机构委托成果、正式课程开发活动，也不是针对任何特定课程或培养项目编写的官方教学文档。

本书手稿源于一个更广泛的独立研究计划，该计划关注现代生成建模的概率与数学基础。其中部分思想的早期简短版本曾以预印本形式发表，先发表于 SSRN，后经修订发表于 arXiv~\cite{chen2025plvm,chen2025unified}。本书在此基础上进一步扩展，形成一部更系统的数学导论，从经典潜变量模型一路讨论到现代生成式人工智能的基础方法。

当前手稿是本书的第二版预印本。与第一版相比，本版在全书范围内加入了更多可视化解释，并对相应正文做了少量澄清。它仍是一份工作中的预印本，后续可能继续进行修订、勘误、解释性补充或增加新的材料。未来亦可能提交经过进一步修订的版本以正式出版。

\section*{再利用说明}
%
%\addcontentsline{toc}{section}{再利用说明}

本手稿以预印本形式公开，供学术阅读、引用、研究参考和学习使用。读者可以在适当署名的前提下引用本书，并将其作为科研、学习和教学中的辅助资源。

除非适用的版权例外另有允许，或未来版本附带的其他许可另有规定，对本书内容进行大规模复制、改编、再分发，或将其转换为衍生教学、培训或课程资源，需获得作者许可。

\section*{致谢}
%
%\addcontentsline{toc}{section}{致谢}

本书源于我持续从第一性原理理解并解释现代生成式人工智能数学基础的努力。在反复尝试把线性代数、概率论、变分推断、随机建模以及现代深度生成模型中的思想连接成一条连贯叙事的过程中，许多解释逐渐变得更加清晰。

感谢更广泛的学术共同体。大量论文、书籍、讲座与讨论共同塑造了本手稿形成时所处的知识背景。文中若仍有错误、遗漏或解释不清之处，责任均在作者本人。

\vspace{1cm}

\begin{flushright}
Tianhua Chen\\
\textit{\bookdate}
\end{flushright}

\tableofcontents
\mainmatter
\chapter[线性代数基础：从 PCA 到自编码器]{线性代数基础：\\从 PCA 到自编码器}
\section*{概述}
\addcontentsline{toc}{section}{概述}

主成分分析（Principal Component Analysis, PCA）是本书一个自然且易于进入的起点，主要有两个原因：
\begin{itemize}
   \item \textbf{搭建通向概率模型与生成模型的桥梁。}
    作为一种广为人知的降维方法，PCA 为本书后续将要展开的若干思想提供了自然的过渡。PCA 的重构表述会自然引向自编码器，在自编码器中，压缩与重构直接从数据中学习。其概率扩展——概率主成分分析（Probabilistic PCA, PPCA）——又从生成建模的角度重新解释了这一熟悉的确定性方法。这样一条从代数投影，到重构，再到概率潜变量建模的路径，预示了现代框架（例如变分自编码器（VAE）和扩散模型）中所采用的技术；在这些更灵活、表达能力更强的框架中，相似的原理会再次出现。

    \item \textbf{重新回顾线性代数中的核心概念。}  
    PCA 能够自然地展示线性代数中的若干关键思想——投影、基变换、特征值与特征向量——这些概念支撑着本书后续许多结构。生成模型运行在高维空间中，数据通过矩阵进行表示与变换。协方差矩阵、方差分析、矩阵分解以及基变换等概念会反复出现，成为理解数据结构的重要工具。
\end{itemize}

基于这些原因，PCA 是本书非常理想的起点。

\section*{概念导图}
\addcontentsline{toc}{section}{概念导图}

本章将按照一条循序渐进的路径展开：从基本线性变换出发，进入 PCA，再由 PCA 过渡到自编码器。每一步的安排都旨在建立直觉、自然连接后续内容，并为后续章节中的概率生成模型做好准备。

\begin{itemize}

    \item[\(\rightarrow\)] \textbf{首先把矩阵理解为一种变换。}  
    我们不会把矩阵--向量乘法只看成机械的计算规则，而是把它理解为向量和坐标系如何被重新塑造。这个几何视角不仅为 PCA 提供必要直觉，也为后续生成模型中反复出现的“基于变换理解模型”的视角打下基础。

    \item[\(\rightarrow\)] \textbf{随后建立矩阵乘法的两个互补视角，并把它们与 PCA 联系起来。}  
    \emph{列组合}视角说明矩阵如何定义一组新的方向，并由这些方向组合得到输出；\emph{行--点积}视角则从投影与对齐的角度解释矩阵乘法。这两种理解共同预示了 PCA 的两个核心作用：改变基，以及通过投影测量变化程度。

    \item[\(\rightarrow\)] \textbf{接下来，把 PCA 表述为寻找最大方差方向的问题，并引入特征值和特征向量作为求解这一问题的关键工具。}  
    将数据投影到候选方向上并测量其离散程度之后，PCA 就转化为一个优化问题。由于协方差矩阵是对称矩阵，其特征向量构成一个标准正交基；在这一基下改写目标函数后，可以看到最大方差方向就是最大特征值对应的特征向量。

    \item[\(\rightarrow\)] \textbf{随后把这一论证扩展到完整的降维与重构。}  
    最大的若干特征值对应的特征向量定义了一组按重要性排序的正交主方向，每一个方向都尽可能捕获剩余方差。由此自然得到投影、重构以及解释方差比，从而同时从几何和实践层面理解 PCA 如何压缩数据。

    \item[\(\rightarrow\)] \textbf{然后进一步揭示 PCA 作为基变换的深层结构。}  
    对协方差矩阵进行对角化，可以理解为一种相似变换：它保持总方差不变，但把方差重新组织到与主方向对齐、彼此去相关的坐标轴上。这一视角也自然引出了更一般的矩阵分解，例如奇异值分解（SVD）。

    \item[\(\rightarrow\)] \textbf{最后，关注点从方差转向重构，从而建立 PCA 与自编码器之间的桥梁。}  
    PCA 可以被重新解释为一个最小化重构误差的编码器--解码器模型；随后可以看到，线性自编码器会恢复同一个最优子空间，而非线性自编码器则把这一思想从线性子空间扩展到弯曲流形。这条路径既揭示了确定性重构模型的能力，也揭示了其局限，并由此引出下一章中的概率潜变量模型。

\end{itemize}

下面的图示对本章路径进行了总结：从基本线性变换，经由 PCA 与重构，最终到达本书后续将使用的自编码器桥梁。

\begin{center}
\begin{tikzpicture}[
  >=Latex,
  every node/.style={font=\small},
  box/.style={draw,rounded corners=2pt,align=center,inner sep=5pt,
             text width=3.05cm,minimum height=1.15cm},
  flow/.style={->,thick}
]
\node[font=\normalsize\bfseries] at (6.0,2.75)
  {线性代数基础概览};

\node[box] (mat) at (1.7,0.95) {
  \textbf{把矩阵理解为变换}\\[3pt]
  {\scriptsize 基向量的像决定\\整个线性映射}
};
\node[box] (proj) at (6.0,0.95) {
  \textbf{内积与投影}\\[3pt]
  {\scriptsize 对齐程度 $\rightarrow$ 所选方向\\上的坐标}
};
\node[box] (eig) at (10.3,0.95) {
  \textbf{特征向量与特征值}\\[3pt]
  {\scriptsize 自然方向\\与缩放强度}
};
\draw[flow] (mat) -- (proj);
\draw[flow] (proj) -- (eig);

\node[box] (pca) at (10.3,-2.30) {
  \textbf{PCA}\\[3pt]
  {\scriptsize 最大方差\\主子空间}
};
\node[box] (basis) at (6.0,-2.30) {
  \textbf{基变换与 SVD}\\[3pt]
  {\scriptsize 对角化协方差\\并揭示结构}
};
\node[box] (ae) at (1.7,-2.30) {
  \textbf{重构桥梁}\\[3pt]
  {\scriptsize 自编码器\\流形\\概率}
};
\draw[flow] (eig) -- (pca);
\draw[flow] (pca) -- (basis);
\draw[flow] (basis) -- (ae);
\end{tikzpicture}
\end{center}


\medskip
\noindent
\section{作为线性变换的矩阵：基方向、投影与复合}

理解矩阵时，最好不要只把它看成一张数字表，而应把它看成作用于向量的变换的表示。矩阵定义了一个\emph{线性变换}：它以一种保持加法与数乘结构的方式移动向量。理解矩阵如何作用于向量，是一项基础能力；它不仅对 PCA 这类线性方法重要，对于由一系列线性与非线性变换构成的更复杂生成模型同样重要。

\subsection{矩阵对基方向的变换}

\paragraph{标准基与恒等变换。}

先从最简单的情形开始。在二维空间中，标准基（也称单位基）向量为
\[
\hat{e}_1 =
\begin{bmatrix}1 \\ 0\end{bmatrix},
\qquad
\hat{e}_2 =
\begin{bmatrix}0 \\ 1\end{bmatrix}.
\]
\(\mathbb{R}^2\) 中的任意向量，都能够唯一地写成这两个方向的线性组合。

将这两个基向量作为矩阵的两列，就得到单位矩阵
\[
I =
\begin{bmatrix}
1 & 0 \\
0 & 1
\end{bmatrix},
\]
它表示一个使所有向量保持不变的变换。从几何上看，恒等变换不会改变任何向量，因此方向、长度和夹角都保持不变。

\paragraph{示例：旋转 \(90^\circ\)。}

为了观察矩阵如何主动改变空间，考虑
\[
R =
\begin{bmatrix}
0 & -1 \\
1 & 0
\end{bmatrix}.
\]
该矩阵的两列告诉我们基向量如何被变换：
\[
\hat{e}_1 \mapsto
\begin{bmatrix}0 \\ 1\end{bmatrix},
\qquad
\hat{e}_2 \mapsto
\begin{bmatrix}-1 \\ 0\end{bmatrix}.
\]
每一个基方向都逆时针旋转了 \(90^\circ\)。由于所有向量都是由这些基向量构成的，因此同样的旋转会一致地作用于整个平面。

这一例子体现了线性代数中的一个核心原则：
\begin{quote}
\emph{一个线性变换完全由它把各个基向量映射到哪里所决定。}
\end{quote}

\paragraph{超越旋转：一般线性变换。}

矩阵并不一定表示旋转。它的各列不要求彼此正交，也不要求具有单位长度。

例如，考虑
\[
A =
\begin{bmatrix}
1 & 1 \\
0 & 1
\end{bmatrix}.
\]
这里，第一个基向量保持不变，而第二个基向量从
\(\begin{bmatrix}0 \\ 1\end{bmatrix}\) 被映射为 \(\begin{bmatrix}1 \\ 1\end{bmatrix}\)。这一变换对应一个剪切（shear）：一个坐标轴发生倾斜，而另一个保持不变。

尽管旋转与剪切在几何效果上不同，但它们的描述方式完全一致：指定各个基向量被映射到哪里。一旦这些映射确定，变换对任意其他向量的作用也就完全确定。

图~\ref{fig:ch1-matrix-transformations} 把这一基向量视角几何化：变换后的基向量决定了整个坐标网格如何被矩阵带到新的位置。

\begin{figure}[H]
\centering
\begin{minipage}[t]{0.4\textwidth}
\vspace{0pt}\centering
\includegraphics[width=\linewidth]{figures/ch01/panels/fig1_1a_rotation.pdf}
\end{minipage}\hfill
\begin{minipage}[t]{0.4\textwidth}
\vspace{0pt}\centering
\includegraphics[width=\linewidth]{figures/ch01/panels/fig1_1b_shear.pdf}
\end{minipage}
\caption{\textbf{基方向决定线性变换。}
浅色虚线表示原始坐标网格，实线表示变换后的网格。箭头 $A\hat e_1$ 与 $A\hat e_2$ 是变换后的基方向，也就是矩阵 $A$ 的两列。旋转保持方形网格，而剪切会使一族网格线发生倾斜；一旦变换后的基方向确定，其他任意向量的变换都可由线性组合得到。}
\label{fig:ch1-matrix-transformations}
\end{figure}

下面转向矩阵乘法，考察这些变换如何作用于任意向量。

\subsection{理解矩阵乘法的两个视角}

\paragraph{矩阵--向量乘法：列组合视角。}

设矩阵 \(A\) 的两列分别为 \(v_1\) 和 \(v_2\)。对于向量 \(x = \begin{bmatrix}a \\ b\end{bmatrix}\)，矩阵--向量乘法给出
\[
Ax = a\,v_1 + b\,v_2.
\]

这个表达式揭示了理解矩阵乘法的第一种基本方式。矩阵的各列定义了一组新的方向，而向量中的各个分量指定这些方向应如何组合。因此，矩阵乘向量可以理解为：用矩阵的列方向作加权组合来构造输出。

这一解释与向量本身的构造方式直接对应。任意向量都可以写成
\[
x = a\,\hat{e}_1 + b\,\hat{e}_2,
\]
即标准基向量的线性组合。矩阵--向量乘法所做的事情，只是把这些标准基方向替换成矩阵 \(A\) 的各列所编码的变换后坐标方向。

在一般情况下，同样有
\[
Ax = \sum_{j=1}^n x_j\,a_j,
\]
其中 \(a_j\) 表示 \(A\) 的第 \(j\) 列。

\paragraph{矩阵--向量乘法：逐行点积视角。}

理解矩阵--向量乘法还有第二种同样重要的方式。此时不再聚焦于矩阵的列，而是考察输出向量的各个分量。

在这一视角下，\(Ax\) 的第 \(i\) 个分量由 \(A\) 的第 \(i\) 行与输入向量做点积得到。这种“行乘列”的计算并不是任意规定的数值规则，而具有明确的几何含义。

\begin{tcolorbox}[
  colback=blue!6,
  colframe=blue!55!black,
  title={为什么点积能够编码投影},
  boxrule=0.8pt,
  arc=2mm,
  left=2mm,
  right=2mm,
  top=1mm,
  bottom=1mm
]

\textbf{余弦相似度。}
两个向量 \(v_1\) 和 \(v_2\) 的余弦相似度为
\[
\cos\theta = \frac{v_1^\top v_2}{\|v_1\|\,\|v_2\|},
\]
其中 \(\theta\) 是两者之间的夹角。

\textbf{标量投影。}
对上式变形，可以得到 \(v_1\) 在 \(v_2\) 方向上的标量投影：
\[
\text{proj}_{v_2}(v_1)
=
\frac{v_1^\top v_2}{\|v_2\|}.
\]
该量衡量 \(v_1\) 沿 \(v_2\) 方向所包含的分量大小。

\textbf{点积作为缩放后的投影。}
因此，点积可以写成
\[
v_1^\top v_2
=
\|v_2\| \times \text{（\(v_1\) 在 \(v_2\) 上的投影）}.
\]
如果 \(v_2\) 是单位向量，那么点积就等于标量投影。

\textbf{向量投影。}
将标量投影乘以 \(v_2\) 的单位方向，就得到相应的\emph{向量投影}：
\[
\mathrm{proj}_{v_2}(v_1)
=
\frac{v_1^\top v_2}{\|v_2\|^2}\, v_2.
\]
它给出了 \(v_1\) 沿 \(v_2\) 方向的分量的向量表示。

\end{tcolorbox}
上述讨论说明，逐行点积计算中出现的点积，在几何上可以理解为对投影或对齐程度的度量。从这一角度看，\(Ax\) 的每一个分量都记录了输入向量与 \(A\) 的某一个行方向之间的对齐程度。

\paragraph{推广到矩阵--矩阵乘法。}

同样的思想可以自然推广到矩阵--矩阵乘法。两个矩阵相乘，对应两个线性变换的复合。

乘积 \(AB\) 可以逐列理解为
\[
AB =
\big[\,A b_1 \;\; A b_2 \;\; \cdots \;\; A b_p\,\big],
\]
即 \(B\) 的每一列分别经过 \(A\) 的变换。等价地，当它作用于向量 \(v\) 时，乘积 \(AB\) 表示两个连续发生的变换：
\[
(AB)v = A(Bv).
\]
向量先经过 \(B\) 的变换，再由 \(A\) 继续变换。每一个阶段仍然都可以使用前面介绍的列组合视角和逐行点积视角来理解。

\medskip
\noindent
综合起来，这些视角构成了 PCA 的直接线性代数基础，并会在我们逐步建立更具表达能力的生成模型时反复出现：
\begin{itemize}
\item 矩阵的列描述一个变换如何重塑空间；
\item 列组合解释变换后的向量如何构造；
\item 逐行点积说明如何提取投影；
\item 矩阵--矩阵乘法则把多个变换串联起来。
\end{itemize}

\section{特征向量、特征值与矩阵的自然坐标系}
在深入 PCA 之前，我们再引入一个理解矩阵如何作用于空间的基本工具。上一节中，我们把矩阵解释为一种线性变换，它通过混合、拉伸与投影不同方向来重塑向量。为了更精确地描述这种作用，现在希望找到一些方向，使变换在这些方向上的行为尽可能简单。

特征向量与特征值恰好提供了这种语言。它们确定了矩阵的一些特殊方向：沿这些方向，变换不会把向量与其他方向混合，而只会把它映射为自身的一个标量倍数。

\paragraph{特征值与特征向量。}
给定方阵 \(A \in \mathbb{R}^{d \times d}\)，若存在非零向量 \(v \in \mathbb{R}^d\) 使得
\[
A v = \lambda v
\]
对某个标量 \(\lambda \in \mathbb{R}\) 成立，则称 \(v\) 为 \(A\) 的一个特征向量，\(\lambda\) 为对应的特征值。这个等式说明，当 \(A\) 作用于 \(v\) 时，输出仍落在 \(v\) 张成的同一条直线上；其长度被缩放 \(|\lambda|\) 倍，如果 \(\lambda<0\)，方向还会发生反转。

对于一般向量，乘以 \(A\) 可能造成旋转、拉伸、剪切或其他形式的方向改变。特征向量则标识了不会与其他方向混合的特殊方向：输出仍位于原来的直线上，仅相差一个标量因子 \(\lambda\)。因此，可以把特征向量理解为矩阵 \(A\) 所定义变换下的\emph{纯缩放方向}。

\paragraph{特征分解：在矩阵的自然坐标系中表示矩阵。}
如果矩阵 \(A \in \mathbb{R}^{d \times d}\) 有 \(d\) 个线性无关的特征向量，就可以把它们按列组成矩阵 \(V\)，并把对应的特征值放在对角矩阵 \(\Lambda\) 的对角线上。

具体地，若 \(v_1, v_2, \dots, v_d\) 是 \(A\) 的特征向量，对应特征值分别为 \(\lambda_1, \lambda_2, \dots, \lambda_d\)，则
\[
A [\, v_1 \; v_2 \; \cdots \; v_d \,]
=
[\, A v_1 \; A v_2 \; \cdots \; A v_d \,]
=
[\, \lambda_1 v_1 \; \lambda_2 v_2 \; \cdots \; \lambda_d v_d \,].
\]

用紧凑的矩阵形式可写为
\[
A V = V \Lambda,
\]
其中
\[
V = [\, v_1 \; v_2 \; \cdots \; v_d \,],
\qquad
\Lambda =
\begin{bmatrix}
\lambda_1 & 0 & \cdots & 0 \\
0 & \lambda_2 & \cdots & 0 \\
\vdots & \vdots & \ddots & \vdots \\
0 & 0 & \cdots & \lambda_d
\end{bmatrix}.
\]

对上式整理可得特征分解
\[
A = V \Lambda V^{-1}.
\]

\begin{tcolorbox}[
  colback=blue!6,
  colframe=blue!55!black,
  title={直觉：把特征分解理解为基变换},
  boxrule=0.8pt,
  arc=2mm,
  left=2mm,
  right=2mm,
  top=1mm,
  bottom=1mm
]

线性代数中的一个核心思想是：同一个向量可以在不同的坐标系中进行描述。特征分解通过引入一个\emph{新的基}把这一点显式化；在这个新基中，矩阵的作用会变得格外简单。

矩阵 \(V\) 的各列由 \(A\) 的特征向量组成。这些特征向量为整个空间定义了一组新的坐标轴，也就是特征向量基。在这一基下表示一个向量，就是用它沿这些优选方向的坐标来描述该向量。

\medskip
\textbf{改变坐标。}
给定在标准基下表示的向量 \(x\)，我们通过
\[
x_{\text{eig}} = V^{-1} x
\]
进行坐标变换。这个操作并不会改变向量本身，只会改变向量的\emph{表示方式}，把它转换为特征向量坐标。

\medskip
\textbf{特征基下的简单作用。}
一旦把向量表示到特征向量基中，\(A\) 的作用就变得非常简单：每一个坐标可以彼此独立地缩放，缩放因子就是对应的特征值，而这些缩放由对角矩阵 \(\Lambda\) 表示：
\[
x_{\text{eig}} \;\mapsto\; \Lambda x_{\text{eig}}.
\]

\medskip
\textbf{回到原始基。}
完成缩放后，再通过
\[
x' = V(\Lambda x_{\text{eig}})
\]
把结果转换回原始坐标系，从而在标准基下重构该向量。

\medskip
综合起来，分解 \(A = V \Lambda V^{-1}\) 可以读作
\[
x
\;\xrightarrow{\;V^{-1}\;}
\text{基变换}
\;\xrightarrow{\;\Lambda\;}
\text{各方向独立缩放}
\;\xrightarrow{\;V\;}
\text{回到原始基}.
\]

因此，特征分解说明：\emph{只要选择合适的坐标系}，一个矩阵的作用就可以化为简单的缩放。
\end{tcolorbox}

\paragraph{对称矩阵这一特殊情形。}
对于一般方阵，实数域上的特征分解可能不存在；也可能出现复特征值，或者没有足够多线性无关的特征向量。然而，当 \(A\) 是实对称矩阵（\(A=A^\top\)）时，情况会大大简化。

此时有若干重要性质成立：
\begin{itemize}
    \item \(A\) 的所有特征值都是实数。
    \item 存在一个\emph{正交}矩阵 \(U\)，其各列为 \(A\) 的特征向量，使得
    \[
    A = U \Lambda U^\top,
    \]
    其中 \(U\) 的各列均为单位长度且两两正交，并满足 \(U^\top U = I\)。
\end{itemize}

这种特殊结构意味着，对称矩阵总可以通过标准正交基变换进行对角化，因此坐标变换本身会保持长度和夹角不变。正如后面将看到的，这一性质对 PCA 至关重要，因为经验协方差矩阵始终是对称矩阵。

至此，我们已经具备了表述并求解 PCA 问题所需的工具。

\section{通过方差最大化理解主成分分析}

主成分分析（PCA）\cite{pearson1901,hotelling1933} 通过引入一个适应数据结构的新坐标系，为数据降维提供了一套系统方法。其核心思想，是识别一组新的坐标轴，使数据沿这些轴的变化程度最大。将坐标轴与数据最为分散的方向对齐后，PCA 就揭示了数据在哪些方向上包含最多信息，并使我们能够在尽量少损失结构的前提下获得低维表示。

PCA 的直觉可以概括为两个核心原则：
\begin{enumerate}
    \item \textbf{把方差视为信息量。}  
    数据在某个方向上的方差越大，数据点通常越容易被区分，因此该方向携带的结构信息也越丰富。相反，方差很小的方向对区分数据的贡献较少。

    \item \textbf{顺序进行正交最大化。}  
    PCA 首先找出能够捕获最大可能方差的单一方向。随后继续选择新的方向，在保持与已选方向正交的约束下，尽可能捕获剩余方差。
\end{enumerate}

\subsection{PCA 的数学表述}

\paragraph{数据中心化与投影方差的定义。}
为了便于表述 PCA，首先把数据中心化，使整个数据云以原点为中心。给定数据点 \(\{x^{(i)}\}_{i=1}^n \subset \mathbb{R}^d\)，假设
\[
\frac{1}{n}\sum_{i=1}^n x^{(i)} = 0.
\]
中心化可以确保我们测量的是数据相对于其中心的变化程度，而不会受到任意平移偏置的影响。

前面在介绍投影和点积时已经讨论过：沿某个方向的变化程度，可以通过考察数据在该方向上的投影强度来理解。设 \(v \in \mathbb{R}^d\) 是一个单位向量，表示我们希望把数据投影到的候选方向，因此 \(\|v\|=1\)。对于数据点 \(x^{(i)}\)，其在 \(v\) 上的标量投影为
\[
v^\top x^{(i)},
\]
因为 \(v\) 的长度为 1。相应的向量投影为
\[
(v^\top x^{(i)})\,v.
\]
这再次说明，当 \(v\) 是单位向量时，点积恰好给出了沿该方向的标量投影。

由于数据已经中心化，投影标量值的均值同样为 0：
\[
\frac{1}{n}\sum_{i=1}^n v^\top x^{(i)}
=
v^\top \left(\frac{1}{n}\sum_{i=1}^n x^{(i)}\right)
=
0.
\]
为了量化这些投影值的离散程度，我们计算它们相对于该均值的方差。因此，数据沿方向 \(v\) 投影后的方差为
\[
\mathrm{Var}(v^\top X)
=
\frac{1}{n}\sum_{i=1}^n \bigl(v^\top x^{(i)}\bigr)^2.
\]
从几何上看，投影越大、分布越分散的方向，对应的方差越高。

\paragraph{协方差矩阵形式。}
投影方差可以用矩阵形式紧凑地写为
\[
\mathrm{Var}(v^\top X)
=
\frac{1}{n}\sum_{i=1}^n v^\top x^{(i)} (x^{(i)})^\top v
=
v^\top \left( \frac{1}{n}\sum_{i=1}^n x^{(i)} (x^{(i)})^\top \right) v.
\]
括号内的矩阵就是经验协方差矩阵
\[
\Sigma = \frac{1}{n} \sum_{i=1}^n x^{(i)} (x^{(i)})^\top.
\]
因此，数据沿方向 \(v\) 投影后的方差可以写成
\[
\mathrm{Var}(v^\top X) = v^\top \Sigma\, v.
\]

于是，PCA 问题就是在单位长度约束下寻找使上述量最大的方向：
\[
\max_{v \in \mathbb{R}^d} \; v^\top \Sigma v
\quad \text{约束为} \quad v^\top v = 1.
\]

\subsection{求解 PCA 优化问题}

协方差矩阵 \(\Sigma\) 按定义就是对称矩阵。如上一节所述，对称矩阵具有正交特征分解：
\[
\Sigma = U \Lambda U^\top,
\qquad
\Lambda = \mathrm{diag}(\lambda_1,\ldots,\lambda_d),
\]
其中 \(U\) 的各列是标准正交特征向量，特征值按
\[
\lambda_1 \ge \lambda_2 \ge \cdots \ge \lambda_d \ge 0
\]
排序。

为了求解优化问题，我们把目标函数表示到这一特征向量坐标系中。

\paragraph{步骤 1：变换到特征向量坐标。}
按照前面介绍的基变换解释，把方向向量 \(v\) 在特征向量基下表示为
\[
\tilde{v} = U^\top v.
\]
由于 \(U\) 为正交矩阵，\(U^\top = U^{-1}\)，因此该操作只是把 \(v\) 的坐标从标准基转换到特征向量基。

正交变换还会保持向量长度：
\[
\tilde{v}^\top \tilde{v}
=
(U^\top v)^\top (U^\top v)
=
v^\top U U^\top v
=
v^\top v.
\]
因此，单位长度约束 \(v^\top v = 1\) 会转化为
\[
\tilde{v}^\top \tilde{v} = 1.
\]

\paragraph{步骤 2：把目标函数对角化。}
将特征分解代入目标函数，有
\[
v^\top \Sigma v
=
v^\top U \Lambda U^\top v
=
(U^\top v)^\top \Lambda (U^\top v)
=
\tilde{v}^\top \Lambda \tilde{v}.
\]
由于 \(\Lambda\) 为对角矩阵，上式化简为
\[
\tilde{v}^\top \Lambda \tilde{v}
=
\sum_{i=1}^d \lambda_i \tilde{v}_i^2.
\]
这明确表明：沿方向 \(v\) 捕获到的方差，是各个特征值的加权和，而权重由 \(v\) 与各特征向量的对齐程度决定。

\paragraph{步骤 3：给出方差的上界。}
由于 \(\tilde{v}_i^2 \ge 0\)，且 \(\sum_{i=1}^d \tilde{v}_i^2 = 1\)，目标函数实际上是各特征值的加权平均：
\[
\sum_{i=1}^d \lambda_i \tilde{v}_i^2
\le
\lambda_1 \sum_{i=1}^d \tilde{v}_i^2
=
\lambda_1.
\]
任何较小特征值的组合都不可能超过最大的特征值。

\paragraph{步骤 4：得到最优方向。}
只有当全部权重都集中到最大特征值上时，才会达到该上界：
\[
\tilde{v} = [1,\,0,\,\ldots,\,0]^\top.
\]
将其变换回原始坐标，得到
\[
v = U \tilde{v} = u_1,
\]
即最大特征值 \(\lambda_1\) 对应的特征向量。\footnote{译注：这里给出的是一个最优解。若最大特征值为单重，则正、负两个单位特征向量都达到上界；若最大特征值为重根，则其特征子空间中的任意单位向量都达到上界。原文的坐标选择不表示最优解唯一。}

因此，使投影方差最大的方向，恰好就是协方差矩阵的第一特征向量，而最大方差等于 \(\lambda_1\)。

\paragraph{后续主成分。}
因此，第一个主成分就是 \(u_1\)。其余主成分通过重复相同的优化过程获得，同时要求新方向与此前已经选择的方向正交。第 \(k\) 个主成分就是第 \(k\) 大特征值对应的特征向量 \(u_k\)，并满足
\[
u_k^\top u_j = 0 \qquad \text{对所有 } j < k \text{ 成立}.
\]

从几何上看，每一个新的主方向都会在保持与此前方向正交的同时，尽可能捕获尚未解释的变化。由于这些方向是协方差矩阵的正交特征向量，由它们得到的主成分坐标彼此不相关。协方差矩阵的特征分解自然地给出了这样一组有序、标准正交的方向。

\subsection{用于降维的 PCA}

一旦确定主成分，PCA 就提供了一套具体的降维过程：先把数据投影到一个低维子空间中；如果需要，还可以再把它重构回原始空间。

设
\[
U_k = [\,u_1\;u_2\;\cdots\;u_k\,] \in \mathbb{R}^{d \times k}
\]
表示由前 \(k\) 个主成分组成的矩阵，也就是协方差矩阵 \(\Sigma\) 中对应最大 \(k\) 个特征值 \(\lambda_1 \ge \lambda_2 \ge \cdots \ge \lambda_k\) 的特征向量。PCA 通过只使用这 \(k\) 个捕获最大方差的方向来表示数据点 \(x \in \mathbb{R}^d\)，从而降低其维度。

\paragraph{投影（编码）。}
\(x\) 的低维表示通过将其投影到主成分子空间得到：
\[
z = U_k^\top x \in \mathbb{R}^k.
\]
向量 \(z\) 包含 \(x\) 沿各主方向的坐标，构成数据的压缩表示。

\paragraph{重构（解码）。}
只使用这 \(k\) 个成分对 \(x\) 进行近似，可通过把低维表示映射回原始空间实现：
\[
\hat{x} = U_k z = U_k U_k^\top x.
\]
这一操作取得 \(x\) 在主成分空间中的坐标，再利用 \(U_k\) 中的主方向进行组合，从而在原始特征空间中重构出一个向量。由于只保留前 \(k\) 个主成分，因此重构过程中不可避免地会损失部分信息。

值得注意的是，当 \(k=d\) 时，有 \(U_k \in \mathbb{R}^{d\times d}\)，并且
\[
U_k U_k^\top = I,
\]
因此重构是精确的。而当 \(k<d\) 时，\(U_k U_k^\top\) 不再是单位矩阵，\(\hat{x}\) 只能是原始数据点的近似。

\paragraph{解释方差比。}
在实践中，需要量化降维后保留了原始总变化量中的多少。最常用的指标就是\emph{解释方差比}（explained variance ratio）。

原始数据集中的总方差为
\[
\sum_{i=1}^d \lambda_i,
\]
而前 \(k\) 个主成分捕获的方差为
\[
\sum_{i=1}^k \lambda_i.
\]
因此，\(k\) 维表示所保留的方差比例为
\[
\text{解释方差比}
=
\frac{\sum_{i=1}^k \lambda_i}{\sum_{i=1}^d \lambda_i},
\]
它为选择主成分数量提供了一个有理论依据、由数据驱动的标准，在实际应用中被广泛采用。

下面进一步仔细考察 PCA 的结构。这将说明为什么解释方差比能够以一种有原则的方式量化降维所保留的信息量，并揭示下一节将讨论的 PCA 更深层本质。

\section{把 PCA 理解为基变换：相似变换与 SVD}
现在从更深的结构视角重新审视 PCA。上一节中，PCA 是作为方差最大化问题的解得到的。本节将说明，同一过程也可以理解为一种\emph{基变换}：它将协方差矩阵对角化，并把方差重新组织到一组正交方向上。这一视角能够解释为什么总方差保持不变、为什么特征值能够自然量化各主方向上捕获的方差，以及为什么 PCA 会产生彼此不相关的坐标。

这一视角还会自然引出奇异值分解（Singular Value Decomposition, SVD）。SVD 将同样的几何思想推广到协方差矩阵之外，也推广到非方阵和非对称矩阵。虽然基本形式的 PCA 并不要求掌握 SVD，但它提供了更一般的线性代数视角，并将在本书后续内容中持续发挥作用。

为了建立这种解释，先介绍两个关键概念：矩阵的迹以及相似变换。

\subsection{迹、相似变换与方差保持}

\paragraph{协方差矩阵的迹。}
对于任意方阵 \(A\in\mathbb{R}^{d\times d}\)，其迹定义为对角元素之和：
\[
\mathrm{Tr}(A)=\sum_{i=1}^d a_{ii}.
\]
当 \(A\) 是协方差矩阵 \(\Sigma\) 时，每个对角元素 \(\sigma_{ii}\) 都是第 \(i\) 个坐标的方差。因此
\[
\mathrm{Tr}(\Sigma)
=
\sum_{i=1}^d \Sigma_{ii}
=
\sum_{i=1}^d \mathrm{Var}(X_i),
\]
这恰好就是数据集沿所有坐标方向的总方差。

这一点对 PCA 很重要。基变换虽然可能改变方差在各个坐标上的分配方式，但总方差不会改变。\footnote{译注：这里关于总方差的结论应限于 PCA 所用的标准正交基变换。一般坐标变换 \(y=A^{-1}x\) 给出 \(\operatorname{Cov}(y)=A^{-1}\Sigma A^{-\top}\)，并非一般的相似变换，总方差通常会改变。对线性算子而言，相似变换保持迹的代数结论仍成立。}矩阵的迹为表达这一总方差提供了紧凑方式。

\paragraph{相似变换。}
前面已经看到，若一个新基由可逆矩阵 \(A\) 定义，则把向量表示到该基中对应坐标变换
\[
v_A = A^{-1}v.
\]
同样的思想也适用于线性变换本身。当线性变换 \(T\) 在这个新基中进行表示时，其矩阵表示变为
\[
T_A = A^{-1}TA,
\]
这称为\emph{相似变换}（similarity transformation）。它表示的仍然是\emph{同一个}线性变换，只是改用由基 \(A\) 决定的另一套坐标系统进行描述，如下推导所示。

\begin{tcolorbox}[
  colback=blue!6,
  colframe=blue!55!black,
  title={为什么 $T_A = A^{-1} T A$},
  boxrule=0.8pt,
  arc=2mm,
  left=2mm,
  right=2mm,
  top=1mm,
  bottom=1mm
]
若 \(v_A=A^{-1}v\) 表示向量 \(v\) 在由 \(A\) 定义的新基中的坐标，那么线性变换 \(T\) 在这个坐标系中也必须具有对应表示 \(T_A\)。两个矩阵描述的是\emph{同一个}几何变换，因此，在新坐标中应用 \(T_A\)，必须与先在原坐标中应用 \(T\)、再改变基得到相同结果：
\[
v'_A = T_A v_A,
\qquad
v'_A = A^{-1}(Tv).
\]
由于这两条路径必须在 \(A\) 坐标中产生同一个向量，代入 \(v_A=A^{-1}v\) 得
\[
T_A A^{-1}v = A^{-1}Tv
\quad\text{对所有 }v\text{ 成立}.
\]
两边左乘 \(A\)，有
\[
A T_A A^{-1}=T,
\qquad\Longrightarrow\qquad
T_A=A^{-1}TA.
\]
\end{tcolorbox}

\paragraph{基变换下迹的不变性。}

迹的一个关键代数性质是循环恒等式
\[
\mathrm{Tr}(BC)=\mathrm{Tr}(CB),
\]
只要矩阵乘积 \(BC\) 与 \(CB\) 均有定义即可。

结合这一性质与相似变换 \(T_A=A^{-1}TA\)，可得
\[
\mathrm{Tr}(T_A)
=
\mathrm{Tr}(A^{-1}TA)
=
\mathrm{Tr}(T).
\]

这说明迹在相似变换下保持不变。因此，虽然基变换可能重新分配方差在各个坐标上的表达方式，但不会改变总方差。

\subsection{把 PCA 理解为相似变换}

现在从基变换视角重新审视 PCA。对于单个数据点 \(x\in\mathbb{R}^d\)，低维表示通过投影到主成分方向获得，即 \(z=U_k^\top x\)。在整个数据集层面，如果把中心化后的数据点按行收集到矩阵 \(X\) 中，那么 PCA 对应线性变换 \(Z=XU\)，其中 \(U\) 是正交矩阵，其各列由协方差矩阵 \(\Sigma\) 的特征向量组成。

\medskip
\noindent
下面考察这种基变换如何改变数据的\emph{协方差结构}。

\paragraph{变换后数据的协方差。}
变换后数据的协方差矩阵为
\[
\tilde{\Sigma} = \mathrm{Cov}(Z)=\mathrm{Cov}(XU).
\]

一般地，随机向量 \(Y\) 的协方差可以写为
\[
\mathrm{Cov}(Y)
=\mathbb{E}\big[(Y-\mathbb{E}[Y])^\top(Y-\mathbb{E}[Y])\big].
\]
在当前情形下，数据已经中心化，因此 \(\mathbb{E}[X]=0\)。由于 \(U\) 是固定的线性变换，有
\[
\mathbb{E}[Z]=\mathbb{E}[XU]=\mathbb{E}[X]U=0.
\]
于是变换后数据的协方差化简为
\[
\tilde{\Sigma}
=\mathbb{E}\big[(Z-\mathbb{E}[Z])^\top(Z-\mathbb{E}[Z])\big]
=\mathbb{E}[Z^\top Z].
\]

代入 \(Z=XU\) 得
\[
\tilde{\Sigma}
=\mathbb{E}[(XU)^\top(XU)]
=U^\top\mathbb{E}[X^\top X]U
=U^\top\Sigma U.
\]

\medskip
\noindent
\textbf{PCA 之前与之后。}
现在可以明确写出这一核心结构结果：
\[
\mathrm{Cov}(X)=\Sigma,
\qquad
\mathrm{Cov}(Z)=U^\top\Sigma U.
\]
也就是说，PCA 按照
\[
\Sigma\;\longmapsto\;U^\top\Sigma U
\]
来变换协方差矩阵，这正是原协方差矩阵的一次\emph{相似变换}。

\paragraph{相似变换的结果。}
由于相似变换保持迹不变，有
\[
\mathrm{Tr}(\Sigma)
=
\mathrm{Tr}(U^\top\Sigma U)
=
\mathrm{Tr}(\Lambda),
\]
其中 \(\Sigma=U\Lambda U^\top\) 是协方差矩阵的特征分解。

这正体现了 PCA 的几何意义。PCA 之前，协方差矩阵 \(\Sigma\) 可能包含非对角项，表示不同坐标之间存在相关性。变换到特征向量基之后，协方差矩阵变成对角形式：
\[
U^\top\Sigma U
=
\Lambda
=
\mathrm{diag}(\lambda_1,\ldots,\lambda_d).
\]
因此，变换后的坐标彼此不相关，而每个主方向上的方差会以对应特征值的形式直接出现在对角线上。

PCA 并不会创造或消除方差。它只是把同样的总方差重新表达在一套与数据自然方向对齐的坐标系中。由迹表示的总方差保持不变，但其分布从相关坐标重新组织到了彼此正交的主方向上。

这也解释了为什么解释方差比具有如此自然的形式。保留前 \(k\) 个主成分，就保留方差 \(\sum_{i=1}^k\lambda_i\)，而总方差为 \(\sum_{i=1}^d\lambda_i\)。因此
\(\frac{\sum_{i=1}^k\lambda_i}{\sum_{i=1}^d\lambda_i}\)
度量的正是舍弃其余
方向之后所保留的总方差比例。

图~\ref{fig:ch1-pca-change-of-basis} 直接展示了这种坐标变换：同一数据云在主成分坐标中变得与坐标轴对齐，同时总方差保持不变。

\begin{figure}[H]
\centering
\begin{minipage}[t]{0.4\textwidth}
\vspace{0pt}\centering
\includegraphics[width=\linewidth]{figures/ch01/panels/fig1_3a_original_coordinates.pdf}
\end{minipage}\hfill
\begin{minipage}[t]{0.42\textwidth}
\vspace{0pt}\centering
\includegraphics[width=\linewidth]{figures/ch01/panels/fig1_3b_pca_coordinates.pdf}
\end{minipage}
\caption{\textbf{PCA 通过改变基来对角化协方差。}
图中分别在原始坐标和主成分坐标下展示同一组数据。在新基中，主方向成为坐标轴，各坐标之间的协方差消失，而总方差保持不变。}
\label{fig:ch1-pca-change-of-basis}
\end{figure}

\subsection{从特征分解到奇异值分解}

前面的讨论表明，PCA 可以理解为一种基变换：它把协方差矩阵对角化，并将方差重新组织到彼此正交的方向上。由此自然引出一个更一般的问题：
\begin{quote}
如果我们想理解的矩阵不是协方差矩阵，不是对称矩阵，甚至不是方阵，该怎么办？
\end{quote}
奇异值分解（Singular Value Decomposition, SVD）提供了相应的一般框架~\cite{strang2016introduction}。

特征分解建立在关系 \(Av=\lambda v\) 上，它寻找那些在矩阵作用后仍保持原方向、仅发生缩放的方向。沿这些方向，变换的行为尤其简单：输出仍处于同一条直线上，只被缩放 \(\lambda\) 倍。

这一视角很强大，但也存在限制。它只适用于方阵，并且通常还需要额外条件才能保证存在完整且性质良好的实数分解。

SVD 对这一思想进行了推广。它不再要求寻找“被矩阵作用后仍保持自身方向”的向量，而提出一个更灵活的问题：\emph{哪些相互正交的输入方向会被映射到哪些相互正交的输出方向？这些方向分别被拉伸多少？}

\paragraph{SVD 视角。}
对于实矩阵 \(A\in\mathbb{R}^{m\times n}\)，奇异值分解写为
\[
A = U S V^\top,
\]
其中，\(V\in\mathbb{R}^{n\times n}\) 包含一组标准正交的输入方向；\(U\in\mathbb{R}^{m\times m}\) 包含一组标准正交的输出方向；而 \(S\in\mathbb{R}^{m\times n}\) 是广义意义上的矩形对角矩阵，其主对角线上为非负数
\(\sigma_1\ge\sigma_2\ge\cdots\ge0\)。

这些量满足
\[
A v_i = \sigma_i u_i,
\qquad
A^\top u_i = \sigma_i v_i,
\]
其中 \(v_i\) 是右奇异向量，\(u_i\) 是对应的左奇异向量，\(\sigma_i\) 是对应的奇异值。

从几何上看，SVD 表明每个正交输入方向 \(v_i\) 会被映射为某个正交输出方向 \(u_i\)，并由奇异值 \(\sigma_i\) 决定缩放倍数。与特征分解不同，输入方向和输出方向不需要重合；这也正是为什么 SVD 对\emph{任意}实矩阵都存在，包括矩形矩阵和非对称矩阵。于是，特征分解可以看作一种特殊情形：当 \(A\) 为对称矩阵时，输入和输出方向重合（\(U=V\)），奇异值等于特征值的绝对值。\footnote{译注：原文此处的 \(U=V\) 还要求 \(A\) 为半正定矩阵。对于含负特征值的实对称矩阵，奇异值仍等于特征值的绝对值，但负特征值对应的左右奇异向量方向相反。例如 \(A=[-1]\) 的奇异值为 1，左右奇异向量不能同时取同一个单位向量。}

SVD 再次强化了本书中一个会反复出现的核心主题：一旦选择了合适的坐标系统或分解方式，复杂的线性变换往往会变得容易理解得多。PCA 对协方差矩阵所做的事情，就是旋转到特征向量基中；SVD 则把同样的思想推广到任意矩阵。

即使后续模型超越了线性变换，这种结构化的思考习惯仍然很重要。我们会不断寻找新的表示，使复杂运算能够分解成更简单、更容易解释的组成部分。从这一意义上说，SVD 不只是一个数值工具，更是 PCA 背后同一套几何思想的概念推广。

\begin{tcolorbox}[
  colback=blue!6,
  colframe=blue!55!black,
  title={秩一结构与矩阵的几何含义},
  boxrule=0.8pt,
  arc=2mm,
  left=2mm,
  right=2mm,
  top=1mm,
  bottom=1mm
]

SVD 最具启发性的解释之一，是它的秩一展开：
\[
A = \sum_{i=1}^r \sigma_i\,u_i v_i^\top,
\]
其中矩阵的秩 \(r\) 表示构造该矩阵所需的独立方向数量。

每一项 \(\sigma_i u_i v_i^\top\) 都是一个\emph{秩一变换}，由两个向量的外积构成。从几何上看，它执行一个简单的三步过程：
\[
\text{投影到 }v_i
\;\rightarrow\;
\text{按 }\sigma_i\text{ 缩放}
\;\rightarrow\;
\text{沿 }u_i\text{ 重新定向}.
\]

由于 \(v_i^\top x\) 是一个标量，输出 \(u_i v_i^\top x\) 总是单一向量 \(u_i\) 的某个标量倍数。因此所有可能输出都位于同一个固定方向上，该变换最多只能张成一维输出空间，所以秩为 1。这说明任意矩阵都可以理解为这些基本外积变换的叠加，其中每一项都捕获一种独立的作用模式。

\end{tcolorbox}

\section{从 PCA 到自编码器：重构、子空间与流形}

前面几节主要从\emph{方差最大化}的角度建立 PCA：我们寻找数据最为分散的方向，并把这些方向作为新的坐标系。这个视角揭示了 PCA 如何重新组织数据，使最有信息量的变化显式呈现出来。

不过，还有另一种同样自然且非常直观的方式来理解同一个方法。我们可以不再问“哪些方向捕获了最大的方差”，而改问一个重构问题：
\begin{quote}
\emph{如果先把数据压缩为一个 \(k\) 维表示，再从中进行重构，那么哪一种 \(k\) 维表示损失的信息最少？}
\end{quote}
这种重新表述把 PCA 从一个特征值问题转变为一个近似问题。几何结构并没有改变，但关注重点由\emph{捕获方差}转向了\emph{最小化重构误差}。这个转变很重要，因为在自编码器以及许多后续表示学习模型中，真正起组织作用的原则将不再是方差最大化，而是重构。

\subsection{从方差最大化到重构}

设 \(U_k=[\,u_1,\dots,u_k\,]\in\mathbb{R}^{d\times k}\) 表示由协方差矩阵前 \(k\) 个特征向量组成的矩阵。对于中心化数据点 \(x_i\in\mathbb{R}^d\)，PCA 给出低维表示
\[
z_i = U_k^\top x_i,
\]
它记录了 \(x_i\) 沿各主方向的坐标。从该表示进行重构得到
\[
\hat{x}_i = U_k z_i = U_k U_k^\top x_i,
\]
这恰好是 \(x_i\) 到由 \(U_k\) 各列张成的 \(k\) 维子空间上的正交投影。

有了原始点 \(x_i\) 和其重构 \(\hat{x}_i\) 后，信息损失最自然的度量就是两者之差。由此得到平方重构误差
\[
\|x_i-\hat{x}_i\|^2
=
\|x_i-U_kU_k^\top x_i\|^2.
\]
对所有数据点求和可得总重构损失
\[
L(U_k)
=
\sum_{i=1}^n \|x_i-U_kU_k^\top x_i\|^2.
\]

从这一视角看，PCA 可以写成
\[
\min_{U_k^\top U_k=I_k}
\;
\sum_{i=1}^n \|x_i-U_kU_k^\top x_i\|^2,
\]
它所回答的问题是：
\begin{quote}
\emph{哪一个具有标准正交基的 \(k\) 维子空间，能够使数据压缩后再重构所造成的失真最小？}
\end{quote}

这一重构形式与前面的 PCA 方差最大化形式完全等价，并与经典的最佳低秩近似视角密切相关~\cite{jolliffe2016pca}。虽然目标函数写法不同，但两者都会得到同一个最优 \(k\) 维子空间，即由协方差矩阵最大的若干特征值对应的特征向量所张成的子空间。

图~\ref{fig:ch1-pca-variance-geometry} 在二维情形下展示了二者的等价性：主方向既使投影后的坐标最为分散，也使正交重构残差最小。

\begin{figure}[H]
\centering
\begin{minipage}[t]{0.42\textwidth}
\vspace{0pt}\centering
\includegraphics[width=\linewidth]{figures/ch01/panels/fig1_2a_bad_direction.pdf}
\end{minipage}\hfill
\begin{minipage}[t]{0.42\textwidth}
\vspace{0pt}\centering
\includegraphics[width=\linewidth]{figures/ch01/panels/fig1_2b_pc_direction.pdf}
\end{minipage}
\caption{\textbf{PCA 将最大投影方差与最小重构误差联系起来。}
同一个中心化数据集分别投影到一个非主方向和第一主成分方向上。主成分方向具有更大的投影方差，也产生更小的正交重构残差。为使图形清晰，仅对部分样本点绘出了残差线段。}
\label{fig:ch1-pca-variance-geometry}
\end{figure}

这种视角尤其重要，因为它把 PCA 重写为一对映射：
\[
x\;\longrightarrow\;z=U_k^\top x
\;\longrightarrow\;\hat{x}=U_k z.
\]
也就是说，PCA 可以被看成一个\emph{编码器}后接一个\emph{解码器}，二者都是线性的，并且都可解析求得。

\paragraph{把自编码器理解为学习得到的编码器--解码器模型。}

当 PCA 被写成重构形式后，一个自然推广就出现了。我们不再根据特征向量解析地构造编码器和解码器，而是通过最小化重构误差，直接从数据中\emph{学习}它们，由此得到自编码器的视角~\cite{baldi1989neural,hinton2006reducing}。

形式上，自编码器引入两个带参数的映射。给定输入 \(x\in\mathbb{R}^d\)：
\begin{itemize}
    \item \textbf{编码器}把输入压缩为潜在表示
    \[
    z=q_\phi(x)\in\mathbb{R}^k,\qquad k<d,
    \]
    其中 \(q_\phi\) 由参数 \(\phi\) 参数化；

    \item \textbf{解码器}根据这一潜在编码重构输入：
    \[
    \hat{x}=p_\theta(z),
    \]
    其中 \(p_\theta\) 由参数 \(\theta\) 参数化。
\end{itemize}

二者共同定义了熟悉的编码器--解码器流程
\[
x\;\longrightarrow\;z=q_\phi(x)
\;\longrightarrow\;\hat{x}=p_\theta(z).
\]
然后通过最小化重构误差来学习参数 \((\phi,\theta)\)，例如
\[
\min_{\phi,\theta}\;
\|p_\theta(q_\phi(x))-x\|^2,
\]
或对整个数据集写成
\[
\min_{\phi,\theta}\;
\sum_{i=1}^n \|p_\theta(q_\phi(x_i))-x_i\|^2.
\]

该目标与 PCA 的重构形式直接对应：
\[
\min_{U_k^\top U_k=I_k}
\sum_{i=1}^n \|U_kU_k^\top x_i-x_i\|^2.
\]
本质区别在于：PCA 中的编码器和解码器是与标准正交特征向量绑定的固定线性映射，而自编码器则把二者视为可学习函数。

\subsection{线性自编码器在什么条件下恢复 PCA？}

线性自编码器与 PCA 之间的联系是经典结果~\cite{baldi1989neural}。不过，线性自编码器本身并不能保证自动恢复 PCA。若想精确恢复 PCA，有两个结构要素十分重要。

\paragraph{1. 线性。}
假设编码器和解码器都是线性的：
\[
q_\phi(x)=W_qx,
\qquad
p_\theta(z)=W_pz,
\]
其中
\(W_q\in\mathbb{R}^{k\times d}\)、\(W_p\in\mathbb{R}^{d\times k}\)。
那么重构为
\[
\hat{x}=W_pW_qx.
\]
此时模型仍处在 PCA 的几何框架中：所有重构结果都位于 \(\mathbb{R}^d\) 中的某个 \(k\) 维线性子空间内。

如果引入非线性激活函数，这一限制就会消失。模型此时可以表示弯曲的非线性结构，而不再局限于单一平坦子空间。因此，如果目标是保持在经典 PCA 框架内，线性是必要条件。

\paragraph{2. 子空间与基的区别。}
PCA 不仅识别一个 \(k\) 维子空间，而且使用由主特征向量构成的标准正交基来描述该子空间。相比之下，线性自编码器并不会被强制学习一个标准正交基。仅仅最小化重构误差，并不会自动施加正交性约束。

这一差异很重要，但也较为细微。对于降维来说，真正关键的几何对象不是某一组特定的基，而是这些基向量所张成的\emph{子空间}。

\paragraph{子空间与列空间。}
给定向量集合 \(\{v_1,\ldots,v_k\}\subset\mathbb{R}^d\)，它们的张成空间定义为
\[
\mathrm{Span}(v_1,\ldots,v_k)
=
\left\{
\sum_{i=1}^k \alpha_i v_i:\alpha_i\in\mathbb{R}
\right\}.
\]
不同向量集合完全可能张成同一个子空间。例如
\[
\mathrm{Span}\!\left(
\begin{bmatrix}1\\0\end{bmatrix},
\begin{bmatrix}0\\1\end{bmatrix}
\right)
=
\mathrm{Span}\!\left(
\begin{bmatrix}1\\1\end{bmatrix},
\begin{bmatrix}1\\-1\end{bmatrix}
\right)
=
\mathbb{R}^2.
\]

当向量作为矩阵 \(A\in\mathbb{R}^{d\times k}\) 的列收集起来时，它们的张成空间称为 \(A\) 的\emph{列空间}（column space，也称 range）：
\[
\mathcal{R}(A)
=
\{Az:z\in\mathbb{R}^k\}.
\]

这正是降维真正关心的几何对象。PCA 使用标准正交基 \(U_k=[\,u_1,\dots,u_k\,]\) 来表示最优的 \(k\) 维子空间；线性自编码器则可能使用另一组通常并不正交的基来表示同一个子空间。

\paragraph{重构损失下的子空间恢复。}
对于线性自编码器，
\[
\hat{x}=W_pW_qx,
\]
任意重构结果都位于 \(W_p\) 的列空间中：
\[
\hat{x} \in \mathcal{R}(W_p).
\]
如果模型在平方重构误差下得到最优训练，那么这个列空间会恢复与 PCA 相同的最优 \(k\) 维子空间，即使具体的基向量并不相同。

这一点非常关键：PCA 为最优子空间选择了一组特定的标准正交基，而线性自编码器只需要恢复子空间本身即可。从重构角度看，这已经足够。只要两组基张成同一组方向，那么无论这些基是否正交，它们表示的都是同一个几何对象。

如果之后希望得到一组标准正交基，可以对学习得到的解码器矩阵进行正交化。更一般地，任何张成某个子空间的满秩向量集合，都可以使用 Gram--Schmidt 正交化、QR 分解或相关矩阵分解方法转换为同一子空间的一组标准正交基。

例如，进行 QR 分解
\[
W_p=QR,
\]
且 \(Q^\top Q=I_k\)，即可得到矩阵 \(Q\)，其各列构成与 \(W_p\) 完全相同子空间的一组标准正交基。关键在于：正交化改变的是用于表示子空间的\emph{基}，而不是子空间本身。

\subsection{从线性子空间到非线性流形}

当编码器和解码器被允许具有非线性时，自编码器真正的表达能力才显现出来。在线性情形下，如果解码器为
\[
p_\theta(z)=W_pz,
\]
那么所有重构结果都位于同一个线性子空间
\[
\{\hat{x}=W_pz:z\in\mathbb{R}^k\}
\]
中。无论潜在编码 \(z\) 如何变化，输出始终被限制在 \(\mathbb{R}^d\) 中一张平坦的 \(k\) 维“平面”上。

如果解码器变为非线性，这一几何结构将发生根本变化。所有可能重构组成的集合变为
\[
\mathcal{M}_\theta
=
\{p_\theta(z):z\in\mathbb{R}^k\},
\]
它不再被限制为一个线性子空间，而可能形成嵌入在高维环境空间中的弯曲低维曲面。这样的曲面称为\emph{流形}（manifold）。

因此，非线性自编码器以一种深刻的方式推广了 PCA。它不再只学习单一的平坦子空间，而可以学习一个弯曲流形，从而更灵活地适应数据的内在结构。

图~\ref{fig:ch1-pca-ae-manifold} 区分了两个容易混淆的概念：线性自编码器可能使用不同的基，但仍张成与 PCA 相同的子空间；而非线性解码器则真正改变了几何结构，使重构结果可以位于弯曲流形上。

\begin{figure}[H]
\centering
\begin{minipage}[t]{0.42\textwidth}
  \vspace{0pt}\centering
  \includegraphics[width=\linewidth]{figures/ch01/panels/fig1_4a_pca_basis.pdf}
\end{minipage}\hfill
\begin{minipage}[t]{0.42\textwidth}
  \vspace{0pt}\centering
  \includegraphics[width=\linewidth]{figures/ch01/panels/fig1_4b_linear_ae_basis.pdf}
\end{minipage}

\vspace{0.8em}

\begin{minipage}[t]{0.42\textwidth}
  \vspace{0pt}\centering
  \includegraphics[width=\linewidth]{figures/ch01/panels/fig1_4c_nonlinear_manifold.pdf}
\end{minipage}\hfill
\begin{minipage}[t]{0.42\textwidth}
  \vspace{0pt}
  \caption{\textbf{从 PCA 子空间到学习得到的流形。}
  PCA 使用一组特定的标准正交基表示线性子空间。线性自编码器可以用另一组通常并不正交的解码器基恢复同一个最优子空间；即使所表示的平面一致，其基向量也不必与 PCA 相同。当解码器变为非线性后，重构集合不再局限于单一平坦子空间，而可以形成弯曲的低维流形。}
  \label{fig:ch1-pca-ae-manifold}
\end{minipage}
\end{figure}

\paragraph{非线性编码器与解码器。}
在实际应用中，非线性自编码器通常由神经网络实现。编码器可以写成
\[
z=q_\phi(x),
\]
其中 \(q_\phi\) 由仿射变换与非线性激活的复合构成；解码器类似地写为
\[
\hat{x}=p_\theta(z).
\]
整体结构依然是
\[
x\;\longrightarrow\;z\;\longrightarrow\;\hat{x},
\]
但允许的变换类别已经远比 PCA 丰富。

这种从线性子空间到非线性流形的转变，是自编码器成为现代表示学习重要工具的关键原因之一。此时，PCA 可以看作一个更大范围的“基于重构的模型家族”中最简单的成员：线性、正交，并且具有解析解。

\paragraph{为什么确定性自编码器还不是生成模型。}
尽管确定性自编码器具有更强的表达能力，但它们仍缺少生成建模的一个关键要素：
\begin{quote}
\emph{对潜变量如何生成数据的概率描述。}
\end{quote}

确定性自编码器把每个输入 \(x\) 映射到一个唯一的潜在编码 \(z\)，再映射到一个唯一的重构 \(\hat{x}\)。这并没有提供有原则的方法来
\begin{itemize}
    \item 为潜在表示分配概率；
    \item 从指定分布中抽取潜变量以生成新样本；
    \item 或者依据基于似然的原则训练模型。
\end{itemize}

为了获得这些能力，我们必须超越确定性重构模型，引入\emph{概率潜变量模型}。这一转变会直接引向概率主成分分析（PPCA），随后再引向变分自编码器（VAE）；在 VAE 中，重构思想被保留下来，但被嵌入到了概率框架之中。

从这一视角看，后续路径已经十分清楚。PCA 通过方差最大化引入线性子空间；PCA 的重构形式自然引出自编码器；线性自编码器恢复与 PCA 相同的最优子空间，而非线性自编码器把这一思想推广到流形。但要真正进入现代生成建模，我们还必须再前进一步：为潜在表示赋予概率分布。

下一章将从概率主成分分析开始完成这一转变。

\section*{本章小结与参考文献}
\addcontentsline{toc}{section}{本章小结与参考文献}

本章介绍了 PCA 所需的线性代数思想，并以此建立从经典降维方法到自编码器的桥梁。

我们首先把矩阵理解为线性变换，并强调基变换、点积和投影是 PCA 背后的几何语言。由此自然引出特征向量与特征值，它们标识矩阵作用最为简单的一组特殊方向。

借助这些工具，我们先把 PCA 建立为方差最大化问题，随后又把它解释为一种将协方差矩阵对角化的基变换。这两个视角说明了为什么主成分就是最大的若干特征值对应的特征向量、为什么特征值能够度量所捕获的方差，以及 PCA 如何通过投影和重构实现降维。

随后，本章把关注点从方差转向重构。把 PCA 重新表述为线性编码器--解码器模型后，它便与自编码器建立了直接联系。线性自编码器能够恢复与 PCA 相同的最优子空间，而非线性自编码器则把这一思想从平坦子空间推广到弯曲流形。

\bigskip

本章内容建立在线性代数、矩阵分解和降维的标准理论之上。关于底层线性代数与矩阵视角，Strang~\cite{strang2016introduction} 是一部广泛适用的参考资料；Jolliffe 和 Cadima~\cite{jolliffe2016pca} 则给出了 PCA 及其解释的现代综述。就历史背景而言，PCA 的经典起源可以追溯到 Pearson~\cite{pearson1901} 和 Hotelling~\cite{hotelling1933}。Baldi 和 Hornik~\cite{baldi1989neural} 讨论了线性自编码器与 PCA 之间的联系，而 Hinton 和 Salakhutdinov~\cite{hinton2006reducing} 则对深度自编码器的发展产生了重要影响。

\chapter[概率主成分分析：通向潜变量生成建模的桥梁]{概率主成分分析：\\通向潜变量生成建模的桥梁}

\section*{概述}
\addcontentsline{toc}{section}{概述}

从确定性表示学习走向概率生成建模之后，概率主成分分析（Probabilistic Principal Component Analysis, PPCA）自然成为继 PCA 和自编码器之后的下一步，主要有两个原因：

\begin{itemize}
    \item \textbf{引入概率潜变量建模。}  
    上一章中的 PCA 和自编码器展示了如何利用低维结构表示数据，但二者从根本上仍是确定性的。PPCA 则引入了本书后续将反复出现的关键概率要素：潜变量的先验、给定潜在原因时观测量的条件分布、用于学习的边缘似然，以及用于推断的后验分布。从这个意义上说，PPCA 是本书中的第一个概率潜变量模型。

    \item \textbf{为现代生成模型提供一个解析可解的桥梁。}  
    由于 PPCA 同时具有线性结构和高斯分布，其所有核心量仍然可以闭式推导。这种少见的可解性使它成为理想的过渡模型：我们可以在每一步都保持数学透明的环境中，研究基于似然的学习、后验推断、Jensen 不等式、证据下界（Evidence Lower Bound, ELBO）以及期望最大化（Expectation Maximisation, EM）算法。后续在表达能力更强、但更难处理的模型（例如变分自编码器）中，这些思想都会再次出现。
\end{itemize}

\section*{概念导图}
\addcontentsline{toc}{section}{概念导图}

本章按照一条从确定性降维走向概率潜变量建模的路径展开。
\begin{itemize}

    \item[\(\rightarrow\)] \textbf{首先把 PPCA 表述为 PCA 的概率扩展。}  
    我们引入一个潜变量，把它视为观测数据背后隐藏的低维原因；随后加入高斯噪声，把确定性的 PCA 重构变成对可能观测值的完整生成分布。

    \item[\(\rightarrow\)] \textbf{随后推导观测数据的边缘似然。}  
    由于 PPCA 是线性高斯模型，可以精确地对潜变量进行积分消去，从而得到观测数据本身也服从高斯分布；其协方差可以明确区分潜在结构与各向同性噪声。

    \item[\(\rightarrow\)] \textbf{有了显式边缘似然后，再比较三种优化视角。}  
    PPCA 可以通过解析最大似然、通用的基于梯度的优化，或者 EM 算法进行学习。这种比较一方面说明 PPCA 为什么具有异常良好的可解性，另一方面也说明为什么从概念上说，EM 才是连接后续内容最重要的路径。

    \item[\(\rightarrow\)] \textbf{随后从精确可解的 PPCA 走向一般潜变量情形。}  
    我们引入 Jensen 不等式，把它作为将难以计算的对数似然替换为可计算下界的关键数学工具，并由此得到 ELBO。

    \item[\(\rightarrow\)] \textbf{接下来把 EM 推导为迭代最大化这一证据下界的方法。}  
    E 步更新潜变量分布，M 步更新模型参数。在 PPCA 中，这一过程会恢复与精确解析推导相同的解，但此时我们可以把它理解为一个更一般的推断原则的特殊情形。

    \item[\(\rightarrow\)] \textbf{最后识别 PPCA 中仍然存在的一个局限，并由此引出下一章。}  
    尽管 PPCA 为潜变量引入了先验，但它并没有显式正则化推断得到的潜在空间，使其与先验保持一致。这个局限会直接引向变分自编码器：在 VAE 中，概率推断与潜在空间正则化被统一到同一个学习目标中。

\end{itemize}

\medskip
\noindent PPCA 的主要概率关系概括如下图所示。
\begin{center}
\begin{tikzpicture}[
  >=Latex,
  every node/.style={font=\small},
  box/.style={
    draw,
    rounded corners=2pt,
    align=center,
    inner sep=7pt,
    minimum height=1.20cm,
    text width=3.75cm
  },
  lowerbox/.style={
    draw,
    rounded corners=2pt,
    align=center,
    inner sep=7pt,
    minimum height=1.12cm,
    text width=3.75cm
  },
  flow/.style={->,thick}
]

\node[box] (prior) at (0,0) {
  \textbf{潜变量先验}\\[2pt]
  $\mathbf{z}\sim\mathcal{N}(\mathbf{0},\mathbf{I})$
};

\node[box] (gen) at (5.15,0) {
  \textbf{条件生成}\\[2pt]
  $p(\mathbf{x}\mid\mathbf{z})
  =\mathcal{N}(\mathbf{W}\mathbf{z}+\boldsymbol{\mu},
  \alpha^2\mathbf{I})$
};

\node[box] (marginal) at (10.30,0) {
  \textbf{边缘数据模型}\\[2pt]
  $p(\mathbf{x})
  =\mathcal{N}\!\bigl(
  \boldsymbol{\mu},
  \mathbf{W}\mathbf{W}^{\top}+\alpha^2\mathbf{I}
  \bigr)$
};

\node[lowerbox] (posterior) at (5.15,-2.75) {
  \textbf{后验推断}\\[2pt]
  $p(\mathbf{z}\mid\mathbf{x})
  \propto p(\mathbf{x}\mid\mathbf{z})p(\mathbf{z})$\\[-1pt]
  在 PPCA 中为\\精确高斯分布
};

\node[lowerbox] (learning) at (10.30,-2.75) {
  \textbf{参数学习}\\[2pt]
  最大化 $\log p(\mathbf{x})$\\[-1pt]
  直接优化或使用 ELBO / EM
};

\draw[flow] (prior) -- (gen);
\draw[flow] (gen) -- (marginal);
\draw[flow] (gen) -- node[right] {\scriptsize Bayes} (posterior);
\draw[flow] (marginal) -- node[right] {\scriptsize 似然} (learning);
\draw[flow] (posterior) -- node[above] {\scriptsize EM} (learning);

\node[font=\normalsize\bfseries] at (5.15,1.45) {PPCA 概览};

\end{tikzpicture}
\end{center}

\medskip

\section{概率主成分分析的形式化}

上一章中，PCA 和自编码器引入了这样一种思想：高维数据可以通过低维潜在结构进行表示。然而，二者从根本上仍然是\emph{确定性的}：一旦选择某个潜在表示，对应的重构也随之固定。迈向生成建模的关键一步，是超越这种确定性图景，使一个潜在点不再只对应单一观测，而是对应一个可能观测值的\emph{分布}。

概率主成分分析（PPCA）~\cite{tipping1999ppca} 是本书中最简单、也最明确地完成这一转变的模型。它保留了 PCA 的线性潜在结构，但把它嵌入概率生成过程中，因此成为我们遇到的第一个具体潜变量生成模型。

\paragraph{把 PPCA 理解为潜变量生成模型。}
在正式给出模型之前，先说明\emph{潜变量}（latent variable）这个术语。“latent”最直接的意思就是“隐藏的”：它指一个无法被直接观测、但被假设存在于数据背后并解释其部分结构的变量。在 PPCA 中，观测向量 \(\mathbf{x}\) 是我们实际看到的数据点，而 \(\mathbf{z}\) 是一个较低维、无法直接观测的隐藏表示。模型假定 \(\mathbf{x}\) 是由 \(\mathbf{z}\) 生成的，而不是把 \(\mathbf{x}\) 看成孤立对象。从“只处理观测数据”转向“观测数据由隐藏原因生成”，正是潜变量生成建模的定义性思想之一。

PPCA 的生成过程包含两步：
\begin{enumerate}
    \item 从标准高斯先验中采样潜向量
    \[
        \mathbf{z}\sim\mathcal{N}(\mathbf{0},\mathbf{I}).
    \]

    \item 对其施加线性变换、加入均值偏置，再叠加高斯噪声 \(\boldsymbol{\epsilon}\) 生成观测：
    \[
        \mathbf{x}=\mathbf{W}\mathbf{z}+\boldsymbol{\mu}+\boldsymbol{\epsilon},
        \qquad
        \boldsymbol{\epsilon}\sim\mathcal{N}(\mathbf{0},\alpha^2\mathbf{I}).
    \]
\end{enumerate}

其中：
\begin{itemize}
    \item \(\mathbf{z}\in\mathbb{R}^k\) 是潜变量；
    \item \(\mathbf{x}\in\mathbb{R}^d\) 是观测数据点；
    \item \(\mathbf{W}\in\mathbb{R}^{d\times k}\) 把潜在坐标映射到观测空间；
    \item \(\boldsymbol{\mu}\in\mathbb{R}^d\) 是数据均值；
    \item \(\alpha^2\) 控制各向同性观测噪声。
\end{itemize}

这个模型已经包含了后续许多生成模型的基本组成要素：潜变量上的先验、从潜变量映射到观测的生成规则，以及对可能输出的概率分布。与 PCA 相比，关键的新意在于：潜在表示不再精确决定一个单点，而是决定一个合理观测分布的\emph{中心}。

\paragraph{观测的条件分布。}

虽然生成方程
\[
\mathbf{x}=\mathbf{W}\mathbf{z}+\boldsymbol{\mu}+\boldsymbol{\epsilon}
\]
看起来与确定性映射相似，但由于包含随机噪声项，其含义从根本上是概率性的。给定潜向量 \(\mathbf{z}\) 后，观测服从高斯分布：
\[
\mathbf{x}\mid\mathbf{z}
\sim
\mathcal{N}\!\bigl(\mathbf{W}\mathbf{z}+\boldsymbol{\mu},\;\alpha^2\mathbf{I}\bigr).
\]

这正是核心的概率化转变：一旦固定 \(\mathbf{z}\)，\(\mathbf{W}\mathbf{z}+\boldsymbol{\mu}\) 决定的是观测分布的\emph{均值}，而不是单一重构结果。图~\ref{fig:ppca-generative-story} 用 \(k=1,d=2\) 的简单情形展示了这一几何结构。随着 \(z\) 变化，条件均值在观测空间中描出一条直线；各向同性高斯噪声则在每个均值周围产生可能观测。因而，潜变量控制沿低维子空间的位置，而噪声允许观测偏离该子空间。

\begin{figure}[H]
\centering
\begin{minipage}[t]{0.45\textwidth}
  \vspace{0pt}\centering
  \includegraphics[width=\linewidth]{figures/ch02/panels/fig2_1a_prior.pdf}
\end{minipage}\hfill
\begin{minipage}[t]{0.45\textwidth}
  \vspace{0pt}\centering
  \includegraphics[width=\linewidth]{figures/ch02/panels/fig2_1b_means.pdf}
\end{minipage}

\vspace{0.8em}

\begin{minipage}[t]{0.45\textwidth}
  \vspace{0pt}\centering
  \includegraphics[width=\linewidth]{figures/ch02/panels/fig2_1c_conditionals.pdf}
\end{minipage}\hfill
\begin{minipage}[t]{0.45\textwidth}
  \vspace{0pt}
  \caption{\textbf{从潜在坐标到概率观测。}
  图中从先验 $p(z)=\mathcal{N}(0,1)$ 中选取若干潜在坐标，将其映射为观测空间中的条件均值 $\mathbf{W}z+\boldsymbol{\mu}$，再以这些均值为中心构造各向同性条件分布 $p(\mathbf{x}\mid z)=\mathcal{N}(\mathbf{W}z+\boldsymbol{\mu},\alpha^2\mathbf{I})$。子图 (c) 展示的是在三个选定潜变量值条件下的样本，而不是边缘分布 $p(\mathbf{x})$。}
  \label{fig:ppca-generative-story}
\end{minipage}
\end{figure}


\paragraph{与 PCA 的联系：零噪声极限。}

此时自然会问：PPCA 与经典 PCA 到底是什么关系？最清晰的理解方式是考察零噪声极限。回顾 PPCA 模型：
\[
\mathbf{x}=\mathbf{W}\mathbf{z}+\boldsymbol{\mu}+\boldsymbol{\epsilon},
\qquad
\mathbf{z}\sim\mathcal{N}(\mathbf{0},\mathbf{I}),
\qquad
\boldsymbol{\epsilon}\sim\mathcal{N}(\mathbf{0},\alpha^2\mathbf{I}).
\]
当观测噪声趋近于零，即 \(\alpha^2\to0\) 时，条件分布
\[
p(\mathbf{x}\mid\mathbf{z})
=
\mathcal{N}\!\bigl(\mathbf{x}\mid\mathbf{W}\mathbf{z}+\boldsymbol{\mu},\;\alpha^2\mathbf{I}\bigr)
\]
会越来越集中在其均值附近。在极限情况下，随机性消失，每个潜在坐标 \(\mathbf{z}\) 都映射为唯一的确定性点：
\[
\mathbf{x}=\mathbf{W}\mathbf{z}+\boldsymbol{\mu}.
\]

从这个意义上说，可以把 PCA 理解为 PPCA 的零噪声极限。潜在子空间仍然存在，但围绕它的概率扩散消失。因而 PPCA 把 PCA 作为一个特殊极限情形包含在内，同时引入了一个本质上新的要素：重构周围的不确定性。

下一节将看到 PPCA 仍然具有异常良好的可解性：由于模型是线性高斯的，其边缘似然 \(p(\mathbf{x})\) 仍可解析推导。

\section{为什么 PPCA 可以解析求解}

把 PPCA 表述成潜变量生成模型后，下一个问题就是如何根据数据拟合它。在 PCA 中，我们寻找最大方差方向；在自编码器中，我们最小化重构误差。相比之下，PPCA 是一个概率模型，因此最自然的拟合原则是最大化模型对观测数据赋予的似然。

\subsection{从似然到边缘化}
假设观测到数据集
\[
\{\mathbf{x}_1,\dots,\mathbf{x}_N\}.
\]
PPCA 下数据集的对数似然为
\[
\sum_{i=1}^N\log p(\mathbf{x}_i).
\]
通常也可以除以 \(N\)，写成平均对数似然
\[
\frac{1}{N}\sum_{i=1}^N\log p(\mathbf{x}_i),
\]
还可以进一步写成关于经验数据分布的期望：
\[
\mathbb{E}_{q(\mathbf{x})}\big[\log p(\mathbf{x})\big].
\]
这里的 \(q(\mathbf{x})\) 仅表示经验分布，它为每一个观测数据点分配相同概率质量。因此，对 \(q(\mathbf{x})\) 取期望，本质上就是对训练样本求平均。后续更一般的概率建模情形中，这种记号会频繁出现。

因此，PPCA 的目标为
\[
\max_{\mathbf{W},\boldsymbol{\mu},\alpha}
\;
\mathbb{E}_{q(\mathbf{x})}\big[\log p(\mathbf{x};\mathbf{W},\boldsymbol{\mu},\alpha)\big].
\]
这里，边缘似然依赖模型参数 \((\mathbf{W},\boldsymbol{\mu},\alpha)\)。不过，为避免记号过于繁重，在含义明确时我们通常会省略这种依赖，仅写成 \(p(\mathbf{x})\)。

这里的关键量是边缘似然 \(p(\mathbf{x})\)。根据全概率公式，它通过对潜变量积分消去得到：
\[
p(\mathbf{x})
=
\int p(\mathbf{x}\mid\mathbf{z})\,p(\mathbf{z})\,d\mathbf{z}.
\]

\begin{tcolorbox}[
  colback=blue!6,
  colframe=blue!55!black,
  title={为什么最大化对数似然而不是似然？},
  boxrule=0.8pt,
  arc=2mm,
  left=2mm,
  right=2mm,
  top=1mm,
  bottom=1mm
]
最大化\emph{对数}似然而非似然本身，既能提高数值稳定性，也能简化代数处理。对于独立数据点，
\[
L=\prod_{i=1}^N p(\mathbf{x}_i),
\qquad
\log L=\sum_{i=1}^N\log p(\mathbf{x}_i).
\]
对数把许多很小概率值的乘积变成求和，从而避免数值下溢，同时得到更易分析和优化的加法形式：
\[
\frac{1}{N}\sum_{i=1}^N\log p(\mathbf{x}_i)
=
\mathbb{E}_{q(\mathbf{x})}\big[\log p(\mathbf{x})\big].
\]
\end{tcolorbox}

乍看之下，这个边缘似然似乎很难计算。代入 PPCA 模型得
\[
p(\mathbf{x})
=
\int
\mathcal{N}\!\bigl(\mathbf{x}\mid\mathbf{W}\mathbf{z}+\boldsymbol{\mu},\alpha^2\mathbf{I}\bigr)
\;
\mathcal{N}\!\bigl(\mathbf{z}\mid\mathbf{0},\mathbf{I}\bigr)
\,d\mathbf{z}.
\]
由于 \(\mathbf{z}\) 遍历连续潜在空间，看起来精确计算该积分应该十分困难。对于大多数潜变量模型而言，这种直觉确实是正确的。PPCA 的特殊之处恰恰在于，它的线性高斯结构使这一边缘化仍然可以解析完成。

\begin{tcolorbox}[
  colback=blue!5!white,
  colframe=blue!75!black,
  title={为什么线性高斯模型易于解析处理？},
  boxrule=0.8pt,
  arc=2mm,
  left=2mm,
  right=2mm,
  top=1mm,
  bottom=1mm
]
\begin{itemize}
    \item \textbf{线性变换保持高斯性。}  
    如果一个随机向量服从高斯分布，那么它的任何仿射变换仍然服从高斯分布。

    \item \textbf{独立高斯变量之和仍为高斯。}  
    如果两个高斯随机变量彼此独立，那么它们的和仍然服从高斯分布。

    \item \textbf{对高斯潜变量积分消去后仍保持高斯。}  
    如果 \((\mathbf{x},\mathbf{z})\) 的联合分布是高斯分布，那么对 \(\mathbf{z}\) 做边缘化后，\(\mathbf{x}\) 的分布仍然是高斯的。
\end{itemize}
\end{tcolorbox}

借助这些事实，我们无需对积分进行暴力计算，就能够直接识别 \(\mathbf{x}\) 的分布。

\subsection{推导并解释边缘分布}

回顾 PPCA 的生成方程：
\[
\mathbf{x}=\mathbf{W}\mathbf{z}+\boldsymbol{\mu}+\boldsymbol{\epsilon},
\qquad
\mathbf{z}\sim\mathcal{N}(\mathbf{0},\mathbf{I}),
\qquad
\boldsymbol{\epsilon}\sim\mathcal{N}(\mathbf{0},\alpha^2\mathbf{I}),
\]
其中假设潜变量 \(\mathbf{z}\) 与噪声项 \(\boldsymbol{\epsilon}\) 相互独立。这说明 \(\mathbf{x}\) 是通过对一个高斯潜变量做仿射变换，再加上独立高斯噪声得到的。由于高斯变量的仿射变换仍为高斯，而独立高斯变量的和仍为高斯，因此 \(\mathbf{x}\) 本身也必然服从高斯分布。

同一个边缘分布也可以通过展开联合密度
\[
p(\mathbf{x},\mathbf{z})=p(\mathbf{x}\mid\mathbf{z})p(\mathbf{z})
\]
再利用所谓的\emph{配方法}（completing the square）对 \(\mathbf{z}\) 积分得到。附录~\ref{app:completing-square} 回顾了这一代数路径，并展示高斯联合分布如何导出高斯边缘分布和后验分布。这里先使用更简单的、基于高斯变量仿射变换的分布论证。

现在可以直接求出其均值与协方差。

\paragraph{均值。}
取期望得到
\[
\mathbb{E}[\mathbf{x}]
=
\mathbb{E}[\mathbf{W}\mathbf{z}+\boldsymbol{\mu}+\boldsymbol{\epsilon}]
=
\mathbf{W}\,\mathbb{E}[\mathbf{z}]+\boldsymbol{\mu}+\mathbb{E}[\boldsymbol{\epsilon}]
=
\boldsymbol{\mu},
\]
因为 \(\mathbb{E}[\mathbf{z}]=\mathbf{0}\)，且 \(\mathbb{E}[\boldsymbol{\epsilon}]=\mathbf{0}\)。

\paragraph{协方差。}
由
\[
\mathbf{x}=\mathbf{W}\mathbf{z}+\boldsymbol{\mu}+\boldsymbol{\epsilon},
\]
并注意常数偏置 \(\boldsymbol{\mu}\) 不影响协方差，有
\[
\operatorname{Cov}[\mathbf{x}]
=
\operatorname{Cov}[\mathbf{W}\mathbf{z}+\boldsymbol{\epsilon}].
\]

更显式地展开时，回顾一般恒等式
\[
\operatorname{Cov}[a+b]
=
\operatorname{Cov}[a]
+
\operatorname{Cov}[b]
+
\operatorname{Cov}[a,b]
+
\operatorname{Cov}[b,a].
\]
令 \(a=\mathbf{W}\mathbf{z}\)、\(b=\boldsymbol{\epsilon}\)，得到
\[
\operatorname{Cov}[\mathbf{x}]
=
\operatorname{Cov}[\mathbf{W}\mathbf{z}]
+
\operatorname{Cov}[\boldsymbol{\epsilon}]
+
\operatorname{Cov}[\mathbf{W}\mathbf{z},\boldsymbol{\epsilon}]
+
\operatorname{Cov}[\boldsymbol{\epsilon},\mathbf{W}\mathbf{z}].
\]

由于 \(\mathbf{z}\) 与 \(\boldsymbol{\epsilon}\) 相互独立，交叉协方差项为零：
\[
\operatorname{Cov}[\mathbf{W}\mathbf{z},\boldsymbol{\epsilon}]=\mathbf{0},
\qquad
\operatorname{Cov}[\boldsymbol{\epsilon},\mathbf{W}\mathbf{z}]=\mathbf{0}.
\]
因此
\[
\operatorname{Cov}[\mathbf{x}]
=
\operatorname{Cov}[\mathbf{W}\mathbf{z}]
+
\operatorname{Cov}[\boldsymbol{\epsilon}].
\]

再使用线性变换下协方差的标准公式：
\[
\operatorname{Cov}[\mathbf{W}\mathbf{z}]
=
\mathbf{W}\operatorname{Cov}[\mathbf{z}]\,\mathbf{W}^\top.
\]
所以
\[
\operatorname{Cov}[\mathbf{x}]
=
\mathbf{W}\operatorname{Cov}[\mathbf{z}]\,\mathbf{W}^\top
+
\operatorname{Cov}[\boldsymbol{\epsilon}].
\]

最后代入
\[
\operatorname{Cov}[\mathbf{z}]=\mathbf{I},
\qquad
\operatorname{Cov}[\boldsymbol{\epsilon}]=\alpha^2\mathbf{I},
\]
得到
\[
\operatorname{Cov}[\mathbf{x}]
=
\mathbf{W}\mathbf{W}^\top+\alpha^2\mathbf{I}.
\]

综合均值与协方差，观测的边缘分布为
\[
p(\mathbf{x})
=
\mathcal{N}\!\bigl(
\mathbf{x}\mid
\boldsymbol{\mu},
\mathbf{W}\mathbf{W}^\top+\alpha^2\mathbf{I}
\bigr).
\]

这是 PPCA 的关键解析结果：潜变量已经被精确积分消去，而得到的观测边缘分布仍然是高斯分布。

\paragraph{协方差结构的解释。}
项 \(\mathbf{W}\mathbf{W}^\top\) 表示潜在变化经映射后在观测空间中贡献的协方差，因为
\[
\operatorname{Cov}(\mathbf{W}\mathbf{z})
=
\mathbf{W}\operatorname{Cov}(\mathbf{z})\mathbf{W}^\top
=
\mathbf{W}\mathbf{W}^\top.
\]
它是低秩矩阵：
\[
\operatorname{rank}(\mathbf{W}\mathbf{W}^\top)\le k,
\]
因此由潜变量诱导的变化最多只分布在 \(d\) 维观测空间中的一个 \(k\) 维子空间内。若
\[
\mathbf{W}=[\,w_1\;\;w_2\;\;\cdots\;\;w_k\,],
\qquad
\mathbf{W}\mathbf{W}^\top=\sum_{j=1}^k w_jw_j^\top,
\]
那么这一结构化协方差可以看成与各个潜在方向对应的一组秩一贡献之和。

相对而言，\(\alpha^2\mathbf{I}\) 在观测空间每个方向上贡献相同方差。因此，PPCA 同时建模两类变化：通过 \(\mathbf{W}\mathbf{W}^\top\) 描述\emph{沿潜在子空间各方向}的变化，通过各向同性噪声描述子空间\emph{周围}的偏离。图~\ref{fig:ppca-covariance-decomposition} 直观总结了这一分解，将潜变量诱导的结构与各向同性观测噪声明确区分开。

\begin{figure}[H]
\centering
\begin{minipage}[t]{0.48\textwidth}
  \vspace{0pt}\centering
  \includegraphics[width=\linewidth]{figures/ch02/panels/fig2_2a_latent_cov.pdf}
\end{minipage}\hfill
\begin{minipage}[t]{0.48\textwidth}
  \vspace{0pt}\centering
  \includegraphics[width=\linewidth]{figures/ch02/panels/fig2_2b_noise_cov.pdf}
\end{minipage}

\vspace{0.8em}

\begin{minipage}[t]{0.48\textwidth}
  \vspace{0pt}\centering
  \includegraphics[width=\linewidth]{figures/ch02/panels/fig2_2c_total_cov.pdf}
\end{minipage}\hfill
\begin{minipage}[t]{0.48\textwidth}
  \vspace{0pt}
  \caption{\textbf{PPCA 中的协方差分解。}
  在图示的 $k=1,d=2$ 情形中，$\mathbf{W}\mathbf{W}^{\top}$ 为秩一矩阵，只沿 $\mathrm{span}(\mathbf{W})$ 贡献方差；$\alpha^2\mathbf{I}$ 则是各向同性且满秩的。二者之和给出满秩的边缘协方差，同时描述沿潜在子空间及其周围的变化。实线与虚线分别表示一个和两个标准差的范围。}
  \label{fig:ppca-covariance-decomposition}
\end{minipage}
\end{figure}


\bigskip

下一节从 \(p(\mathbf{x})\) 的形式转向学习问题。既然边缘似然已经可以显式写出，我们就可以讨论如何根据数据估计参数 \((\mathbf{W},\boldsymbol{\mu},\alpha)\)。

\section{学习 PPCA：三种优化视角}

我们已经推导出边缘分布
\[
p(\mathbf{x})
=
\mathcal{N}\!\bigl(
\mathbf{x}\mid
\boldsymbol{\mu},
\mathbf{W}\mathbf{W}^\top+\alpha^2\mathbf{I}
\bigr),
\]
现在可以从“表示”转向“学习”。PPCA 的参数 \((\mathbf{W},\boldsymbol{\mu},\alpha)\) 通过最大化观测数据集的对数似然进行拟合：
\[
\max_{\mathbf{W},\boldsymbol{\mu},\alpha}
\;
\mathbb{E}_{q(\mathbf{x})}\big[\log p(\mathbf{x})\big].
\]

这时自然会产生一个问题：
\begin{quote}
\emph{既然似然已经能够显式计算，应当怎样优化它？}
\end{quote}
有三种有用的理解方式。第一种利用 PPCA 特有的解析结构；第二种使用通用的基于梯度的优化；第三种是期望最大化（EM），它尤其重要，因为它揭示了潜变量模型更深层的原理，而这些原理随后会通向变分推断与 VAE。

\subsection{直接解析最大似然}

由于 \(\mathbf{x}\) 的边缘分布是高斯分布，PPCA 具有一条异常简洁的学习路径。原则上，可以把对数似然显式写出，对各参数求导，把导数设为零，再求解驻点：
\[
\nabla_{\mathbf{W}}\log p(\mathbf{x})=0,
\qquad
\nabla_{\boldsymbol{\mu}}\log p(\mathbf{x})=0,
\qquad
\frac{\partial}{\partial\alpha}\log p(\mathbf{x})=0.
\]

这就是经典的最大似然策略。它的吸引力很直观：如果这些方程可以精确求解，那么模型参数就能够直接写成数据经验协方差结构的函数。

对 PPCA 而言，这一路径之所以可行，是因为从始至终模型都保持线性高斯结构。最终解可以闭式写出，并与样本协方差矩阵的特征结构紧密相关。这也是 PPCA 具有重要教学价值的原因之一：它是一个真正的潜变量生成模型，却仍然拥有精确的最大似然解。

与此同时，这条路径也暴露出一个局限。直接解析优化虽然优雅，却高度依赖于模型足够简单，导数方程能够精确求解。一旦潜变量模型变为非线性或非高斯，这种闭式求解方式通常就会失效。而且，即使解析解存在，获得它也可能需要矩阵求逆、特征分解等计算开销较大的操作，尤其是在高维情形下。

\subsection{基于梯度的优化}

第二种方式，是把对数似然直接看成目标函数，再利用梯度上升进行数值优化。此时参数按如下方式迭代更新：
\[
\theta
\leftarrow
\theta+\eta\,\nabla_\theta\mathbb{E}_{q(\mathbf{x})}\big[\log p(\mathbf{x};\theta)\big],
\]
其中 \(\theta=\{\mathbf{W},\boldsymbol{\mu},\alpha\}\)，\(\eta\) 为步长。

这是最通用的优化视角。与直接解析最大似然不同，它不要求最优参数本身可以闭式求得，而是把 PPCA 当作任意可微概率模型，通过梯度逐步爬升似然曲面。

在 PPCA 中，边缘似然 \(p(\mathbf{x})\) 可以显式计算，因此其梯度也可以显式计算。与解析路线的区别，并不在于梯度法绕开了似然本身，而在于它不必把整个优化问题精确地闭式求解。更一般地，即使参数没有精确闭式解，只要目标函数或其梯度能够计算——或者至少能够近似——仍然可以使用基于梯度的方法。

这一点值得强调，因为它预示了现代大多数生成模型的训练方式。后续章节中，一旦闭式解不再存在，基于梯度的优化就会成为默认的计算工具。

不过，对 PPCA 这个特例而言，梯度上升并不是最具启发性的路径。它当然可用，但没有利用模型特殊的潜变量结构，也没有揭示为什么处理潜变量模型时，人们往往会显式进行后验推断，而不只是直接优化目标。要理解这一点，需要第三种视角。

\subsection{期望最大化（EM）}

第三条路径是期望最大化（Expectation--Maximisation, EM）算法~\cite{dempster1977em}。与直接解析优化不同，EM 并不试图一步求解似然；与通用梯度上升不同，它也不会把潜变量仅仅当成黑箱目标函数中的隐藏干扰项。相反，它显式利用模型的\emph{潜变量结构}。

基本思想很简单。边缘对数似然
\[
\log p(\mathbf{x})
=
\log\int p(\mathbf{x},\mathbf{z})\,d\mathbf{z}
\]
可能很难直接优化，但联合分布
\[
p(\mathbf{x},\mathbf{z})
\]
往往具有简单得多的形式。EM 通过在以下两个任务之间交替来利用这一点：
\begin{itemize}
    \item 在当前参数下估计潜变量；
    \item 利用这些估计更新模型参数。
\end{itemize}

对 PPCA 而言，EM 并非严格必需，因为模型本身已经能够解析求解。但它仍然极其重要：它揭示了适用于潜变量模型的一般优化原则，而即使精确解不再存在，这一原则仍然成立。

从更广的角度看，PPCA 之所以特别有价值，是因为它属于少数几个所有核心概率量都能显式写出的潜变量模型：先验 \(p(\mathbf{z})\)、条件分布 \(p(\mathbf{x}\mid\mathbf{z})\)、边缘似然 \(p(\mathbf{x})\) 和后验 \(p(\mathbf{z}\mid\mathbf{x})\)。在更复杂模型中，这些量仍然处于核心位置，但不再都具有解析可解性。尤其是，学习由边缘似然 \(p(\mathbf{x})\) 驱动，而推断依赖后验 \(p(\mathbf{z}\mid\mathbf{x})\)。这正是 EM 如此重要的原因：它提供了一种结构化方法，在潜变量问题的这两个方面之间交替。

对后续内容更重要的是，EM 会自然引出两个更深层思想：在潜变量上引入辅助分布，以及用一个可以迭代优化的下界替代难处理的对数似然。这些思想并不属于 PPCA 专有，而是潜变量建模一般理论的一部分；后续讨论变分推断时，它们会以更一般的形式再次出现。

因此，下面转向 EM 背后的数学基础：Jensen 不等式与证据下界（ELBO）。

\section{Jensen 不等式与证据下界}

为了理解 EM 为什么如此重要，现在需要暂时从 PPCA 退一步，提出一个更一般的问题。当边缘似然
\[
p(\mathbf{x};\theta)=\int p(\mathbf{x},\mathbf{z};\theta)\,d\mathbf{z}
\]
不再具有解析可解性时，该怎么办？PPCA 的特殊之处在于，这个积分仍然能够精确计算。但大多数潜变量模型并没有这种便利。在更一般的情形下，潜变量可能通过非线性变换、非高斯条件分布，或者二者共同进入模型，此时边缘似然就无法再写出闭式结果。

这正是 Jensen 不等式发挥作用的地方。它提供了一种系统方法，把困难的对数似然替换成一个可计算的下界。这个下界就是\emph{证据下界}（Evidence Lower Bound, ELBO），也是 EM 以及后续章节中更多生成模型背后的核心思想之一。

\subsection{从边缘似然到 Jensen 不等式}

PPCA 的潜变量结构并非 PPCA 独有。更一般地，许多生成模型都假设数据通过两个阶段产生：
\begin{enumerate}
    \item 从先验 \(p(\mathbf{z})\) 中采样潜变量 \(\mathbf{z}\)；
    \item 再从条件分布 \(p(\mathbf{x}\mid\mathbf{z};\theta)\) 中生成观测 \(\mathbf{x}\)。
\end{enumerate}
相应的边缘似然为
\[
p(\mathbf{x};\theta)
=
\int p(\mathbf{x}\mid\mathbf{z};\theta)\,p(\mathbf{z})\,d\mathbf{z}
=
\int p(\mathbf{x},\mathbf{z};\theta)\,d\mathbf{z}.
\]

对于数据集 \(\{\mathbf{x}^{(i)}\}\)，最大似然学习变为
\[
\max_\theta\sum_i\log p(\mathbf{x}^{(i)};\theta)
=
\max_\theta\sum_i\log\int p(\mathbf{x}^{(i)},\mathbf{z};\theta)\,d\mathbf{z}.
\]
困难正来自这一表达式的结构：必须对潜变量进行积分，而对数却位于积分外部。在 PPCA 中，线性高斯结构使我们可以绕开这一障碍；一般情况下，则需要一个更系统的工具。

这个工具从 Jensen 不等式开始。回顾一下，\emph{凸函数}向上弯曲，因此连接其图像上任意两点的直线位于函数图像上方；\emph{凹函数}向下弯曲，因此相应直线位于函数图像下方。对于随机变量 \(X\) 和凸函数 \(f\)，Jensen 不等式给出
\[
f(\mathbb{E}[X])\le\mathbb{E}[f(X)].
\]
如果 \(f\) 是凹函数，不等号方向反转：
\[
f(\mathbb{E}[X])\ge\mathbb{E}[f(X)].
\]

在这里，我们真正关心的函数是对数。由于 \(\log\) 是凹函数，
\[
\log\mathbb{E}[X]\ge\mathbb{E}[\log X].
\]

这恰好就是我们需要的形式。潜变量模型中，难处理的量通常形如
\[
\log\int(\cdots)\,d\mathbf{z},
\]
等价地也可以看成 \(\log\mathbb{E}[\cdot]\)，即对数位于积分或期望之外。Jensen 不等式允许我们把对数移入期望内部，代价是从精确等式变成一个下界，但这个下界要容易处理得多。

\subsection{推导证据下界}

为了应用 Jensen 不等式，先把边缘似然改写成期望。令 \(Q(\mathbf{z})\) 是潜变量上的任意概率分布，满足
\[
\int Q(\mathbf{z})\,d\mathbf{z}=1.
\]
于是
\[
p(\mathbf{x};\theta)
=
\int p(\mathbf{x},\mathbf{z};\theta)\,d\mathbf{z}
=
\int Q(\mathbf{z})
\frac{p(\mathbf{x},\mathbf{z};\theta)}{Q(\mathbf{z})}\,d\mathbf{z}
=
\mathbb{E}_{\mathbf{z}\sim Q}
\!\left[
\frac{p(\mathbf{x},\mathbf{z};\theta)}{Q(\mathbf{z})}
\right].
\]
取对数得到
\[
\log p(\mathbf{x};\theta)
=
\log
\mathbb{E}_{\mathbf{z}\sim Q}
\!\left[
\frac{p(\mathbf{x},\mathbf{z};\theta)}{Q(\mathbf{z})}
\right].
\]

乍看之下，这可能只是一个纯粹的代数技巧，但它其实是关键步骤。分布 \(Q(\mathbf{z})\) 是我们自由选择的辅助分布，引入它的目的，就是把边缘似然写成期望。一旦做到这一点，就可以使用 Jensen 不等式。

现在对改写后的边缘似然应用 Jensen 不等式：
\[
\log p(\mathbf{x};\theta)
=
\log
\mathbb{E}_{\mathbf{z}\sim Q}
\!\left[
\frac{p(\mathbf{x},\mathbf{z};\theta)}{Q(\mathbf{z})}
\right].
\]
由于对数函数是凹函数，Jensen 不等式给出
\[
\log
\mathbb{E}_{\mathbf{z}\sim Q}
\!\left[
\frac{p(\mathbf{x},\mathbf{z};\theta)}{Q(\mathbf{z})}
\right]
\ge
\mathbb{E}_{\mathbf{z}\sim Q}
\!\left[
\log
\frac{p(\mathbf{x},\mathbf{z};\theta)}{Q(\mathbf{z})}
\right].
\]
因此
\[
\log p(\mathbf{x};\theta)
\ge
\mathbb{E}_{\mathbf{z}\sim Q}
\!\left[
\log
\frac{p(\mathbf{x},\mathbf{z};\theta)}{Q(\mathbf{z})}
\right].
\]

这个下界称为\emph{证据下界}（Evidence Lower Bound），简称\emph{ELBO}：
\[
\mathcal{L}(Q,\theta;\mathbf{x})
=
\mathbb{E}_{\mathbf{z}\sim Q}
\!\left[
\log
\frac{p(\mathbf{x},\mathbf{z};\theta)}{Q(\mathbf{z})}
\right].
\]
把分式中的对数展开，可得到等价形式
\[
\mathcal{L}(Q,\theta;\mathbf{x})
=
\mathbb{E}_{Q}\big[\log p(\mathbf{x},\mathbf{z};\theta)\big]
-
\mathbb{E}_{Q}\big[\log Q(\mathbf{z})\big].
\]

这一表达式更易处理，因为此时对数直接作用在期望内部的联合密度上，而不是位于一个难以计算的积分之外。因而，它更适合做迭代优化。

这就是本节的核心结果。与其直接最大化通常难以计算的 \(\log p(\mathbf{x};\theta)\)，不如最大化一个可处理的下界；该下界依赖两个对象：模型参数 \(\theta\) 和辅助分布 \(Q(\mathbf{z})\)。

\subsection{什么时候下界会变成精确等式？}

只有理解它何时取紧，下界才真正有用。对于 Jensen 不等式，当凹函数的输入——这里就是对数的输入——相对于求平均的变量保持常数时，等号成立。这里被平均的变量是 \(\mathbf{z}\)，所以等号成立当且仅当
\[
\frac{p(\mathbf{x},\mathbf{z};\theta)}{Q(\mathbf{z})}
=
c
\qquad
\text{对所有 }\mathbf{z}\text{ 成立},
\]
其中 \(c\) 与 \(\mathbf{z}\) 无关。

整理得
\[
Q(\mathbf{z})\propto p(\mathbf{x},\mathbf{z};\theta).
\]
为了确定归一化常数，对 \(\mathbf{z}\) 积分。由于 \(Q(\mathbf{z})\) 是合法概率分布，
\[
\int Q(\mathbf{z})\,d\mathbf{z}=1.
\]
因此
\[
p(\mathbf{x};\theta)
=
\int p(\mathbf{x},\mathbf{z};\theta)\,d\mathbf{z}
=
c\int Q(\mathbf{z})\,d\mathbf{z}
=
c.
\]
所以常数必须为
\[
c=p(\mathbf{x};\theta).
\]
将其代回等号条件得
\[
Q(\mathbf{z})
=
\frac{p(\mathbf{x},\mathbf{z};\theta)}{p(\mathbf{x};\theta)}
=
p(\mathbf{z}\mid\mathbf{x};\theta),
\]
这正是真实的潜变量后验分布。

因此，ELBO 恰好等于真实对数似然，当且仅当
\[
Q(\mathbf{z})=p(\mathbf{z}\mid\mathbf{x};\theta),
\]
也就是辅助分布与真实潜变量后验完全一致。这一点非常关键：ELBO 并不是任意构造的替代目标，它与真实对数似然的接近程度由 \(Q(\mathbf{z})\) 对真实后验的逼近程度所控制。

这个视角也进一步说明了 PPCA 在本章中的作用。PPCA 的特殊之处并不在于它绕开了潜变量问题，而在于它恰好能够把这个问题精确解决。其后验 \(p(\mathbf{z}\mid\mathbf{x})\) 可解析计算，边缘似然 \(p(\mathbf{x})\) 也可精确求得，因此 PPCA 并不严格需要 ELBO。

然而，ELBO 框架揭示的是 PPCA 所体现的一般原则。相同的潜变量结构在更复杂模型中仍然存在，只是精确后验和精确边缘似然不再可得。ELBO 正是当精确高斯计算失效之后，让我们能够继续进行学习的工具。

\bigskip

现在已经可以引入期望最大化算法。ELBO 依赖两个耦合的要素：辅助分布 \(Q(\mathbf{z})\) 与模型参数 \(\theta\)。这自然暗示出一种交替策略：先改进 \(Q(\mathbf{z})\)，使下界尽可能紧；然后改进 \(\theta\)，使下界本身变得更大。

这种交替逻辑正是 EM 所形式化的内容。下一节将说明 EM 如何作为一个迭代最大化 ELBO 的过程自然产生，并展示在 PPCA 这个特殊情形下，它如何恢复此前利用精确高斯计算得到的同一个解。

\section{期望最大化：最大化下界}

潜变量模型中的根本困难是边缘对数似然
\[
\log p(\mathbf{x};\theta)
=
\log\int p(\mathbf{x},\mathbf{z};\theta)\,d\mathbf{z},
\]
因为对数位于隐藏变量 \(\mathbf{z}\) 的边缘化积分之外。这使直接优化变得困难，即使联合分布 \(p(\mathbf{x},\mathbf{z};\theta)\) 本身可能具有相对简单的形式。

ELBO 通过引入下界
\[
\mathcal{L}(Q,\theta;\mathbf{x})
=
\mathbb{E}_{Q}\big[\log p(\mathbf{x},\mathbf{z};\theta)\big]
-
\mathbb{E}_{Q}\big[\log Q(\mathbf{z})\big]
\]
来解决这一问题。它依赖两个相互耦合的对象：辅助分布 \(Q(\mathbf{z})\) 与模型参数 \(\theta\)。于是自然得到一个交替优化策略：
\begin{quote}
\emph{如果固定其中一个，能否改进另一个？}
\end{quote}
这正是 EM 所形式化的思想。直观上，当二者中的一个固定时，ELBO 就变成只关于另一个对象的优化问题。EM 利用这种可分结构，在“改进潜变量分布”和“改进模型参数”之间交替进行。

\subsection{EM 的两个步骤}

每一轮 EM 迭代都由两个阶段组成。

\paragraph{E 步（Expectation step）。}
固定当前参数 \(\theta^{(t)}\)，选择使 ELBO 尽可能紧的分布 \(Q(\mathbf{z})\)。上一节已经知道，当
\[
Q(\mathbf{z})=p(\mathbf{z}\mid\mathbf{x};\theta^{(t)})
\]
时，下界最紧。因此，E 步设定
\[
Q^{(t)}(\mathbf{z})
=
p(\mathbf{z}\mid\mathbf{x};\theta^{(t)}).
\]
换句话说，E 步在当前参数值下计算潜变量的后验分布。

这也澄清了一个重要问题。在标准 EM 中，E 步假设后验 \(p(\mathbf{z}\mid\mathbf{x};\theta^{(t)})\) 足够容易处理，可以精确计算。PPCA 正是如此，因此 EM 可以直接应用。但更一般地，后验往往没有闭式形式。在这种情况下，我们就不能把 \(Q\) 设为精确后验，而必须用某种近似来代替它。这正是从 EM 走向\emph{变分推断}的关键一步。

还需要注意的是，“后验可计算”并不自动意味着整个优化问题都有闭式解。对 EM 来说，关键要求是：在当前参数值下，后验 \(p(\mathbf{z}\mid\mathbf{x};\theta)\)，或者至少关于该后验所需的期望，能够被计算出来。这就足以完成 E 步。然而，随后的参数更新仍然可能需要通过 M 步迭代求解，而不是一次性直接解析得到。

PPCA 的特殊性在于，它的若干核心量——包括边缘似然和后验——都可以解析处理，而且模型还具有闭式最大似然解。在许多其他模型中，即使后验足够容易处理、可以执行 EM，完整的参数优化问题仍然可能必须依赖迭代过程。

\paragraph{M 步（Maximisation step）。}
固定 \(Q^{(t)}\) 后，通过关于 \(\theta\) 最大化 ELBO 来更新模型参数：
\[
\theta^{(t+1)}
=
\arg\max_{\theta}
\mathcal{L}(Q^{(t)},\theta;\mathbf{x}).
\]
由于熵项
\[
-\mathbb{E}_{Q^{(t)}}[\log Q^{(t)}(\mathbf{z})]
\]
不依赖 \(\theta\)，这等价于最大化
\[
\theta^{(t+1)}
=
\arg\max_{\theta}
\mathbb{E}_{Q^{(t)}}\big[\log p(\mathbf{x},\mathbf{z};\theta)\big].
\]
这个量通常称为\emph{完整数据对数似然的期望}（expected complete-data log-likelihood）。

因此，EM 的两步更新模式可以写成
\[
\boxed{
\begin{aligned}
\text{E 步：}\quad
&Q^{(t)}(\mathbf{z})=p(\mathbf{z}\mid\mathbf{x};\theta^{(t)}),\\[4pt]
\text{M 步：}\quad
&\theta^{(t+1)}
=
\arg\max_{\theta}
\mathbb{E}_{Q^{(t)}}\big[\log p(\mathbf{x},\mathbf{z};\theta)\big].
\end{aligned}}
\]

\paragraph{为什么 EM 不会降低似然。}
EM 的单调性可以直接由这两个优化步骤以及 ELBO 的下界性质推出。首先，E 步之后，下界在当前参数处是紧的：
\[
\mathcal{L}(Q^{(t)},\theta^{(t)};\mathbf{x})
=
\log p(\mathbf{x};\theta^{(t)}).
\]
第二，在固定 \(Q^{(t)}\) 时，M 步选择不会使该下界下降的新参数：
\[
\mathcal{L}(Q^{(t)},\theta^{(t+1)};\mathbf{x})
\ge
\mathcal{L}(Q^{(t)},\theta^{(t)};\mathbf{x}).
\]
最后，由于 ELBO 是对数似然的下界，
\[
\log p(\mathbf{x};\theta^{(t+1)})
\ge
\mathcal{L}(Q^{(t)},\theta^{(t+1)};\mathbf{x}).
\]
将这些关系合并，就得到一次完整 EM 迭代不会降低观测数据的对数似然。图~\ref{fig:em-elbo-geometry} 从几何角度展示了同一个过程：算法交替地“收紧”下界并把下界“抬高”。

\begin{figure}[H]
\centering
\begin{tikzpicture}[x=1cm,y=1cm,>=Latex,font=\small]
\def\thetacur{2.0}
\def\thetanext{2.714}

% Common schematic functions:
% log likelihood: 3 - 0.25(theta-3)^2
% post-E ELBO: likelihood - 0.10(theta-2)^2
% loose pre-E ELBO: likelihood - 0.60 - 0.15(theta-1.5)^2

% Panel (a): E-step
\begin{scope}[shift={(0,0)}]
  \draw[->] (0,0) -- (5.0,0) node[right] {$\theta$};
  \draw[->] (0,0) -- (0,3.35);

  \draw[thick] plot[domain=0.3:4.7,samples=100]
    (\x,{3 - 0.25*(\x-3)^2});
  \node[anchor=west] at (3.25,3.02) {$\log p(\mathbf{x};\theta)$};

  \draw[dashed] plot[domain=0.3:4.7,samples=100]
    (\x,{3 - 0.25*(\x-3)^2 - 0.60 - 0.15*(\x-1.5)^2});
  \node[anchor=west] at (0.55,0.92) {$\mathcal{L}(Q,\theta;\mathbf{x})$};

  \draw[thick,dash dot] plot[domain=0.3:4.7,samples=100]
    (\x,{3 - 0.25*(\x-3)^2 - 0.10*(\x-2)^2});
  \node[anchor=west] at (3.15,2.12)
    {$\mathcal{L}(Q^{(t)},\theta;\mathbf{x})$};

  \draw[densely dotted] (\thetacur,0) -- (\thetacur,2.75);
  \fill (\thetacur,2.75) circle (1.4pt);
  \node[below] at (\thetacur,0) {$\theta^{(t)}$};

  \draw[->] (1.45,1.28) .. controls (1.35,1.80) and (1.55,2.25) .. (1.90,2.62);
  \node[align=center] at (0.95,1.80) {\scriptsize 更新\\[-1pt]\scriptsize $Q$};

  \node[font=\small\bfseries] at (2.5,3.72) {(a) E 步：收紧下界};
\end{scope}

% Panel (b): M-step
\begin{scope}[shift={(7.0,0)}]
  \draw[->] (0,0) -- (5.0,0) node[right] {$\theta$};
  \draw[->] (0,0) -- (0,3.35);

  \draw[thick] plot[domain=0.3:4.7,samples=100]
    (\x,{3 - 0.25*(\x-3)^2});
  \node[anchor=west] at (3.25,3.02) {$\log p(\mathbf{x};\theta)$};

  % Exactly the same post-E-step bound as in panel (a)
  \draw[thick,dash dot] plot[domain=0.3:4.7,samples=100]
    (\x,{3 - 0.25*(\x-3)^2 - 0.10*(\x-2)^2});
  \node[anchor=west] at (3.15,2.12)
    {$\mathcal{L}(Q^{(t)},\theta;\mathbf{x})$};

  \draw[densely dotted] (\thetacur,0) -- (\thetacur,2.75);
  \draw[densely dotted] (\thetanext,0) -- (\thetanext,2.929);
  \fill (\thetacur,2.75) circle (1.4pt);
  \fill (\thetanext,2.929) circle (1.4pt);

  \node[below,xshift=-4pt] at (\thetacur,0) {$\theta^{(t)}$};
  \node[below,xshift=7pt] at (\thetanext,0) {$\theta^{(t+1)}$};

  \draw[->,thick] (2.08,2.78) -- (2.62,2.92);

  \node[font=\small\bfseries] at (2.5,3.72) {(b) M 步：抬高下界};
\end{scope}

\node[align=center] at (6.0,-0.92) {
$\displaystyle
\log p(\mathbf{x};\theta^{(t+1)})
\ge
\mathcal{L}(Q^{(t)},\theta^{(t+1)};\mathbf{x})
\ge
\mathcal{L}(Q^{(t)},\theta^{(t)};\mathbf{x})
=
\log p(\mathbf{x};\theta^{(t)})$};

\end{tikzpicture}
\caption{\textbf{把 EM 理解为对 ELBO 的交替优化。}
在 E 步中，把 $Q$ 设为精确后验，使 ELBO 在当前参数处与对数似然相切。固定这一界后，M 步更新参数以提高该下界。下方不等式说明完整的一次迭代不会降低观测数据的对数似然。曲线仅为示意。}
\label{fig:em-elbo-geometry}
\end{figure}

\subsection{把 EM 应用于 PPCA}

在 PPCA 中，潜变量后验是可解析处理的，因此可以显式执行 EM。这使 PPCA 成为特别有价值的示例：它展示了一个一般的、基于 ELBO 的算法如何在线性高斯模型中退化为闭式更新。

对于 PPCA，潜变量的后验分布是高斯分布。这仍然来自此前推导边缘分布时所利用的线性高斯结构，而配方法推导见附录~\ref{app:completing-square}。后验为
\[
p(\mathbf{z}\mid\mathbf{x})
=
\mathcal{N}\!\bigl(
\mathbf{z}\mid
\mathbf{M}^{-1}\mathbf{W}^\top(\mathbf{x}-\boldsymbol{\mu}),
\alpha^2\mathbf{M}^{-1}
\bigr),
\qquad
\mathbf{M}=\mathbf{W}^\top\mathbf{W}+\alpha^2\mathbf{I}.
\]
图~\ref{fig:ppca-four-probabilistic-views} 把 PPCA 中至此已经出现的四个概率对象放在一起展示：先验 \(p(\mathbf{z})\)、条件分布 \(p(\mathbf{x}\mid\mathbf{z})\)、边缘分布 \(p(\mathbf{x})\) 和后验 \(p(\mathbf{z}\mid\mathbf{x})\)。

\begin{figure}[H]
\centering
\begin{minipage}[t]{0.45\textwidth}
\centering\includegraphics[width=\linewidth]{figures/ch02/panels/fig2_3a_prior.pdf}
\end{minipage}\hfill
\begin{minipage}[t]{0.45\textwidth}
\centering\includegraphics[width=\linewidth]{figures/ch02/panels/fig2_3b_conditional.pdf}
\end{minipage}

\vspace{0.4em}

\begin{minipage}[t]{0.45\textwidth}
\centering\includegraphics[width=\linewidth]{figures/ch02/panels/fig2_3c_marginal.pdf}
\end{minipage}\hfill
\begin{minipage}[t]{0.45\textwidth}
\centering\includegraphics[width=\linewidth]{figures/ch02/panels/fig2_3d_posterior.pdf}
\end{minipage}
\caption{\textbf{同一个 PPCA 模型的四种概率视角。}
先验 $p(z)$ 描述潜在坐标；固定 $z^\ast$ 得到 $p(\mathbf{x}\mid z^\ast)$；对 $z$ 积分得到边缘分布 $p(\mathbf{x})$；观察到 $\mathbf{x}^\ast$ 后得到后验 $p(z\mid\mathbf{x}^\ast)$。后验子图中标出了真实生成用的 $z^\ast$；由于存在观测噪声，它不必等于后验均值。在 PPCA 中，这四种分布均为高斯分布且都可解析求得。}
\label{fig:ppca-four-probabilistic-views}
\end{figure}


根据这一后验，E 步可以计算
\[
\mathbb{E}[\mathbf{z}]
=
\mathbf{M}^{-1}\mathbf{W}^\top(\mathbf{x}-\boldsymbol{\mu}),
\]
以及
\[
\mathbb{E}[\mathbf{z}\mathbf{z}^\top]
=
\alpha^2\mathbf{M}^{-1}
+
\mathbb{E}[\mathbf{z}]\,\mathbb{E}[\mathbf{z}]^\top.
\]

随后，M 步把这些期望代入完整数据对数似然的期望，并关于参数 \((\mathbf{W},\boldsymbol{\mu},\alpha)\) 最大化，从而得到 PPCA 参数的闭式更新方程。

这里真正重要的并不只是“PPCA 可以用 EM 来拟合”，而是“PPCA 的精确解可以被重新解释为 EM 过程的固定点”。换句话说，EM 并没有得到不同的模型，它提供的是对\emph{同一个模型如何学习}的另一种理解。

更广泛地看，EM 告诉我们，潜变量模型中的学习可以理解为以下两个过程的交替：
\begin{itemize}
    \item \textbf{推断}：给定数据，估计隐藏变量；
    \item \textbf{学习}：利用这些推断得到的潜变量更新模型参数。
\end{itemize}

\subsection{从 PPCA 到正则化潜在表示}
尽管 PPCA 引入了潜变量先验 \(p(\mathbf{z})=\mathcal{N}(\mathbf{0},\mathbf{I})\)，并定义了合法的生成模型，但它还没有提供现代潜变量模型（例如 VAE）中将成为核心的那种灵活、正则化的推断框架。

对于任意单个观测 \(\mathbf{x}\)，后验 \(p(\mathbf{z}\mid\mathbf{x})\) 描述了在模型下哪些潜在坐标可以合理解释该数据点。如果对经验数据分布上的这些后验取平均，就得到\emph{聚合后验}（aggregate posterior）：
\[
\mathbb{E}_{q(\mathbf{x})}\!\left[p(\mathbf{z}\mid\mathbf{x})\right].
\]
该分布描述整个数据集如何被表示在潜在空间中。

在 PPCA 中，由于模型是线性高斯的，这个后验可以精确计算。这当然是 PPCA 的优点，但也恰好揭示了其局限的来源：干净的数学结果依赖很强的假设。潜变量到观测的映射必须是线性的，观测噪声必须是高斯的，而且后验推断必须能够解析完成。

当我们离开线性高斯世界时，这一点就变得非常重要。现代生成模型希望保留同样的潜变量思想，同时让潜变量到观测的映射变得更加灵活。一旦引入这种灵活性，使 PPCA 如此透明的简单闭式计算通常就不复存在。因此，我们需要一种新的框架：它能够学习表达能力更强的生成映射、从数据中推断潜在表示，并使潜在空间保持一种有利于生成的组织结构。

变分自编码器正是为此而生。可以把 VAE 看成 PPCA 叙事的非线性、变分扩展：它保留“潜变量生成观测”的基本思想，却用一个可学习的推断--生成框架替代精确高斯计算。关于这一框架的细节，包括变分后验的准确作用以及潜在空间正则化项，将在下一章展开。

因此，从 PPCA 过渡到 VAE，并不是从“非生成模型”转向“生成模型”；PPCA 本身已经是生成模型。真正发生的，是从解析可解的线性高斯潜变量模型，过渡到一个更灵活的非线性潜变量模型，在后者中，推断与生成必须被共同学习。

下一章中，变分自编码器将把重构、概率推断和潜在空间正则化统一进一个变分目标，从而完整展开这一思想。

\section*{本章小结与参考文献}
\addcontentsline{toc}{section}{本章小结与参考文献}

本章把概率主成分分析（PPCA）作为本书中的第一个概率潜变量模型进行了介绍。它在 PCA 的确定性图景之上，引入潜变量先验、概率解码器、用于学习的边缘似然以及用于推断的后验分布。

由于 PPCA 是线性高斯模型，其所有核心量都保持解析可解。这使我们能够显式推导边缘似然，把其协方差解释为结构化低秩成分与各向同性噪声之和，并把 PPCA 看作一个少见的模型：其先验、似然、边缘分布和后验都能够写成闭式。

随后，我们比较了三种模型学习视角：直接解析最大似然、基于梯度的优化以及 EM。其中 EM 尤其重要，因为它揭示了潜变量学习更深层的结构：在隐藏变量推断和模型参数优化之间交替。由此进一步引出 Jensen 不等式和 ELBO，它们把 PPCA 的思想推广到无法进行精确边缘化的模型。

最后，我们识别出 PPCA 的一个重要局限：虽然它为潜变量引入了先验，却并没有显式正则化聚合后的推断潜在空间，使其与该先验对齐。这直接引向下一章的变分自编码器（VAE）；VAE 会把重构、概率推断以及潜在空间正则化统一到一个变分目标中。

\bigskip

本章内容建立在概率主成分分析、潜变量模型和最大似然推断的经典工作之上。PPCA 的原始形式由 Tipping 和 Bishop~\cite{tipping1999ppca} 提出，而 EM 算法的经典工作来自 Dempster、Laird 和 Rubin~\cite{dempster1977em}。希望进一步学习这些思想，尤其是高斯潜变量模型、EM 及其概率解释的读者，可以参考 Bishop~\cite{bishop2006prml} 与 Murphy~\cite{murphy2012ml} 的教材。

\chapter[变分自编码器：从概率潜变量到变分推断]{变分自编码器：\\从概率潜变量到变分推断}

\section*{概述}
\addcontentsline{toc}{section}{概述}

变分自编码器（Variational Autoencoder, VAE）是现代深度生成学习的基础模型之一。它在本书中占据重要位置，因为它把前几章分别发展的两条脉络汇合起来：一条是作为基于重构的表示学习器的自编码器，另一条是作为潜变量生成模型的概率主成分分析。

在确定性自编码器中，每个输入被映射到一个单一潜在编码，然后再进行重构。在 PPCA 中，潜变量被概率化处理，但模型仍然保持线性高斯结构。VAE 把这两种思想结合到更灵活的非线性框架中：编码器输出潜变量上的分布，解码器定义观测上的条件分布，而模型学习则通过优化变分下界完成，而不是直接优化精确边缘似然。

这使 VAE 成为从经典潜变量模型通向现代深度生成建模的自然桥梁。它保留了 PPCA 的概率结构——先验、似然、后验和边缘似然——但用一个学习得到的近似后验替代精确高斯推断。上一章引入的证据下界（ELBO）在这里成为实际训练目标，在重构准确性与潜在空间正则化之间取得平衡。

尽管 VAE 已不再是最新一代生成模型中最占主导地位的范式，但从概念上看，它依然处于核心位置。VAE 展示了如何把概率推断、神经网络、潜在表示与基于梯度的优化统一到一个端到端生成模型中。因此，本章是从解析可解的潜变量模型迈向本书后续更广泛深度生成框架的第一个重要步骤。

\section*{概念导图}
\addcontentsline{toc}{section}{概念导图}

本章追踪从确定性自编码器走向变分潜变量生成模型的过程。

\begin{itemize}

    \item[\(\rightarrow\)] \textbf{首先把确定性自编码器改造成概率模型。}  
    编码器不再把每个观测映射为单一潜在编码，而是输出分布 \(q_{\phi}(\mathbf{z}\mid\mathbf{x})\)；与此同时，解码器定义条件分布 \(p_{\theta}(\mathbf{x}\mid\mathbf{z})\)。这样，模型既能表达不确定性，也真正具备生成能力。

    \item[\(\rightarrow\)] \textbf{随后从潜变量模型本身推导学习目标。}  
    从难以计算的边缘似然 \(p_{\theta}(\mathbf{x})\) 出发，我们使用上一章建立的 ELBO 框架，并将其具体化到 VAE。由此得到熟悉的“重构项 + 正则项”分解，揭示模型如何在拟合数据与组织潜在空间之间取得平衡。

    \item[\(\rightarrow\)] \textbf{最后解释 VAE 如何训练，以及这种训练实现了什么。}  
    变分后验 \(q_{\phi}(\mathbf{z}\mid\mathbf{x})\) 充当对真实后验的学习近似，而重参数化技巧使这一随机推断过程变得可微，从而能够用基于梯度的方法进行训练。

\end{itemize}

\medskip
\noindent VAE 的主要概率关系概括如下图所示。
\begin{center}
\begin{tikzpicture}[
  >=Latex,
  every node/.style={font=\small},
  box/.style={
    draw,
    rounded corners=2pt,
    align=center,
    inner sep=6pt,
    minimum height=1.02cm,
    text width=3.0cm
  },
  objective/.style={
    draw,
    rounded corners=2pt,
    align=center,
    inner sep=6pt,
    minimum height=1.08cm,
    text width=4.2cm
  },
  smallbox/.style={
    draw,
    rounded corners=2pt,
    align=center,
    inner sep=5pt,
    minimum height=0.88cm,
    text width=2.4cm
  },
  flow/.style={->,thick},
  secondary/.style={->,densely dashed}
]

\node[smallbox] (x) at (0,0) {
  \textbf{观测数据}\\[1pt]
  $\mathbf{x}$
};

\node[box] (enc) at (4.1,0) {
  \textbf{变分编码器}\\[1pt]
  $q_{\phi}(\mathbf{z}\mid\mathbf{x})$
};

\node[smallbox] (z) at (8.1,0) {
  \textbf{潜变量样本}\\[1pt]
  $\mathbf{z}$
};

\node[box] (dec) at (12.2,0) {
  \textbf{概率解码器}\\[1pt]
  $p_{\theta}(\mathbf{x}\mid\mathbf{z})$
};

\node[smallbox] (prior) at (8.1,-2.45) {
  \textbf{潜变量先验}\\[1pt]
  $p(\mathbf{z})=\mathcal{N}(\mathbf{0},\mathbf{I})$
};

\node[objective] (elbo) at (3.0,-2.45) {
  \textbf{ELBO 目标}\\[2pt]
  $\mathbb{E}_{q_{\phi}(\mathbf{z}\mid\mathbf{x})}
   [\log p_{\theta}(\mathbf{x}\mid\mathbf{z})]$
  $-\;
  \mathrm{KL}\!\left(
  q_{\phi}(\mathbf{z}\mid\mathbf{x})
  \,\|\,p(\mathbf{z})\right)$
};

\draw[flow] (x) -- (enc);
\draw[flow] (enc) -- (z);
\draw[flow] (z) -- (dec);

% Prior supplies latent samples when generating new observations.
\draw[flow] (prior) -- node[right] {\scriptsize 生成} (z);

% Dashed links indicate distributions entering the ELBO terms.
\draw[secondary] (prior) -- (elbo);
\draw[secondary] (enc) -- (elbo);
\draw[secondary] (dec.south west) to[out=-120,in=0] (elbo.east);

\node[font=\normalsize\bfseries] at (6.1,1.45) {VAE 概览};

\end{tikzpicture}
\end{center}
\medskip

\section{从自编码器到变分自编码器}

前几章并行发展了两种思想。一方面，自编码器说明数据可以被压缩为较低维的潜在表示，也就是一种紧凑的内部描述或特征编码，再由此重构输入。另一方面，概率主成分分析（PPCA）说明如何把这样的表示变成真正的概率生成过程。变分自编码器（VAE）~\cite{kingma2014autoencoding,rezende2014stochastic} 把这两条路径结合起来。

标准自编码器具有熟悉的确定性形式
\[
\mathbf{x}
\;\xrightarrow{\text{编码器}}\;
\mathbf{z}
\;\xrightarrow{\text{解码器}}\;
\hat{\mathbf{x}}.
\]
给定输入 \(\mathbf{x}\)，编码器产生一个唯一潜在编码 \(\mathbf{z}\)，解码器再把该编码映射为一个唯一重构 \(\hat{\mathbf{x}}\)。这种结构对于压缩和表示学习很有用，但有一个重要局限：潜在编码只是一个确定的点，而不是概率对象。因此，模型无法自然表达不确定性，也没有提供一个有原则的潜变量生成框架。

VAE 通过把确定性的潜在编码替换为潜变量上的\emph{概率分布}来解决这一问题。这与此前从 PCA 过渡到 PPCA 的过程相呼应：在那里，我们也是通过引入不确定性，把确定性的低维描述转化为概率潜变量模型。VAE 不再把 \(\mathbf{x}\) 编码为潜在空间中的一个精确点，而是编码为一个概率分布，用于表达“哪些潜在表示最可能解释该观测”这一不确定性。这就是核心概念转变：
\begin{quote}
\emph{潜在表示不再是单一编码，而是一个可以从中采样编码的分布。}
\end{quote}

\subsection{让编码器概率化}

在确定性自编码器中，编码器是函数
\[
f_{\text{enc}}:\mathbf{x}\mapsto\mathbf{z}.
\]
对于每个观测 \(\mathbf{x}\)，它只返回一个潜向量 \(\mathbf{z}\)。

在 VAE 中，编码器则定义一个易于处理的潜变量条件分布
\[
q_{\phi}(\mathbf{z}\mid\mathbf{x}),
\]
它作为对真实后验的学习近似。这个分布通常称为\emph{变分后验}（variational posterior）或\emph{近似后验}（approximate posterior）。其作用是描述：给定观测数据点 \(\mathbf{x}\) 后，哪些潜变量是合理的。

这里使用“后验”一词是有意的。上一章的 PPCA 中，我们看到潜变量模型的真实后验为
\[
p_{\theta}(\mathbf{z}\mid\mathbf{x}),
\]
即观察到 \(\mathbf{x}\) 之后，潜变量的条件分布。在 VAE 中，这个真实后验通常难以解析计算。主要原因是编码器和解码器通常由多层非线性神经网络建模，因此后验不再保留使 PPCA 可解析求解的简单线性高斯形式。所以，我们改为引入可处理的近似 \(q_{\phi}(\mathbf{z}\mid\mathbf{x})\)。此时最关键的结构变化很简单：编码器不再输出单一向量，而是输出一个分布的参数。

实际中，编码器通常由神经网络实现，把 \(\mathbf{x}\) 映射为某个高斯分布的参数。最常见的假设为
\[
q_{\phi}(\mathbf{z}\mid\mathbf{x})
=
\mathcal{N}\!\bigl(
\mathbf{z}\mid
\boldsymbol{\mu}_{\phi}(\mathbf{x}),
\operatorname{diag}(\boldsymbol{\sigma}_{\phi}^2(\mathbf{x}))
\bigr),
\]
或者在简化的各向同性表述中，
\[
q_{\phi}(\mathbf{z}\mid\mathbf{x})
=
\mathcal{N}\!\bigl(
\mathbf{z}\mid
\boldsymbol{\mu}_{\phi}(\mathbf{x}),
\boldsymbol{\sigma}_{\phi}^2(\mathbf{x})\mathbf{I}
\bigr).
\]
因此，编码器网络并不直接输出 \(\mathbf{z}\)，而是输出潜在空间中高斯分布的均值和方差参数。

从概念上看，编码器现在具有如下形式：
\[
\mathbf{x}
\;\longrightarrow\;
\bigl(\boldsymbol{\mu}_{\phi}(\mathbf{x}),
\boldsymbol{\sigma}_{\phi}(\mathbf{x})\bigr)
\;\longrightarrow\;
q_{\phi}(\mathbf{z}\mid\mathbf{x}).
\]
随后通过采样得到潜在编码：
\[
\mathbf{z}\sim q_{\phi}(\mathbf{z}\mid\mathbf{x}).
\]

因此，对于固定输入 \(\mathbf{x}\)，确定性自编码器与 VAE 的差别十分直接：前者返回一个潜在点，而后者的编码器定义一组可能的潜在值所构成的分布，可以从中采样不同编码。图~\ref{fig:vae-ae-versus-posterior} 从几何上展示了这种“点到分布”的转变。

\begin{figure}[H]
\centering
\includegraphics[width=0.94\textwidth]{figures/ch03/figure3_1_ae_vs_vae_latent_representation.pdf}
\caption{\textbf{从确定性潜在编码到变分后验。}
对于固定输入 $\mathbf{x}$，标准自编码器只产生一个潜在点，而 VAE 编码器给出合理潜变量值的分布 $q_{\phi}(\mathbf{z}\mid\mathbf{x})$。图中用样本、均值和椭圆等高线展示常见的对角高斯选择。两者使用相同中心仅是为了更直观地突出“单点”与“分布”之间的差别。}
\label{fig:vae-ae-versus-posterior}
\end{figure}

\begin{tcolorbox}[
  colback=blue!6,
  colframe=blue!55!black,
  title={为什么选择易处理的高斯近似后验？},
  boxrule=0.8pt,
  arc=2mm,
  left=2mm,
  right=2mm,
  top=1mm,
  bottom=1mm
]

VAE 中的真实后验 \(p_{\theta}(\mathbf{z}\mid\mathbf{x})\) 通常复杂到无法精确计算，因此我们用一个更简单的分布族 \(q_{\phi}(\mathbf{z}\mid\mathbf{x})\) 替代它，使其能够在实际中表示和学习。选择高斯分布尤其方便，主要有两个原因。

\begin{itemize}
    \item 它容易表示。神经网络只需输出均值和方差，就可以紧凑、可微地描述潜在分布。
    \item 它便于计算。我们可以轻松从中采样，而且后续学习目标中将出现的许多数学量，在采用高斯分布时都更容易求值。
\end{itemize}
这里所谓“易处理”，只是指近似后验被选得足够简单，可以在学习过程中进行操作，并不意味着整个 VAE 又重新变成解析可解模型。与 PPCA 不同，VAE 的解码器仍然是非线性的，精确数据似然仍难以计算，我们仍需要替代目标与基于梯度的优化来训练模型。

\end{tcolorbox}

\subsection{让解码器概率化}

解码端也会发生类似转变。在确定性自编码器中，解码器是函数
\[
f_{\text{dec}}:\mathbf{z}\mapsto\hat{\mathbf{x}},
\]
它对每个潜在输入给出一个唯一重构 \(\hat{\mathbf{x}}\)。

在 VAE 中，解码器则定义条件分布
\[
p_{\theta}(\mathbf{x}\mid\mathbf{z}),
\]
用于描述如何由潜变量生成观测。这就是模型的生成部分。

同样，解码器通常由神经网络参数化。给定 \(\mathbf{z}\) 后，解码器输出数据空间中某个分布的参数。对于连续数据，一种常见选择是高斯分布，例如
\[
p_{\theta}(\mathbf{x}\mid\mathbf{z})
=
\mathcal{N}\!\bigl(
\mathbf{x}\mid
\boldsymbol{\mu}_{\theta}(\mathbf{z}),
\sigma^2\mathbf{I}
\bigr),
\]
而对于二值数据则经常使用 Bernoulli 似然。具体选择取决于所建模的数据类型，但底层原则始终相同：
\begin{quote}
\emph{解码器输出的不是一个确定性重构点，而是给定潜变量时可能观测值的概率分布。}
\end{quote}

实践中，人们通常会用该分布的均值或众数作为模型的代表性重构，但从概念上说，解码器定义的是完整条件概率律，而不是单点预测。因此，编码器把数据映射为潜在空间中的分布，而解码器则把潜变量映射为数据空间中的分布。

\subsection{先验与学习困难}

\paragraph{潜变量先验。}
为了完成概率潜变量模型，还需要引入潜变量的先验分布：
\[
p(\mathbf{z}).
\]
与 PPCA 类似，标准 VAE 通常从简单的标准高斯先验开始：
\[
p(\mathbf{z})=\mathcal{N}(\mathbf{0},\mathbf{I}).
\]

这个先验有两个作用。

第一，它提供生成的起点。训练完成后，可以采样
\[
\mathbf{z}\sim p(\mathbf{z}),
\]
再从解码器分布中采样生成观测：
\[
\mathbf{x}\sim p_{\theta}(\mathbf{x}\mid\mathbf{z}).
\]

第二，它提供一个参考分布，后续会用它来正则化推断得到的潜变量分布。这是相对于 PPCA 的一个重要推进：在 PPCA 中，先验虽然存在，但模型并不会显式鼓励推断得到的潜在表示与先验对齐；在 VAE 中，这种对齐将直接成为学习目标的一部分。

因此，先验赋予潜在空间全局结构，而解码器把该空间中的点转换为数据。

\paragraph{为什么这会带来新的困难。}

把上述部分合在一起，VAE 包含三个概率组件：
\[
p(\mathbf{z}),
\qquad
q_{\phi}(\mathbf{z}\mid\mathbf{x}),
\qquad
p_{\theta}(\mathbf{x}\mid\mathbf{z}),
\]
分别对应潜变量先验、用于推断的近似后验，以及用于生成的解码器似然。

这使 VAE 的结构既类似自编码器，也类似 PPCA，但比二者都更加丰富。像自编码器一样，它具有编码器--解码器架构；像 PPCA 一样，它是概率潜变量模型。但与 PPCA 不同，编码器和解码器现在通常都是非线性神经网络，这大幅增强了建模能力，同时也破坏了我们此前依赖的解析可解性。

在 PPCA 中，线性高斯结构使得边缘似然 \(p(\mathbf{x})\) 和后验 \(p(\mathbf{z}\mid\mathbf{x})\) 等量都能显式推导。在 VAE 中，这些量通常都不再具有闭式形式。因此，虽然 VAE 具有清晰的概率结构，但其学习问题已经不能像 PPCA 那样解析求解。这就直接带来了本章的核心问题：
\begin{quote}
\emph{如果真实边缘似然和后验都难以计算，我们应当转而优化什么目标？}
\end{quote}

答案就是上一章介绍的同一条一般原则：证据下界（ELBO）。下一节将把这一原则具体化到 VAE，并推导使模型能够训练的目标函数。

\section{从难以计算似然到 ELBO}

建立了 VAE 的概率结构之后，现在可以讨论如何学习模型。与 PPCA 一样，最自然的原则仍然是最大似然：希望选择模型参数，使模型对实际观测数据赋予较高概率。

对于单个观测 \(\mathbf{x}\)，相关量是边缘似然
\[
p_{\theta}(\mathbf{x})
=
\int p_{\theta}(\mathbf{x},\mathbf{z})\,d\mathbf{z}
=
\int p_{\theta}(\mathbf{x}\mid\mathbf{z})\,p(\mathbf{z})\,d\mathbf{z},
\]
对数据集 \(\{\mathbf{x}^{(i)}\}\)，学习目标为
\[
\max_{\theta}\sum_i\log p_{\theta}(\mathbf{x}^{(i)}).
\]

它在形式上与上一章 PPCA 的目标完全相同，但现在存在一个关键差异。PPCA 的线性高斯结构允许我们精确计算关于 \(\mathbf{z}\) 的积分；VAE 中的解码器 \(p_{\theta}(\mathbf{x}\mid\mathbf{z})\) 通常由非线性神经网络参数化，因此此前支持精确边缘化的简单线性高斯结构不再存在。结果就是边缘似然
\[
\log p_{\theta}(\mathbf{x})
=
\log\int p_{\theta}(\mathbf{x}\mid\mathbf{z})\,p(\mathbf{z})\,d\mathbf{z}
\]
通常难以计算。

同样的困难也立刻出现在后验中：
\[
p_{\theta}(\mathbf{z}\mid\mathbf{x})
=
\frac{p_{\theta}(\mathbf{x}\mid\mathbf{z})\,p(\mathbf{z})}
{p_{\theta}(\mathbf{x})}.
\]
分子可以写出，但要将其归一化就需要同一个难以计算的边缘似然 \(p_{\theta}(\mathbf{x})\)。因此，在 VAE 中，边缘似然和真实后验通常都没有闭式形式。

这正是上一章的变分机制成为必需工具的地方。我们不直接优化难以计算的对数似然，而是优化一个可处理的下界：证据下界（ELBO）。

\subsection{引入变分后验}

第 2 章推导 ELBO 时，引入了潜变量上的辅助分布 \(Q(\mathbf{z})\)。在 VAE 中，这个辅助分布现在变成一个学习得到的、参数化的真实后验近似：
\[
Q(\mathbf{z})
\;\longrightarrow\;
q_{\phi}(\mathbf{z}\mid\mathbf{x}),
\]
其中 \(q_{\phi}(\mathbf{z}\mid\mathbf{x})\) 就是上一节引入的编码器网络。

这一替换，是把一般 ELBO 框架转化为 VAE 目标的关键一步。我们不再尝试计算难以处理的真实后验 \(p_{\theta}(\mathbf{z}\mid\mathbf{x})\)，而是使用易处理的近似 \(q_{\phi}(\mathbf{z}\mid\mathbf{x})\)，并优化 \(\log p_{\theta}(\mathbf{x})\) 的下界。

应用 ELBO 构造得到
\[
\log p_{\theta}(\mathbf{x})
\geq
\mathbb{E}_{q_{\phi}(\mathbf{z}\mid\mathbf{x})}
\!\left[
\log\frac{p_{\theta}(\mathbf{x},\mathbf{z})}
{q_{\phi}(\mathbf{z}\mid\mathbf{x})}
\right].
\]
因此，单个数据点的 VAE 目标为
\[
\mathcal{L}(\mathbf{x};\theta,\phi)
=
\mathbb{E}_{q_{\phi}(\mathbf{z}\mid\mathbf{x})}
\!\left[
\log\frac{p_{\theta}(\mathbf{x},\mathbf{z})}
{q_{\phi}(\mathbf{z}\mid\mathbf{x})}
\right].
\]

这就是针对 VAE 具体化之后的 ELBO。它依赖两组参数：
\begin{itemize}
    \item \(\theta\)：解码器网络 \(p_{\theta}(\mathbf{x}\mid\mathbf{z})\) 的参数；
    \item \(\phi\)：编码器网络 \(q_{\phi}(\mathbf{z}\mid\mathbf{x})\) 的参数。
\end{itemize}
训练 VAE，就是同时关于模型的这两个部分最大化该下界。

下面把同一个目标改写成三种等价形式。每种形式突出学习问题的一个不同侧面：潜变量结构、与精确似然之间的差距，以及实践中重构与正则化之间的平衡。

\subsection{ELBO 的三种等价形式}

\paragraph{ELBO 的第一种形式：联合密度形式。}
ELBO 最紧凑的形式为
\[
\mathcal{L}(\mathbf{x};\theta,\phi)
=
\mathbb{E}_{q_{\phi}(\mathbf{z}\mid\mathbf{x})}
\!\left[
\log\frac{p_{\theta}(\mathbf{x},\mathbf{z})}
{q_{\phi}(\mathbf{z}\mid\mathbf{x})}
\right].
\]
展开分式的对数，可以更清楚地看到其结构：
\[
\mathcal{L}(\mathbf{x};\theta,\phi)
=
\mathbb{E}_{q_{\phi}(\mathbf{z}\mid\mathbf{x})}
[\log p_{\theta}(\mathbf{x},\mathbf{z})]
-
\mathbb{E}_{q_{\phi}(\mathbf{z}\mid\mathbf{x})}
[\log q_{\phi}(\mathbf{z}\mid\mathbf{x})].
\]

这一形式在概念上很有用，因为它与上一章的一般 ELBO 推导展示的是完全相同的潜变量结构。第一项要求采样得到的潜变量与联合模型 \(p_{\theta}(\mathbf{x},\mathbf{z})\) 相容；第二项则来自变分分布本身。

它也清楚说明了 VAE 在做什么。模型希望通过潜变量 \(\mathbf{z}\) 解释观测 \(\mathbf{x}\)，但由于精确后验不可得，只能转而优化依赖近似后验 \(q_{\phi}(\mathbf{z}\mid\mathbf{x})\) 的下界。

\paragraph{ELBO 的第二种形式：精确似然减去后验失配。}

利用 贝叶斯公式
\[
p_{\theta}(\mathbf{x},\mathbf{z})
=
p_{\theta}(\mathbf{z}\mid\mathbf{x})\,p_{\theta}(\mathbf{x}),
\]
可将 ELBO 重写为
\[
\begin{aligned}
\mathcal{L}(\mathbf{x};\theta,\phi)
&=
\mathbb{E}_{q_{\phi}(\mathbf{z}\mid\mathbf{x})}
\!\left[
\log\frac{p_{\theta}(\mathbf{z}\mid\mathbf{x})\,p_{\theta}(\mathbf{x})}
{q_{\phi}(\mathbf{z}\mid\mathbf{x})}
\right] \\[4pt]
&=
\log p_{\theta}(\mathbf{x})
-
\mathbb{E}_{q_{\phi}(\mathbf{z}\mid\mathbf{x})}
\!\left[
\log\frac{q_{\phi}(\mathbf{z}\mid\mathbf{x})}
{p_{\theta}(\mathbf{z}\mid\mathbf{x})}
\right].
\end{aligned}
\]

\begin{tcolorbox}[
  colback=blue!6,
  colframe=blue!55!black,
  title={KL 散度},
  boxrule=0.8pt,
  arc=2mm,
  left=2mm,
  right=2mm,
  top=1mm,
  bottom=1mm
]
对于两个概率分布 \(q\) 和 \(p\)，从 \(q\) 到 \(p\) 的\emph{Kullback--Leibler（KL）散度}定义为
\[
\mathrm{KL}(q\,\|\,p)
=
\mathbb{E}_{q}\!\left[\log\frac{q}{p}\right].
\]
它衡量在关于 \(q\) 取期望时，\(q\) 与 \(p\) 之间有多大差异。

这里尤其重要的两个基本事实是：
\begin{itemize}
    \item \(\mathrm{KL}(q\,\|\,p)\ge0\)；
    \item 当且仅当 \(q=p\) 几乎处处成立时，\(\mathrm{KL}(q\,\|\,p)=0\)。
\end{itemize}

因此，KL 散度可以看成两个分布之间非负的失配度量。严格来说它并不是真正的几何距离，因为它不对称，但在概率建模中极其有用。
\end{tcolorbox}

识别出第二项正是 KL 散度后，可得
\[
\boxed{
\mathcal{L}(\mathbf{x};\theta,\phi)
=
\log p_{\theta}(\mathbf{x})
-
\mathrm{KL}\!\bigl(
q_{\phi}(\mathbf{z}\mid\mathbf{x})
\,\|\,\,
p_{\theta}(\mathbf{z}\mid\mathbf{x})
\bigr).
}
\]

这一形式表明，ELBO 与真实对数似然之间只相差近似后验与真实后验之间的散度。由于 KL 散度始终非负，
\[
\mathrm{KL}\!\bigl(
q_{\phi}(\mathbf{z}\mid\mathbf{x})
\,\|\,\,
p_{\theta}(\mathbf{z}\mid\mathbf{x})
\bigr)\ge0,
\]
立即得到
\[
\mathcal{L}(\mathbf{x};\theta,\phi)\le\log p_{\theta}(\mathbf{x}),
\]
这也解释了为什么 ELBO 的确是一个下界。

这一形式还清楚说明 ELBO 最大化试图实现什么：一方面提高模型证据的下界，另一方面同时减小近似后验与真实后验之间的失配。

\paragraph{ELBO 的第三种形式：重构加正则化。}

为了得到实践中最常用的形式，把联合分布分解为
\[
p_{\theta}(\mathbf{x},\mathbf{z})
=
p_{\theta}(\mathbf{x}\mid\mathbf{z})\,p(\mathbf{z}).
\]
将其代入 ELBO，得到
\[
\begin{aligned}
\mathcal{L}(\mathbf{x};\theta,\phi)
&=
\mathbb{E}_{q_{\phi}(\mathbf{z}\mid\mathbf{x})}
\!\left[
\log\frac{p_{\theta}(\mathbf{x}\mid\mathbf{z})\,p(\mathbf{z})}
{q_{\phi}(\mathbf{z}\mid\mathbf{x})}
\right] \\[4pt]
&=
\mathbb{E}_{q_{\phi}(\mathbf{z}\mid\mathbf{x})}
\!\left[\log p_{\theta}(\mathbf{x}\mid\mathbf{z})\right]
+
\mathbb{E}_{q_{\phi}(\mathbf{z}\mid\mathbf{x})}
\!\left[
\log p(\mathbf{z})-\log q_{\phi}(\mathbf{z}\mid\mathbf{x})
\right].
\end{aligned}
\]
第二个期望同样是一个负 KL 散度，因此得到
\[
\boxed{
\mathcal{L}(\mathbf{x};\theta,\phi)
=
\mathbb{E}_{q_{\phi}(\mathbf{z}\mid\mathbf{x})}
\!\left[\log p_{\theta}(\mathbf{x}\mid\mathbf{z})\right]
-
\mathrm{KL}\!\bigl(
q_{\phi}(\mathbf{z}\mid\mathbf{x})
\,\|\,\,
p(\mathbf{z})
\bigr).
}
\]

这就是最熟悉的 VAE 目标。它包含作用不同的两个部分：
\begin{itemize}
    \item \textbf{重构项}
    \[
    \mathbb{E}_{q_{\phi}(\mathbf{z}\mid\mathbf{x})}
    [\log p_{\theta}(\mathbf{x}\mid \mathbf{z})],
    \]
    它鼓励解码器在给定采样潜编码时为观测数据赋予较高概率，因此起到重构目标的作用；

    \item \textbf{正则化项}
    \[
    \mathrm{KL}\!\bigl(
    q_{\phi}(\mathbf{z}\mid \mathbf{x})
    \,\|\,\,
    p(\mathbf{z})
    \bigr),
    \]
    它鼓励近似后验保持接近先验。
\end{itemize}

第一项通过解码器促进准确重构。
第二项则防止编码器将每个数据点任意放置在潜空间中。
相反，它鼓励推断得到的潜变量分布与所假定的先验保持相容，从而使潜空间具有更加连贯的全局结构。

正是这种正则化使生成式采样具有实际意义。
训练完成后，我们可以采样
\[
\mathbf{z}\sim p(\mathbf{z}),
\qquad
\mathbf{x}\sim p_{\theta}(\mathbf{x}\mid \mathbf{z}),
\]
并期望生成的样本与数据分布相似。

这也回应了 PPCA 一章末尾指出的局限。
在 PPCA 中，模型虽然引入了先验，但并没有显式地对推断得到的潜空间施加正则约束，使其与先验对齐。
而在 VAE 中，这种对齐被直接纳入学习目标之中。

上面推导出的三种形式在代数上完全等价，只是分别强调了同一目标的不同方面。图~\ref{fig:vae-elbo-three-views}
在我们总结 ELBO 最大化的作用之前，将这三种视角汇集在一起。

\begin{figure}[H]
\centering
\begin{tikzpicture}[
  >=Latex,
  every node/.style={font=\small},
  elbo/.style={
    draw,
    rounded corners=2pt,
    align=center,
    inner sep=7pt,
    text width=12.2cm
  },
  view/.style={
    draw,
    rounded corners=2pt,
    align=center,
    inner sep=7pt,
    text width=5.75cm,
    minimum height=2.35cm
  },
  flow/.style={->,thick}
]

\node[elbo] (joint) at (0,2.25) {
  \textbf{联合密度形式}\\[4pt]
  $\displaystyle
  \mathcal{L}(\mathbf{x};\theta,\phi)
  =
  \mathbb{E}_{q_{\phi}(\mathbf{z}\mid\mathbf{x})}
  \!\left[
  \log
  \frac{p_{\theta}(\mathbf{x},\mathbf{z})}
       {q_{\phi}(\mathbf{z}\mid\mathbf{x})}
  \right]$
};

\node[view] (inf) at (-3.45,-1.2) {
  \textbf{推断视角}\\[4pt]
  $\displaystyle
  \begin{aligned}
  \mathcal{L}
  &=
  \log p_{\theta}(\mathbf{x})\\[-1pt]
  &\quad-
  \mathrm{KL}\!\left(
  q_{\phi}(\mathbf{z}\mid\mathbf{x})
  \,\|\,p_{\theta}(\mathbf{z}\mid\mathbf{x})
  \right)
  \end{aligned}$\\[5pt]
  精确对数似然与 ELBO 之间的差距来自
  近似后验与真实后验的不匹配。
};

\node[view] (train) at (3.45,-1.2) {
  \textbf{训练视角}\\[4pt]
  $\displaystyle
  \begin{aligned}
  \mathcal{L}
  &=
  \mathbb{E}_{q_{\phi}(\mathbf{z}\mid\mathbf{x})}
  [\log p_{\theta}(\mathbf{x}\mid\mathbf{z})]\\[-1pt]
  &\quad-
  \mathrm{KL}\!\left(
  q_{\phi}(\mathbf{z}\mid\mathbf{x})
  \,\|\,p(\mathbf{z})
  \right)
  \end{aligned}$\\[5pt]
  期望重构对数似然与
  朝先验的正则化相互权衡。
};

\draw[flow] (joint) -- node[left] {\scriptsize Bayes 规则} (inf);
\draw[flow] (joint) -- node[right] {\scriptsize 分解联合分布} (train);

\node[font=\normalsize\bfseries] at (0,3.45)
  {同一个 ELBO 的三种视角};

\end{tikzpicture}
\caption{\textbf{VAE ELBO 的三种等价视角。}
联合密度形式是一般变分目标。利用 Bayes 规则可得到推断视角，即“精确对数似然减去后验不匹配”；对联合分布进行分解则得到训练视角，即“期望重构对数似然减去朝潜变量先验的正则化”。三种形式在代数上完全等价。}
\label{fig:vae-elbo-three-views}
\end{figure}

\paragraph{ELBO 实现了什么。}

现在，我们可以从两个互补的视角理解 VAE 的目标。

从推断视角看，ELBO 表达的是：
\begin{quote}
\emph{学习一个既接近真实后验、又能够很好解释数据的近似后验。}
\end{quote}

从表征学习视角看，它表达的是：
\begin{quote}
\emph{学习能够准确重构数据的潜编码，同时使潜空间保持有组织的结构，并与一个简单先验相容。}
\end{quote}

这就是为什么 VAE 并不是“加入噪声的自编码器”那么简单。
它的目标并非仅仅完成重构，而是在潜变量生成模型内部执行近似概率推断。

与此同时，这一目标也带来了新的计算挑战。
重构项是关于随机潜变量的期望，而该潜变量是从
$q_{\phi}(\mathbf{z}\mid\mathbf{x})$ 中采样得到的。
为了关于编码器参数 $\phi$ 优化 ELBO，并由此将近似后验与解码器联合训练，我们需要一种能够让梯度穿过这一采样过程的方法。

这正是重参数化技巧的用途，也是下一节的出发点。

\section{学习变分自编码器：重参数化与优化}

上一节已经说明，VAE 通过最大化 ELBO 来训练：
\[
\mathcal{L}(\mathbf{x};\theta,\phi)
=
\mathbb{E}_{q_{\phi}(\mathbf{z}\mid\mathbf{x})}
\!\left[
\log p_{\theta}(\mathbf{x}\mid\mathbf{z})
\right]
-
\mathrm{KL}\!\bigl(
q_{\phi}(\mathbf{z}\mid\mathbf{x})
\,\|\,\,
p(\mathbf{z})
\bigr).
\]
剩下的问题是：这个目标中包含对编码器分布所生成随机潜样本的期望，那么在实践中应如何对它进行优化？

ELBO 的重构项包含对潜变量样本的期望：
\[
\mathbb{E}_{q_{\phi}(\mathbf{z}\mid\mathbf{x})}
\!\left[
\log p_{\theta}(\mathbf{x}\mid\mathbf{z})
\right].
\]
这意味着在训练过程中，我们必须先从编码器分布中采样 $\mathbf{z}$，再将该样本送入解码器。

乍看之下，这似乎很直接：先采样
\[
\mathbf{z}\sim q_{\phi}(\mathbf{z}\mid\mathbf{x}),
\]
计算重构项，然后进行优化。
困难在于，采样所依据的分布本身依赖编码器参数 $\phi$。
如果把 $\mathbf{z}\sim q_{\phi}(\mathbf{z}\mid\mathbf{x})$
视为一个黑盒随机操作，那么样本 $\mathbf{z}$ 对 $\phi$ 的依赖就被隐藏在采样步骤内部，普通反向传播因此无法穿过编码器中的这一过程。

换句话说，一旦给定 $\mathbf{z}$，对解码器参数 $\theta$ 求导并不困难；真正棘手的是编码器参数 $\phi$，因为它决定了 $\mathbf{z}$ 是\emph{如何}被采样出来的。
因此，如果希望利用基于梯度的优化高效训练 VAE，就需要把这个随机采样步骤改写成可微计算。

\paragraph{重参数化技巧。}
关键思想是把随机性的来源与对编码器参数的依赖分离开来。
对于最常见的对角协方差高斯近似后验，
\[
q_{\phi}(\mathbf{z}\mid\mathbf{x})
=
\mathcal{N}\!\bigl(
\mathbf{z}\mid
\boldsymbol{\mu}_{\phi}(\mathbf{x}),
\operatorname{diag}(\boldsymbol{\sigma}_{\phi}^2(\mathbf{x}))
\bigr),
\]
我们可以先从标准正态分布采样噪声
\[
\boldsymbol{\epsilon}\sim\mathcal{N}(\mathbf{0},\mathbf{I}),
\]
再定义
\[
\mathbf{z}
=
\boldsymbol{\mu}_{\phi}(\mathbf{x})
+
\boldsymbol{\sigma}_{\phi}(\mathbf{x})\odot \boldsymbol{\epsilon}.
\]
这就是重参数化技巧~\cite{kingma2014autoencoding,rezende2014stochastic}。
图~\ref{fig:vae-reparameterisation} 对比了直接随机采样视角与重参数化后的计算方式。

\begin{figure}[H]
\centering
\begin{tikzpicture}[
  >=Latex,
  every node/.style={font=\small},
  box/.style={
    draw,
    rounded corners=2pt,
    align=center,
    inner sep=5pt,
    minimum height=0.92cm,
    text width=2.55cm
  },
  tinybox/.style={
    draw,
    rounded corners=2pt,
    align=center,
    inner sep=4pt,
    minimum height=0.82cm,
    text width=2.0cm
  },
  flow/.style={->,thick},
  stochastic/.style={->,thick,dashed}
]

% ---------------- Direct sampling ----------------
\begin{scope}[shift={(-6.15,0)}]
\node[font=\normalsize\bfseries] at (4.75,2.05)
  {(a) 直接的参数依赖采样};

\node[tinybox] (x1) at (0,0.55) {$\mathbf{x}$};
\node[box] (enc1) at (3.0,0.55) {
  编码器\\[-1pt]
  $\boldsymbol{\mu}_{\phi}(\mathbf{x}),
   \boldsymbol{\sigma}_{\phi}(\mathbf{x})$
};
\node[box] (draw1) at (6.25,0.55) {
  随机采样\\[-1pt]
  $\mathbf{z}\sim q_{\phi}(\mathbf{z}\mid\mathbf{x})$
};
\node[box] (dec1) at (9.5,0.55) {
  解码器\\[-1pt]
  $p_{\theta}(\mathbf{x}\mid\mathbf{z})$
};

\draw[flow] (x1) -- (enc1);
\draw[stochastic] (enc1) -- (draw1);
\draw[flow] (draw1) -- (dec1);

\node[align=center,text width=8.8cm] at (4.75,-0.95) {
把采样视为黑盒随机操作，会遮蔽通往编码器参数 $\phi$ 的常规路径式反向传播通路。
};
\end{scope}

% ---------------- Reparameterised sampling ----------------
\begin{scope}[shift={(-6.15,-4.55)}]
\node[font=\normalsize\bfseries] at (4.75,2.05)
  {(b) 重参数化采样};

\node[tinybox] (x2) at (0,0.55) {$\mathbf{x}$};
\node[box] (enc2) at (3.0,0.55) {
  编码器\\[-1pt]
  $\boldsymbol{\mu}_{\phi}(\mathbf{x}),
   \boldsymbol{\sigma}_{\phi}(\mathbf{x})$
};
\node[tinybox] (eps) at (3.0,-1.05) {
  $\boldsymbol{\epsilon}\sim
  \mathcal{N}(\mathbf{0},\mathbf{I})$
};
\node[box] (combine) at (6.25,0.55) {
  确定性变换\\[-1pt]
  $\mathbf{z}
  =
  \boldsymbol{\mu}_{\phi}
  +
  \boldsymbol{\sigma}_{\phi}
  \odot\boldsymbol{\epsilon}$
};
\node[box] (dec2) at (9.5,0.55) {
  解码器\\[-1pt]
  $p_{\theta}(\mathbf{x}\mid\mathbf{z})$
};

\draw[flow] (x2) -- (enc2);
\draw[flow] (enc2) -- (combine);
\draw[flow] (eps) -- (combine);
\draw[flow] (combine) -- (dec2);

\node[align=center,text width=8.8cm] at (4.75,-2.15) {
在固定 $\boldsymbol{\epsilon}$ 时，$\mathbf{z}$ 是编码器输出的确定性可微函数，因此显式提供了通往 $\phi$ 的路径式梯度。
};
\end{scope}

\end{tikzpicture}
\caption{\textbf{重参数化技巧把随机性与编码器参数分离。}
直接从 $\mathbf{z}\sim q_{\phi}(\mathbf{z}\mid\mathbf{x})$ 采样，会使穿过采样步骤的常规路径式梯度变得不明显。改为采样与参数无关的 $\boldsymbol{\epsilon}\sim\mathcal{N}(\mathbf{0},\mathbf{I})$，并令 $\mathbf{z}=\boldsymbol{\mu}_{\phi}(\mathbf{x})+\boldsymbol{\sigma}_{\phi}(\mathbf{x})\odot\boldsymbol{\epsilon}$，既保持同一个高斯分布，又使固定噪声样本下的 $\mathbf{z}$ 对编码器输出可微。}
\label{fig:vae-reparameterisation}
\end{figure}

这一技巧并没有消除随机性，而是把随机性转移到了与编码器参数无关的辅助噪声 $\boldsymbol{\epsilon}$ 上。
对于固定的一次噪声采样，潜变量就是一个确定性的可微函数
\[
\mathbf{z}=\mathbf{z}(\mathbf{x},\phi,\boldsymbol{\epsilon}),
\]
因此普通的路径导数可以通过标准反向传播穿过编码器进行传播。

\paragraph{编码器与解码器的联合优化。}
有了重参数化技巧之后，就可以通过同时关于两组参数最大化 ELBO，对 VAE 进行端到端训练：
\[
(\theta,\phi)
\quad\longmapsto\quad
\max \; \mathcal{L}(\mathbf{x};\theta,\phi).
\]
实践中，对每个输入，编码器首先输出
$\boldsymbol{\mu}_{\phi}(\mathbf{x})$ 和
$\boldsymbol{\sigma}_{\phi}(\mathbf{x})$；随后通过重参数化构造一个潜变量样本，计算 ELBO，再将梯度同时反向传播经过编码器和解码器。
这种优化通常使用小批量随机梯度方法或其现代变体完成。

这是 VAE 最重要的概念性成果之一。
它把近似后验推断与生成模型学习转化为同一个可微计算图。
从这一点开始，编码器和解码器不再被作为两个彼此分离的对象训练：它们通过 ELBO 相互耦合，并被联合学习。

从这个意义上说，训练 VAE 本身就是一种变分推断。
我们不再精确计算真实后验 $p_{\theta}(\mathbf{z}\mid\mathbf{x})$，而是学习一个易处理的近似
$q_{\phi}(\mathbf{z}\mid\mathbf{x})$，并同时关于解码器参数 $\theta$ 与编码器参数 $\phi$ 优化 ELBO。

\paragraph{一种变体：$\beta$-VAE。}

重构与正则化之间的平衡可以采用不同方式调节。其中一个重要扩展是
$\beta$-VAE~\cite{higgins2017beta}，它在 KL 项前引入一个权重系数：
\[
\mathcal{L}_{\beta}
=
\mathbb{E}_{q_{\phi}(\mathbf{z}\mid \mathbf{x})}
\!\left[
\log p_{\theta}(\mathbf{x}\mid \mathbf{z})
\right]
-
\beta\,
\mathrm{KL}\!\bigl(
q_{\phi}(\mathbf{z}\mid \mathbf{x})
\,\|\,\,
p(\mathbf{z})
\bigr).
\]
当 $\beta>1$ 时，潜空间受到更强的正则化，往往会得到更有结构、或更具解耦性的表征，但通常要以重构质量下降为代价。

这一扩展尤其清楚地揭示了变分潜变量模型中的一个核心张力：
\begin{quote}
\emph{潜空间既必须足够规整以支持生成，又必须具有足够的信息量以支持有用的表征与重构。}
\end{quote}

从这一视角来看，VAE 标志着现代生成建模发展中的重要一步。
PPCA 引入了一个可解析处理的概率潜变量模型；VAE 则通过以学习得到的近似取代精确推断，并把 ELBO 转化为可实际优化的训练目标，将同样的潜变量思想推进到非线性神经网络场景。

这些思想在后续章节中仍将发挥作用。
尤其是，尽管扩散模型经常从其他角度引入，但基于 ELBO 的推理方式和潜变量视角同样为理解它们提供了有用框架。
下一章中，我们将看到这些概率思想如何进一步扩展到现代生成式人工智能最重要的框架之一。

\section*{本章小结与参考文献}
\addcontentsline{toc}{section}{本章小结与参考文献}

本章将变分自编码器（VAE）介绍为自编码器的概率扩展，同时也是通过 PPCA 建立起来的潜变量框架在非线性情形下的延续。

由于模型是非线性的，边缘似然
$p_{\theta}(\mathbf{x})$ 与真实后验
$p_{\theta}(\mathbf{z}\mid\mathbf{x})$ 一般都不存在闭式表达。
随后我们看到，ELBO 有若干种等价形式，其中最常见的一种可分解为重构项与正则化项。
由此可以清楚地看到，VAE 并非只是重构输入，而是在学习一种与简单先验保持相容的潜在表征，从而支持有意义的生成。

最后，本章说明了重参数化技巧如何把随机潜变量采样改写为对不依赖模型参数的噪声所施加的可微变换，使编码器和解码器能够利用基于梯度的优化进行联合训练。

最初的 VAE 形式由 Kingma 和 Welling 提出~\cite{kingma2014autoencoding}，Rezende、Mohamed 与 Wierstra 也发展了密切相关的随机梯度变分方法~\cite{rezende2014stochastic}。
有关更广泛的概率建模背景，可参见 Bishop~\cite{bishop2006prml} 与 Murphy~\cite{murphy2012ml}。
关于 VAE 的系统性介绍以及较新进展的综述，可参见 Kingma 和 Welling~\cite{kingma2019introduction}。

\chapter{去噪扩散概率模型}
\section*{概述}
\addcontentsline{toc}{section}{概述}

扩散模型是现代生成模型中最重要的模型族之一，在广泛的应用中取得了出色性能，尤其是在图像生成领域。
其核心思想简单而强大：不是一步生成一个样本，而是学习通过一系列去噪操作，将噪声逐渐转化为具有结构的数据。

本章聚焦于离散时间的\emph{去噪扩散概率模型}（Denoising Diffusion Probabilistic Model，DDPM）框架~\cite{ho2020denoising}，它为这一思想提供了最清晰的概率建模形式之一。
从历史上看，基于扩散的生成建模由 Sohl-Dickstein 等人~\cite{sohl2015deep} 开创，其灵感来自物理学中的一些过程，尤其是随机性随时间逐渐累积的过程。
例如在物理世界中，一滴墨水进入水中后会缓慢扩散，直到原有结构与周围环境难以区分。
扩散模型借用了这种直觉，并将过程反转：我们不再观察有序结构如何消解为噪声，而是学习如何从噪声出发，逐步重建结构。

不过在本书中，我们首先从 PPCA 与 VAE 已经建立的概率建模和潜变量视角来理解扩散模型。
这为后续引入连续时间动力学视角提供了一个自然入口。
随后在第~5 章中，我们将重新讨论相关的微积分、动力系统以及密度演化视角，并在第~6 章逐步构建更连续的解释。

从方法论角度看，扩散模型与此前的潜变量模型仍具有重要的结构联系。
关键区别在于：扩散模型并不只使用一个潜变量或一次随机编码，而是引入一整条随时间相连的潜状态序列。
实践中，这使扩散模型包含两个彼此衔接的阶段：一个固定的前向过程，将数据逐步扰动为噪声；以及一个学习得到的反向去噪过程，从噪声逐步恢复结构。
这种时间化设计赋予模型更强的灵活性，同时也改变了推断与学习的形式。

\section*{概念导图}
\addcontentsline{toc}{section}{概念导图}

本章延续前面章节的建模逻辑，利用潜变量、后验推断和变分学习来解释序贯加噪与去噪过程。

\begin{itemize}

    \item[\(\rightarrow\)] \textbf{首先从前向扩散过程开始。}
    干净样本 $\mathbf{x}_0$ 在 $T$ 个步骤中逐渐受到高斯噪声污染，形成潜在轨迹
    $\mathbf{x}_1,\dots,\mathbf{x}_T$。
    这定义了一条固定的马尔可夫链，并给出了任意噪声状态 $\mathbf{x}_t$ 与原始样本 $\mathbf{x}_0$ 之间的闭式关系。

    \item[\(\rightarrow\)] \textbf{随后把反向过程引入为需要学习的生成模型。}
    模型从时刻 $T$ 的简单高斯先验出发，学习一条反向去噪链，以恢复数据分布。
    这条反向链是扩散模型中的未知部分，也是学习的主要对象。

    \item[\(\rightarrow\)] \textbf{接着，把扩散模型改写为 ELBO 形式。}
    观测变量是干净样本 $\mathbf{x}_0$，潜变量则是整条轨迹 $\mathbf{x}_{1:T}$。
    已知的前向过程在 ELBO 中扮演辅助分布的角色，从而把难以直接处理的似然转化为可处理的下界。

    \item[\(\rightarrow\)] \textbf{然后，将 ELBO 分解为局部去噪项。}
    由此可以看出，每个时间步上真正需要训练的量，是学习得到的反向转移
    $p_\theta(\mathbf{x}_{t-1}\mid \mathbf{x}_t)$ 与真实逆后验
    $q(\mathbf{x}_{t-1}\mid \mathbf{x}_t,\mathbf{x}_0)$ 之间的差异。

    \item[\(\rightarrow\)] \textbf{随后推导真实逆后验，并在噪声域中重新表达。}
    由于前向过程是高斯过程，真实逆后验同样是高斯分布。
    通过噪声变量对其均值进行重参数化，可以消除对 $\mathbf{x}_0$ 的显式依赖，并得到标准 DDPM 参数化。

    \item[\(\rightarrow\)] \textbf{最后得到实际使用的 DDPM 目标与生成过程。}
    学习问题可简化为噪声预测，因此训练变成对高斯噪声的回归任务；而采样则成为一个迭代去噪过程，将纯噪声逐步转化为数据样本。

\end{itemize}

\medskip
\noindent 下图对 DDPM 的核心概率结构进行了可视化概括。
\begin{center}
\begin{tikzpicture}[
  >=Latex,
  every node/.style={font=\small},
  state/.style={draw,rounded corners=2pt,minimum width=1.65cm,minimum height=0.82cm,align=center},
  wide/.style={draw,rounded corners=2pt,text width=4.2cm,minimum height=1.02cm,align=center,inner sep=6pt},
  flow/.style={->,thick},
  rev/.style={->,thick,dashed}
]
\node[font=\normalsize\bfseries] at (6.3,2.0) {DDPM 概览};

\node[state] (x0) at (0,0.65) {$\mathbf{x}_0$\\干净数据};
\node[state] (x1) at (3.0,0.65) {$\mathbf{x}_1$};
\node at (5.0,0.65) {$\cdots$};
\node[state] (xt) at (7.0,0.65) {$\mathbf{x}_t$};
\node at (9.0,0.65) {$\cdots$};
\node[state] (xT) at (12.0,0.65) {$\mathbf{x}_T$\\高斯噪声};

\draw[flow] (x0) -- node[above] {\scriptsize 固定的 $q$} (x1);
\draw[flow] (x1) -- (4.45,0.65);
\draw[flow] (5.55,0.65) -- (xt);
\draw[flow] (xt) -- (8.45,0.65);
\draw[flow] (9.55,0.65) -- node[above] {\scriptsize 加噪} (xT);

\draw[rev] (xT.south) to[out=-135,in=-45] node[midway,above] {\scriptsize 学习到的反向过程 $p_\theta$} (xt.south);
\draw[rev] (xt.south) to[out=-135,in=-45] (x1.south);
\draw[rev] (x1.south) to[out=-135,in=-45] node[below] {\scriptsize 去噪} (x0.south);

\node[wide] (shortcut) at (3.25,-2.15) {
\textbf{训练捷径}\\[2pt]
采样 $t,\boldsymbol{\epsilon}$，直接跳到
$\displaystyle
\mathbf{x}_t=
\sqrt{\bar{\alpha}_t}\mathbf{x}_0+
\sqrt{1-\bar{\alpha}_t}\boldsymbol{\epsilon}$
};

\node[wide] (learn) at (9.2,-2.15) {
\textbf{学习}\\[2pt]
由 $(\mathbf{x}_t,t)$ 预测 $\boldsymbol{\epsilon}$，使
$p_\theta(\mathbf{x}_{t-1}\mid\mathbf{x}_t)$ 匹配局部反向目标
};

\draw[flow] (shortcut) -- (learn);
\end{tikzpicture}
\end{center}
\medskip

\section{前向扩散过程：从数据到噪声}

扩散模型与此前学习的生成框架相比，有一个重要不同。
在 PPCA 和 VAE 中，每个观测通常通过一个相对直接的潜在表示来描述。
相比之下，在扩散模型中，一个数据点对应于一个\emph{潜状态序列}，它记录了随时间逐渐发生的变换。
因此，生成不再是一次性的解码步骤，而是一条轨迹。

从整体上看，DDPM 包含两个互补过程。
\textbf{前向过程}逐渐破坏干净数据样本 $\mathbf{x}_0$，直到它几乎成为纯高斯噪声。
\textbf{反向过程}，也就是\textbf{去噪过程}，试图逐步撤销这种破坏：从噪声出发，一步步恢复具有结构的样本。

我们先从模型中较简单的一半——前向过程——开始。
在扩散模型文献中，这一过程通常记为 $q$。
它是固定且已知的，而不是从数据中学习得到的。
它的作用是规定如何在 $T$ 个离散时间步中逐渐扰动干净样本 $\mathbf{x}_0$，得到一系列噪声越来越强的潜状态
\[
\mathbf{x}_1,\mathbf{x}_2,\dots,\mathbf{x}_T.
\]

\subsection{高斯加噪与前向过程的闭式形式}

\paragraph{逐步高斯加噪过程。}
设 $\mathbf{x}_0$ 表示一个观测数据样本。
前向扩散过程在每一步加入少量高斯噪声，从而产生上述序列。
形式上，从 $\mathbf{x}_{t-1}$ 到 $\mathbf{x}_t$ 的转移定义为
\[
q(\mathbf{x}_t\mid \mathbf{x}_{t-1})
=
\mathcal{N}\!\bigl(
\mathbf{x}_t \mid
\sqrt{1-\beta_t}\,\mathbf{x}_{t-1},
\beta_t\mathbf{I}
\bigr),
\]
其中 $\beta_t\in(0,1)$ 控制第 $t$ 步注入多少噪声。

等价地，同一步骤可以写成
\[
\mathbf{x}_t
=
\sqrt{1-\beta_t}\,\mathbf{x}_{t-1}
+
\sqrt{\beta_t}\,\mathbf{z}_t,
\qquad
\mathbf{z}_t\sim\mathcal{N}(\mathbf{0},\mathbf{I}).
\]
这个更新有一个直观解释：
$\sqrt{1-\beta_t}$ 保留上一状态中的一部分信号，而
$\sqrt{\beta_t}$ 对当前步骤新加入的高斯扰动进行缩放。
因此，每个新状态 $\mathbf{x}_t$ 都比前一个状态稍微更嘈杂。

这里，$\beta_t$ 控制第 $t$ 步加入多少新的高斯噪声，因而也决定了随着过程推进，原始数据结构被破坏的速度。

\paragraph{$\mathbf{x}_t$ 与 $\mathbf{x}_0$ 之间的闭式联系。}
上述单步规则描述了 $\mathbf{x}_t$ 如何依赖 $\mathbf{x}_{t-1}$，但在很多情况下，更有用的是直接建立 $\mathbf{x}_t$ 与原始干净样本 $\mathbf{x}_0$ 的关系。
事实证明，这是前向扩散过程最有用的性质之一，因为它允许我们直接跳到任意时间步 $t$，而无需模拟整条链。

从
\[
\mathbf{x}_t
=
\sqrt{1-\beta_t}\,\mathbf{x}_{t-1}
+
\sqrt{\beta_t}\,\mathbf{z}_t,
\qquad
\mathbf{z}_t\sim\mathcal{N}(\mathbf{0},\mathbf{I}),
\]
出发，我们可以递归地把较早时间步代入后面的时间步。

例如，
\[
\begin{aligned}
\mathbf{x}_1
&=
\sqrt{1-\beta_1}\,\mathbf{x}_0
+
\sqrt{\beta_1}\,\mathbf{z}_1,\\[4pt]
\mathbf{x}_2
&=
\sqrt{1-\beta_2}\,\mathbf{x}_1
+
\sqrt{\beta_2}\,\mathbf{z}_2.
\end{aligned}
\]
将第一式代入第二式可得
\[
\mathbf{x}_2
=
\sqrt{(1-\beta_2)(1-\beta_1)}\,\mathbf{x}_0
+
\sqrt{1-\beta_2}\sqrt{\beta_1}\,\mathbf{z}_1
+
\sqrt{\beta_2}\,\mathbf{z}_2.
\]

此时定义
\[
\alpha_t = 1-\beta_t,
\qquad
\bar{\alpha}_t = \prod_{s=1}^{t}\alpha_s.
\]
其中，$\alpha_t$ 是第 $t$ 步的信号保留因子，而 $\bar{\alpha}_t$ 是截至时刻 $t$ 的累积信号保留因子。

现在规律已经很清楚。
每增加一个时间步，剩余信号都会再乘上一个因子
$\sqrt{1-\beta_t}=\sqrt{\alpha_t}$，同时加入新的高斯扰动。
经过 $t$ 步之后，从 $\mathbf{x}_0$ 继承下来的信号因此被缩放为
\[
\prod_{s=1}^{t}\sqrt{1-\beta_s}
=
\prod_{s=1}^{t}\sqrt{\alpha_s}
=
\sqrt{\bar{\alpha}_t}.
\]

与此同时，许多彼此独立的高斯噪声项可以合并成一个整体高斯变量。
由于独立高斯变量的加权和仍然是高斯变量，累积噪声可以写成一个标准高斯变量，再由总方差因子 $1-\bar{\alpha}_t$ 进行缩放。
于是得到闭式关系
\[
\mathbf{x}_t
=
\sqrt{\bar{\alpha}_t}\,\mathbf{x}_0
+
\sqrt{1-\bar{\alpha}_t}\,\boldsymbol{\epsilon},
\qquad
\boldsymbol{\epsilon}\sim\mathcal{N}(\mathbf{0},\mathbf{I}).
\]

以概率分布形式写为
\[
q(\mathbf{x}_t\mid \mathbf{x}_0)
=
\mathcal{N}\!\left(
\mathbf{x}_t \mid
\sqrt{\bar{\alpha}_t}\,\mathbf{x}_0,\,
(1-\bar{\alpha}_t)\mathbf{I}
\right).
\]

这里的记号很重要：
\begin{itemize}
    \item $\mathbf{z}_t$ 表示第 $t$ 步新加入的\emph{局部}高斯噪声；
    \item $\boldsymbol{\epsilon}$ 表示\emph{整体累积}噪声，它概括了截至时刻 $t$ 的所有加噪步骤的综合作用。
\end{itemize}

上述闭式表达非常有用。
它告诉我们，经过 $t$ 步之后，样本 $\mathbf{x}_t$ 中仍保留原始信号的
$\sqrt{\bar{\alpha}_t}$ 比例，而剩余部分已经被高斯噪声取代。
当 $t=0$ 时，$\bar{\alpha}_0=1$，因此 $\mathbf{x}_t=\mathbf{x}_0$。
当 $t=T$ 足够大，使得 $\bar{\alpha}_T\approx 0$ 时，反复的高斯扰动已经冲淡原始结构，使样本接近各向同性高斯噪声。

这种 $\mathbf{x}_t$ 与 $\mathbf{x}_0$ 之间的直接关系会使后续训练高效得多，因为我们可以在任意时间步直接采样噪声状态，而无需模拟整条轨迹。

图~\ref{fig:ddpm-forward-closed-form} 对这两种效应进行了可视化总结：随着噪声系数增大，信号系数逐渐减小，而边缘分布也逐渐失去其原始结构。

\begin{figure}[H]
\centering
\includegraphics[width=0.98\textwidth]{figures/ch04/figure4_1_forward_diffusion_closed_form.pdf}
\caption{\textbf{前向扩散及其闭式加噪关系。}
上排展示在示例线性噪声日程下若干扩散时刻的边缘样本。随着 $t$ 增大，数据结构逐渐被高斯噪声淹没；下方曲线则显示信号系数下降、噪声系数上升。各快照直接从 $q(\mathbf{x}_t\mid\mathbf{x}_0)$ 采样，因此表示不同时刻的分布，而不是同一条已实现的 Markov 轨迹。}
\label{fig:ddpm-forward-closed-form}
\end{figure}

\subsection{前向链与马尔可夫假设}

由于扩散过程引入了一整条潜状态序列，完整的前向轨迹可以由联合分布
\[
q(\mathbf{x}_{0:T}) = q(\mathbf{x}_0)\,q(\mathbf{x}_{1:T}\mid \mathbf{x}_0)
\]
描述。
根据概率的链式法则，
\[
q(\mathbf{x}_{0:T})
=
q(\mathbf{x}_0)
\prod_{t=1}^{T}
q(\mathbf{x}_t\mid \mathbf{x}_{t-1},\mathbf{x}_{t-2},\dots,\mathbf{x}_0).
\]

关键的简化假设是前向过程具有\textbf{马尔可夫性}。
这意味着每个新状态只依赖于紧邻的前一状态：
\[
q(\mathbf{x}_t\mid \mathbf{x}_{t-1},\mathbf{x}_{t-2},\dots,\mathbf{x}_0)
=
q(\mathbf{x}_t\mid \mathbf{x}_{t-1}).
\]
因此，完整的前向分布可分解为
\[
q(\mathbf{x}_{0:T})
=
q(\mathbf{x}_0)\prod_{t=1}^{T}q(\mathbf{x}_t\mid \mathbf{x}_{t-1}).
\]

这使前向过程在数学上易于处理和分析。
但它并不是主要的学习对象。
一旦数据样本 $\mathbf{x}_0$ 和噪声调度
$\{\beta_t\}_{t=1}^{T}$ 给定，前向过程 $q$ 就完全已知。

\begin{tcolorbox}[
  colback=blue!6,
  colframe=blue!55!black,
  title={前向过程小结},
  boxrule=0.8pt,
  arc=2mm,
  left=2mm,
  right=2mm,
  top=1mm,
  bottom=1mm
]
\[
q(\mathbf{x}_t\mid \mathbf{x}_{t-1})
=
\mathcal{N}\!\bigl(
\mathbf{x}_t\mid
\sqrt{1-\beta_t}\,\mathbf{x}_{t-1},
\beta_t\mathbf{I}
\bigr),
\qquad
\alpha_t=1-\beta_t,
\qquad
\bar{\alpha}_t=\prod_{s=1}^{t}\alpha_s.
\]

\[
\mathbf{x}_t
=
\sqrt{\bar{\alpha}_t}\,\mathbf{x}_0
+
\sqrt{1-\bar{\alpha}_t}\,\boldsymbol{\epsilon},
\qquad
\boldsymbol{\epsilon}\sim\mathcal{N}(\mathbf{0},\mathbf{I}),
\]
因此等价地有
\[
q(\mathbf{x}_t\mid \mathbf{x}_0)
=
\mathcal{N}\!\left(
\mathbf{x}_t\mid
\sqrt{\bar{\alpha}_t}\,\mathbf{x}_0,\,
(1-\bar{\alpha}_t)\mathbf{I}
\right).
\]

\[
q(\mathbf{x}_{0:T})
=
q(\mathbf{x}_0)\prod_{t=1}^{T}q(\mathbf{x}_t\mid \mathbf{x}_{t-1}).
\]
\end{tcolorbox}

然而，对生成而言真正重要的是\textbf{反向}过程，也就是\textbf{去噪}过程：
从一个随机高斯点出发，逐渐重建出类似真实数据的样本。
这一反向过程没有已知的解析形式，必须学习，通常由神经网络来参数化。
为了理解它，我们现在重新引入前面发展出的 ELBO 与变分推断工具，并将其适配到这一随时间变化的潜变量序列，从而使反向过程能够被表达、分析并最终训练。

\section{反向过程与 ELBO 形式化}

前向过程为我们提供了一种固定且易处理的方法：把干净数据样本 $\mathbf{x}_0$ 变换成噪声逐渐增大的潜状态序列
$\mathbf{x}_1,\dots,\mathbf{x}_T$，最终接近高斯噪声。
但这一过程并不是我们真正想学习的部分。
对于生成而言，重要的是相反的方向：从噪声出发逐步恢复结构。
这就是\textbf{反向过程}的作用，通常记为 $p$。

从整体上看，反向过程试图撤销前向链引入的扰动。
它从一个简单的高斯样本 $\mathbf{x}_T$ 开始，逐步生成
$\mathbf{x}_{T-1},\mathbf{x}_{T-2},\dots,\mathbf{x}_0$，目标是使最终状态 $\mathbf{x}_0$ 看起来像真实的数据样本。
换句话说，如果前向过程沿着
\[
\mathbf{x}_0 \longrightarrow \mathbf{x}_1 \longrightarrow \cdots \longrightarrow \mathbf{x}_T
\]
演化，那么反向过程要建模的就是相反方向
\[
\mathbf{x}_T \longrightarrow \mathbf{x}_{T-1} \longrightarrow \cdots \longrightarrow \mathbf{x}_0.
\]

\subsection{需要学习的反向链}

为了从概率角度描述生成，我们引入一个参数化的反向过程。
根据一般的概率链式法则，整条反向轨迹的联合分布原则上可以写成
\[
\begin{aligned}
p_\theta(\mathbf{x}_{0:T})
&= p_\theta(\mathbf{x}_0,\mathbf{x}_{1:T}) \\[4pt]
&= p_\theta(\mathbf{x}_0\mid \mathbf{x}_{1:T})\,p_\theta(\mathbf{x}_{1:T}) \\[4pt]
&= p_\theta(\mathbf{x}_0\mid \mathbf{x}_{1:T})\,
   p_\theta(\mathbf{x}_1\mid \mathbf{x}_{2:T})\,
   p_\theta(\mathbf{x}_{2:T}) \\[4pt]
&= p_\theta(\mathbf{x}_0\mid \mathbf{x}_{1:T})\,
   p_\theta(\mathbf{x}_1\mid \mathbf{x}_{2:T})\,
   p_\theta(\mathbf{x}_2\mid \mathbf{x}_{3:T})\,
   \cdots\,
   p_\theta(\mathbf{x}_{T-1}\mid \mathbf{x}_T)\,
   p(\mathbf{x}_T).
\end{aligned}
\]
这是最一般的展开形式：每个更早的状态都可能依赖所有更晚的状态。

不过在扩散模型中，我们与前向过程一样采用简化假设，认为反向链同样具有\textbf{马尔可夫性}。
即
\[
p_\theta(\mathbf{x}_{t-1}\mid \mathbf{x}_{t:T})
=
p_\theta(\mathbf{x}_{t-1}\mid \mathbf{x}_t).
\]
在这一一阶马尔可夫假设下，反向联合分布简化为
\[
p_\theta(\mathbf{x}_{0:T})
=
p(\mathbf{x}_T)\prod_{t=1}^{T} p_\theta(\mathbf{x}_{t-1}\mid \mathbf{x}_t).
\]
这个表达式应当从右向左阅读。
我们先从最终时间步的一个简单先验开始，通常取
\[
p(\mathbf{x}_T)=\mathcal{N}(\mathbf{0},\mathbf{I}),
\]
然后连续应用学习得到的反向转移
$p_\theta(\mathbf{x}_{t-1}\mid \mathbf{x}_t)$，直到最终得到时刻 $0$ 的样本。

这就是模型中关键的未知对象。
前向过程 $q$ 在选定噪声调度 $\{\beta_t\}$ 后就完全确定，而反向转移并没有已知的闭式形式，因此必须从数据中学习。
实践中，参数 $\theta$ 属于一个神经网络，该网络将被训练来建模这些反向去噪步骤。

\subsection{为什么反向过程很困难}

乍看之下，我们似乎只需要把前向规则直接反转即可。
但事情并没有这么简单。
尽管前向转移
\[
q(\mathbf{x}_t\mid \mathbf{x}_{t-1})
\]
是高斯分布且有明确解析定义，相应的反向转移
\[
q(\mathbf{x}_{t-1}\mid \mathbf{x}_t)
\]
却不能直接以可用的形式得到。

原因可以通过贝叶斯公式看出：
\[
q(\mathbf{x}_{t-1}\mid \mathbf{x}_t)
=
\frac{
q(\mathbf{x}_t\mid \mathbf{x}_{t-1})\,q(\mathbf{x}_{t-1})
}{
q(\mathbf{x}_t)
}.
\]
前向条件分布 $q(\mathbf{x}_t\mid \mathbf{x}_{t-1})$ 是按构造已知的。
困难在于边缘分布 $q(\mathbf{x}_{t-1})$ 与 $q(\mathbf{x}_t)$。
它们并不是人为指定的简单固定高斯分布，而是由未知数据分布 $q(\mathbf{x}_0)$ 经过前向加噪链推送后诱导得到的。
因此，尽管单个前向步骤很简单，真实反向条件分布却隐式依赖底层数据分布，无法作为一个可直接计算或采样的简单规则使用。

当考虑基于似然的学习时，还会出现第二个困难。
模型定义了整条反向轨迹 $\mathbf{x}_{0:T}$ 上的联合分布，但我们最终关心的是观测到的干净数据的边缘似然：
\[
p_\theta(\mathbf{x}_0)
=
\int p_\theta(\mathbf{x}_{0:T})\,d\mathbf{x}_{1:T}.
\]
这正是此前潜变量模型中边缘似然在扩散模型中的对应形式。
观测变量是干净样本 $\mathbf{x}_0$，而整条轨迹 $\mathbf{x}_{1:T}$ 扮演潜变量的角色，必须被积分消去。

$p_\theta(\mathbf{x}_0)$ 是我们希望学习和优化的关键量，因为它衡量模型给真实数据样本赋予了多大概率。
如果一个生成模型能够为观测数据赋予较高概率，就意味着它对数据域上的分布学得较好。
然而在扩散模型中，这个边缘似然只通过反向链
\[
p_\theta(\mathbf{x}_{0:T})
=
p(\mathbf{x}_T)\prod_{t=1}^{T} p_\theta(\mathbf{x}_{t-1}\mid \mathbf{x}_t)
\]
间接定义。
所以，把反向转移 $p_\theta(\mathbf{x}_{t-1}\mid \mathbf{x}_t)$ 学好，正是我们构建一个能够给 $\mathbf{x}_0$ 赋予高似然模型的方式。

困难在于，计算 $p_\theta(\mathbf{x}_0)$ 需要对整条潜在轨迹 $\mathbf{x}_{1:T}$ 积分。
由于这条轨迹既是高维的，又带有时间索引，这个积分通常不可处理。

因此，正如 VAE 中一样，核心学习问题变成：
\begin{quote}
\emph{当精确边缘似然不存在闭式形式时，我们应如何优化模型？}
\end{quote}
答案仍然是使用\textbf{证据下界（ELBO）}。

\subsection{扩散模型的 ELBO}
\paragraph{在扩散模型中重新表达 ELBO。}
此时，把扩散模型重新与前面章节建立的一般 ELBO 框架联系起来会很有帮助。

在 PPCA 和 VAE 中，我们把边缘似然写成一般的潜变量形式
\[
p(\mathbf{x}) = \int p(\mathbf{x},\mathbf{z})\,d\mathbf{z},
\]
其中 $\mathbf{x}$ 是观测变量，$\mathbf{z}$ 是潜变量。
在扩散模型中，同样的结构仍然存在，只是记号具体化为
\[
\mathbf{x}\;\longrightarrow\;\mathbf{x}_0,
\qquad
\mathbf{z}\;\longrightarrow\;\mathbf{x}_{1:T}.
\]
也就是说，观测变量是干净数据样本 $\mathbf{x}_0$，而潜变量不再是单一编码，而是整条噪声轨迹 $\mathbf{x}_{1:T}$。

因此，边缘似然写成
\[
p_\theta(\mathbf{x}_0)
=
\int p_\theta(\mathbf{x}_{0:T})\,d\mathbf{x}_{1:T}.
\]

在一般的 ELBO 推导中，我们还会在潜变量上引入一个辅助分布 $Q(\mathbf{z})$。
在扩散模型中，这一角色自然由
\[
q(\mathbf{x}_{1:T}\mid \mathbf{x}_0)
\]
承担，即\emph{以观测数据为条件的已知前向过程}。

这一点值得强调。
在 VAE 中，由于真实后验难以处理且没有闭式编码器，辅助分布必须以学习得到的近似
$q_\phi(\mathbf{z}\mid \mathbf{x})$ 的形式引入。
而扩散模型的情况不同：
一旦给定 $\mathbf{x}_0$，前向加噪过程已经为整条潜在轨迹提供了完整的条件分布。
也就是说，对给定的干净数据点 $\mathbf{x}_0$，模型已经精确规定了如何采样其噪声版本
$\mathbf{x}_1,\mathbf{x}_2,\dots,\mathbf{x}_T$。
所以，ELBO 所需要的辅助分布并不是额外叠加在模型之上的学习结构；它已经作为前向扩散过程本身内置在模型中。

这是扩散模型非常简洁优雅的性质之一。
“变分分布”已经直接可得：
\[
Q(\mathbf{Z})
\;\longrightarrow\;
q(\mathbf{x}_{1:T}\mid \mathbf{x}_0).
\]
换言之，扩散模型不需要像 VAE 那样额外学习一个用于潜变量的独立编码器网络。
已知前向过程本身已经提供了 ELBO 所需的潜状态条件分布。

\begin{tcolorbox}[
  colback=blue!6,
  colframe=blue!55!black,
  title={ELBO 形式下的扩散模型各量},
  boxrule=0.8pt,
  arc=2mm,
  left=2mm,
  right=2mm,
  top=1mm,
  bottom=1mm
]
\[
\text{观测变量： }\mathbf{x}\rightarrow \mathbf{x}_0,
\qquad
\text{潜变量： }\mathbf{z}\rightarrow \mathbf{x}_{1:T}.
\]

\[
\text{学习得到的生成联合分布：}\;
p(\mathbf{x},\mathbf{z})
\rightarrow
p_\theta(\mathbf{x}_{0:T})
=
p(\mathbf{x}_T)\prod_{t=1}^{T} p_\theta(\mathbf{x}_{t-1}\mid \mathbf{x}_t).
\]

\[
\text{辅助/变分分布：}\;
Q(\mathbf{z})
\rightarrow
q(\mathbf{x}_{1:T}\mid \mathbf{x}_0),
\]
它由前向扩散过程直接给出。

\[
\text{边缘似然：}\;
p_\theta(\mathbf{x}_0)
=
\int p_\theta(\mathbf{x}_{0:T})\,d\mathbf{x}_{1:T},
\]
由于需要边缘化整条潜在轨迹，这一量通常难以直接计算。
\end{tcolorbox}

\paragraph{扩散模型的 ELBO。}
完成上述对应之后，标准 ELBO 构造给出
\[
\log p_\theta(\mathbf{x}_0)
\ge
\mathbb{E}_{q(\mathbf{x}_{1:T}\mid \mathbf{x}_0)}
\!\left[
\log
\frac{
p_\theta(\mathbf{x}_{0:T})
}{
q(\mathbf{x}_{1:T}\mid \mathbf{x}_0)
}
\right].
\]
这就是专门适用于扩散模型的 ELBO。

为了在整个数据集上训练模型，我们还需要对经验数据分布 $q(\mathbf{x}_0)$ 取外层期望，也就是对干净训练样本本身取期望。
由此得到总体目标
\[
\mathbb{E}_{q(\mathbf{x}_0)}
\left[
\log p_\theta(\mathbf{x}_0)
\right]
\ge
\mathbb{E}_{q(\mathbf{x}_0)}
\mathbb{E}_{q(\mathbf{x}_{1:T}\mid \mathbf{x}_0)}
\!\left[
\log
\frac{
p_\theta(\mathbf{x}_{0:T})
}{
q(\mathbf{x}_{1:T}\mid \mathbf{x}_0)
}
\right].
\]

因此，从本书前面建立的潜变量视角看，扩散模型很自然地落入同一个一般概率框架：
它仍然是生成式潜变量模型，只不过潜变量由单个潜编码变成了带时间索引的潜变量序列。
我们并不直接优化难以处理的边缘似然，而是利用两个组成部分构造并优化一个可处理的下界：学习得到的反向生成模型
$p_\theta(\mathbf{x}_{0:T})$ 与已知前向过程
$q(\mathbf{x}_{1:T}\mid \mathbf{x}_0)$。

因此，下一步的关键是展开并化简这个 ELBO，从而看清学习问题如何分解到各个扩散时间步。
这正是下一节的目标：我们将更明确地展开反向过程，并推导逐步目标，最终得到实际使用的 DDPM 训练损失。

\section{从扩散 ELBO 到逐步目标}

现在，我们把上一节推导的扩散 ELBO 改写为能够直接揭示模型究竟在学习什么的形式。
关键思想是：尽管似然涉及整条潜在轨迹 $\mathbf{x}_{1:T}$，最终目标却可以分解为一系列局部去噪步骤。

\subsection{展开扩散 ELBO}

\paragraph{代入前向链与反向链。}

回顾已经建立的两个分解：
\[
p_\theta(\mathbf{x}_{0:T})
=
p(\mathbf{x}_T)\prod_{t=1}^{T} p_\theta(\mathbf{x}_{t-1}\mid \mathbf{x}_t),
\]
以及
\[
q(\mathbf{x}_{1:T}\mid \mathbf{x}_0)
=
\prod_{t=1}^{T} q(\mathbf{x}_t\mid \mathbf{x}_{t-1}).
\]
如前所述，由于前向过程是以干净数据样本 $\mathbf{x}_0$ 为条件定义的，也可以等价地显式保留这一条件，写成
\[
q(\mathbf{x}_{1:T}\mid \mathbf{x}_0)
=
\prod_{t=1}^{T} q(\mathbf{x}_t\mid \mathbf{x}_{t-1},\mathbf{x}_0).
\]
根据马尔可夫性质，在单步前向转移中对 $\mathbf{x}_0$ 的条件实际上是冗余的；但暂时把它保留下来，在稍后利用贝叶斯公式改写 ELBO 时会很有用。

把前向链与反向链代入 ELBO，得到
\[
\begin{aligned}
\mathcal{L}_{\mathrm{ELBO}}
&=
\mathbb{E}_{q(\mathbf{x}_0)}
\mathbb{E}_{q(\mathbf{x}_{1:T}\mid \mathbf{x}_0)}
\!\left[
\log
\frac{
p(\mathbf{x}_T)\prod_{t=1}^{T} p_\theta(\mathbf{x}_{t-1}\mid \mathbf{x}_t)
}{
\prod_{t=1}^{T} q(\mathbf{x}_t\mid \mathbf{x}_{t-1},\mathbf{x}_0)
}
\right].
\end{aligned}
\]
利用对数把乘积变为求和，可写为
\[
\begin{aligned}
\mathcal{L}_{\mathrm{ELBO}}
&=
\mathbb{E}_{q(\mathbf{x}_0)}
\mathbb{E}_{q(\mathbf{x}_{1:T}\mid \mathbf{x}_0)}
\!\left[
\log p(\mathbf{x}_T)
+
\sum_{t=1}^{T}
\log
\frac{
p_\theta(\mathbf{x}_{t-1}\mid \mathbf{x}_t)
}{
q(\mathbf{x}_t\mid \mathbf{x}_{t-1},\mathbf{x}_0)
}
\right].
\end{aligned}
\]

乍看之下，这似乎已经是一个逐步目标。
但它还不是最有用的形式。
这里的分母仍然是\emph{前向}加噪分布，而模型真正试图学习的是一个\emph{反向}去噪过程。

\paragraph{学习到的去噪过程应当逼近什么？}

学习得到的反向转移
\[
p_\theta(\mathbf{x}_{t-1}\mid \mathbf{x}_t)
\]
就是模型的去噪规则：给定一个噪声状态 $\mathbf{x}_t$，它应当输出稍微更干净的状态 $\mathbf{x}_{t-1}$ 的分布。
因此，自然的学习目标并不是前向扰动规则
$q(\mathbf{x}_t\mid \mathbf{x}_{t-1})$，而是由前向过程诱导出的相应\emph{真实逆后验}。

这里必须区分两个密切相关的对象：

\begin{itemize}
    \item \textbf{学习得到的反向过程}，或者说\textbf{学习得到的去噪过程}
    \[
    p_\theta(\mathbf{x}_{t-1}\mid \mathbf{x}_t),
    \]
    它是我们参数化并训练的模型；

    \item \textbf{由前向过程诱导的真实逆后验}
    \[
    q(\mathbf{x}_{t-1}\mid \mathbf{x}_t,\mathbf{x}_0),
    \]
    它是在给定干净样本 $\mathbf{x}_0$ 后，由已知加噪机制所蕴含的精确单步去噪分布。
\end{itemize}

为什么希望学习得到的去噪过程与这一后验匹配？
因为前向过程精确规定了真实数据是如何被扰动的。
所以在训练过程中，它也同时告诉我们，这种扰动在单步意义下“正确的逆过程”应该是什么。
如果学习得到的反向转移在每个时间步都能够模拟真实逆后验，那么整条学习得到的反向链就有能力撤销前向加噪过程并恢复数据分布。

一个细微但重要的点是：我们并不直接使用
\[
q(\mathbf{x}_{t-1}\mid \mathbf{x}_t).
\]
这个对象是由完整数据分布诱导出的无条件反向分布，并不能直接解析处理。
然而在训练期间，干净数据样本 $\mathbf{x}_0$ 是已知的。
因此，我们可以改用
\[
q(\mathbf{x}_{t-1}\mid \mathbf{x}_t,\mathbf{x}_0),
\]
它是针对某个具体观测样本、由已知前向过程诱导出的条件后验。
这正是我们显式保留 $\mathbf{x}_0$ 的原因：对观测数据进行条件化之后，相关的反向分布变得解析可处理。

因此，学习反向过程的任务可以更精确地表述为：
\begin{quote}
\emph{在每个扩散时间步上，训练建模的去噪转移
$p_\theta(\mathbf{x}_{t-1}\mid \mathbf{x}_t)$
去匹配真实逆后验
$q(\mathbf{x}_{t-1}\mid \mathbf{x}_t,\mathbf{x}_0)$。}
\end{quote}

\subsection{利用贝叶斯公式改写 ELBO}\begin{tcolorbox}[
  colback=blue!6,
  colframe=blue!55!black,
  title={显式得到真实逆后验的一个实用技巧},
  boxrule=0.8pt,
  arc=2mm,
  left=2mm,
  right=2mm,
  top=1mm,
  bottom=1mm
]
技巧是让所有项都显式以观测到的干净样本 $\mathbf{x}_0$ 为条件，然后应用贝叶斯公式。从前向转移
$q(\mathbf{x}_t\mid \mathbf{x}_{t-1},\mathbf{x}_0)$ 出发，在给定 $\mathbf{x}_0$ 的条件下，贝叶斯公式给出：
\[
\begin{aligned}
q(\mathbf{x}_t\mid \mathbf{x}_{t-1},\mathbf{x}_0)
&=
\frac{
q(\mathbf{x}_{t-1},\mathbf{x}_t\mid \mathbf{x}_0)
}{
q(\mathbf{x}_{t-1}\mid \mathbf{x}_0)
}=
\frac{
q(\mathbf{x}_{t-1}\mid \mathbf{x}_t,\mathbf{x}_0)\,
q(\mathbf{x}_t\mid \mathbf{x}_0)
}{
q(\mathbf{x}_{t-1}\mid \mathbf{x}_0)
}.
\end{aligned}
\]
重新整理后得到真实逆后验：
\[
q(\mathbf{x}_{t-1}\mid \mathbf{x}_t,\mathbf{x}_0)
=
\frac{
q(\mathbf{x}_t\mid \mathbf{x}_{t-1},\mathbf{x}_0)\,
q(\mathbf{x}_{t-1}\mid \mathbf{x}_0)
}{
q(\mathbf{x}_t\mid \mathbf{x}_0)
}.
\]
\end{tcolorbox}

为了把这个真实逆后验引入 ELBO，我们对前向转移应用贝叶斯公式：
\[
q(\mathbf{x}_t\mid \mathbf{x}_{t-1},\mathbf{x}_0)
=
\frac{
q(\mathbf{x}_{t-1}\mid \mathbf{x}_t,\mathbf{x}_0)\,
q(\mathbf{x}_t\mid \mathbf{x}_0)
}{
q(\mathbf{x}_{t-1}\mid \mathbf{x}_0)
}.
\]
把它代入 ELBO 的分母，得到
\[
\begin{aligned}
\mathcal{L}_{\mathrm{ELBO}}
&=
\mathbb{E}_{q(\mathbf{x}_0)}
\mathbb{E}_{q(\mathbf{x}_{1:T}\mid \mathbf{x}_0)}
\!\left[
\log p(\mathbf{x}_T)
+
\sum_{t=1}^{T}
\log
\frac{
p_\theta(\mathbf{x}_{t-1}\mid \mathbf{x}_t)\,
q(\mathbf{x}_{t-1}\mid \mathbf{x}_0)
}{
q(\mathbf{x}_{t-1}\mid \mathbf{x}_t,\mathbf{x}_0)\,
q(\mathbf{x}_t\mid \mathbf{x}_0)
}
\right].
\end{aligned}
\]
现在把乘积的对数拆成和：
\[
\begin{aligned}
\mathcal{L}_{\mathrm{ELBO}}
&=
\mathbb{E}_{q(\mathbf{x}_0)}
\mathbb{E}_{q(\mathbf{x}_{1:T}\mid \mathbf{x}_0)}
\!\left[
\log p(\mathbf{x}_T)
+
\sum_{t=1}^{T}
\log
\frac{
p_\theta(\mathbf{x}_{t-1}\mid \mathbf{x}_t)
}{
q(\mathbf{x}_{t-1}\mid \mathbf{x}_t,\mathbf{x}_0)
}
+
\sum_{t=1}^{T}
\log
\frac{
q(\mathbf{x}_{t-1}\mid \mathbf{x}_0)
}{
q(\mathbf{x}_t\mid \mathbf{x}_0)
}
\right].
\end{aligned}
\]

第二个求和是一个望远镜求和。
显式展开即可看出：
\[
\begin{aligned}
\sum_{t=1}^{T}
\log
\frac{
q(\mathbf{x}_{t-1}\mid \mathbf{x}_0)
}{
q(\mathbf{x}_t\mid \mathbf{x}_0)
}
&=
\log
\frac{q(\mathbf{x}_0\mid \mathbf{x}_0)}{q(\mathbf{x}_1\mid \mathbf{x}_0)}
+
\log
\frac{q(\mathbf{x}_1\mid \mathbf{x}_0)}{q(\mathbf{x}_2\mid \mathbf{x}_0)}
+\cdots+
\log
\frac{q(\mathbf{x}_{T-1}\mid \mathbf{x}_0)}{q(\mathbf{x}_T\mid \mathbf{x}_0)}
\\[6pt]
&=
\log\!\left[
\frac{q(\mathbf{x}_0\mid \mathbf{x}_0)}{q(\mathbf{x}_1\mid \mathbf{x}_0)}
\cdot
\frac{q(\mathbf{x}_1\mid \mathbf{x}_0)}{q(\mathbf{x}_2\mid \mathbf{x}_0)}
\cdots
\frac{q(\mathbf{x}_{T-1}\mid \mathbf{x}_0)}{q(\mathbf{x}_T\mid \mathbf{x}_0)}
\right]
\\[6pt]
&=
\log\!\left[
\frac{q(\mathbf{x}_0\mid \mathbf{x}_0)}{q(\mathbf{x}_T\mid \mathbf{x}_0)}
\right] =
\log q(\mathbf{x}_0\mid \mathbf{x}_0)
-\log q(\mathbf{x}_T\mid \mathbf{x}_0).
\end{aligned}
\]
因为所有中间项都会两两抵消。
严格来说，$q(\mathbf{x}_0\mid \mathbf{x}_0)$ 是集中在观测值 $\mathbf{x}_0$ 处的退化分布，对应于狄拉克 delta 而不是普通概率密度。
由于它不包含任何可学习参数，因此对目标函数只贡献一个加性常数，在优化时可以忽略。
于是 ELBO 变为
\[
\begin{aligned}
\mathcal{L}_{\mathrm{ELBO}}
&=
\mathbb{E}_{q(\mathbf{x}_0)}
\mathbb{E}_{q(\mathbf{x}_{1:T}\mid \mathbf{x}_0)}
\!\left[
\log p(\mathbf{x}_T)
-
\log q(\mathbf{x}_T\mid \mathbf{x}_0)
+
\sum_{t=1}^{T}
\log
\frac{
p_\theta(\mathbf{x}_{t-1}\mid \mathbf{x}_t)
}{
q(\mathbf{x}_{t-1}\mid \mathbf{x}_t,\mathbf{x}_0)
}
\right]
+\text{常数}.
\end{aligned}
\]

\subsection{从全局 ELBO 到局部学习项}

现在，我们把 ELBO 重新组织为标准的逐步分解。
主要需要处理两部分：
\begin{itemize}
    \item 与终点 $\mathbf{x}_T$ 有关的\textbf{边界项}；
    \item 与学习得到的去噪转移
    $p_\theta(\mathbf{x}_{t-1}\mid \mathbf{x}_t)$ 有关的\textbf{时间步项}。
\end{itemize}

从上面的表达式出发，先分离出只涉及终点 $\mathbf{x}_T$ 的部分：
\[
\mathbb{E}_{q(\mathbf{x}_{1:T}\mid \mathbf{x}_0)}
\!\big[
\log p(\mathbf{x}_T)-\log q(\mathbf{x}_T\mid \mathbf{x}_0)
\big].
\]
由于被积函数只依赖 $\mathbf{x}_T$，在前向链下可以把其他潜变量
$\mathbf{x}_1,\dots,\mathbf{x}_{T-1}$ 全部边缘化，只留下边缘分布
$q(\mathbf{x}_T\mid \mathbf{x}_0)$：
\[
\begin{aligned}
&\mathbb{E}_{q(\mathbf{x}_{1:T}\mid \mathbf{x}_0)}
\!\big[
\log p(\mathbf{x}_T)-\log q(\mathbf{x}_T\mid \mathbf{x}_0)
\big] \\[4pt]
&\qquad=
\mathbb{E}_{q(\mathbf{x}_T\mid \mathbf{x}_0)}
\!\big[
\log p(\mathbf{x}_T)-\log q(\mathbf{x}_T\mid \mathbf{x}_0)
\big] \\[4pt]
&\qquad=
-\mathbb{E}_{q(\mathbf{x}_T\mid \mathbf{x}_0)}
\!\left[
\log \frac{q(\mathbf{x}_T\mid \mathbf{x}_0)}{p(\mathbf{x}_T)}
\right] \\[4pt]
&\qquad=
-\mathrm{KL}\!\bigl(
q(\mathbf{x}_T\mid \mathbf{x}_0)\,\|\,p(\mathbf{x}_T)
\bigr).
\end{aligned}
\]
这就是\textbf{先验匹配项}：它比较前向过程到达最终时间步时的分布，与学习到的反向过程作为起点的简单高斯先验。

接下来处理关于 $t$ 的求和。
这里必须谨慎处理内部期望
$\mathbb{E}_{q(\mathbf{x}_{1:T}\mid \mathbf{x}_0)}$。
整个求和都位于期望内部，但每个求和项
\[
\log
\frac{
p_\theta(\mathbf{x}_{t-1}\mid \mathbf{x}_t)
}{
q(\mathbf{x}_{t-1}\mid \mathbf{x}_t,\mathbf{x}_0)
}
\]
只依赖局部变量 $\mathbf{x}_{t-1}$ 和 $\mathbf{x}_t$，以及观测样本 $\mathbf{x}_0$。
因此，对固定的 $t$，其余潜变量同样可以被边缘化掉。

对于特殊情形 $t=1$，有
\[
\begin{aligned}
&\mathbb{E}_{q(\mathbf{x}_{1:T}\mid \mathbf{x}_0)}
\!\left[
\log
\frac{
p_\theta(\mathbf{x}_0\mid \mathbf{x}_1)
}{
q(\mathbf{x}_0\mid \mathbf{x}_1,\mathbf{x}_0)
}
\right] =
\mathbb{E}_{q(\mathbf{x}_1\mid \mathbf{x}_0)}
\!\left[
\log
\frac{
p_\theta(\mathbf{x}_0\mid \mathbf{x}_1)
}{
q(\mathbf{x}_0\mid \mathbf{x}_1,\mathbf{x}_0)
}
\right].
\end{aligned}
\]
一旦以 $\mathbf{x}_0$ 为条件，$q(\mathbf{x}_0\mid \mathbf{x}_1,\mathbf{x}_0)$ 就是退化分布，对应于观测值处的点质量，并且不包含可学习参数。
因此它只贡献一个加性常数，留下\textbf{重构项}
\[
\mathbb{E}_{q(\mathbf{x}_1\mid \mathbf{x}_0)}
\big[
\log p_\theta(\mathbf{x}_0\mid \mathbf{x}_1)
\big].
\]

对于剩余时间步 $t=2,\dots,T$，我们以相同方式将内部期望局部化：
\[
\begin{aligned}
&\mathbb{E}_{q(\mathbf{x}_{1:T}\mid \mathbf{x}_0)}
\!\left[
\log
\frac{
p_\theta(\mathbf{x}_{t-1}\mid \mathbf{x}_t)
}{
q(\mathbf{x}_{t-1}\mid \mathbf{x}_t,\mathbf{x}_0)
}
\right] \\[6pt]
&\qquad=
\mathbb{E}_{q(\mathbf{x}_{t-1},\mathbf{x}_t\mid \mathbf{x}_0)}
\!\left[
\log
\frac{
p_\theta(\mathbf{x}_{t-1}\mid \mathbf{x}_t)
}{
q(\mathbf{x}_{t-1}\mid \mathbf{x}_t,\mathbf{x}_0)
}
\right] \\[6pt]
&\qquad=
\mathbb{E}_{q(\mathbf{x}_t\mid \mathbf{x}_0)}
\mathbb{E}_{q(\mathbf{x}_{t-1}\mid \mathbf{x}_t,\mathbf{x}_0)}
\!\left[
\log
\frac{
p_\theta(\mathbf{x}_{t-1}\mid \mathbf{x}_t)
}{
q(\mathbf{x}_{t-1}\mid \mathbf{x}_t,\mathbf{x}_0)
}
\right].
\end{aligned}
\]
现在可以识别出，内层期望就是某个 KL 散度的负值：
\[
\begin{aligned}
&\mathbb{E}_{q(\mathbf{x}_{t-1}\mid \mathbf{x}_t,\mathbf{x}_0)}
\!\left[
\log
\frac{
p_\theta(\mathbf{x}_{t-1}\mid \mathbf{x}_t)
}{
q(\mathbf{x}_{t-1}\mid \mathbf{x}_t,\mathbf{x}_0)
}
\right] \\[4pt]
&\qquad=
-\mathbb{E}_{q(\mathbf{x}_{t-1}\mid \mathbf{x}_t,\mathbf{x}_0)}
\!\left[
\log
\frac{
q(\mathbf{x}_{t-1}\mid \mathbf{x}_t,\mathbf{x}_0)
}{
p_\theta(\mathbf{x}_{t-1}\mid \mathbf{x}_t)
}
\right] \\[4pt]
&\qquad=
-\mathrm{KL}\!\big(
q(\mathbf{x}_{t-1}\mid \mathbf{x}_t,\mathbf{x}_0)
\,\|\,\,
p_\theta(\mathbf{x}_{t-1}\mid \mathbf{x}_t)
\big).
\end{aligned}
\]
因此，每个中间时间步的贡献都变为
\[
-
\mathbb{E}_{q(\mathbf{x}_t\mid \mathbf{x}_0)}
\Big[
\mathrm{KL}\!\big(
q(\mathbf{x}_{t-1}\mid \mathbf{x}_t,\mathbf{x}_0)
\,\|\,\,
p_\theta(\mathbf{x}_{t-1}\mid \mathbf{x}_t)
\big)
\Big].
\]

把边界项和各时间步项合并起来，负 ELBO 可以写成
\[
\begin{aligned}
-\mathcal{L}_{\mathrm{ELBO}}
=
\mathbb{E}_{q(\mathbf{x}_0)}
\Big[
&\underbrace{
\mathrm{KL}\!\bigl(
q(\mathbf{x}_T\mid \mathbf{x}_0)\,\|\,p(\mathbf{x}_T)
\bigr)
}_{\text{先验匹配项}}
\\[4pt]
&+
\sum_{t=2}^{T}
\underbrace{
\mathbb{E}_{q(\mathbf{x}_t\mid \mathbf{x}_0)}
\Big[
\mathrm{KL}\!\big(
q(\mathbf{x}_{t-1}\mid \mathbf{x}_t,\mathbf{x}_0)
\,\|\,\,
p_\theta(\mathbf{x}_{t-1}\mid \mathbf{x}_t)
\big)
\Big]
}_{\text{逐步去噪项}}
\\[4pt]
&-
\underbrace{
\mathbb{E}_{q(\mathbf{x}_1\mid \mathbf{x}_0)}
\big[
\log p_\theta(\mathbf{x}_0\mid \mathbf{x}_1)
\big]
}_{\text{重构项}}
\Big].
\end{aligned}
\]

这一分解把三种作用清晰区分开来：噪声终点处的先验匹配、中间时间步的局部反向转移匹配，以及数据端的重构。
图~\ref{fig:ddpm-elbo-local-decomposition} 展示了这些项如何分布在前向链和反向链上。

\begin{figure}[H]
\centering
\begin{tikzpicture}[
  >=Latex,
  every node/.style={font=\small},
  state/.style={draw,circle,minimum size=0.76cm,inner sep=1pt},
  term/.style={draw,rounded corners=2pt,align=center,inner sep=5pt},
  flow/.style={->,thick},
  reverse/.style={->,thick,dashed}
]
\node[font=\normalsize\bfseries] at (6.1,3.0)
  {扩散 ELBO 沿链分解};

\node[state] (x0) at (0,1.35) {$x_0$};
\node[state] (x1) at (2.15,1.35) {$x_1$};
\node at (3.75,1.35) {$\cdots$};
\node[state] (xtm1) at (5.7,1.35) {$x_{t-1}$};
\node[state] (xt) at (8.15,1.35) {$x_t$};
\node at (9.8,1.35) {$\cdots$};
\node[state] (xT) at (12.0,1.35) {$x_T$};

% Forward chain
\draw[flow] (x0) -- node[above] {\scriptsize 固定的 $q$} (x1);
\draw[flow] (x1) -- (3.2,1.35);
\draw[flow] (4.3,1.35) -- (xtm1);
\draw[flow] (xtm1) -- node[above] {\scriptsize $q(x_t\mid x_{t-1})$} (xt);
\draw[flow] (xt) -- (9.25,1.35);
\draw[flow] (10.35,1.35) -- (xT);

% Learned reverse chain
\draw[reverse] (x1) -- node[below] {\scriptsize $p_\theta$} (x0);
\draw[reverse] (xt) -- node[below] {\scriptsize $p_\theta(x_{t-1}\mid x_t)$} (xtm1);
\draw[reverse] (xT) -- node[below] {\scriptsize 学习到的反向过程} (10.35,1.35);

\node[term,text width=3.55cm] (recon) at (1.1,-1.15) {
\textbf{重构}\\[2pt]
$\displaystyle
-\mathbb E_{q(x_1\mid x_0)}
[\log p_\theta(x_0\mid x_1)]$
};

\node[term,text width=5.15cm] (local) at (7.0,-1.35) {
\textbf{中间步骤，$t=2,\ldots,T$}\\[2pt]
$\displaystyle
\begin{aligned}
&\mathbb E_{q(x_t\mid x_0)}\Big[
\mathrm{KL}\!\big(\\[-1pt]
&q(x_{t-1}\mid x_t,x_0)
\,\|\,p_\theta(x_{t-1}\mid x_t)
\big)\Big]
\end{aligned}$
};

\node[term,text width=3.85cm] (prior) at (10.55,-3.65) {
\textbf{端点／先验匹配}\\[2pt]
$\displaystyle
\mathrm{KL}\!\big(
q(x_T\mid x_0)\,\|\,p(x_T)
\big)$
};

\draw[flow] (x0.south) to[out=-90,in=90] (recon.north west);
\draw[flow] (xtm1.south) to[out=-90,in=90] (local.north west);
\draw[flow] (xt.south) to[out=-90,in=90] (local.north east);
\draw[flow] (xT.south) to[out=-90,in=90] (prior.north);

\node[font=\scriptsize,align=center,text width=11.5cm] at (5.8,-5.0) {
全局潜轨迹目标可化为一个边界项、一个重构项，以及若干局部反向转移匹配项之和。
};

\end{tikzpicture}
\caption{\textbf{从扩散 ELBO 到局部去噪项。}
固定的前向链 $q$ 产生潜在轨迹，而 $p_\theta$ 建模反向转移。重新整理负 ELBO 后，可得到端点先验匹配、数据重构，以及一系列中间 KL 项；这些项把每个学习到的反向转移与精确后验 $q(\mathbf{x}_{t-1}\mid\mathbf{x}_t,\mathbf{x}_0)$ 进行比较。}
\label{fig:ddpm-elbo-local-decomposition}
\end{figure}

\noindent
关键结论是：学习可以归结为在每个时间步上，让学习得到的去噪转移
$p_\theta(\mathbf{x}_{t-1}\mid \mathbf{x}_t)$
匹配真实逆后验
$q(\mathbf{x}_{t-1}\mid \mathbf{x}_t,\mathbf{x}_0)$。
因此，下一步就是求出这个后验的解析形式。

\section{从真实逆后验到 DDPM 学习}

前面几节已经说明，扩散 ELBO 中核心的可训练项比较的是学习得到的去噪转移
\[
p_\theta(\mathbf{x}_{t-1}\mid \mathbf{x}_t)
\]
与真实逆后验
\[
q(\mathbf{x}_{t-1}\mid \mathbf{x}_t,\mathbf{x}_0).
\]
本节的目标是把这两个视角连接起来：首先确定前向过程所蕴含的精确高斯逆后验，然后说明该后验如何导出实际使用的 DDPM 学习目标与采样过程。

\subsection{由已知前向量得到真实逆后验}

上一节给出了恒等式
\[
q(\mathbf{x}_{t-1}\mid \mathbf{x}_t,\mathbf{x}_0)
=
\frac{
q(\mathbf{x}_t\mid \mathbf{x}_{t-1},\mathbf{x}_0)\,
q(\mathbf{x}_{t-1}\mid \mathbf{x}_0)
}{
q(\mathbf{x}_t\mid \mathbf{x}_0)
}.
\]
这个形式很有用，因为右侧每一项都由已经完全已知的前向过程决定。特别地，
\[
q(\mathbf{x}_t\mid \mathbf{x}_{t-1},\mathbf{x}_0)
=
q(\mathbf{x}_t\mid \mathbf{x}_{t-1})
=
\mathcal{N}\!\bigl(
\mathbf{x}_t \mid
\sqrt{\alpha_t}\,\mathbf{x}_{t-1},
(1-\alpha_t)\mathbf{I}
\bigr),
\]
其中 $\alpha_t=1-\beta_t$。由前面推导的前向闭式关系，我们还知道
\[
q(\mathbf{x}_{t-1}\mid \mathbf{x}_0)
=
\mathcal{N}\!\bigl(
\mathbf{x}_{t-1}\mid
\sqrt{\bar{\alpha}_{t-1}}\,\mathbf{x}_0,
(1-\bar{\alpha}_{t-1})\mathbf{I}
\bigr),
\]
以及
\[
q(\mathbf{x}_t\mid \mathbf{x}_0)
=
\mathcal{N}\!\bigl(
\mathbf{x}_t\mid
\sqrt{\bar{\alpha}_t}\,\mathbf{x}_0,
(1-\bar{\alpha}_t)\mathbf{I}
\bigr).
\]

因此，真实逆后验完全由前向扩散过程中的高斯量表示。
由于上面的各项都是高斯分布，关于 $\mathbf{x}_{t-1}$ 的后验也必然是高斯分布。
所以可以写成
\[
q(\mathbf{x}_{t-1}\mid \mathbf{x}_t,\mathbf{x}_0)
=
\mathcal{N}\!\bigl(
\mathbf{x}_{t-1}\mid
\tilde{\boldsymbol{\mu}}_t(\mathbf{x}_t,\mathbf{x}_0),
\tilde{\beta}_t\mathbf{I}
\bigr),
\]
其中均值为 $\tilde{\boldsymbol{\mu}}_t(\mathbf{x}_t,\mathbf{x}_0)$，方差系数为 $\tilde{\beta}_t$。

利用前面介绍的配方法，并结合附录~\ref{app:diffusion-posterior} 中的完整推导，真实逆后验具有如下高斯形式，其均值为
\[
\tilde{\boldsymbol{\mu}}_t(\mathbf{x}_t,\mathbf{x}_0)
=
\frac{
\sqrt{\bar{\alpha}_{t-1}}\,\beta_t
}{
1-\bar{\alpha}_t
}\,\mathbf{x}_0
+
\frac{
\sqrt{\alpha_t}(1-\bar{\alpha}_{t-1})
}{
1-\bar{\alpha}_t
}\,\mathbf{x}_t,
\]
方差系数为
\[
\tilde{\beta}_t
=
\frac{
1-\bar{\alpha}_{t-1}
}{
1-\bar{\alpha}_t
}\,\beta_t.
\]

这个结果值得解释。
后验均值
$\tilde{\boldsymbol{\mu}}_t(\mathbf{x}_t,\mathbf{x}_0)$
是两部分信息的加权组合：
干净样本 $\mathbf{x}_0$ 告诉我们轨迹从何处出发，而当前噪声状态 $\mathbf{x}_t$ 告诉我们在时刻 $t$ 位于何处。
因此，真实逆后验通过结合干净起点与当前噪声观测的信息，确定最可能的、稍微更干净的上一状态。

后验方差 $\tilde{\beta}_t\mathbf{I}$ 同样由前向扩散调度固定。
也就是说，一旦选择了前向加噪过程，对应反向步骤的不确定性也随之确定。
所以从 ELBO 视角看，反向模型并不是要匹配任意去噪分布，而是要匹配由已知前向过程诱导出的特定高斯目标。

图~\ref{fig:ddpm-true-reverse-posterior} 将这一目标可视化为两个已知前向过程因子的高斯融合。

\begin{figure}[H]
\centering
\includegraphics[width=0.8\textwidth]{figures/ch04/figure4_3_true_reverse_posterior_gaussian_fusion.pdf}
\caption{\textbf{把精确反向后验理解为高斯融合。}
在固定 $\mathbf{x}_0$ 与 $\mathbf{x}_t$ 时，前向边缘分布 $q(\mathbf{x}_{t-1}\mid\mathbf{x}_0)$ 与转移似然 $q(\mathbf{x}_t\mid\mathbf{x}_{t-1})$（把后者视为 $\mathbf{x}_{t-1}$ 的函数）给出两个高斯因子。二者乘积得到 $q(\mathbf{x}_{t-1}\mid\mathbf{x}_t,\mathbf{x}_0)$。一维图展示了各向同性多变量情形中同样的逐分量计算。}
\label{fig:ddpm-true-reverse-posterior}
\end{figure}

尽管这个后验是精确的，但它仍然显式依赖干净数据样本 $\mathbf{x}_0$。
训练时 $\mathbf{x}_0$ 是已知的，因此没有问题；但在生成时，$\mathbf{x}_0$ 恰恰是我们希望创建的未知对象，所以不能直接使用这一后验。

\subsection{通过噪声变量重参数化后验}

关键观察是，前向过程已经给出了干净样本 $\mathbf{x}_0$、噪声样本 $\mathbf{x}_t$ 与标准高斯噪声变量之间的直接关系：
\[
\mathbf{x}_t
=
\sqrt{\bar{\alpha}_t}\,\mathbf{x}_0
+
\sqrt{1-\bar{\alpha}_t}\,\boldsymbol{\epsilon},
\qquad
\boldsymbol{\epsilon}\sim \mathcal{N}(\mathbf{0},\mathbf{I}).
\]

解出 $\mathbf{x}_0$，得到
\[
\mathbf{x}_0
=
\frac{1}{\sqrt{\bar{\alpha}_t}}
\Bigl(
\mathbf{x}_t-\sqrt{1-\bar{\alpha}_t}\,\boldsymbol{\epsilon}
\Bigr).
\]
现在把这个表达式代入后验均值
$\tilde{\boldsymbol{\mu}}_t(\mathbf{x}_t,\mathbf{x}_0)$，以消除其中对 $\mathbf{x}_0$ 的显式依赖。

从
\[
\tilde{\boldsymbol{\mu}}_t(\mathbf{x}_t,\mathbf{x}_0)
=
\frac{
\sqrt{\bar{\alpha}_{t-1}}\,\beta_t
}{
1-\bar{\alpha}_t
}\,\mathbf{x}_0
+
\frac{
\sqrt{\alpha_t}(1-\bar{\alpha}_{t-1})
}{
1-\bar{\alpha}_t
}\,\mathbf{x}_t
\]
出发，将 $\mathbf{x}_0$ 的表达式代入可得
\[
\begin{aligned}
\tilde{\boldsymbol{\mu}}_t(\mathbf{x}_t,\boldsymbol{\epsilon})
&=
\frac{
\sqrt{\bar{\alpha}_{t-1}}\,\beta_t
}{
1-\bar{\alpha}_t
}
\cdot
\frac{1}{\sqrt{\bar{\alpha}_t}}
\Bigl(
\mathbf{x}_t-\sqrt{1-\bar{\alpha}_t}\,\boldsymbol{\epsilon}
\Bigr)
+
\frac{
\sqrt{\alpha_t}(1-\bar{\alpha}_{t-1})
}{
1-\bar{\alpha}_t
}\,\mathbf{x}_t.
\end{aligned}
\]

化简系数后，得到
\[
\tilde{\boldsymbol{\mu}}_t(\mathbf{x}_t,\boldsymbol{\epsilon})
=
\frac{1}{\sqrt{\alpha_t}}
\left(
\mathbf{x}_t
-
\frac{\beta_t}{\sqrt{1-\bar{\alpha}_t}}\,
\boldsymbol{\epsilon}
\right).
\]

\begin{tcolorbox}[
  colback=blue!6,
  colframe=blue!55!black,
  title={系数化简背后的代数过程},
  boxrule=0.8pt,
  arc=2mm,
  left=2mm,
  right=2mm,
  top=1mm,
  bottom=1mm
]
利用
\[
\bar{\alpha}_t=\alpha_t\bar{\alpha}_{t-1},
\qquad
\frac{\sqrt{\bar{\alpha}_{t-1}}}{\sqrt{\bar{\alpha}_t}}
=
\frac{1}{\sqrt{\alpha_t}},
\]
$\mathbf{x}_t$ 的系数变为
\[
\frac{\beta_t}{\sqrt{\alpha_t}(1-\bar{\alpha}_t)}
+
\frac{\sqrt{\alpha_t}(1-\bar{\alpha}_{t-1})}{1-\bar{\alpha}_t}
=
\frac{
\beta_t+\alpha_t(1-\bar{\alpha}_{t-1})
}{
\sqrt{\alpha_t}(1-\bar{\alpha}_t)
}.
\]
由于
\[
\beta_t=1-\alpha_t,
\qquad
\bar{\alpha}_t=\alpha_t\bar{\alpha}_{t-1},
\]
可得
\[
\beta_t+\alpha_t(1-\bar{\alpha}_{t-1})
=
(1-\alpha_t)+\alpha_t-\alpha_t\bar{\alpha}_{t-1}
=
1-\bar{\alpha}_t,
\]
因此化简为
\[
\frac{1-\bar{\alpha}_t}{\sqrt{\alpha_t}(1-\bar{\alpha}_t)}
=
\frac{1}{\sqrt{\alpha_t}}.
\]
类似地，$\boldsymbol{\epsilon}$ 的系数变为
\[
\frac{
\sqrt{\bar{\alpha}_{t-1}}\,\beta_t
}{
\sqrt{\bar{\alpha}_t}\sqrt{1-\bar{\alpha}_t}
}
=
\frac{\beta_t}{\sqrt{\alpha_t}\sqrt{1-\bar{\alpha}_t}}.
\]
\end{tcolorbox}

因此，真实逆后验可以改写为
\[
\boxed{
\begin{aligned}
q(\mathbf{x}_{t-1}\mid \mathbf{x}_t,\boldsymbol{\epsilon})
&=
\mathcal{N}\!\left(
\mathbf{x}_{t-1}\mid
\frac{1}{\sqrt{\alpha_t}}
\left(
\mathbf{x}_t
-
\frac{\beta_t}{\sqrt{1-\bar{\alpha}_t}}\,
\boldsymbol{\epsilon}
\right),
\tilde{\beta}_t\mathbf{I}
\right).
\end{aligned}
}
\]

这一形式极其重要。
它说明，反向步骤可以仅通过当前噪声样本 $\mathbf{x}_t$ 与生成该样本的累积噪声 $\boldsymbol{\epsilon}$ 来刻画。
因此，模型不必直接恢复 $\mathbf{x}_0$，而可以转而预测底层噪声。

\subsection{参数化学习得到的反向过程，并推导实际 DDPM 目标}

这就导出了标准 DDPM 参数化。
真实噪声 $\boldsymbol{\epsilon}$ 对模型而言是未知的，因此我们引入神经网络
\[
\boldsymbol{\epsilon}_\theta(\mathbf{x}_t,t)
\]
根据噪声样本 $\mathbf{x}_t$ 和时间步 $t$ 来预测它。

学习得到的反向转移定义为
\[
p_\theta(\mathbf{x}_{t-1}\mid \mathbf{x}_t)
=
\mathcal{N}\!\bigl(
\mathbf{x}_{t-1}\mid
\boldsymbol{\mu}_\theta(\mathbf{x}_t,t),
\tilde{\beta}_t\mathbf{I}
\bigr),
\]
其中均值为
\[
\boldsymbol{\mu}_\theta(\mathbf{x}_t,t)
=
\frac{1}{\sqrt{\alpha_t}}
\left(
\mathbf{x}_t
-
\frac{\beta_t}{\sqrt{1-\bar{\alpha}_t}}\,
\boldsymbol{\epsilon}_\theta(\mathbf{x}_t,t)
\right).
\]

其解释非常直接：
如果网络能够准确预测真实噪声，那么学习得到的反向均值就会与精确逆后验均值一致，因此学习到的去噪过程也就匹配了正确的局部反向步骤。

\bigskip

回顾 ELBO 中核心的可训练项
\[
\mathbb{E}_{q(\mathbf{x}_t\mid \mathbf{x}_0)}
\Big[
\mathrm{KL}\!\big(
q(\mathbf{x}_{t-1}\mid \mathbf{x}_t,\mathbf{x}_0)
\,\|\,\,
p_\theta(\mathbf{x}_{t-1}\mid \mathbf{x}_t)
\big)
\Big].
\]
现在两个分布都是具有相同方差
$\tilde{\beta}_t\mathbf{I}$ 的高斯分布。
因此，它们之间的 KL 散度会化简为均值之差的平方：
\[
\mathrm{KL}\!\big(
q(\mathbf{x}_{t-1}\mid \mathbf{x}_t,\mathbf{x}_0)
\,\|\,\,
p_\theta(\mathbf{x}_{t-1}\mid \mathbf{x}_t)
\big)
\propto
\bigl\|
\tilde{\boldsymbol{\mu}}_t(\mathbf{x}_t,\boldsymbol{\epsilon})
-
\boldsymbol{\mu}_\theta(\mathbf{x}_t,t)
\bigr\|^2.
\]
代入两个均值表达式，得到
\[
\tilde{\boldsymbol{\mu}}_t(\mathbf{x}_t,\boldsymbol{\epsilon})
-
\boldsymbol{\mu}_\theta(\mathbf{x}_t,t)
=
\frac{\beta_t}{\sqrt{\alpha_t}\sqrt{1-\bar{\alpha}_t}}
\Bigl(
\boldsymbol{\epsilon}
-
\boldsymbol{\epsilon}_\theta(\mathbf{x}_t,t)
\Bigr).
\]
因此，忽略一个依赖时间步的缩放因子后，KL 项就变成噪声预测的均方误差：
\[
\mathrm{KL}\!\big(
q(\mathbf{x}_{t-1}\mid \mathbf{x}_t,\mathbf{x}_0)
\,\|\,\,
p_\theta(\mathbf{x}_{t-1}\mid \mathbf{x}_t)
\big)
\propto
\bigl\|
\boldsymbol{\epsilon}
-
\boldsymbol{\epsilon}_\theta(\mathbf{x}_t,t)
\bigr\|^2.
\]

这就是 DDPM 中最核心的化简：局部变分项在相差一个时间相关系数的意义下，归结为对所加入高斯噪声的回归问题。
因此，精确 ELBO 会导出带时间权重的噪声预测均方误差。
实践中，Ho 等人~\cite{ho2020denoising} 常用的简化目标去掉了这个依赖时间步的权重，采用
\[
\boxed{
\mathcal{L}_{\mathrm{simple}}
=
\mathbb{E}_{t,\mathbf{x}_0,\boldsymbol{\epsilon}}
\Bigl[
\bigl\|
\boldsymbol{\epsilon}
-
\boldsymbol{\epsilon}_\theta(\mathbf{x}_t,t)
\bigr\|^2
\Bigr]
}
\]
其中
$\mathbf{x}_t
=
\sqrt{\bar{\alpha}_t}\,\mathbf{x}_0
+
\sqrt{1-\bar{\alpha}_t}\,\boldsymbol{\epsilon}$。
图~\ref{fig:ddpm-noise-prediction-bridge} 总结了从局部后验匹配到这一实际噪声预测损失的代数桥梁。

\begin{figure}[H]
\centering
\begin{tikzpicture}[
  >=Latex,
  every node/.style={font=\small},
  box/.style={
    draw,
    rounded corners=2pt,
    align=center,
    inner sep=6pt,
    text width=5.45cm,
    minimum height=1.38cm
  },
  narrow/.style={
    draw,
    rounded corners=2pt,
    align=center,
    inner sep=6pt,
    text width=4.15cm,
    minimum height=1.15cm
  },
  flow/.style={->,thick}
]

\node[box] (true) at (2.9,1.45) {
\textbf{精确反向后验均值}\\[3pt]
$\displaystyle
\tilde{\boldsymbol{\mu}}_t
=
\frac{1}{\sqrt{\alpha_t}}
\left(
\mathbf{x}_t-
\frac{\beta_t}{\sqrt{1-\bar{\alpha}_t}}
\,\boldsymbol{\epsilon}
\right)$
};

\node[box] (learned) at (9.1,1.45) {
\textbf{学习到的反向均值}\\[3pt]
$\displaystyle
\boldsymbol{\mu}_{\theta}
=
\frac{1}{\sqrt{\alpha_t}}
\left(
\mathbf{x}_t-
\frac{\beta_t}{\sqrt{1-\bar{\alpha}_t}}
\,\boldsymbol{\epsilon}_{\theta}(\mathbf{x}_t,t)
\right)$
};

\draw[flow] (true) -- (learned);

\node[narrow] (samevar) at (2.9,-0.95) {
\textbf{共享的固定方差}\\[2pt]
$q:\ \tilde{\beta}_t\mathbf I
\qquad
p_\theta:\ \tilde{\beta}_t\mathbf I$
};

\node[narrow] (meanmse) at (9.1,-0.95) {
\textbf{高斯 KL 化为\\均值匹配}\\[2pt]
$\displaystyle
\mathrm{KL}(q\,\|\,p_\theta)
\propto
\|\tilde{\boldsymbol{\mu}}_t-\boldsymbol{\mu}_\theta\|^2$
};

\draw[flow] (true.south) -- (samevar.north);
\draw[flow] (learned.south) -- (meanmse.north);
\draw[flow] (samevar) -- (meanmse);

\node[box] (noise) at (6.0,-3.45) {
\textbf{共享系数只留下时间依赖的缩放}\\[3pt]
$\displaystyle
\|\tilde{\boldsymbol{\mu}}_t-\boldsymbol{\mu}_\theta\|^2
\propto
\|\boldsymbol{\epsilon}
-\boldsymbol{\epsilon}_{\theta}(\mathbf{x}_t,t)\|^2$
};

\draw[flow] (meanmse) -- (noise);

\node[narrow] (simple) at (6.0,-5.95) {
\textbf{实用 DDPM 目标}\\[3pt]
$\displaystyle
\mathcal L_{\mathrm{simple}}
=
\mathbb E
\left[
\|\boldsymbol{\epsilon}
-\boldsymbol{\epsilon}_{\theta}(\mathbf{x}_t,t)\|^2
\right]$
};

\draw[flow] (noise) -- (simple);

\end{tikzpicture}

\caption{\textbf{从精确后验匹配到噪声预测。}
把精确反向均值与学习到的反向均值写成同一形式，并使用相同的固定方差，局部高斯 KL 就转化为均值平方误差匹配。由于两者共享系数，该目标进一步化为真实噪声 $\boldsymbol{\epsilon}$ 与其预测 $\boldsymbol{\epsilon}_\theta(\mathbf{x}_t,t)$ 之间带时间权重的平方误差。常用的 $\mathcal{L}_{\mathrm{simple}}$ 目标会去掉这一时间依赖权重。}
\label{fig:ddpm-noise-prediction-bridge}
\end{figure}

\subsection{训练与采样如何工作}

现在，训练循环在概念上已经非常简单。
对于每个干净数据样本 $\mathbf{x}_0$，我们：
\begin{enumerate}
    \item 采样一个时间步 $t$；
    \item 采样高斯噪声 $\boldsymbol{\epsilon}\sim\mathcal{N}(\mathbf{0},\mathbf{I})$；
    \item 构造噪声状态
    \[
    \mathbf{x}_t
    =
    \sqrt{\bar{\alpha}_t}\,\mathbf{x}_0
    +
    \sqrt{1-\bar{\alpha}_t}\,\boldsymbol{\epsilon};
    \]
    \item 训练网络根据 $(\mathbf{x}_t,t)$ 预测 $\boldsymbol{\epsilon}$。
\end{enumerate}
由于可以直接从 $\mathbf{x}_0$ 采样得到 $\mathbf{x}_t$，训练时不需要逐步模拟完整前向链。

训练完成后，模型通过迭代运行学习到的反向过程来生成数据。
我们从
\[
\mathbf{x}_T\sim\mathcal{N}(\mathbf{0},\mathbf{I})
\]
开始，然后对 $t=T,T-1,\dots,1$ 依次采样
\[
\mathbf{x}_{t-1}\sim p_\theta(\mathbf{x}_{t-1}\mid \mathbf{x}_t).
\]
在每一步，模型利用预测噪声
$\boldsymbol{\epsilon}_\theta(\mathbf{x}_t,t)$ 计算反向均值，并去除少量扰动。
经过所有时间步反复执行后，最终状态 $\mathbf{x}_0$ 就是生成样本。

\bigskip
\noindent
至此，我们完成了离散时间 DDPM 的主要概率建模逻辑。
前向过程固定，反向过程通过学习获得，ELBO 化简为局部去噪目标，而最终实际使用的目标进一步化为噪声预测。

\section*{本章小结与参考文献}
\addcontentsline{toc}{section}{本章小结与参考文献}

本章从概率潜变量视角介绍了去噪扩散概率模型（DDPM）。
与 PPCA 和 VAE 中由单一潜变量发挥核心作用不同，扩散模型处理的是完整的带时间索引潜在轨迹
$\mathbf{x}_{1:T}$；该轨迹由固定的高斯马尔可夫过程逐步扰动数据而产生。

我们首先说明了前向过程如何把干净样本 $\mathbf{x}_0$ 变换为噪声逐渐增强的状态，并推导了 $\mathbf{x}_t$ 与 $\mathbf{x}_0$ 之间的闭式关系。
随后转向必须学习的反向过程，并利用 ELBO 框架把扩散模型表达为一个具有难以直接计算边缘似然的潜变量生成模型。

接着，我们把 ELBO 分解为局部去噪项，从而揭示每个时间步上真正需要训练的量，是学习得到的反向转移
$p_\theta(\mathbf{x}_{t-1}\mid \mathbf{x}_t)$ 与真实逆后验
$q(\mathbf{x}_{t-1}\mid \mathbf{x}_t,\mathbf{x}_0)$ 之间的差异。
由于前向过程是高斯过程，这个后验同样是高斯分布，并可以闭式推导。
然而，由于它显式依赖干净样本 $\mathbf{x}_0$，在生成时不能直接使用。

最后，通过利用噪声变量对后验进行重参数化，我们说明了扩散学习如何归结为噪声预测。
这就得到实际 DDPM 目标：训练神经网络预测在给定时间步生成噪声样本时所加入的高斯噪声，然后利用这个学习到的预测器，把纯高斯噪声迭代去噪为数据样本。
以这种方式，DDPM 将前面章节逐步建立的概率潜变量思想扩展成一种基于渐进扰动与渐进重构的序贯生成框架。

\medskip
\noindent

基于扩散的生成建模框架由 Sohl-Dickstein 等人~\cite{sohl2015deep} 提出，他们把生成建模表述为对渐进加噪过程的学习式逆转，并从非平衡热力学中获得启发。
这项工作确立了后来 DDPM 仍然沿用的多个核心组成：固定的前向扰动过程、学习得到的反向过程，以及连接数据与噪声的一系列中间潜状态。

现代 DDPM 形式由 Ho、Jain 和 Abbeel~\cite{ho2020denoising} 建立，他们说明了这一序贯潜变量模型可以通过变分目标进行有效训练，而该目标进一步简化为噪声预测。
Nichol 与 Dhariwal~\cite{nichol2021improved} 等后续工作又改进了 DDPM 的实际训练与采样表现。
本书后续介绍的得分模型与连续时间视角会为同一去噪原理提供互补解释，但本章重点是 DDPM 的离散潜变量构造。

\chapter{连续时间生成建模的微积分基础}

\section*{概述}
\addcontentsline{toc}{section}{概述}

上一章研究了离散时间下的去噪扩散概率模型（DDPM），其中生成过程被描述为一系列有限的加噪与去噪步骤。这一视角十分强大，但也自然引出了下一个问题：当步长越来越小时，会发生什么？此时，演化过程更适合被视为时间上的连续运动。

这个问题并非扩散模型所独有。更一般地，机器学习与生成建模中的许多方法都可以理解为反复执行如下局部更新：
\[
x_{t+\Delta t}\approx x_t+\text{小幅更新}.
\]
一个更具体的例子是
\[
x_{t+\Delta t}=x_t+\alpha f(x_t)\Delta t+\beta\,\text{噪声},
\]
其中下一状态由当前状态加上一个有结构的变化得到，而在某些情形下还会再加入随机扰动。
在离散时间中，这类规则被逐步写出；而在连续时间中，同样的思想要通过无穷小运动来表达：状态在任意短的时间区间内如何变化，以及随机性如何进入这种运动。
这就是为什么微积分语言在这里变得不可或缺。

因此，本章从\emph{运动与轨迹}开始。
在研究完整概率分布之前，先理解单个状态如何在局部更新规则下移动、这种运动如何在连续时间中变成微分方程，以及确定性运动和随机运动如何在统一框架下表示，会更有帮助。

但对生成建模而言，仅有运动还不够。
我们最终关心的不只是某一个粒子移动到哪里，而是一整团粒子如何演化，使其概率密度变成我们希望的形状。
因此，下一步将转向\emph{密度演化}：粒子运动如何重塑概率质量，先研究确定性输运，再研究漂移加噪声的动力学，最终得到 Fokker--Planck 方程。

所以，本章按照下面的路线展开：
\[
\text{离散更新}
\;\longrightarrow\;
\text{连续运动}
\;\longrightarrow\;
\text{密度演化}.
\]
这些内容构成下一章的微积分基础；下一章会把它们与基于得分的生成建模相联系，并进一步进入完全连续时间的扩散模型。

\section*{概念导图}
\addcontentsline{toc}{section}{概念导图}

本章建立从离散更新规则过渡到连续时间生成动力学，以及从粒子运动过渡到完整概率密度演化所需的微积分基础。

\begin{itemize}
    \item[\(\rightarrow\)] \textbf{首先通过收缩研究确定性连续时间运动。}
    在一维情形下，收缩展示微分方程如何显式决定一条轨迹；在多维情形下，同一思想通过矩阵与矩阵指数自然推广，并揭示不同方向可以以不同速率收缩。

    \item[\(\rightarrow\)] \textbf{随后通过布朗运动引入随机运动。}
    布朗运动是连续时间随机波动的典型模型，其高斯增量的方差随时间线性增长。它提供了后续扩散过程所依赖的基本随机成分。

    \item[\(\rightarrow\)] \textbf{接着，在 Ornstein--Uhlenbeck 过程中把漂移与噪声结合起来。}
    OU 过程展示了确定性收缩与随机扰动如何共存于同一个系统中，形成均值回归的“漂移+噪声”动力学，并预示后续连续时间生成模型的结构。

    \item[\(\rightarrow\)] \textbf{然后，从粒子运动转向确定性密度演化。}
    从流映射下的概率守恒出发，我们先推导有限时间的变量变换公式，再令时间步趋于无穷小，得到一维 Liouville 方程及沿轨迹的对数密度变化规律。

    \item[\(\rightarrow\)] \textbf{进一步把确定性密度演化推广到多维。}
    这里会引入 Jacobian 作为流映射的局部线性化，行列式用于衡量局部体积变化，散度则刻画净局部流出。这些量共同导出多维 Liouville 方程及其沿轨迹形式。

    \item[\(\rightarrow\)] \textbf{最后加入布朗扩散并得到 Fokker--Planck 方程。}
    纯布朗运动通过高斯平均使密度平滑，由此得到扩散方程；再加入确定性漂移，就得到完整的 Fokker--Planck 方程，说明输运与扩散如何共同控制连续时间随机密度演化。
\end{itemize}

\medskip
\noindent 本章建立的主要关系总结如下图所示。
\begin{center}
\resizebox{0.98\textwidth}{!}{%
\begin{tikzpicture}[
  >=Latex,
  every node/.style={font=\small},
  stage/.style={
    draw, rounded corners=2pt, align=center, inner sep=6pt,
    text width=3.0cm, minimum height=1.35cm
  },
  next/.style={
    draw, rounded corners=2pt, align=center, inner sep=6pt,
    text width=2.65cm, minimum height=1.35cm
  },
  flow/.style={->,thick}
]
\node[font=\normalsize\bfseries] at (6.15,2.55)
  {连续时间基础概览};

\node[stage] (disc) at (0,0) {
  \textbf{小步离散}\\[-1pt]\textbf{更新}\\[3pt]
  $x_{t+\Delta t}\approx x_t+$ 局部变化
};

\node[stage] (运动) at (4.1,0) {
  \textbf{粒子层面}\\[-1pt]\textbf{运动}\\[3pt]
  {\scriptsize ODE: $dX_t=f(X_t,t)\,dt$}\\
  {\scriptsize SDE:}\\[-1pt]{\scriptsize $dX_t=f\,dt+\sigma\,dW_t$}
};

\node[stage] (density) at (8.2,0) {
  \textbf{密度层面}\\[-1pt]\textbf{演化}\\[3pt]
  Liouville：输运\\
  Fokker--Planck：输运 $+$ 扩散
};

\node[next] (得分) at (12.3,0) {
  \textbf{下一章}\\[3pt]
  密度演化\\
  $\longrightarrow$ 得分\\
  $\nabla_x\log p_t(x)$
};

\draw[flow] (disc) -- node[above] {\scriptsize 取极限} (运动);
\draw[flow] (运动) -- node[above] {\scriptsize 系综} (density);
\draw[flow] (density) -- (得分);
\end{tikzpicture}%
}
\end{center}

\medskip

\section{从离散更新到连续运动}

本节引入本章后续将反复使用的连续时间运动语言。
我们先建立状态、速度、微分方程和局部近似等核心概念，使离散更新规则能够被重新解释为连续运动规律。
随后研究三个基本例子：确定性收缩，它以有结构的方式把状态拉向原点；布朗运动，它引入连续时间随机波动；以及 Ornstein--Uhlenbeck 过程，它把漂移和噪声统一在一个随机系统中。
这些例子共同构成后续密度演化与连续时间生成动力学的基本运动模式。

\subsection{连续时间运动的基本语言}

首先引入描述连续时间运动所需的基本语言。
目标是精确定义状态“平滑演化”意味着什么、局部变化如何度量，以及离散更新规则如何转写成无穷小形式。

设 $x(t)$ 表示时刻 $t$ 的某个量的位置。
在一维中，$x(t)$ 只是时间的标量函数；在多维中，它变成一个所有坐标都随时间演化的向量。
一阶导数 $\frac{dx}{dt}$ 描述 $x$ 相对于时间的瞬时变化率，通常解释为速度，也常写作点记号 $\dot{x}(t)$。
类似地，二阶导数 $\frac{d^2x}{dt^2}=\ddot{x}(t)$ 描述速度的变化率，通常解释为加速度。

符号 $dx$ 与 $dt$ 常被非正式地理解为位置的无穷小变化和时间的无穷小变化。
基本直觉是：如果 $\frac{dx}{dt}$ 是变化率，那么在极短的时间区间内，对应的位置变化可表示为
\[
dx=\frac{dx}{dt}\,dt.
\]
因此在无穷小形式下，局部变化仍可理解为
\[
\text{变化}=\text{变化率}\times\text{持续时间}.
\]

一旦我们写出描述位置如何随时间变化的方程，就得到了\emph{微分方程}。
在最简单的情形下，这种方程规定状态在每个可能位置上的速度，速度是其当前位置的函数，有时还显式依赖时间。
在高维中，这种局部运动规则可从几何上理解为：为空间中的每一点指定一个运动方向。
因此，可以把微分方程理解为“运动规则”，把向量场理解为由所有局部箭头构成的图景；这些箭头表示如果一个粒子被放在不同位置，它将向哪个方向移动。

这些思想在这里尤其有用，因为它们允许我们把迭代更新规则重新解释为连续运动规律。
在一个很小的时间间隔 $\Delta t$ 内，一阶近似给出
\[
x(t+\Delta t)\approx x(t)+\frac{dx}{dt}\,\Delta t.
\]
这意味着稍后时刻的值约等于当前值加上“速度乘以一个很小的时间增量”。
因此，它正是离散更新规则在连续时间中的对应形式。

\begin{tcolorbox}[
  colback=blue!4,
  colframe=blue!50!black,
  title={一阶 Taylor 展开},
  boxrule=0.6pt,
  arc=2mm,
  left=2mm,
  right=2mm,
  top=1mm,
  bottom=1mm
]
Taylor 展开利用某一点处的导数来近似该点附近的光滑函数。
在局部形式下，对于很小的位移 $\Delta x$，可以写成
\[
f(x+\Delta x)
=
\sum_{k=0}^{\infty}
\frac{f^{(k)}(x)}{k!}(\Delta x)^k.
\]
写出前几项得到
\[
f(x+\Delta x)
=
f(x)
+
f'(x)\Delta x
+
\frac{f''(x)}{2!}(\Delta x)^2
+
\frac{f^{(3)}(x)}{3!}(\Delta x)^3
+\cdots.
\]

\emph{一阶} Taylor 近似只保留常数项和线性项：
\[
f(x+\Delta x)
\approx
f(x)+f'(x)\Delta x.
\]
被舍弃的项包含 $(\Delta x)^2$、$(\Delta x)^3$ 等更高次幂。
当 $\Delta x$ 很小时，这些高阶项会比线性项小得多。

在当前语境中，函数是轨迹 $x(t)$，展开变量是时间，小位移则是 $\Delta t$。
因此
\[
x(t+\Delta t)
=
x(t)
+
\left.\frac{dx}{dt}\right|_{t}\Delta t
+
O(\Delta t^2).
\]
只保留一阶项就得到
$x(t+\Delta t)
\approx
x(t)
+
\left.\frac{dx}{dt}\right|_{t}\Delta t$。
\end{tcolorbox}

\subsection{确定性收缩}

首先考虑确定性运动：一旦给定当前状态，未来演化就被完全确定。
最简单的例子是\emph{收缩}，即状态被持续拉回原点的运动。
我们先从一维开始，再自然推广到多维，此时不同坐标可以具有不同的收缩速率。

\paragraph{一维收缩。}
在一维中，向原点收缩可表示为
\[
\frac{dx}{dt}=-\alpha x,
\qquad \alpha>0.
\]
这意味着速度始终与当前位置成正比，但符号相反。
若 $x>0$，则 $\dot{x}<0$，状态向左移动、靠近原点；若 $x<0$，则 $\dot{x}>0$，状态向右移动、同样靠近原点。
无论哪种情况，运动都指向内部，而且状态离原点越远，受到的拉回作用越强。

这个方程可以直接求解：
\[
\begin{aligned}
\frac{dx}{dt}=-\alpha x
&\;\Longrightarrow\;
\frac{1}{x}\frac{dx}{dt}=-\alpha
\;\Longrightarrow\;
\frac{d}{dt}\log|x|=-\alpha \\[4pt]
&\;\Longrightarrow\;
\log|x(t)|=-\alpha t + C
\;\Longrightarrow\;
x(t)=Ce^{-\alpha t}.
\end{aligned}
\]
利用初始条件 $x(0)=x_0$，得到 $C=x_0$，因此
\[
\boxed{
x(t)=x_0e^{-\alpha t}.
}
\]

这使系统行为一目了然。
在 $t=0$ 时，状态为 $x_0$；随着时间增加，指数因子 $e^{-\alpha t}$ 逐渐衰减到零，因此 $x(t)$ 的绝对值逐渐趋近原点。
参数 $\alpha$ 控制收缩速度：$\alpha$ 越大，衰减越快；$\alpha$ 越小，衰减越慢。

因此，即便如此简单的例子也已经体现连续时间运动的基本原则：微分方程规定局部运动的方向和速度，而其解则给出随时间形成的完整轨迹。

\paragraph{推广到多维。}
同样的思想可以自然推广到更高维。
现在设 $x(t)\in\mathbb{R}^2$ 为二维状态向量
\[
x(t)=
\begin{bmatrix}
x_1(t)\\
x_2(t)
\end{bmatrix}.
\]
一个简单的多维收缩模型为
\[
\frac{d}{dt}x(t)=Ax(t),
\]
其中 $A$ 是常数矩阵。
若希望两个坐标轴方向彼此独立地收缩，可以取
\[
A=
\begin{bmatrix}
-\alpha_1 & 0\\
0 & -\alpha_2
\end{bmatrix},
\qquad \alpha_1,\alpha_2>0.
\]
于是系统变成
\[
\frac{d}{dt}
\begin{bmatrix}
x_1(t)\\
x_2(t)
\end{bmatrix}
=
\begin{bmatrix}
-\alpha_1 & 0\\
0 & -\alpha_2
\end{bmatrix}
\begin{bmatrix}
x_1(t)\\
x_2(t)
\end{bmatrix},
\]
等价于两个标量方程
\[
\dot{x}_1(t)=-\alpha_1x_1(t),
\qquad
\dot{x}_2(t)=-\alpha_2x_2(t).
\]
每个坐标都按照自身的速率独立收缩：
\[
x_1(t)=x_1(0)e^{-\alpha_1 t},
\qquad
x_2(t)=x_2(0)e^{-\alpha_2 t}.
\]
紧凑地写为
\[
x(t)=e^{At}x(0),
\]
其中 $e^{At}$ 是\emph{矩阵指数}。
在当前例子中，$A$ 为对角矩阵，因此矩阵指数只需对每个对角元素应用普通标量指数：
\[
e^{At}
=
\begin{bmatrix}
e^{-\alpha_1 t} & 0\\
0 & e^{-\alpha_2 t}
\end{bmatrix}.
\]

\begin{tcolorbox}[
  colback=blue!4,
  colframe=blue!50!black,
  title={通过 Maclaurin 级数理解矩阵指数},
  boxrule=0.6pt,
  arc=2mm,
  left=2mm,
  right=2mm,
  top=1mm,
  bottom=1mm
]
光滑函数在零点附近的 Taylor 展开称为\emph{Maclaurin 级数}。对于标量函数 $f$，其形式为\footnote{译注：原文把光滑函数与其无穷 Taylor 级数相等作为一般表述；仅有光滑性并不足以保证这一点，还需函数在所讨论邻域内解析。这里的指数函数满足该条件。后文推导扩散方程时，只需在适当正则条件下使用带余项的有限阶 Taylor 展开。}
\[
f(u)
=
\sum_{k=0}^{\infty}\frac{f^{(k)}(0)}{k!}u^k
=
f(0)
+
f'(0)u
+
\frac{f''(0)}{2!}u^2
+
\frac{f^{(3)}(0)}{3!}u^3
+\cdots.
\]
对于指数函数，其任意阶导数仍然是 $e^u$，所以在 $u=0$ 处都等于 $1$。因此
\[
e^u
=
1+u+\frac{u^2}{2!}+\frac{u^3}{3!}+\cdots
=
\sum_{k=0}^{\infty}\frac{u^k}{k!}.
\]

矩阵指数使用同一个级数来定义，只需把标量 $u$ 替换成矩阵 $At$：
\[
\boxed{
e^{At}
=
I
+
At
+
\frac{(At)^2}{2!}
+
\frac{(At)^3}{3!}
+\cdots.
}
\]
这里，标量 $1$ 被单位矩阵 $I$ 取代，因为 $I$ 是矩阵乘法中的乘法单位元。
对于当前使用的对角矩阵
$A=
\begin{bmatrix}
-\alpha_1 & 0\\
0 & -\alpha_2
\end{bmatrix}$，
有
\[
At=
\begin{bmatrix}
-\alpha_1 t & 0\\
0 & -\alpha_2 t
\end{bmatrix}.
\]
对角矩阵的幂仍然是对角矩阵：
\[
(At)^k
=
\begin{bmatrix}
(-\alpha_1 t)^k & 0\\
0 & (-\alpha_2 t)^k
\end{bmatrix}.
\]
将这些幂代入级数，得到
\[
e^{At}
=
\begin{bmatrix}
\sum_{k=0}^{\infty}\frac{(-\alpha_1 t)^k}{k!} & 0\\
0 & \sum_{k=0}^{\infty}\frac{(-\alpha_2 t)^k}{k!}
\end{bmatrix}
=
\begin{bmatrix}
e^{-\alpha_1 t} & 0\\
0 & e^{-\alpha_2 t}
\end{bmatrix}.
\]
\end{tcolorbox}

这个多维例子引入了一个重要的新特征：运动在不同方向上可以以不同速度收缩。
如果 $\alpha_1>\alpha_2$，那么第一个坐标比第二个坐标衰减得更快，因此状态在 $x_1$ 方向上的收缩速度高于 $x_2$ 方向。
所以运动仍然趋向原点，但各方向上的收缩并不需要均匀一致。

对角形式也说明坐标之间仍不存在相互作用：每个分量都独立演化。
虽然为了清晰起见这里写成二维情形，但同一思想可以直接推广到 $\mathbb{R}^d$。
对角矩阵
\[
A=\mathrm{diag}(-\alpha_1,-\alpha_2,\dots,-\alpha_d)
\]
会使每个坐标方向独立地进行指数收缩。

因此，对角收缩是一个很有用的第一个多维例子。
它以最简单的方式引入基于矩阵的连续运动，也自然架起了从标量微分方程到向量值连续动力学的桥梁。

\subsection{布朗运动}

到目前为止，我们的讨论完全是确定性的：一旦固定初始状态，未来轨迹就由控制它的微分方程完全决定。
然而，许多连续时间生成模型还包含随机性。
在每一个很短的时间区间内，状态都可能受到随机扰动，因此仅凭初始条件无法完全预测后续运动。

描述这类连续时间随机波动的基本模型就是\emph{布朗运动}（Brownian motion），也称为\emph{Wiener 过程}~\cite{oksendal2003sde}。
通常记为 $\{W_t\}_{t\ge 0}$，可以通过三个基本性质刻画：
\begin{enumerate}
    \item $W_0=0$；
    \item 不相交时间区间上的增量彼此独立；
    \item 对任意 $\Delta t>0$，
    \[
    W_{t+\Delta t}-W_t \sim \mathcal{N}(0,\Delta t).
    \]
\end{enumerate}

第三个性质尤其重要。
它表示在长度为 $\Delta t$ 的短时间区间内，随机增量均值为零、方差为 $\Delta t$。
因此，布朗运动引入的波动平均而言没有偏置，但其不确定性会随经历的时间增加而增长。

高斯假设也很自然。
如果把一个短区间内的运动想象为大量微小且近似独立的随机扰动的累积结果，那么它们的总和可以很好地用高斯随机变量建模。
这与中心极限定理背后的直觉密切相关：许多微小独立效应相加时，其总和趋向于高斯形式。

一个有用的理解方法是把区间 $[0,t]$ 划分为许多小段，并把 $W_t$ 写成独立高斯增量之和。
若把 $[0,t]$ 等分为 $n$ 段，每段长度
$\Delta t=\frac{t}{n}$，则
\[
W_t-W_0
=
\sum_{k=0}^{n-1}\bigl(W_{(k+1)\Delta t}-W_{k\Delta t}\bigr)
=
\sum_{k=0}^{n-1}\Delta W_k,
\qquad
\Delta W_k:=W_{(k+1)\Delta t}-W_{k\Delta t}.
\]

由于 $W_0=0$，得到
\[
W_t=\sum_{k=0}^{n-1}\Delta W_k.
\]

\begin{tcolorbox}[
  colback=blue!4,
  colframe=blue!50!black,
  title={时刻 $t$ 的布朗运动均值与方差},
  boxrule=0.6pt,
  arc=2mm,
  left=2mm,
  right=2mm,
  top=1mm,
  bottom=1mm
]
每个增量满足
\[
\Delta W_k\sim \mathcal{N}(0,\Delta t),
\qquad
\mathbb{E}[\Delta W_k]=0,
\qquad
\mathrm{Var}(\Delta W_k)=\Delta t=\frac{t}{n}.
\]
因此
\[
\mathbb{E}[W_t]
=
\sum_{k=0}^{n-1}\mathbb{E}[\Delta W_k]
=
0,
\]
并且利用独立性，
\[
\mathrm{Var}(W_t)
=
\sum_{k=0}^{n-1}\mathrm{Var}(\Delta W_k)
=
\sum_{k=0}^{n-1}\frac{t}{n}
=
t.
\]
所以
\[
\mathbb{E}[W_t]=0,
\qquad
\mathrm{Var}(W_t)=t.
\]
\end{tcolorbox}

因此，在任意固定时刻 $t$，布朗运动都以零为中心，但其离散程度随时间线性增长。
这意味着布朗运动不会把状态拉向任何偏好的位置；它描述的是一种随机游走，其不确定性随着时间推移不断增大。

一个有用的启发式写法是
\[
W_{t+\Delta t}-W_t \approx \sqrt{\Delta t}\,\varepsilon_t,
\qquad
\varepsilon_t\sim\mathcal{N}(0,1).
\]
这个形式特别清楚地展示了缩放规律。
由于增量方差为 $\Delta t$，其典型大小是 $\sqrt{\Delta t}$ 阶，而不是 $\Delta t$ 阶。
这种平方根缩放是随机连续时间运动最典型的特征之一，后续也会在扩散模型中再次出现。

布朗运动常被非正式地写成微分形式 $dW_t$，表示一个无穷小随机增量。
虽然 $dW_t$ 不是确定性微积分意义下的普通微分，但这种记号非常有用，因为它使随机运动能够用与常微分方程平行的形式书写。

\subsection{Ornstein--Uhlenbeck 过程}

Ornstein--Uhlenbeck（OU）过程是一个特别重要的随机连续时间运动例子。
它把确定性收缩与布朗噪声结合起来：
\[
dX_t=-\alpha X_t\,dt+\sigma\,dW_t,
\qquad
\alpha>0,\ \sigma>0.
\]

这个方程包含两部分。
漂移项 $-\alpha X_t\,dt$ 像前面介绍的一维收缩模型一样，把状态拉回原点；噪声项 $\sigma\,dW_t$ 则注入由布朗运动驱动的随机扰动。
因此，OU 过程平衡了两种相互竞争的趋势：确定性地拉向原点，以及随机地偏离原点。

在一个很小的时间区间 $\Delta t$ 内，可以启发式地写成更新式
\[
X_{t+\Delta t}
\approx
X_t-\alpha X_t\,\Delta t+\sigma\sqrt{\Delta t}\,\varepsilon_t,
\qquad
\varepsilon_t\sim\mathcal{N}(0,1).
\]
这个离散时间视角很有用。
它说明下一状态由当前状态加上两部分得到：
\begin{itemize}
    \item 大小为 $\Delta t$ 阶的确定性修正；
    \item 大小为 $\sqrt{\Delta t}$ 阶的随机扰动。
\end{itemize}
因此，OU 过程已经具有后续 Langevin 动力学和扩散模型中会出现的基本“漂移+噪声”结构。

OU 过程还存在显式解：
\[
X_t
=
X_0e^{-\alpha t}
+
\sigma \int_0^t e^{-\alpha(t-s)}\,dW_s.
\]
这个公式清楚揭示了其结构。
第一项 $X_0e^{-\alpha t}$ 正是初始状态的确定性指数收缩；第二项累积所有过去的随机扰动，但每个扰动都乘以衰减因子 $e^{-\alpha(t-s)}$。
因此，更早发生的噪声影响会逐渐减弱，而更近期的扰动影响更大。

这就是为什么 OU 过程常被描述为一种\emph{均值回归}随机过程：它始终具有随机性，同时保持稳定的恢复趋势。
图~\ref{fig:continuous-three-motion-laws} 将这一行为与确定性收缩和布朗运动并列展示，清楚说明 OU 过程如何把两种基本效应结合起来。

不过，到目前为止，我们仍然只是在单个粒子的层面描述运动。
下一步要理解的是，这类运动如何随时间重塑整个概率密度。

\begin{figure}[H]
\centering
\begin{minipage}[t]{0.45\textwidth}
  \vspace{0pt}\centering
  \includegraphics[width=\linewidth]{figures/ch05/panels/fig5_1a_contraction.pdf}\\[-0.2em]
  {\small $\displaystyle \frac{dX_t}{dt}=-\alpha X_t$}
\end{minipage}\hfill
\begin{minipage}[t]{0.45\textwidth}
  \vspace{0pt}\centering
  \includegraphics[width=\linewidth]{figures/ch05/panels/fig5_1b_brownian.pdf}\\[-0.2em]
  {\small $\displaystyle W_{t+\Delta t}-W_t\sim\mathcal N(0,\Delta t)$}
\end{minipage}

%\vspace{0.8em}

\begin{minipage}[t]{0.45\textwidth}
  \vspace{0pt}\centering
  \includegraphics[width=\linewidth]{figures/ch05/panels/fig5_1c_ou.pdf}\\[-0.2em]
  {\small $\displaystyle dX_t=-\alpha X_t\,dt+\sigma\,dW_t$}
\end{minipage}\hfill
\begin{minipage}[t]{0.45\textwidth}
  \vspace{0pt}
  \caption{\textbf{三种基本的连续时间运动规律。}
  确定性收缩产生平滑轨迹并逐渐衰减到原点。Brownian 运动具有零均值的独立高斯增量，因此轨迹随机游走，其典型扩散尺度与 $\sqrt{t}$ 成正比；虚线表示单位 Brownian 运动的 $\pm2\sqrt{t}$。Ornstein--Uhlenbeck 过程把随机波动与指向原点的回复漂移结合起来；虚线表示平稳标准差的两倍。}
  \label{fig:continuous-three-motion-laws}
\end{minipage}
\end{figure}


\section{确定性密度演化}

上一节关注的是单个粒子的运动。
这一视角已经很有用，因为许多连续时间模型确实由作用于状态或轨迹的微分方程来描述。
然而，对生成建模而言，一个粒子的运动只是故事的一部分。
最终我们关心的是一整组粒子的\emph{概率密度}如何随时间变化。

因此，本节完成一个重要过渡：从\emph{轨迹动力学}转向\emph{密度动力学}。
我们从最简单的确定性情形开始，在这种情形下，粒子由一个映射输运，并且概率质量守恒。
这会先导出有限时间的变量变换公式，再导出无穷小密度演化规律，最后通过散度得到其多维形式。
整个过程中，底层思想始终相同：
\[
\text{粒子运动}
\;\longrightarrow\;
\text{概率质量输运}
\;\longrightarrow\;
\text{确定性密度演化}.
\]

\subsection{有限时间密度变换}

先从最基本的密度原则开始：\emph{概率在确定性输运下守恒}。
这里的含义是：动力系统以完全确定的方式把粒子从一个位置移动到另一个位置，而与粒子关联的概率质量只是随之被携带，不会被创造或消失。

设时刻 $0$ 的粒子位置为 $x_0$，其密度为 $p_0(x_0)$。
再假设确定性动力系统把这个初始点输运到稍后的
\[
x_t=\phi_t(x_0),
\]
其中 $\phi_t$ 表示把粒子从时刻 $0$ 移动到时刻 $t$ 的轨迹映射或流映射。

为了理解密度如何变化，考虑 $x_0$ 附近一个长度为 $dx_0$ 的极小区间。
该微小区间中包含的概率质量近似为
\[
p_0(x_0)\,dx_0.
\]
演化到时刻 $t$ 后，同一小块概率质量被输运到 $x_t=\phi_t(x_0)$ 附近的新区间。
设输运后区间长度为 $dx_t$。
由于确定性运动只是在搬运粒子而不创造或销毁概率质量，总概率必须保持不变：
\[
p_0(x_0)\,dx_0=p_t(x_t)\,dx_t.
\]

因此，关键问题变成：小区间 $dx_0$ 在流映射下如何变换？
在一维中，这种局部拉伸或压缩由映射的导数度量。
推导如下。

\begin{tcolorbox}[
  colback=blue!4,
  colframe=blue!50!black,
  title={一维变量变换公式的推导},
  boxrule=0.6pt,
  arc=2mm,
  left=2mm,
  right=2mm,
  top=1mm,
  bottom=1mm
]
不仅跟踪点 $x_0$，同时跟踪其附近的点 $x_0+dx_0$。
在同一个流映射下，
\[
x_t=\phi_t(x_0), \qquad x_t+dx_t=\phi_t(x_0+dx_0),
\]
所以
$dx_t=\phi_t(x_0+dx_0)-\phi_t(x_0)$。
由于 $dx_0$ 很小，对 $\phi_t$ 在 $x_0$ 附近作一阶 Taylor 展开，得到
\[
\phi_t(x_0+dx_0)\approx \phi_t(x_0)+\frac{d\phi_t(x_0)}{dx_0}\,dx_0.
\]
因此
$dx_t\approx \frac{d\phi_t(x_0)}{dx_0}\,dx_0
= \frac{dx_t}{dx_0}\,dx_0$。
代入概率守恒关系，
\[
p_0(x_0)\,dx_0
=
p_t(x_t)\left|\frac{dx_t}{dx_0}\right|dx_0,
\]
约去 $dx_0$ 后得到
\[
\boxed{
p_t(x_t)=p_0(x_0)\left|\frac{dx_t}{dx_0}\right|^{-1}.
}
\]
\end{tcolorbox}

这就是由确定性流诱导的一维变量变换公式。
它的直观含义非常清楚：如果变换把一个区间拉长，那么同样数量的概率会分摊到更大的区域中，因此密度降低；如果变换把区间压缩，那么同样的概率质量被挤压到更小区域中，因此密度增大。
所以，导数 $\frac{dx_t}{dx_0}$ 衡量局部长度改变了多少，而其倒数则告诉我们密度应如何重新缩放。

同样的思想可以自然推广到多维。
在一维中，我们跟踪一个微小\emph{长度}如何在映射下变化；二维中要跟踪微小\emph{面积}，三维中要跟踪微小\emph{体积}，一般情况下则是一个 $d$ 维微小体积元。
衡量这种局部缩放的量不再是普通导数，而是变换 Jacobian 矩阵的行列式。
因此，多维变量变换公式变成
\[
p_t(x_t)
=
p_0(x_0)
\left|
\det\!\left(\frac{\partial x_t}{\partial x_0}\right)
\right|^{-1}.
\]
所以在高维中，行列式扮演一维导数 $\frac{dx_t}{dx_0}$ 的角色，用来度量一个微小局部区域被流拉伸或压缩了多少。

需要注意的是，这个公式描述的是\emph{有限时间}确定性变换。
它告诉我们密度如何在时刻 $0$ 与时刻 $t$ 之间变化，但尚未描述密度变化的无穷小速率。
为此，下面推导局部时间意义下的连续性方程。

\subsection{一维无穷小密度演化：Liouville 方程}

有限时间变量变换公式告诉我们密度在时刻 $0$ 与时刻 $t$ 之间如何变换。
然而在连续时间建模中，更重要的往往是理解密度在一个\emph{无穷小}时间区间内如何变化，因为连续动力学正是由这种局部更新构成的。
这就导出了\emph{Liouville 方程}；在确定性情形下，它也称为连续性方程。

设粒子服从一维常微分方程
\[
\frac{dx}{dt}=f(x,t),
\]
其中 $f(x,t)$ 是局部速度场。
为了理解密度如何变化，现在采用固定空间视角：不再跟随某一个粒子，而是站在空间中的一个小区间处，观察概率质量如何穿过该区间。

考虑短区间 $[x,x+\Delta x]$。
时刻 $t$ 区间内的总概率质量近似为
\[
p(x,t)\,\Delta x.
\]
在一个无穷小时间区间内，这部分局部概率质量只能通过左边界流入，或通过右边界流出而发生改变。

相关的流量称为\emph{概率通量}，即密度乘以速度：
\[
p(x,t)f(x,t).
\]
它表示单位时间内穿过某一位置的概率质量。
因此，左边界流入量为
\[
p(x,t)f(x,t),
\]
而右边界流出量为
\[
p(x+\Delta x,t)f(x+\Delta x,t).
\]

概率守恒意味着，区间内概率质量的变化率必须等于流入减去流出：
\[
\frac{\partial}{\partial t}p(x,t)\,\Delta x
=
p(x,t)f(x,t)-p(x+\Delta x,t)f(x+\Delta x,t).
\]
两边除以 $\Delta x$，得到
\[
\frac{\partial}{\partial t}p(x,t)
=
\frac{p(x,t)f(x,t)-p(x+\Delta x,t)f(x+\Delta x,t)}{\Delta x}.
\]
整理为
\[
\frac{\partial}{\partial t}p(x,t)
=
-
\frac{p(x+\Delta x,t)f(x+\Delta x,t)-p(x,t)f(x,t)}{\Delta x}.
\]
令 $\Delta x\to 0$，右侧变为空间导数，因此得到
\[
\boxed{
\partial_t p(x,t)
=
-\partial_x\bigl(p(x,t)f(x,t)\bigr).
}
\]
这就是一维 Liouville 方程。

该方程表示密度的变化由概率通量的净流出决定。
若某个区域流出的概率多于流入，密度就会下降；若流入多于流出，密度就会上升。
利用乘积法则展开右侧往往很有帮助：
\[
\partial_t p
=
-\bigl(f\,\partial_x p + p\,\partial_x f\bigr).
\]
由此可以看出两个不同的贡献。

项 $-f\,\partial_x p$ 依赖密度自身的空间斜率，因此描述\emph{输运}：已有密度轮廓被速度场携带着穿过空间。
如果流只是把粒子从一个位置移动到另一个位置，而不改变它们之间的相对间距，那么主要效果就是整个密度轮廓发生平移。

相比之下，项 $-p\,\partial_x f$ 依赖速度场的空间斜率，因此描述\emph{形变}：相邻粒子的速度不再完全相同，于是它们会彼此远离或相互压缩。
如果 $\partial_x f>0$，速度随位置增大而增大，相邻粒子彼此分开，密度下降；如果 $\partial_x f<0$，速度随位置增大而减小，相邻粒子相互压缩，密度上升。

\subsection{沿轨迹的密度变化}

Liouville 方程描述固定空间位置上的密度如何变化。
但还有另一个同样有用的视角：沿着一个运动粒子的轨迹，密度如何变化？
这一视角把注意力从固定位置上发生的变化，转向粒子随流移动时自身经历的变化。

设粒子服从确定性动力学
\[
\frac{dx(t)}{dt}=f(x(t),t).
\]
现在考虑沿该粒子路径取值的密度
\[
p(x(t),t).
\]
由于位置和时间都在变化，我们利用链式法则对其求导。
等价地，这就是\emph{全导数}：它同时包含对时间的显式变化，以及由于粒子本身在空间中移动而引起的隐式变化。
因此
\[
\frac{d}{dt}p(x(t),t)
=
\partial_t p(x(t),t)
+
\frac{dx(t)}{dt}\,\partial_x p(x(t),t).
\]
由于 $\dot{x}(t)=f(x(t),t)$，得到
\[
\frac{d}{dt}p(x(t),t)
=
\partial_t p + f\,\partial_x p.
\]

代入 Liouville 方程
\[
\partial_t p
=
-\partial_x(pf)
=
-(f\,\partial_x p + p\,\partial_x f),
\]
于是
\[
\frac{d}{dt}p(x(t),t)
=
-(f\,\partial_x p + p\,\partial_x f)+f\,\partial_x p.
\]
输运项相互抵消，只剩
\[
\frac{d}{dt}p(x(t),t)
=
-p(x(t),t)\,\partial_x f(x(t),t).
\]
两边除以 $p(x(t),t)$，得到
\[
\frac{d}{dt}\log p(x(t),t)
=
-\partial_x f(x(t),t).
\]
因此
\[
\boxed{
\frac{d}{dt}\log p(x(t),t)
=
-\partial_x f(x(t),t).
}
\]

这是一个重要结果。
上一小节把密度演化分解成输运项和形变项。
而这里可以看到，一旦我们随粒子一起移动，输运项就消失了：从粒子自身的视角看，它不再观察密度轮廓被流携带穿过空间，因为它自己就在随流移动。
此时剩下的只有局部膨胀或压缩效应。

因此，当 $\partial_x f>0$ 时，相邻粒子趋于分离，沿轨迹的密度下降；当 $\partial_x f<0$ 时，相邻粒子趋于压缩，密度上升。
所以，固定空间视角和随粒子移动视角强调了同一动力学的不同方面。
Liouville 方程描述整个密度场如何在空间中改变，而上面的公式告诉我们密度沿单条轨迹如何变化。

\subsection{散度与多维 Liouville 方程}

现在把密度演化从一维推广到完整的多维状态空间。
主要的新特征是运动与概率通量都变成向量值。
在一维中，标量导数 $\partial_x$ 足以描述概率如何流入或流出一个小区间；而在多维中，对应的量是\emph{散度}，它衡量一个向量场从空间微小区域中向外流出的净量。
本节将说明散度如何从局部体积变化中自然出现，以及它如何导出多维 Liouville 方程。

设 $x(t)\in\mathbb{R}^d$，粒子按照
\[
\frac{dx}{dt}=v(x,t)
\]
演化，其中速度场为
\[
v(x,t)=
\begin{bmatrix}
v_1(x,t)\\
v_2(x,t)\\
\vdots\\
v_d(x,t)
\end{bmatrix}.
\]

与一维情形一样，基本思想仍然是局部概率守恒。
不过在写出密度方程之前，需要先理解空间中的一个微小区域如何在流作用下变化。
在一维中，这由映射的导数度量，它告诉我们一个小区间如何被拉伸或压缩；在多维中，对应对象是\emph{Jacobian 矩阵}及其行列式。

在很小的时间步 $t\to t+dt$ 上，运动可以近似为
\[
x(t+dt)=x(t)+v(x(t),t)\,dt.
\]
所以无穷小变换是
\[
x\mapsto x+v(x,t)\,dt.
\]

为了理解一个微小体积如何变化，我们考察更新后的点如何依赖起始位置。
这种依赖通过对映射关于 $x$ 求导来刻画，从而得到该变换的\emph{Jacobian 矩阵}：
\[
\frac{\partial x(t+dt)}{\partial x}
=
I+\frac{\partial v}{\partial x}dt.
\]
若把速度场的 Jacobian 记为
\[
J_v(x,t)=\frac{\partial v}{\partial x},
\]
则
\[
\frac{\partial x(t+dt)}{\partial x}=I+J_v(x,t)\,dt.
\]

Jacobian 告诉我们微小位移在流作用下如何变化。
要从位移过渡到体积，就取其行列式。

\begin{tcolorbox}[
  colback=blue!4,
  colframe=blue!50!black,
  title={局部体积缩放与散度},
  boxrule=0.6pt,
  arc=2mm,
  left=2mm,
  right=2mm,
  top=1mm,
  bottom=1mm
]
局部体积缩放因子为
\[
\det\!\bigl(I+J_v(x,t)\,dt\bigr).
\]
当 $dt$ 很小时，一阶展开给出
\[
\det(I+A\,dt)=1+\operatorname{tr}(A)\,dt+O(dt^2).
\]
令 $A=J_v(x,t)$，得到
\[
\det\!\bigl(I+J_v(x,t)\,dt\bigr)
=
1+\operatorname{tr}(J_v(x,t))\,dt+O(dt^2).
\]
Jacobian 的迹有一个专门名称：
\[
\operatorname{tr}(J_v(x,t))
=
\nabla\cdot v(x,t)
=
\sum_{i=1}^d \frac{\partial v_i}{\partial x_i}(x,t).
\]
因此
\[
\det\!\bigl(I+J_v(x,t)\,dt\bigr)
=
1+(\nabla\cdot v(x,t))\,dt+O(dt^2),
\]
所以一个微小体积元的一阶演化为
\[
dV(t+dt)=\bigl(1+(\nabla\cdot v(x,t))\,dt\bigr)dV(t).
\]
\end{tcolorbox}

这就给出了散度的关键几何解释：它是一个微小体积在流作用下瞬时膨胀或收缩的速率。
若 $\nabla\cdot v>0$，流在局部使体积膨胀；若 $\nabla\cdot v<0$，则在局部压缩体积。

有了这个几何图景，就可以推导密度演化方程。
一个随流移动的小体积中所包含的概率质量必须保持不变：
\[
p(x,t)\,dV(t)
=
p(x+v(x,t)dt,t+dt)\,dV(t+dt).
\]
代入体积缩放关系
\[
dV(t+dt)=\bigl(1+(\nabla\cdot v)\,dt\bigr)dV(t),
\]
得到
\[
p(x,t)
=
p(x+vdt,t+dt)\bigl(1+(\nabla\cdot v)\,dt\bigr).
\]

接下来，对密度应用多变量一阶 Taylor 展开。
\begin{tcolorbox}[
  colback=blue!4,
  colframe=blue!50!black,
  title={多维 Liouville 方程的推导},
  boxrule=0.6pt,
  arc=2mm,
  left=2mm,
  right=2mm,
  top=1mm,
  bottom=1mm
]
对于依赖位置与时间的函数 $p(x,t)$，多变量一阶 Taylor 展开为
\[
p(x+\Delta x,t+\Delta t)
\approx
p(x,t)+\nabla p(x,t)\cdot \Delta x+\partial_t p(x,t)\Delta t.
\]
在这里的流中，$\Delta x=v(x,t)dt$ 且 $\Delta t=dt$，因此
\[
\begin{aligned}
p(x+vdt,t+dt)
&=
p(x,t)+\nabla p(x,t)\cdot(vdt)+\partial_t p(x,t)dt+O(dt^2)\\
&=
p(x,t)+(v\cdot\nabla p)dt+(\partial_t p)dt+O(dt^2).
\end{aligned}
\]

把它代入概率守恒关系
\[
p(x,t)
=
p(x+vdt,t+dt)\bigl(1+(\nabla\cdot v)dt\bigr),
\]
得到
\[
\begin{aligned}
p
&=
\left[p+(\partial_t p)dt+(v\cdot\nabla p)dt\right]
\left[1+(\nabla\cdot v)dt\right]
+O(dt^2)\\
&=
p+(\partial_t p)dt+(v\cdot\nabla p)dt+p(\nabla\cdot v)dt+O(dt^2).
\end{aligned}
\]
两边减去 $p$、除以 $dt$ 并令 $dt\to0$，得到
\[
0=\partial_t p+v\cdot\nabla p+p(\nabla\cdot v).
\]
利用恒等式 $\nabla\cdot(pv)=v\cdot\nabla p+p(\nabla\cdot v)$，可写为
\[
0=\partial_t p+\nabla\cdot(pv),
\]
等价地，
\[
\boxed{
\partial_t p(x,t)=-\nabla\cdot\bigl(p(x,t)v(x,t)\bigr).
}
\]
\end{tcolorbox}

这就是多维 Liouville 方程，它是一维连续性方程的自然推广。
原来的一维标量导数 $\partial_x(pf)$ 现在变成了向量值概率通量 $pv$ 的散度。
因此，密度的变化由概率从空间微小区域中净流出的程度决定。

沿轨迹的视角也可以自然推广。
沿着满足
\[
\dot{x}(t)=v(x(t),t)
\]
的路径，对数密度按照
\[
\boxed{
\frac{d}{dt}\log p(x(t),t)
=
-\nabla\cdot v(x(t),t).
}
\]
演化。
所以在多维中，与一维完全类似，一旦我们随粒子本身一起移动，输运部分就消失，剩余的密度变化由流的局部膨胀或压缩决定。
图~\ref{fig:continuous-flow-density-jacobian} 总结了有限时间体积变化、密度重新缩放和无穷小散度关系之间的几何联系。

\begin{figure}[H]
\centering
\resizebox{0.76\textwidth}{!}{%
\begin{tikzpicture}[
  >=Latex,
  every node/.style={font=\small},
  panel/.style={draw,rounded corners=2pt,minimum width=4.0cm,minimum height=3.25cm,align=center},
  eq/.style={draw,rounded corners=2pt,align=center,inner sep=5pt,text width=5.4cm},
  flow/.style={->,thick}
]
\node[font=\normalsize\bfseries] at (6.0,3.05)
  {确定性流：局部体积变化控制密度};

\node[panel] (left) at (2.0,0.65) {};
\node[panel] (right) at (10.0,0.65) {};

\node[font=\bfseries] at (2.0,1.85) {初始局部区域};
\node[font=\bfseries] at (10.0,1.85) {经过流映射 $F_t$ 后};

\draw[thick] (1.05,0.05) rectangle (2.95,1.10);
\node at (2.0,0.58) {$\Delta V_0$};
\node at (2.0,-0.45) {$p_0(x_0)$};

\draw[thick] (9.25,0.25) rectangle (10.75,0.92);
\node at (10.0,0.58) {$\Delta V_t$};
\node at (10.0,-0.45) {$p_t(x_t)$};

\draw[flow] (4.0,0.75) -- node[above,align=center]
  {$x_t=F_t(x_0)$\\[-1pt]
   $J_t=\partial F_t/\partial x_0$} (8.0,0.75);

\node[eq] (volume) at (3.0,-2.05) {
  \textbf{局部体积缩放}\\[3pt]
  $\displaystyle \Delta V_t\approx |\det J_t|\,\Delta V_0$
};

\node[eq] (density) at (9.0,-2.05) {
  \textbf{概率守恒}\\[3pt]
  $\displaystyle p_t(F_t(x_0))=\frac{p_0(x_0)}{|\det J_t|}$
};

\draw[flow] (volume) -- (density);

\node[eq,text width=7.0cm] (instant) at (6.0,-4.05) {
  \textbf{沿轨迹的无穷小极限}\\[3pt]
  $\displaystyle \frac{d}{dt}\log p(x(t),t)=-\nabla\!\cdot v(x(t),t)$
};


\end{tikzpicture}%
}
\caption{\textbf{概率守恒把局部体积变化与密度变化联系起来。}
确定性流映射 $F_t$ 会移动粒子及其周围的体积元。Jacobian 行列式度量局部体积缩放，因此概率守恒要求密度按照体积因子的倒数重新缩放。在无穷小极限下，这给出沿轨迹的 Liouville 关系：局部膨胀降低密度，而局部压缩提高密度。}
\label{fig:continuous-flow-density-jacobian}
\end{figure}


\section{随机密度演化与 Fokker--Planck 方程}

到目前为止，密度演化完全是确定性的：概率质量由流输运，其局部变化由散度控制。
然而，许多连续时间生成模型还包含随机扰动。
一旦加入这种扩散，密度就不再只通过输运和局部膨胀/压缩而变化，还会在空间中扩散开来。
这就把我们从 Liouville 方程带到\emph{Fokker--Planck 方程}~\cite{risken1996fokker}。

\subsection{纯布朗扩散作为高斯平滑}

为了建立直觉，先考虑一维纯布朗扩散：
\[
dX_t=\sigma\,dW_t.
\]
在很短的时间区间 $\Delta t$ 内，增量服从高斯分布：
\[
X_{t+\Delta t}-X_t \sim \mathcal{N}(0,\sigma^2\Delta t).
\]
因此，一个原本位于 $y$ 的粒子，在这一短时间步之后有一定概率到达 $x$ 附近，而这一概率由以 $y$ 为中心的高斯核描述。
所以更新后的密度可由所有邻近旧位置的贡献累加得到：
\[
p(x,t+\Delta t)
=
\int p(y,t)\,
\frac{1}{\sqrt{2\pi \sigma^2\Delta t}}
\exp\!\left(
-\frac{(x-y)^2}{2\sigma^2\Delta t}
\right)dy.
\]

换句话说，短时间后的 $x$ 处密度是先前 $x$ 附近密度的加权平均，权重由短时间高斯转移核给出。
这就是与窄高斯核的卷积。
直观上，概率质量从邻近位置扩散进入 $x$，同时原本位于 $x$ 附近的一些概率也扩散出去；净结果是密度被平滑并逐渐展宽。

为了更方便分析，记高斯增量为 $\eta$，其中
\[
\eta\sim \mathcal{N}(0,\sigma^2\Delta t).
\]
作变量替换
\[
\eta=x-y,
\qquad\text{因此}\qquad
y=x-\eta.
\]
于是密度更新可以写成
\[
p(x,t+\Delta t)
=
\int p(x-\eta,t)\,G_{\Delta t}(\eta)\,d\eta,
\]
其中
\[
G_{\Delta t}(\eta)
=
\frac{1}{\sqrt{2\pi \sigma^2\Delta t}}
\exp\!\left(
-\frac{\eta^2}{2\sigma^2\Delta t}
\right)
\]
是增量的高斯密度。
由于
\[
\int G_{\Delta t}(\eta)\,d\eta=1,
\]
也可以把它理解成期望：
\[
p(x,t+\Delta t)=\mathbb{E}_{\eta}\bigl[p(x-\eta,t)\bigr].
\]

\subsection{从布朗扩散到 Fokker--Planck 方程}

现在利用 Taylor 级数在 $x$ 附近展开 $p(x-\eta,t)$。
一般而言，对于足够光滑的一元函数，
$f(x+h)=\sum_{k=0}^{\infty}\frac{f^{(k)}(x)}{k!}h^k$。
令 $f(\cdot)=p(\cdot,t)$、$h=-\eta$，得到
\[
\begin{aligned}
p(x-\eta,t)
&=
\sum_{k=0}^{\infty}
\frac{\partial_x^{\,k}p(x,t)}{k!}(-\eta)^k \\[4pt]
&=
p(x,t)
-\eta\,\partial_x p(x,t)
+\frac{(-\eta)^2}{2!}\partial_{xx}p(x,t)
+\frac{(-\eta)^3}{3!}\partial_{xxx}p(x,t)
+\cdots \\[4pt]
&=
p(x,t)-\eta\,\partial_x p(x,t)
+\frac{\eta^2}{2}\partial_{xx}p(x,t)
+O(\eta^3).
\end{aligned}
\]
对 $\eta$ 取期望，得到
\[
p(x,t+\Delta t)
=
\mathbb{E}_{\eta}\bigl[p(x-\eta,t)\bigr]
=
p(x,t)
-
\mathbb{E}[\eta]\partial_x p(x,t)
+
\frac{1}{2}\mathbb{E}[\eta^2]\partial_{xx}p(x,t)
+\cdots.
\]
由于高斯增量均值为零，
\[
\mathbb{E}[\eta]=0,
\]
一阶项消失。
其方差为
\[
\mathbb{E}[\eta^2]=\sigma^2\Delta t,
\]
所以二阶项保留下来：
\[
p(x,t+\Delta t)
=
p(x,t)+\frac{\sigma^2\Delta t}{2}\partial_{xx}p(x,t)+o(\Delta t).
\]
两边减去 $p(x,t)$，除以 $\Delta t$，再令 $\Delta t\to 0$，得到
\[
\boxed{
\partial_t p(x,t)=\frac{\sigma^2}{2}\partial_{xx}p(x,t).
}
\]
因此，纯布朗运动导出了纯扩散方程。

这就是随机扩散在密度层面的特征。
Liouville 方程通过确定性输运移动密度，而扩散方程则使密度在空间中变得更加平滑并展宽。

现在再加入确定性漂移：
\[
dX_t=f(X_t,t)\,dt+\sigma\,dW_t.
\]
在短时间区间内，粒子会同时经历两种效应：
\begin{itemize}
    \item 由漂移 $f$ 引起的 $\Delta t$ 阶确定性位移；
    \item 由布朗项引起的 $\sqrt{\Delta t}$ 阶随机高斯位移。
\end{itemize}
最终的密度演化把确定性的 Liouville 输运项与上面推导出的扩散项结合起来：
\[
\boxed{
\partial_t p(x,t)
=
-\partial_x\bigl(f(x,t)p(x,t)\bigr)
+
\frac{\sigma^2}{2}\partial_{xx}p(x,t).
}
\]
这就是扩散系数 $\sigma$ 为常数时的一维\textbf{Fokker--Planck 方程}。
图~\ref{fig:continuous-liouville-fokker-planck} 对比了仅有漂移、仅有扩散以及二者结合时对应的密度演化。

\begin{figure}[H]
\centering
\begin{minipage}[t]{0.48\textwidth}
  \vspace{0pt}\centering
  \includegraphics[width=\linewidth]{figures/ch05/panels/fig5_3a_drift_only.pdf}\\[-0.2em]
  {\small $\displaystyle \partial_t p=-\partial_x(fp)$}
\end{minipage}\hfill
\begin{minipage}[t]{0.48\textwidth}
  \vspace{0pt}\centering
  \includegraphics[width=\linewidth]{figures/ch05/panels/fig5_3b_diffusion_only.pdf}\\[-0.2em]
  {\small $\displaystyle \partial_t p=\frac{\sigma^2}{2}\partial_{xx}p$}
\end{minipage}

\vspace{0.8em}

\begin{minipage}[t]{0.48\textwidth}
  \vspace{0pt}\centering
  \includegraphics[width=\linewidth]{figures/ch05/panels/fig5_3c_ou_density.pdf}\\[-0.2em]
  {\small $\displaystyle \partial_t p=-\partial_x(fp)+\frac{\sigma^2}{2}\partial_{xx}p$}
\end{minipage}\hfill
\begin{minipage}[t]{0.48\textwidth}
  \vspace{0pt}
  \caption{\textbf{漂移输运概率，扩散使概率铺展。}
  从同一个高斯密度出发，确定性收缩使均值朝原点移动，同时减小方差；纯 布朗扩散则保持均值不变而增大方差。Ornstein--Uhlenbeck 过程结合了两种效应：其均值回归零点，而方差趋近于平稳值 $\sigma^2/(2\alpha)$。图中曲线为精确的高斯密度解。}
  \label{fig:continuous-liouville-fokker-planck}
\end{minipage}
\end{figure}


\subsection{多维形式与解释}

在多维中，确定性部分变为
\[
-\nabla\cdot\bigl(p(x,t)v(x,t)\bigr),
\]
扩散项变为
\[
\frac{\sigma^2}{2}\Delta p(x,t),
\]
其中 Laplacian 算子为
\[
\Delta p
=
\nabla\cdot(\nabla p)
=
\sum_{i=1}^d \frac{\partial^2 p}{\partial x_i^2}.
\]
因此，对于常数各向同性扩散系数 $\sigma$，多维 Fokker--Planck 方程为
\[
\boxed{
\partial_t p(x,t)
=
-\nabla\cdot\bigl(p(x,t)v(x,t)\bigr)
+
\frac{\sigma^2}{2}\Delta p(x,t).
}
\]

这是连续时间概率建模中最重要的密度演化定律之一。
Liouville 方程告诉我们密度如何在确定性流下移动；Fokker--Planck 方程在此基础上加入随机扩散。
二者共同构成从基本连续时间运动过渡到完全连续扩散模型底层密度动力学的主要桥梁。

下一章将直接建立在这些基础之上。
关键的下一步是理解密度演化如何与对数密度的梯度相联系，以及这一视角如何自然导向基于得分的学习和完全连续时间扩散模型。

\section*{本章小结与参考文献}
\addcontentsline{toc}{section}{本章小结与参考文献}

本章建立了从离散更新过渡到连续时间生成动力学所需的微积分基础。
首先通过导数、微分方程、无穷小变化与一阶 Taylor 近似，引入连续运动的基本语言。

随后研究了三种核心运动模式。
确定性收缩展示了微分方程如何确定轨迹，先在一维中说明，再借助矩阵和矩阵指数推广到多维。
布朗运动通过高斯增量引入连续时间随机波动；Ornstein--Uhlenbeck 过程则把漂移和噪声统一在一个均值回归系统中。

接着，本章从粒子运动转向密度演化。
对于确定性流，我们推导了有限时间变量变换公式以及 Liouville 方程，并给出了沿轨迹的对数密度形式。
将其推广到多维后，引入了 Jacobian、行列式和散度，它们是描述局部体积变化与密度流动的关键工具。

最后加入布朗扩散，得到 Fokker--Planck 方程，它把确定性输运与随机扩散结合起来。
这构成了通往下一章的主要桥梁；下一章会把密度演化与对数密度梯度、得分学习和完全连续时间扩散模型联系起来。

\bigskip

本章内容建立在微积分、微分方程、随机过程以及连续时间概率论的标准思想之上。
关于机器学习背景和概率建模视角，可参考 Bishop~\cite{bishop2006prml} 与 Murphy~\cite{murphy2012ml}。
关于随机微分方程与布朗运动，Øksendal~\cite{oksendal2003sde} 是经典入门教材。
关于连续性方程、输运及相关偏微分方程视角，可参考 Evans~\cite{evans2010pde}。
更具体地讨论 Fokker--Planck 方程与随机扩散的经典参考文献则是 Risken~\cite{risken1996fokker}。

\chapter[基于得分的生成建模与连续时间扩散]{基于得分的生成建模\\与连续时间扩散}
\section*{概述}
\addcontentsline{toc}{section}{概述}

前面的 DDPM 章节把扩散模型介绍为一种离散潜变量构造：数据逐渐被扰动成噪声，而学习得到的反向过程再逐步重建结构。
这一视角十分强大，因为它把扩散模型与本书前面发展出的概率建模思想联系起来，包括潜变量、高斯转移、后验推断以及变分学习。

但扩散模型还有另一种解释方式，并且这种视角已经成为现代生成建模的核心之一。
我们不再只问“如何写出一条反向马尔可夫链”，而是提出一个更加局部且几何化的问题：如果一个样本当前位于 $x$，应该朝哪个方向移动，才能使它在我们希望生成的目标分布下变得更可能？
答案由\emph{得分函数}给出，即
$\nabla_x \log p(x)$，它指向对数密度局部增长最快的方向。

这种视角变化非常深刻。
概率密度描述概率质量集中在哪里，而得分则描述在这个密度定义的空间中应该如何移动。
它把一个分布转化为向量场。
一旦知道这个场，就可以用它把带噪样本引导回高概率区域。
从这个意义上说，基于得分的建模把生成建模从“显式计算密度”的问题，转化为“学习密度局部几何结构”的问题。

这也是基于得分的生成建模成为理解扩散模型重要方式的原因之一。
与通过一次直接变换生成数据的模型不同，现代扩散系统学习在多个扰动水平上逐步或连续地修正噪声。
在每一个噪声水平，模型都需要局部信息，告诉它应该怎样轻微推动当前带噪样本，使其更像数据。
得分恰好提供了这种信息。

上一章已经给出了理解这一视角所需的微积分工具。
我们研究了粒子在确定性和随机动力学下如何运动，以及整个概率密度如何通过 Liouville 方程与 Fokker--Planck 方程演化。
本章现在将利用这些工具，解释得分场、随机采样和密度演化如何共同构成基于得分的扩散模型。

本章的目的并不是取代 DDPM 视角，而是对它进行补充。
DDPM 说明扩散模型如何组织成一个离散潜变量模型；基于得分的视角则说明，去噪为什么也可以从几何和动力学角度理解：生成由沿一系列带噪分布学习得到的局部方向引导，把样本从随机性逐渐带回有结构的数据区域。

\section*{概念导图}
\addcontentsline{toc}{section}{概念导图}

本章逐步建立扩散模型的得分解释：从局部密度几何，到得分学习，再到多尺度去噪，最后进入连续时间反向扩散。

\begin{itemize}

    \item[\(\rightarrow\)] \textbf{首先从得分函数开始。}
    得分 $\nabla_x \log p(x)$ 被解释为一个局部概率方向场，它指向对数密度增长最快的区域。

    \item[\(\rightarrow\)] \textbf{随后通过 Langevin 动力学把得分与采样联系起来。}
    确定性的得分流把样本推向高密度区域，而布朗噪声防止塌缩并维持探索。
    Fokker--Planck 方程说明，为什么目标密度在这种“得分+噪声”动力学下是平稳分布。

    \item[\(\rightarrow\)] \textbf{接着讨论如何学习得分。}
    Fisher 散度为把学习到的得分 $s_\theta(x)$ 与真实数据得分匹配提供了自然目标。
    得分匹配通过分部积分改写这一目标，从而不再需要计算未知的数据密度。

    \item[\(\rightarrow\)] \textbf{随后通过去噪来稳定得分学习。}
    与其直接学习原始数据得分，不如加入高斯噪声，学习平滑后分布 $p_\sigma$ 的得分。
    这样，得分估计就转化为对去噪方向的回归问题。

    \item[\(\rightarrow\)] \textbf{再从单一噪声尺度扩展到多个噪声尺度。}
    大噪声水平揭示粗粒度全局结构，小噪声水平揭示精细局部细节。
    因此，多尺度得分学习是在不同平滑尺度上学习一族得分场。

    \item[\(\rightarrow\)] \textbf{把多尺度得分与扩散联系起来。}
    当一族带噪分布被组织为同一个前向加噪过程的边缘分布时，一族带噪分布的得分场就形成扩散模型。
    在连续时间中，这对应一条密度路径 $p_t(x)$，$0\leq t\leq T$。

    \item[\(\rightarrow\)] \textbf{利用得分反转扩散。}
    前向过程把数据推向噪声。
    反向时间动力学使用随时间变化的得分
    $\nabla_x \log p_t(x)$
    作为局部修正，把样本从噪声引导回数据。

    \item[\(\rightarrow\)] \textbf{最后重新连接 DDPM。}
    前面曾把 DDPM 作为离散潜变量模型引入；这里重新解释其噪声预测目标，说明它携带了中间带噪分布的类得分信息，从而揭示 DDPM 视角与得分视角其实是互补的。

\end{itemize}

\medskip
\noindent 本章基于得分建模的主要路径总结如下图所示。
\begin{center}
\includegraphics[width=0.95\textwidth]{figures/ch06/chapter6_score_at_a_glance_figure.pdf}
\end{center}

\medskip

\section{从得分函数到 Langevin 采样}

上一章解释了概率密度如何在确定性输运与随机扩散下演化。
我们看到，常微分方程会诱导 Liouville 型密度演化，而含布朗噪声的随机微分方程会导出 Fokker--Planck 方程。
这些结果很重要，因为它们把讨论从单个粒子的运动推进到整个分布的演化。

由此自然出现下一个问题：在所有可能的密度演化中，哪些真正适合生成建模？
在生成建模中，我们并不需要任意运动，而是希望得到一种与某个感兴趣的概率密度相关联的动力学，通常这个密度就是数据分布，或者它的某个平滑版本。
这正是得分函数发挥作用的地方。

得分提供一个局部概率方向场；当它与噪声结合时，就得到 Langevin 动力学——连接密度建模与采样最重要的桥梁之一。
本节会仔细建立这座桥梁：先把得分解释为概率漂移场，再说明如何通过 Langevin 动力学加入噪声，最后解释为什么得到的动力学通过平稳性与目标密度绑定在一起。

因此，本节还不会完整引入现代扩散模型。
它首先建立一个更简单、更基础的原则：如果目标密度的得分可得，那么可以把它用作随机动力学中的漂移场，而这一动力学的平稳分布正是目标分布本身。
这构成从“描述目标密度”到“构造能够从该密度生成样本的动力学”的主要概念桥梁之一。

\subsection{作为概率漂移场的得分函数}

设 $p(x)$ 是 $\mathbb{R}^d$ 上的概率密度。
其\emph{得分函数}定义为
\[
s(x)=\nabla_x \log p(x).
\]
如果密度还依赖时间，那么随时间变化的得分写为
\[
s(x,t)=\nabla_x \log p(x,t).
\]

乍看之下，对\emph{对数}密度求梯度为何如此重要似乎有些意外。
原因在于，得分以一种特别有用的形式编码了局部方向信息。
密度 $p(x)$ 告诉我们在 $x$ 附近集中了多少概率质量，而得分则告诉我们沿哪个方向移动会使对数密度增长得最快。
因此，它提供了一个指向更高概率区域的局部向量场。

在一维中，得分变成
\[
s(x)
=
\frac{d}{dx}\log p(x)
=
\frac{p'(x)}{p(x)}.
\]
这个表达已经很有启发性。
分子 $p'(x)$ 衡量密度的局部变化，而除以 $p(x)$ 会根据密度自身的局部尺度对这种变化进行归一化。
因此，得分衡量的是相对斜率，而不仅仅是绝对斜率。
这使它更适合作为一种概率运动场。

一个简单的高斯例子可以把这种解释具体化。
考虑一维高斯密度
\[
p(x)
=
\frac{1}{\sqrt{2\pi\sigma^2}}
\exp\left(
-\frac{(x-\mu)^2}{2\sigma^2}
\right).
\]
取对数得到
\[
\log p(x)
=
-\frac{(x-\mu)^2}{2\sigma^2}
-\frac{1}{2}\log(2\pi\sigma^2).
\]
关于 $x$ 求导，得到
\[
\frac{d}{dx}\log p(x)
=
-\frac{x-\mu}{\sigma^2}.
\]
因此得分为
\[
\boxed{
s(x)
=
-\frac{x-\mu}{\sigma^2}.
}
\]

这个高斯例子立刻揭示了得分的几何含义。
如果 $x>\mu$，得分为负，方向指回均值；如果 $x<\mu$，得分为正，同样指回均值。
$x$ 离均值越远，得分的绝对值越大。
因此，对于高斯密度，得分就像一个恢复力场，把粒子引向高密度区域。

在多元高斯情形
\[
p(x)=\mathcal{N}(\mu,\Sigma)
\]
中，得分变为
\[
\boxed{
\nabla_x \log p(x)
=
-\Sigma^{-1}(x-\mu).
}
\]
因此，得分仍然是指向高密度中心的向量场，但其形状由协方差结构决定。
方差较小的方向产生更强的恢复作用，方差较大的方向则更弱。
图~\ref{fig:score-field-geometry} 展示了同一思想如何从单个高斯推广到多峰分布：对于多峰密度，得分会形成随空间变化的向量场，其局部方向取决于周围的密度几何结构。

\begin{figure}[H]
\centering
\includegraphics[width=0.98\textwidth]{figures/ch06/figure6_1_score_field_geometry.pdf}
\caption{\textbf{得分把多模态密度转化为局部方向场。}
左图展示精确双分量高斯混合的密度等高线，右图展示对应的得分方向 $\nabla_x\log p(x)$。箭头指向对数密度增大的方向，并弯向局部影响更强的模态。为便于阅读，箭头长度经过归一化，因此图中强调的是方向而非幅度。}
\label{fig:score-field-geometry}
\end{figure}

这启发了一个自然想法。
如果定义确定性流
\[
\frac{dx}{dt}
=
\nabla_x \log p(x),
\]
那么粒子会沿着对数密度增长最快的方向移动。
由于对数函数单调，这也会使粒子朝密度更高的区域移动。
从这个意义上说，得分可以解释为一个\emph{概率漂移场}：它告诉我们，如果想把粒子引导到目标分布下更可能的区域，应该怎样移动。

这种视角非常有用，但也暴露出一个重要局限。
纯确定性的得分流会把粒子沿上坡方向推向众数，因此可能使轨迹塌缩到高度集中的高密度区域，而无法探索完整目标分布。
这种性质对于优化很有用，但对采样还不够。
为了得到真正的采样动力学，必须加入噪声。
这就直接引出了 Langevin 动力学。

\subsection{Langevin 动力学：得分加噪声}

Langevin 动力学结合两个成分：
\begin{enumerate}
    \item 由得分给出的漂移项，把粒子拉向高密度区域；
    \item 布朗噪声项，注入随机性并维持探索。
\end{enumerate}

在连续时间中，Langevin 随机微分方程写作
\[
\boxed{
dX_t
=
\nabla_x \log p^\ast(X_t)\,dt
+
\sqrt{2}\,dW_t,
}
\]
其中 $p^\ast(x)$ 表示目标密度。
这里使用 $p^\ast$ 是为了强调：这是我们希望从中采样的目标分布。

在离散时间中，可以启发式地写成更新式
\[
\boxed{
x_{k+1}
=
x_k
+
\epsilon\,\nabla_x \log p^\ast(x_k)
+
\sqrt{2\epsilon}\,\xi_k,
\qquad
\xi_k\sim\mathcal{N}(0,I),
}
\]
其中 $\epsilon>0$ 是较小步长。

这个更新有非常直接的解释。
当前状态 $x_k$ 首先经过得分项
\[
\epsilon\,\nabla_x \log p^\ast(x_k)
\]
移动，使粒子朝目标密度下更高概率的区域偏移。
随后加入高斯扰动
\[
\sqrt{2\epsilon}\,\xi_k.
\]
这个随机项阻止粒子简单地塌缩到某个众数，并让它继续探索状态空间。
在连续时间极限下，确定性得分漂移与布朗扰动结合成上面的 Langevin SDE。

这一结构与上一章介绍的形式相似。
Langevin 动力学是一个“漂移+噪声”系统，因此其密度按照 Fokker--Planck 方程演化。
对于
\[
dX_t
=
\nabla_x \log p^\ast(X_t)\,dt
+
\sqrt{2}\,dW_t,
\]
漂移场是
\[
v(x)=\nabla_x \log p^\ast(x),
\]
扩散系数为 $\sigma=\sqrt{2}$。
由于常数各向同性扩散下的 Fokker--Planck 方程为
\[
\partial_t p(x,t)
=
-\nabla\cdot\bigl(p(x,t)v(x,t)\bigr)
+
\frac{\sigma^2}{2}\Delta p(x,t),
\]
代入 $v(x)=\nabla_x\log p^\ast(x)$ 与 $\sigma^2=2$，得到
\[
\boxed{
\partial_t p(x,t)
=
-\nabla\cdot\bigl(p(x,t)\nabla_x \log p^\ast(x)\bigr)
+
\Delta p(x,t).
}
\]

该方程结合两种效应。
第一项
$-\nabla\cdot\bigl(p\nabla_x \log p^\ast\bigr)$
描述由得分场诱导的确定性输运；第二项 $\Delta p$ 描述布朗噪声导致的扩散。
漂移项倾向于把概率集中到目标的高密度区域，而扩散项倾向于把概率向外铺开。
重要的是，正是在这种精确组合下，目标密度会成为平稳分布。
这正是 Langevin 动力学能够作为采样机制使用的关键原因。

把它与前面介绍的 Ornstein--Uhlenbeck 过程比较也很有帮助。
OU 过程具有指向原点的线性漂移和高斯平稳分布。
Langevin 动力学推广了同一原则：不再使用固定的线性恢复力，而是把所选目标密度的得分作为恢复场。
从这个意义上说，OU 过程可以视为更一般得分采样机制的高斯特例。

为了显式看出这一点，假设目标密度为标准高斯
$p^\ast(x)=\mathcal{N}(0,I)$。
那么
\[
\nabla_x\log p^\ast(x)=-x,
\]
Langevin 动力学变为
\[
dX_t=-X_t\,dt+\sqrt{2}\,dW_t,
\]
这正是一个特定噪声缩放下的 Ornstein--Uhlenbeck 过程。
因此，当目标密度为高斯时，Langevin 方程恢复 OU 过程；而对于非高斯目标，它通过得分
$\nabla_x\log p^\ast(x)$ 对该原则进行推广。

\subsection{目标分布与平稳分布}

下面解释为什么 Langevin 动力学不只是一个形式化方程，而是真正有意义的采样机制。

假设某个密度 $p(x,t)$ 按照密度演化方程
\[
\partial_t p(x,t)
=
\mathcal{L}[p(x,t)]
\]
演化。
如果系统一旦具有密度 $p^\ast(x)$ 就会一直保持该密度，那么 $p^\ast(x)$ 就称为\emph{平稳分布}。
在 PDE 语言中，这意味着
\[
\mathcal{L}[p^\ast(x)]=0.
\]
等价地，
\[
\partial_t p(x,t)=0
\qquad
\text{当}
\qquad
p(x,t)=p^\ast(x).
\]

这个概念在生成建模中非常重要。
如果目标密度是某个动力学的平稳分布，那么该动力学就与目标分布相一致。
如果进一步地，从其他初始状态出发的轨迹或密度会随时间趋近这一平稳分布，那么该过程就成为从目标分布中采样的机制。
完整的收敛理论还需要额外假设，超出了本教程范围，但仅仅平稳性计算就已经揭示了核心思想。

现在验证 Langevin 动力学的平稳性。
回顾其 Fokker--Planck 方程
\[
\partial_t p(x,t)
=
-\nabla\cdot\bigl(p(x,t)\nabla_x \log p^\ast(x)\bigr)
+
\Delta p(x,t).
\]
令
\[
p(x,t)=p^\ast(x).
\]
右侧变为
\[
-\nabla\cdot\bigl(p^\ast(x)\nabla_x \log p^\ast(x)\bigr)
+
\Delta p^\ast(x).
\]
利用恒等式
\[
p^\ast(x)\nabla_x \log p^\ast(x)
=
\nabla_x p^\ast(x),
\]
因为
\[
\nabla_x \log p^\ast(x)
=
\frac{\nabla_x p^\ast(x)}{p^\ast(x)}.
\]
于是上式化简为
\[
-\nabla\cdot(\nabla_x p^\ast(x))
+
\Delta p^\ast(x).
\]
而根据定义，
\[
\nabla\cdot(\nabla_x p^\ast)
=
\Delta p^\ast.
\]
因此
\[
-\Delta p^\ast(x)+\Delta p^\ast(x)=0.
\]
所以确实有
\[
\boxed{
\partial_t p(x,t)=0
\qquad
\text{当}
\qquad
p(x,t)=p^\ast(x).
}
\]

因此，$p^\ast$ 是 Langevin 动力学
\[
\boxed{
dX_t
=
\nabla_x \log p^\ast(X_t)\,dt
+
\sqrt{2}\,dW_t
}
\]
的一个平稳分布。

这是基于得分的生成建模背后最核心的概念结果之一。
它说明，如果把目标密度的得分作为漂移场，并配合适当的布朗噪声，那么得到的密度演化会保持该目标密度不变。
因此，得分不只是密度的一个有趣导数；它恰好就是使目标密度在自然随机动力学下保持平稳的漂移场。

同时，也必须准确理解这里真正证明了什么。
上面的论证建立的是\emph{平稳性}，而不是“任意初始密度都必然收敛到 $p^\ast$”的完整证明。
后一个问题更深入，涉及随机过程的长期行为。
尽管如此，平稳性结果已经填补了一个重要的概念缺口：它解释了为什么得分漂移配合布朗噪声能够为从感兴趣的分布中采样提供有原则的路径。

在适当条件下，如果让 Langevin 过程演化足够长时间，其分布可能逐渐接近平稳分布 $p^\ast$。
从密度层面看，这意味着概率质量逐渐形成目标分布的形状；从粒子层面看，则意味着不断用 Langevin 动力学更新样本，最终可得到统计行为类似于目标分布抽样的粒子。

因此，密度层面与粒子层面的视角彼此一致。
Fokker--Planck 方程解释分布如何演化，Langevin 更新解释单个样本如何移动，而得分把二者连接起来：局部上，它作为概率漂移场；全局上，它是使目标密度成为平稳分布的关键成分。

\section{学习得分：Fisher 散度与得分匹配}

上一节说明了得分函数为什么对采样重要。
如果目标密度 $p^\ast(x)$ 的得分可得，那么可以把它作为 Langevin 动力学
\[
dX_t
=
\nabla_x \log p^\ast(X_t)\,dt
+
\sqrt{2}\,dW_t
\]
中的漂移场，而该动力学的平稳分布就是 $p^\ast$。
这在概率密度与随机采样动力学之间建立了有原则的联系。

然而在生成建模中，目标密度通常并不存在可直接使用的闭式形式。
我们只有来自某个未知数据生成分布的样本，却不知道该密度本身。
记这个未知分布为 $p_{\text{data}}(x)$。
因此，核心学习问题变成近似其得分
\[
s_{\text{data}}(x)
=
\nabla_x \log p_{\text{data}}(x).
\]
如果能从数据中学到这个向量场，就可以把它用于 Langevin 型采样动力学。

我们用参数化函数 $s_\theta(x)$ 建模得分，通常由参数为 $\theta$ 的神经网络实现。
一个自然目标是在数据分布下最小化学习得分与真实数据得分之间的差异。
这就得到\emph{Fisher 散度}：
\[
\mathcal{J}(\theta)
=
\mathbb{E}_{x\sim p_{\text{data}}}
\left[
\left\|
s_\theta(x)-s_{\text{data}}(x)
\right\|_2^2
\right].
\]
等价地，学习问题是 $\min_\theta \mathcal{J}(\theta)$。
这里，散度定义在两个向量场之间，而优化则针对参数 $\theta$ 进行。
如果 $\mathcal{J}(\theta)=0$，则
\[
s_\theta(x)=s_{\text{data}}(x)
\]
在 $p_{\text{data}}$ 下几乎处处成立。
此时，学习到的向量场就匹配了数据分布的局部几何结构。

然而立刻会遇到一个困难。
该目标似乎需要真实得分 $s_{\text{data}}(x)$，而它依赖未知密度 $p_{\text{data}}(x)$。
得分匹配（score matching）~\cite{hyvarinen2005estimation} 通过改写目标解决了这一问题，使未知真实得分不再显式出现。

\subsection{展开 Fisher 散度}

先展开平方范数：
\[
\begin{aligned}
\mathcal{J}(\theta)
&=
\mathbb{E}_{p_{\text{data}}}
\left[
\bigl(s_\theta(x)-s_{\text{data}}(x)\bigr)^\top
\bigl(s_\theta(x)-s_{\text{data}}(x)\bigr)
\right] \\[4pt]
&=
\mathbb{E}_{p_{\text{data}}}
\left[
\|s_\theta(x)\|_2^2
-
2s_\theta(x)^\top s_{\text{data}}(x)
+
\|s_{\text{data}}(x)\|_2^2
\right].
\end{aligned}
\]
最后一项
\[
\mathbb{E}_{p_{\text{data}}}
\left[
\|s_{\text{data}}(x)\|_2^2
\right]
\]
与 $\theta$ 无关，因此不会影响优化。
所以，最小化 $\mathcal{J}(\theta)$ 等价于最小化
\[
\tilde{\mathcal{J}}(\theta)
=
\mathbb{E}_{p_{\text{data}}}
\left[
\|s_\theta(x)\|_2^2
\right]
-
2
\mathbb{E}_{p_{\text{data}}}
\left[
s_\theta(x)^\top s_{\text{data}}(x)
\right].
\]

第一项可以直接从样本计算，因为只需要在数据点上计算模型 $s_\theta(x)$。
困难在交叉项
\[
\mathbb{E}_{p_{\text{data}}}
\left[
s_\theta(x)^\top s_{\text{data}}(x)
\right],
\]
因为它仍然包含未知真实得分。

得分匹配的关键思想是利用分部积分改写这一交叉项。
我们先在一维中仔细推导，因为此时计算最透明，再把结果推广到多维。

\subsection{一维情形}

为清楚展示核心思想，先考虑一维场景，此时真实数据得分为
\[
s_{\text{data}}(x)
=
\frac{d}{dx}\log p_{\text{data}}(x)
=
\frac{1}{p_{\text{data}}(x)}
\frac{d}{dx}p_{\text{data}}(x).
\]
代入交叉项，得到
\[
\begin{aligned}
\mathbb{E}_{p_{\text{data}}}
\left[
s_\theta(x)s_{\text{data}}(x)
\right]
&=
\int
s_\theta(x)
\frac{d}{dx}\log p_{\text{data}}(x)
p_{\text{data}}(x)
\,dx \\[4pt]
&=
\int
s_\theta(x)
\frac{1}{p_{\text{data}}(x)}
\frac{d}{dx}p_{\text{data}}(x)
p_{\text{data}}(x)
\,dx \\[4pt]
&=
\int
s_\theta(x)
\frac{d}{dx}p_{\text{data}}(x)
\,dx.
\end{aligned}
\]

此时，导数作用在未知密度 $p_{\text{data}}(x)$ 上。
我们的目标是通过分部积分把导数转移到模型 $s_\theta(x)$ 上。

\begin{tcolorbox}[
  colback=blue!4,
  colframe=blue!50!black,
  title={一维分部积分},
  boxrule=0.6pt,
  arc=2mm,
  left=2mm,
  right=2mm,
  top=1mm,
  bottom=1mm
]
乘积法则给出
\[
\frac{d}{dx}\bigl(u(x)v(x)\bigr)
=
u'(x)v(x)+u(x)v'(x).
\]
两边积分，得到
\[
\int u(x)v'(x)\,dx
=
\bigl[u(x)v(x)\bigr]_{-\infty}^{\infty}
-
\int u'(x)v(x)\,dx.
\]
这个恒等式以产生一个边界项为代价，把导数从 $v$ 转移到了 $u$。
\end{tcolorbox}

现在取
\[
u(x)=s_\theta(x),
\qquad
v'(x)=\frac{d}{dx}p_{\text{data}}(x),
\qquad
v(x)=p_{\text{data}}(x).
\]
于是
\[
\int
s_\theta(x)
\frac{d}{dx}p_{\text{data}}(x)
\,dx
=
\bigl[
s_\theta(x)p_{\text{data}}(x)
\bigr]_{-\infty}^{\infty}
-
\int
\frac{d}{dx}s_\theta(x)
p_{\text{data}}(x)
\,dx.
\]

在标准假设下，$p_{\text{data}}(x)$ 在 $|x|\to\infty$ 时衰减得足够快，同时 $s_\theta(x)$ 不会增长得过快，因此乘积 $s_\theta(x)p_{\text{data}}(x)$ 在边界处消失。
所以边界项为零，从而得到
\[
\int
s_\theta(x)
\frac{d}{dx}p_{\text{data}}(x)
\,dx
=
-
\int
\frac{d}{dx}s_\theta(x)
p_{\text{data}}(x)
\,dx.
\]
把右侧识别为 $p_{\text{data}}$ 下的期望，可得
\[
\mathbb{E}_{p_{\text{data}}}
\left[
s_\theta(x)s_{\text{data}}(x)
\right]
=
-
\mathbb{E}_{p_{\text{data}}}
\left[
\frac{d}{dx}s_\theta(x)
\right].
\]

这样，导数就从未知密度 $p_{\text{data}}(x)$ 转移到了模型 $s_\theta(x)$ 上，不再需要显式知道真实数据得分。
把这一结果代入约化后的 Fisher 散度，得到
\[
\begin{aligned}
\tilde{\mathcal{J}}(\theta)
&=
\mathbb{E}_{p_{\text{data}}}
\left[
s_\theta(x)^2
\right]
-
2
\mathbb{E}_{p_{\text{data}}}
\left[
s_\theta(x)s_{\text{data}}(x)
\right] \\[4pt]
&=
\mathbb{E}_{p_{\text{data}}}
\left[
s_\theta(x)^2
\right]
+
2
\mathbb{E}_{p_{\text{data}}}
\left[
\frac{d}{dx}s_\theta(x)
\right].
\end{aligned}
\]

由于把目标函数乘以正的常数不会改变最优解，一维得分匹配目标通常写为
\[
\boxed{
\mathcal{L}_{\text{SM}}(\theta)
=
\mathbb{E}_{p_{\text{data}}}
\left[
\frac{1}{2}s_\theta(x)^2
+
\frac{d}{dx}s_\theta(x)
\right].
}
\]

\subsection{推广到多维}

现在把同样的论证推广到一般情形 $x\in\mathbb{R}^d$。
真实数据得分为
\[
s_{\text{data}}(x)
=
\nabla_x\log p_{\text{data}}(x),
\]
Fisher 散度中的交叉项变为
\[
\begin{aligned}
\mathbb{E}_{p_{\text{data}}}
\left[
s_\theta(x)^\top s_{\text{data}}(x)
\right]
&=
\int
s_\theta(x)^\top
\nabla_x\log p_{\text{data}}(x)
p_{\text{data}}(x)
\,dx \\[4pt]
&=
\int
s_\theta(x)^\top
\nabla_x p_{\text{data}}(x)
\,dx.
\end{aligned}
\]

在一维中由普通导数承担的角色，现在由散度承担。
对于向量场
\[
v(x)=
\begin{bmatrix}
v_1(x)\\
\vdots\\
v_d(x)
\end{bmatrix},
\]
散度定义为
\[
\nabla_x\cdot v(x)
=
\sum_{i=1}^d
\frac{\partial v_i(x)}{\partial x_i}.
\]
正如上一章介绍的，散度衡量向量场的局部净流出。
这里之所以出现散度，是因为模型得分 $s_\theta(x)$ 本身就是一个向量场。

多维分部积分的对应形式来自散度定理。
相关乘积法则为
\[
\nabla_x\cdot
\bigl(p_{\text{data}}(x)s_\theta(x)\bigr)
=
s_\theta(x)^\top\nabla_x p_{\text{data}}(x)
+
p_{\text{data}}(x)\nabla_x\cdot s_\theta(x).
\]
重新整理得
\[
s_\theta(x)^\top\nabla_x p_{\text{data}}(x)
=
\nabla_x\cdot
\bigl(p_{\text{data}}(x)s_\theta(x)\bigr)
-
p_{\text{data}}(x)\nabla_x\cdot s_\theta(x).
\]
在 $\mathbb{R}^d$ 上积分，并假设满足合适的边界条件，使散度项没有边界通量贡献，可得
\[
\int
s_\theta(x)^\top\nabla_x p_{\text{data}}(x)
\,dx
=
-
\int
p_{\text{data}}(x)
\nabla_x\cdot s_\theta(x)
\,dx.
\]
因此
\[
\boxed{
\mathbb{E}_{p_{\text{data}}}
\left[
s_\theta(x)^\top s_{\text{data}}(x)
\right]
=
-
\mathbb{E}_{p_{\text{data}}}
\left[
\nabla_x\cdot s_\theta(x)
\right].
}
\]

这就是一维结果
\[
\mathbb{E}_{p_{\text{data}}}
\left[
s_\theta(x)s_{\text{data}}(x)
\right]
=
-
\mathbb{E}_{p_{\text{data}}}
\left[
\frac{d}{dx}s_\theta(x)
\right]
\]
在多维中的直接对应。
唯一的结构变化是：标量得分的普通导数被向量得分场的散度所替代。

把交叉项恒等式代入约化后的 Fisher 散度，得到
\[
\tilde{\mathcal{J}}(\theta)
=
\mathbb{E}_{p_{\text{data}}}
\left[
\|s_\theta(x)\|_2^2
+
2\nabla_x\cdot s_\theta(x)
\right].
\]
等价地，乘以 $1/2$ 后，一般的得分匹配目标为
\[
\boxed{
\mathcal{L}_{\text{SM}}(\theta)
=
\mathbb{E}_{p_{\text{data}}}
\left[
\frac{1}{2}
\|s_\theta(x)\|_2^2
+
\nabla_x\cdot s_\theta(x)
\right].
}
\]

这就是经典得分匹配目标。
它最初来自学习得分与真实数据得分之间的 Fisher 散度，但最终表达式不再显式包含 $s_{\text{data}}(x)$。
一旦 $s_\theta(x)$ 能够逼近 $s_{\text{data}}(x)$，就可以把它代入前面介绍的 Langevin 动力学。
在适当条件下，这种学习得到的动力学可以用于生成分布近似数据分布的样本。

\section{稳定学习：通过扰动实现去噪得分匹配}

上一节推导了得分匹配目标，它允许我们在不知道数据密度显式形式的情况下，从数据中学习得分函数 $s_\theta(x)$。
虽然理论上非常优雅，但直接得分匹配在高维生成建模中可能难以应用。

核心困难来自真实数据分布的几何结构。
经验数据往往集中在嵌入高维空间的低维流形上或其附近。
因此，数据密度可能高度集中、不规则，或者在数据支撑集之外表现很差。
直接学习这样的分布得分可能非常不稳定，因为在数据覆盖很少或完全没有覆盖的区域，诱导出的向量场可能行为异常。

一个自然解决办法是引入噪声。
高斯扰动会把经验分布或高度集中的数据分布变成一个平滑后的分布，其得分场通常表现更良好。
我们不再直接学习原始数据分布的得分，而是学习其带噪版本的得分。
这就是\emph{去噪得分匹配}（denoising score matching）~\cite{vincent2011connection} 的基础思想。

\subsection{高斯扰动与平滑得分}

首先用加性高斯噪声扰动数据。
给定干净样本
\[
x\sim p_{\text{data}}(x),
\]
构造带噪观测
\[
\tilde{x}
=
x+\varepsilon,
\qquad
\varepsilon\sim\mathcal{N}(0,\sigma^2 I).
\]
等价地，给定干净样本时，扰动后样本的条件分布为
\[
p(\tilde{x}\mid x)
=
\mathcal{N}(\tilde{x};x,\sigma^2 I).
\]

$\tilde{x}$ 的边缘分布通过对干净变量 $x$ 积分得到：
\[
p_\sigma(\tilde{x})
=
\int
p_{\text{data}}(x)
p(\tilde{x}\mid x)
\,dx.
\]
这就是原始数据分布与高斯核的卷积。
其作用是平滑：概率质量被扩散到邻近区域，尖锐结构得到软化，而原本密度几乎为零的区域在扰动后分布下也可能获得非零概率。

现在不再学习原始数据分布的得分，而是学习这个平滑密度的得分：
\[
s_\sigma(\tilde{x})
=
\nabla_{\tilde{x}}\log p_\sigma(\tilde{x}).
\]

下面考察这一带噪分布的得分的结构。
根据对数梯度恒等式，
\[
\nabla_{\tilde{x}}\log p_\sigma(\tilde{x})
=
\frac{1}{p_\sigma(\tilde{x})}
\nabla_{\tilde{x}}p_\sigma(\tilde{x}).
\]
平滑密度的梯度为
\[
\begin{aligned}
\nabla_{\tilde{x}}p_\sigma(\tilde{x})
&=
\nabla_{\tilde{x}}
\int
p_{\text{data}}(x)
p(\tilde{x}\mid x)
\,dx \\[4pt]
&=
\int
p_{\text{data}}(x)
\nabla_{\tilde{x}}p(\tilde{x}\mid x)
\,dx,
\end{aligned}
\]
其中在标准正则条件下可以交换求导和积分。
因此
\[
s_\sigma(\tilde{x})
=
\frac{
\int
p_{\text{data}}(x)
\nabla_{\tilde{x}}p(\tilde{x}\mid x)
\,dx
}{
\int
p_{\text{data}}(x)
p(\tilde{x}\mid x)
\,dx
}.
\]

这个表达具有一个很有用的概率解释。
由贝叶斯公式，给定扰动观测时干净变量的后验为
\[
p(x\mid\tilde{x})
=
\frac{
p_{\text{data}}(x)p(\tilde{x}\mid x)
}{
p_\sigma(\tilde{x})
}.
\]
利用这个后验，带噪分布的得分可以写成
\[
s_\sigma(\tilde{x})
=
\int
p(x\mid\tilde{x})
\nabla_{\tilde{x}}\log p(\tilde{x}\mid x)
\,dx,
\]
等价地，
\[
s_\sigma(\tilde{x})
=
\mathbb{E}_{x\sim p(x\mid\tilde{x})}
\left[
\nabla_{\tilde{x}}\log p(\tilde{x}\mid x)
\right].
\]

对于高斯扰动，
\[
p(\tilde{x}\mid x)
=
\mathcal{N}(\tilde{x};x,\sigma^2 I).
\]
忽略与 $\tilde{x}$ 无关的常数，其对数为
\[
\log p(\tilde{x}\mid x)
=
-\frac{1}{2\sigma^2}
\|\tilde{x}-x\|_2^2
+\text{常数}.
\]
因此
\[
\nabla_{\tilde{x}}\log p(\tilde{x}\mid x)
=
-\frac{1}{\sigma^2}(\tilde{x}-x)
=
\frac{1}{\sigma^2}(x-\tilde{x}).
\]
代入条件期望可得
\[
\boxed{
s_\sigma(\tilde{x})
=
\frac{1}{\sigma^2}
\mathbb{E}_{x\sim p(x\mid\tilde{x})}
\left[
x-\tilde{x}
\right].
}
\]

这个恒等式给出了平滑得分的直观解释。
在扰动点 $\tilde{x}$ 处，得分指向可能生成该点的干净样本的条件平均位置。
换句话说，带噪分布的得分提供一个局部去噪方向：它告诉我们平均而言应如何移动 $\tilde{x}$，才能重新靠近数据分布。

\subsection{为何难以直接计算带噪分布的得分}

上面的恒等式在概念上非常有用，但还没有给出可以直接计算的训练目标。
前向扰动过程很简单：我们可以采样
\[
x\sim p_{\text{data}}(x),
\qquad
\tilde{x}=x+\varepsilon,
\qquad
\varepsilon\sim\mathcal{N}(0,\sigma^2 I).
\]
因此可以很容易生成 $(x,\tilde{x})$ 样本对。

问题在于，带噪分布的得分包含关于后验分布 $p(x\mid\tilde{x})$ 的期望：
\[
s_\sigma(\tilde{x})
=
\frac{1}{\sigma^2}
\mathbb{E}_{x\sim p(x\mid\tilde{x})}
\left[
x-\tilde{x}
\right].
\]
根据贝叶斯公式，
\[
p(x\mid\tilde{x})
=
\frac{
p(\tilde{x}\mid x)p_{\text{data}}(x)
}{
p_\sigma(\tilde{x})
},
\]
其中
\[
p_\sigma(\tilde{x})
=
\int
p_{\text{data}}(x)p(\tilde{x}\mid x)\,dx.
\]
虽然似然 $p(\tilde{x}\mid x)$ 已知，但数据分布 $p_{\text{data}}(x)$ 没有闭式形式，我们只有它的样本。
因此，边缘密度 $p_\sigma(\tilde{x})$ 和后验 $p(x\mid\tilde{x})$ 都无法直接获得。

这阻止我们直接最小化带噪分布下的 Fisher 散度
\[
\mathbb{E}_{\tilde{x}\sim p_\sigma}
\left[
\left\|
s_\theta(\tilde{x})
-
s_\sigma(\tilde{x})
\right\|_2^2
\right],
\]
因为计算该目标需要显式访问真实带噪分布的得分 $s_\sigma(\tilde{x})$。
因此，尽管前向扰动过程可处理，对应得分仍不能以闭式直接得到。

所以我们再次面对一个熟悉局面：得分估计问题定义清楚，但目标无法直接计算。
关键问题在于，能否利用已知扰动过程来避开对后验的显式计算？

\subsection{去噪得分匹配作为回归}

关键观察是，扰动过程定义了 $(x,\tilde{x})$ 上一个可处理的联合分布：
\[
p_{\text{data}}(x)p(\tilde{x}\mid x).
\]
我们也许不知道边缘密度 $p_\sigma(\tilde{x})$ 或后验 $p(x\mid\tilde{x})$ 的显式形式，但可以精确地从这个联合分布采样：先抽取一个干净数据点 $x$，再加入高斯噪声得到 $\tilde{x}$。

去噪得分匹配利用一个基本的最小二乘回归性质来使用这一联合采样过程。

\begin{tcolorbox}[
  colback=blue!4,
  colframe=blue!50!black,
  title={最小二乘回归恒等式},
  boxrule=0.6pt,
  arc=2mm,
  left=2mm,
  right=2mm,
  top=1mm,
  bottom=1mm
]
对于随机变量 $(U,Y)$，使
\[
\mathbb{E}
\left[
\|f(U)-Y\|_2^2
\right]
\]
最小的函数为
\[
f^\star(U)=\mathbb{E}[Y\mid U].
\]
因此，在平方误差损失下，最佳预测就是给定输入后目标变量的条件均值。
\end{tcolorbox}

现在把这一恒等式应用于去噪场景。
令输入变量为扰动观测
\[
U=\tilde{x},
\]
回归目标为缩放后的残差
\[
Y=
\frac{1}{\sigma^2}(x-\tilde{x}).
\]
考虑目标
\[
\boxed{
\mathcal{L}_{\text{DSM}}(\theta)
=
\mathbb{E}_{x\sim p_{\text{data}},\,
\tilde{x}\sim p(\tilde{x}\mid x)}
\left[
\left\|
s_\theta(\tilde{x})
-
\frac{1}{\sigma^2}(x-\tilde{x})
\right\|_2^2
\right].
}
\]

根据最小二乘回归恒等式，最优函数满足
\[
s_\theta^\star(\tilde{x})
=
\mathbb{E}
\left[
\frac{1}{\sigma^2}(x-\tilde{x})
\;\middle|\;
\tilde{x}
\right].
\]
而这个条件期望恰好就是前面推导出的带噪分布的得分：
\[
s_\sigma(\tilde{x})
=
\frac{1}{\sigma^2}
\mathbb{E}
\left[
x-\tilde{x}
\;\middle|\;
\tilde{x}
\right].
\]
因此
\[
\boxed{
s_\theta^\star(\tilde{x})=s_\sigma(\tilde{x}).
}
\]

这就是去噪得分匹配的核心洞见。
单个训练样本对提供可处理的条件目标，而平方误差回归会在给定 $\tilde{x}$ 的条件下对这些目标求平均，从而恢复平滑密度的得分。
图~\ref{fig:score-dsm-regression} 将这种平均几何关系直观展示出来。

\begin{figure}[H]
\centering
\includegraphics[width=0.80\textwidth]{figures/ch06/figure6_3_dsm_regression_geometry.pdf}
\caption{\textbf{把去噪得分匹配理解为最小二乘回归。}
在这个有限示例中，细箭头表示各个可能干净点对应的可处理条件目标 $(x-\tilde{x})/\sigma^2$，其视觉强度反映相应的后验责任度。按后验加权后的平均值由粗箭头表示，终点位于 $\mathbb{E}[x\mid\tilde{x}]$，并等于边缘得分 $\nabla_{\tilde{x}}\log p_\sigma(\tilde{x})$。因此，平方误差回归能够从单个去噪目标恢复带噪分布的得分。}
\label{fig:score-dsm-regression}
\end{figure}

等价地，由于
\[
\tilde{x}=x+\varepsilon,
\]
有
\[
x-\tilde{x}=-\varepsilon.
\]
所以去噪得分匹配目标也可以写成
\[
\mathcal{L}_{\text{DSM}}(\theta)
=
\mathbb{E}_{x\sim p_{\text{data}},\,
\varepsilon\sim\mathcal{N}(0,\sigma^2 I)}
\left[
\left\|
s_\theta(x+\varepsilon)
+
\frac{\varepsilon}{\sigma^2}
\right\|_2^2
\right].
\]
这个形式使学习目标尤其具体：模型被训练去预测负的噪声方向，并乘以 $1/\sigma^2$ 进行缩放。

因此，去噪得分匹配把一个难以直接处理的得分估计问题转化为完全由已知扰动过程定义的监督回归问题。
目标函数只依赖干净数据样本和高斯噪声，不需要计算 $p_{\text{data}}(x)$、$p_\sigma(\tilde{x})$ 或后验 $p(x\mid\tilde{x})$。

需要注意，对于固定的 $\sigma$，模型学习的是扰动后密度 $p_\sigma$ 的得分，而不是直接学习 $p_{\text{data}}$ 的得分。
因此，噪声水平控制了模型正在学习其几何结构的那个分布。
然而，只学习单一噪声水平的得分又引出进一步问题：应该如何选择噪声尺度 $\sigma$？

\section{从多尺度得分学习到连续时间扩散}

到目前为止，本章自底向上建立了基于得分的生成建模。
我们先把得分解释为概率漂移场，再说明 Langevin 动力学如何利用得分进行采样，并进一步推导得分匹配和去噪得分匹配，使得分场能够从数据中学习。

现在进行最后一个概念性过渡：从单一噪声水平的得分学习，过渡到沿一整条带噪分布路径学习得分。这种多尺度视角已经成为现代基于得分生成建模的核心~\cite{song2019generative,song2021scorebased}。
也正是在这里，得分学习开始与扩散模型直接连接起来。

这种联系需要谨慎理解。
在前面的 DDPM 章节中，我们从离散潜变量视角引入扩散模型：前向加噪马尔可夫链、学习得到的反向链，以及变分训练目标。
这里则从另一个方向接近同一大框架。
从得分学习出发，我们询问一族带噪分布的得分场如何把样本从噪声移动回数据。
这会先导出多尺度去噪得分匹配，再导出退火 Langevin 动力学，最后进入连续时间扩散过程。

因此，本节不会重复 DDPM 的推导，而是展示得分视角如何提供一条通往扩散的连续时间路径：
\[
\begin{aligned}
\text{单一带噪分布的得分}
&\;\longrightarrow\;
\text{多个带噪分布的得分}
\\
&\;\longrightarrow\;
\text{退火反向采样}
\;\longrightarrow\;
\text{连续时间得分引导扩散}.
\end{aligned}
\]

\subsection{多尺度去噪得分匹配}

上一节已经说明，对于固定噪声水平 $\sigma$，去噪得分匹配学习的是平滑数据分布 $p_\sigma$ 的得分。
我们通过
\[
\tilde{x}=x+\varepsilon,
\qquad
\varepsilon\sim\mathcal{N}(0,\sigma^2I)
\]
扰动数据，并训练
\[
s_\theta(\tilde{x},\sigma)
\approx
\nabla_{\tilde{x}}\log p_\sigma(\tilde{x}).
\]
对应的去噪目标为
\[
\frac{1}{\sigma^2}(x-\tilde{x})
=
-\frac{\varepsilon}{\sigma^2}.
\]

然而，单一噪声水平只提供数据分布的一个平滑视图。
如果 $\sigma$ 很大，密度非常平滑，得分场能稳定地提供指向整体数据区域的全局引导，但精细细节会被模糊掉。
如果 $\sigma$ 很小，密度更接近原始数据分布，得分场能够捕捉更尖锐的局部细节，但在远离数据支撑集的区域又可能难以可靠估计。
因此，单一噪声水平无法同时提供稳健的全局引导和精确的局部细化。

图~\ref{fig:score-multiscale-fields} 展示了同一个多峰密度在三个高斯平滑尺度下的这一效应：大噪声揭示粗粒度全局几何，小噪声则恢复更尖锐的局部结构。

\begin{figure}[H]
\centering
\begin{minipage}[t]{0.43\textwidth}
  \vspace{0pt}\centering
  \includegraphics[width=\linewidth]{figures/ch06/panels/fig6_4_scale_1.pdf}\\[-0.2em]
  {\small 大噪声：粗尺度几何}
\end{minipage}\hfill
\begin{minipage}[t]{0.43\textwidth}
  \vspace{0pt}\centering
  \includegraphics[width=\linewidth]{figures/ch06/panels/fig6_4_scale_2.pdf}\\[-0.2em]
  {\small 中间尺度}
\end{minipage}

\vspace{0.8em}

\begin{minipage}[t]{0.43\textwidth}
  \vspace{0pt}\centering
  \includegraphics[width=\linewidth]{figures/ch06/panels/fig6_4_scale_3.pdf}\\[-0.2em]
  {\small 小噪声：局部细节}
\end{minipage}\hfill
\begin{minipage}[t]{0.43\textwidth}
  \vspace{0pt}
  \caption{\textbf{不同噪声水平对应不同的得分场。}
  同一个双分量分布在三个尺度上用高斯噪声进行平滑。大噪声揭示粗尺度的全局几何；随着噪声减小，分离的峰和更尖锐的局部结构逐渐显现。因此，多尺度得分模型以噪声水平为条件，而不是让所有尺度共享同一个得分场。为便于阅读，箭头长度经过归一化。}
  \label{fig:score-multiscale-fields}
\end{minipage}
\end{figure}


这说明噪声尺度不应只选择一次。
相反，应在一系列噪声水平上学习一族得分场。

令
\[
\sigma_1>\sigma_2>\cdots>\sigma_K>0
\]
为递减的噪声尺度序列。
对每个 $\sigma_k$，定义扰动后密度
\[
p_{\sigma_k}(\tilde{x})
=
\int
p_{\text{data}}(x)
\mathcal{N}(\tilde{x};x,\sigma_k^2I)
\,dx.
\]
我们不再学习一个得分场，而是训练一个噪声条件模型
\[
s_\theta(\tilde{x},\sigma_k)
\approx
\nabla_{\tilde{x}}\log p_{\sigma_k}(\tilde{x}).
\]

去噪得分匹配目标可以直接从单尺度情形扩展：
\[
\boxed{
\mathcal{L}_{\text{MS-DSM}}(\theta)
=
\mathbb{E}_{k}
\mathbb{E}_{x\sim p_{\text{data}},\,
\tilde{x}\sim\mathcal{N}(x,\sigma_k^2I)}
\left[
\left\|
s_\theta(\tilde{x},\sigma_k)
-
\frac{1}{\sigma_k^2}(x-\tilde{x})
\right\|_2^2
\right].
}
\]
等价地，写 $\tilde{x}=x+\varepsilon$，其中
$\varepsilon\sim\mathcal{N}(0,\sigma_k^2I)$，得到
\[
\mathcal{L}_{\text{MS-DSM}}(\theta)
=
\mathbb{E}_{k}
\mathbb{E}_{x\sim p_{\text{data}},\,
\varepsilon\sim\mathcal{N}(0,\sigma_k^2I)}
\left[
\left\|
s_\theta(x+\varepsilon,\sigma_k)
+
\frac{\varepsilon}{\sigma_k^2}
\right\|_2^2
\right].
\]

对于每一个固定噪声水平，前面同样的回归论证都成立。
最优预测器会恢复相应平滑密度的得分：
\[
s_\theta^\star(\tilde{x},\sigma_k)
=
\nabla_{\tilde{x}}\log p_{\sigma_k}(\tilde{x}).
\]
区别在于，现在模型学习的是一族得分场，而不是单个得分场。
因此，我们不再只对一个密度建模，而是在从高度加噪、简单的分布到轻度加噪、细节丰富的分布之间，对一整套层次化平滑密度建模。

\subsection{退火 Langevin 动力学}

一旦在多个噪声水平上学到得分，就可以用于采样。
这就得到\emph{退火 Langevin 动力学}，它按照从大噪声到小噪声的顺序执行 Langevin 更新。

假设
\[
s_\theta(x,\sigma_k)
\approx
\nabla_x\log p_{\sigma_k}(x),
\qquad
k=1,\dots,K,
\]
并且
\[
\sigma_1>\sigma_2>\cdots>\sigma_K.
\]
我们从一个宽的噪声分布初始化，例如
\[
x^{(0)}\sim\mathcal{N}(0,\sigma_1^2I),
\]
然后对每个尺度 $\sigma_k$ 执行若干次更新：
\[
\boxed{
x
\leftarrow
x
+
\epsilon_k s_\theta(x,\sigma_k)
+
\sqrt{2\epsilon_k}\,\xi,
\qquad
\xi\sim\mathcal{N}(0,I).
}
\]
在尺度 $\sigma_k$ 上迭代若干步后，再移动到更小的下一尺度 $\sigma_{k+1}$。

其逻辑很简单。
在大噪声水平下，得分提供指向整体数据区域的粗略引导；在更小噪声水平下，得分提供更尖锐的局部修正。
因此，采样过程表现为逐步细化：
\[
\text{随机噪声}
\;\longrightarrow\;
\text{粗略数据区域}
\;\longrightarrow\;
\text{局部结构}
\;\longrightarrow\;
\text{精细细节}.
\]

退火 Langevin 动力学把多尺度得分学习转化成明确的生成过程。
它也自然启发了下一步：如果噪声水平的数量越来越多，相邻水平之间的间隔越来越小，那么这种离散细化过程就开始类似连续时间扩散过程。

退火 Langevin 动力学已经是一种生成过程，但还不是 DDPM 意义上的扩散模型。
它使用一组带噪分布 $p_{\sigma_k}$，并在每个尺度上执行 Langevin 细化。
而扩散模型会把这些中间带噪分布组织成同一个前向加噪过程的边缘分布。
在 DDPM 中，这个过程是一条离散马尔可夫链；在连续时间的得分视角中，它由随机微分方程描述。

\subsection{从离散噪声水平到连续扩散过程}

多尺度构造使用有限多个噪声水平来近似从数据到噪声的平滑演进。
当噪声水平数量变大、相邻水平越来越接近时，自然可以用连续时间变量 $t$ 取代离散索引 $\sigma_k$。
我们不再考虑彼此分开的平滑密度
\[
p_{\sigma_1},p_{\sigma_2},\dots,p_{\sigma_K},
\]
而是考虑一族连续密度
\[
p_t(x),
\qquad
0\leq t\leq T.
\]
这一连续时间描述应理解为 DDPM 章节中离散加噪路径的对应形式：有限个中间带噪变量序列被连续索引的一族带噪分布替代。

在一端，$p_0$ 是数据分布，或其理想化版本；在另一端，$p_T$ 被设计成简单的类噪声分布，通常近似高斯。

得分模型也从噪声条件场
\[
s_\theta(x,\sigma_k)
\approx
\nabla_x\log p_{\sigma_k}(x)
\]
变为时间依赖场
\[
s_\theta(x,t)
\approx
\nabla_x\log p_t(x).
\]
因此，多尺度得分学习图景变成
\[
\text{有限噪声尺度}
\quad\longrightarrow\quad
\text{连续的一族带噪密度}
\quad\longrightarrow\quad
\text{前向扩散过程}.
\]

正是在这里，多尺度得分学习与扩散模型真正绑定起来。
在多个噪声水平上学习得分，会得到一族向量场，但仅凭这一点还不能定义完整扩散模型。
只有当这些带噪密度被组织成同一个前向加噪过程在不同时间的边缘分布时，才得到扩散模型。
换言之，必须存在一个随机动力学，使其随时间演化的分布正是 $p_t$。
此时，学习得到的得分 $s_\theta(x,t)$ 不再只是孤立的去噪场，而成为把这一加噪过程反向运行所需的修正项。

连续加噪过程可以由如下随机微分方程描述：
\[
\boxed{
dX_t
=
f(X_t,t)\,dt
+
g(t)\,dW_t.
}
\]
其中 $f(X_t,t)\,dt$ 是确定性漂移，$g(t)\,dW_t$ 注入布朗噪声，函数 $g(t)$ 控制时刻 $t$ 的噪声强度。

这正是上一章建立的“漂移+噪声”结构。
如果 $p_t(x)$ 表示 $X_t$ 的密度，那么在这里的标量、各向同性扩散系数 $g(t)$ 情形下，其演化由 Fokker--Planck 方程控制：
\[
\partial_t p_t(x)
=
-\nabla_x\cdot\bigl(p_t(x)f(x,t)\bigr)
+
\frac{1}{2}g(t)^2\Delta p_t(x).
\]
漂移项在空间中输运概率质量，而扩散项使密度扩散并平滑。

在生成建模中，前向过程并不是直接用于生成新数据。
它的作用是定义一条从数据到噪声的受控路径：
\[
p_0(x)
\approx
p_{\text{data}}(x),
\qquad
p_T(x)
\approx
\text{简单噪声分布}.
\]
生成问题则是把这一过程反向运行：从噪声开始，恢复出有结构的样本。

\subsection{反向时间动力学与得分的作用}

前向扩散过程很容易模拟。
从一个数据样本出发，按照所选随机微分方程不断加入噪声，就能得到噪声越来越强的状态。
但生成需要相反方向：我们希望从一个简单噪声样本开始，反向移动到数据分布。

乍看之下，反转带噪过程似乎不可能，因为扩散看起来会破坏信息。
然而，在适当的正则条件下，扩散过程的反向时间演化同样可以由随机微分方程描述。
关键事实是：反向漂移依赖中间密度 $p_t$ 的得分。

对于前向过程
\[
dX_t
=
f(X_t,t)\,dt
+
g(t)\,dW_t,
\]
在这里考虑的标量扩散情形下，反向时间动力学具有形式
\[
\boxed{
dX_t
=
\left[
f(X_t,t)
-
g(t)^2\nabla_x\log p_t(X_t)
\right]dt
+
g(t)\,d\bar{W}_t.
}
\]
这里采用常见的反向时间记号：过程从 $t=T$ 向 $t=0$ 模拟，所以时间增量应理解为沿反方向。
布朗运动 $\bar{W}_t$ 表示与该反向过程相联系的噪声。

这个方程正是学习得分真正进入生成过程的地方。
前向 SDE 定义一族带噪密度 $p_t$，而反向漂移则需要依赖得分的修正
\[
-g(t)^2\nabla_x\log p_t(X_t).
\]
由于真实得分未知，我们学习一个时间条件模型：
\[
s_\theta(x,t)
\approx
\nabla_x\log p_t(x),
\qquad
0\leq t\leq T.
\]
图~\ref{fig:score-forward-reverse-sde} 总结了这些组成部分如何相互配合。

\begin{figure}[H]
\centering
\begin{tikzpicture}[
  >=Latex,
  every node/.style={font=\small},
  state/.style={
    draw, rounded corners=2pt, align=center, inner sep=5pt,
    text width=2.20cm, minimum height=0.82cm
  },
  eqbox/.style={
    draw, rounded corners=2pt, align=center, inner sep=6pt,
    text width=10.8cm, minimum height=1.12cm
  },
  flow/.style={->,thick},
  reverse/.style={->,thick,dashed}
]
\node[font=\normalsize\bfseries] at (6.1,3.25)
  {得分提供反转扩散所需的修正};

% Forward row
\node[state] (f0) at (0,1.75) {$p_0$\\数据};
\node[state] (ft) at (6.1,1.75) {$p_t$\\中间状态};
\node[state] (fT) at (12.2,1.75) {$p_T$\\简单噪声};

\draw[flow] (f0) -- node[above] {\scriptsize 前向加噪} (ft);
\draw[flow] (ft) -- (fT);

\node[eqbox] (forward) at (6.1,0.15) {
\textbf{前向 SDE}\\[3pt]
$\displaystyle
dX_t=f(X_t,t)\,dt+g(t)\,dW_t
$
\qquad
\textbf{沿其边缘分布学习：}
$\displaystyle
s_\theta(x,t)\approx\nabla_x\log p_t(x)
$
};

% Reverse row
\node[state] (r0) at (0,-1.65) {$p_0$\\数据};
\node[state] (rt) at (6.1,-1.65) {$p_t$\\中间状态};
\node[state] (rT) at (12.2,-1.65) {$p_T$\\简单噪声};

\draw[reverse] (rT) -- node[above] {\scriptsize 反向时间生成} (rt);
\draw[reverse] (rt) -- (r0);

\node[eqbox] (reversebox) at (6.1,-3.45) {
\textbf{反向时间 SDE：从 $T$ 向 $0$ 模拟}\\[3pt]
$\displaystyle
dX_t=
\left[
f(X_t,t)-g(t)^2s_\theta(X_t,t)
\right]dt
+
g(t)\,d\bar W_t
$
};

\draw[flow] (forward) -- (reversebox);

\end{tikzpicture}
\caption{\textbf{得分提供反转扩散所需的修正。}
前向 SDE 定义带时间索引的边缘分布 $p_t$，而以时间为条件的模型学习 $s_\theta(x,t)\approx\nabla_x\log p_t(x)$。对于这里采用的标量、与状态无关的扩散系数，反向漂移与前向漂移相差 $-g(t)^2\nabla_x\log p_t(x)$。用 $s_\theta(x,t)$ 替代未知得分，就得到实际使用的反向时间 SDE，并从 $T$ 向 $0$ 进行模拟。}
\label{fig:score-forward-reverse-sde}
\end{figure}

因此，模型学习的是中间带噪分布的得分，而不只是数据分布本身的得分。
这些局部场提供从简单终端噪声分布回到结构化数据所需的修正信息。

一旦这些得分已知或能够被良好近似，就可以通过模拟反向时间动力学生成样本。
首先从
\[
X_T\sim p_T
\]
开始，其中 $p_T$ 是简单噪声分布；然后利用学习到的得分 $s_\theta(x,t)$，从 $T$ 向 $0$ 反向演化。
最终状态就被视为来自数据分布的近似样本。

至此，我们完成了从得分学习到基于得分扩散的概念过渡。
去噪得分匹配解释如何学习带噪分布的得分；多尺度得分学习把它扩展到多个噪声水平；前向扩散过程把这些带噪分布组织成一条统一的时间路径；反向时间 SDE 则说明如何利用学习到的得分沿这条路径返回。

因此，扩散模型并不只是得分估计器，而仅有得分匹配也还不是完整的生成模型。
前向扩散过程定义从数据到噪声的结构化路径，而得分学习提供反转这条路径所需的局部信息。
生成过程由二者结合得到：从一个简单噪声分布开始，沿学习得到的、由得分引导的反向动力学回到数据分布。

\section{从得分视角理解离散扩散模型}

上面的讨论建立了连续时间的得分扩散视角。
现在把它重新连接到本书前面介绍的离散 DDPM 形式。
目的不是重复潜变量推导，而是从得分视角重新解释 DDPM 中一个熟悉的部分：噪声预测目标。

在 DDPM 章节中，扩散模型通过前向加噪马尔可夫链、学习得到的反向马尔可夫链以及最终化简为噪声预测损失的变分训练目标来介绍。
这里说明为什么噪声预测损失也可以理解为携带中间带噪分布的类得分信息。

\subsection{噪声预测与条件得分}

回顾 DDPM 前向过程，它通过加入高斯噪声逐步扰动干净数据样本 $x_0$。
该过程一个关键性质是，可以直接从 $x_0$ 采样 $x_t$：
\[
x_t
=
\sqrt{\bar{\alpha}_t}\,x_0
+
\sqrt{1-\bar{\alpha}_t}\,\varepsilon,
\qquad
\varepsilon\sim\mathcal{N}(0,I),
\]
其中 $\bar{\alpha}_t$ 控制时刻 $t$ 剩余的信号强度。
等价地，
\[
q(x_t\mid x_0)
=
\mathcal{N}
\left(
x_t;
\sqrt{\bar{\alpha}_t}\,x_0,
(1-\bar{\alpha}_t)I
\right).
\]

从得分视角看，相关对象是关于带噪变量 $x_t$ 的对数密度梯度。
对于条件高斯 $q(x_t\mid x_0)$，有
\[
\log q(x_t\mid x_0)
=
-\frac{1}{2(1-\bar{\alpha}_t)}
\left\|
x_t-\sqrt{\bar{\alpha}_t}x_0
\right\|_2^2
+
\text{常数}.
\]
关于 $x_t$ 求导，得到
\[
\nabla_{x_t}\log q(x_t\mid x_0)
=
-\frac{1}{1-\bar{\alpha}_t}
\left(
x_t-\sqrt{\bar{\alpha}_t}x_0
\right).
\]
利用重参数化关系
\[
x_t-\sqrt{\bar{\alpha}_t}x_0
=
\sqrt{1-\bar{\alpha}_t}\,\varepsilon,
\]
可进一步写成
\[
\boxed{
\nabla_{x_t}\log q(x_t\mid x_0)
=
-\frac{1}{\sqrt{1-\bar{\alpha}_t}}\,
\varepsilon.
}
\]

这个恒等式解释了噪声预测为什么与得分估计密切相关。
如果模型能够预测加入的噪声 $\varepsilon$，那么乘上已知缩放因子
\[
-\frac{1}{\sqrt{1-\bar{\alpha}_t}}
\]
之后，它也就预测了加噪后的高斯分布 $q(x_t\mid x_0)$ 的条件得分。

必须准确理解这一含义。
上述表达是\emph{条件得分}，其中干净样本 $x_0$ 被视为已知。
而反向过程真正需要的是边缘得分
\[
\nabla_{x_t}\log q_t(x_t),
\]
其中 $q_t(x_t)$ 是时刻 $t$ 所有带噪样本的完整分布。
与去噪得分匹配一样，边缘得分可以理解为对所有可能产生同一带噪观测 $x_t$ 的干净样本 $x_0$ 进行平均后得到。

因此，DDPM 的噪声预测为获得得分模型所需的同类信息提供了一个实际回归路径：它给出局部方向，告诉一个带噪样本应该如何被修正，才能变得更像数据。

\subsection{基于得分扩散的统一视角}

现在可以把不同视角并列起来。

从 DDPM 的潜变量视角看，扩散模型由离散前向加噪过程
\[
x_0 \rightarrow x_1 \rightarrow \cdots \rightarrow x_T
\]
与学习得到的反向过程
\[
x_T \rightarrow x_{T-1} \rightarrow \cdots \rightarrow x_0
\]
共同构成。
训练可以通过变分目标表达，而在标准参数化下，最终得到一个预测每个时间步所加入噪声的神经网络。

从得分视角看，同一组带噪变量定义了一族中间带噪分布
\[
q_t(x_t),
\qquad
t=0,\dots,T.
\]
在每个时间步，反向过程都需要局部信息，告诉它应如何把带噪样本移动到对应中间分布下更高概率的区域。
这种局部信息由得分
\[
\nabla_{x_t}\log q_t(x_t)
\]
编码。

因此，两种视角是互补的。
DDPM 视角强调概率潜变量结构与离散反向马尔可夫链；得分视角强调中间带噪分布的几何结构，以及反转扰动所需的局部去噪方向。

在连续时间中，这种几何图景尤其清晰。
离散的一族
\[
q_0,q_1,\dots,q_T
\]
被连续密度路径
\[
p_t(x),
\qquad
0\leq t\leq T
\]
取代，而反向动力学显式包含得分
\[
\nabla_x\log p_t(x).
\]
因此，基于得分的连续时间形式可以看作 DDPM 中同一去噪原则的平滑对应。

这里最重要的概念信息并不是所有形式在细节上完全相同。
DDPM、退火 Langevin 动力学和连续时间得分扩散以不同方式组织生成过程，但都共享同一核心思想：
\[
\text{学习带噪样本在局部应当如何被修正，}
\]
并利用这些学习到的修正把噪声转化为数据。

从这个意义上说，扩散模型是结构化的时间依赖去噪模型。
前向过程创建一系列层次化带噪分布；学习问题估计沿这一层次反向移动所需的信息；采样则从简单噪声开始，反复应用学习到的修正，直到获得类似数据的样本。

至此，本章完成了离散 DDPM 视角与得分视角之间的连接。
DDPM 章节说明扩散模型如何作为潜变量生成模型推导出来；本章则说明它的去噪行为为什么也可以从几何角度理解：模型在不同噪声水平上学习类得分信息，并利用这些信息引导从噪声到数据的反向变换。

\section*{本章小结与参考文献}
\addcontentsline{toc}{section}{本章小结与参考文献}

本章建立了生成建模的得分视角，并利用它构建扩散模型的连续时间解释。
此前 DDPM 章节通过离散潜变量和变分学习引入扩散，而本章从局部概率几何出发理解同一去噪思想。

核心对象是得分函数 $\nabla_x \log p(x)$，它把密度转化为向量场。
我们先把得分解释为概率漂移方向，并说明它与布朗噪声结合后如何形成 Langevin 动力学。
利用 Fokker--Planck 方程进行的平稳性计算解释了：为什么目标密度的得分可以支撑一个以该目标分布为平稳分布的采样动力学。

随后转向学习问题。
由于真实数据密度未知，其得分无法直接计算。
得分匹配通过最小化 Fisher 散度目标并利用分部积分，把未知数据得分从训练损失中消去。
去噪得分匹配进一步通过高斯噪声扰动数据，使这一思想更适用于实际高维建模，并把平滑密度得分的学习转化为回归问题。
在这一视角下，模型学习的是一个带噪样本应如何朝干净数据分布方向移动。

接着，本章从单一噪声水平扩展到一族带噪分布。
多尺度得分学习在大噪声水平提供粗粒度全局引导，在小噪声水平提供精细局部修正。
退火 Langevin 动力学把这些得分场转化为采样过程，而连续时间形式则把这些带噪分布组织成前向随机微分方程的边缘分布。
相应的反向时间动力学说明，时间依赖得分
$\nabla_x \log p_t(x)$
为何是把噪声重新转化为数据所需的关键修正项。

最后，我们把这一得分视角重新连接到 DDPM。
DDPM 的噪声预测目标可以理解为学习中间带噪分布的类得分信息。
因此，DDPM 与得分视角并不是相互竞争的解释，而是互补的：DDPM 强调离散潜变量构造，而得分扩散强调反转噪声过程的几何结构与动力学。

\bigskip

基于得分扩散的发展融合了若干重要研究脉络。
Hyv\"arinen~\cite{hyvarinen2005estimation} 提出了得分匹配原则，说明模型可以通过匹配对数密度梯度进行估计，而不需要密度的归一化常数。
Vincent~\cite{vincent2011connection} 随后把得分匹配与去噪联系起来，说明去噪目标可以恢复扰动后数据分布的得分信息。
这些思想构成本章得分学习与去噪得分匹配推导的基础。

Song 和 Ermon~\cite{song2019generative} 把基于得分的生成建模带入现代深度生成模型场景：他们在多个高斯噪声水平上学习得分，并利用退火 Langevin 动力学采样。
这项工作明确展示了多尺度得分场的重要性，尤其是在数据集中于低维结构附近时。
随后，Song 等人~\cite{song2021scorebased} 通过随机微分方程、反向时间 SDE 及相关采样过程，统一并扩展了连续时间形式。
这些工作与 Ho 等人~\cite{ho2020denoising} 的 DDPM 形式共同构成了现代扩散模型的主要概念基础：扩散既可以理解为概率去噪模型，也可以理解为由得分引导的生成动力学。

\chapter{精确密度模型：归一化流与自回归分解}
\section*{概述}
\addcontentsline{toc}{section}{概述}

前面几章已经展示了进入现代生成建模的多条主要路径。
变分自编码器引入了潜变量建模与近似推断；去噪扩散概率模型把视角转向随机加噪与去噪过程；连续时间与得分章节又进一步通过密度演化、随机微分方程和学习得到的得分场发展了这一思想。

本章转向另一种范式：\emph{精确密度建模}。
这里的目标是构造这样一类生成模型：对于一个观测数据点，其概率密度在模型设计上就可以直接计算。
这改变了建模哲学。
我们不再依赖难以处理的潜变量边缘似然，也不再学习去反转随机扰动过程，而是主动施加某种数学结构，使似然计算从一开始就保持可处理。

第一条路线是\textbf{归一化流}（normalising flows）。
归一化流从一个简单基础分布（例如标准高斯）出发，依次应用一系列可微、可逆变换。
变量变换公式通过 Jacobian 行列式刻画局部拉伸或压缩，从而精确告诉我们密度在每一次变换下如何变化。
只要这些变换可逆，且其 Jacobian 行列式可处理，那么似然计算和采样都是模型中的精确操作。

第二条路线是\textbf{自回归模型}。
它不通过可逆映射对基础分布进行变换，而是直接利用概率链式法则构造联合密度。
一个高维数据点被拆解成有序分量序列，每个分量都以此前分量为条件进行建模。
这种方式无需潜变量、变分下界或 Jacobian 行列式，也能得到精确似然，但采样通常必须顺序执行。

这两种方法在现代生成建模中仍然十分重要，因为它们展示了精确似然建模能够达到怎样的程度。
归一化流为可逆密度估计、不确定性建模以及概率分布的可处理变换提供了清晰框架。
与此同时，自回归模型构成许多有影响力的序列、音频、语言以及像素级图像生成系统的基础。
即使较新的系统会组合多种建模思想，本章涉及的原则——变量变换、可处理 Jacobian、条件分解以及有序生成——仍然会作为核心构件不断出现。

\section*{概念导图}
\addcontentsline{toc}{section}{概念导图}

本章把精确密度建模作为与此前潜变量、扩散和得分方法相对照的另一种范式展开。
核心问题是：灵活的高维生成模型如何同时保留精确似然计算能力？

\begin{itemize}
    \item[\(\rightarrow\)] \textbf{首先通过变量变换公式引入归一化流。}
    归一化流利用可微、可逆变换把简单基础变量 $z$ 映射为数据变量 $x=f(z)$。
    变量变换公式通过 Jacobian 行列式跟踪相应的密度变化。

    \item[\(\rightarrow\)] \textbf{随后把平面流作为第一种非线性流层来研究。}
    平面流通过秩一残差扰动引入非线性，在保持 Jacobian 行列式解析可处理的同时，也展示了可逆性约束的重要性。

    \item[\(\rightarrow\)] \textbf{接着转向仿射耦合层这一实用流构造。}
    仿射耦合层只变换输入的一部分，而以另一部分为条件。
    这种设计使逆映射具有闭式形式，且对数行列式能够高效计算，因此适用于实际流模型。

    \item[\(\rightarrow\)] \textbf{然后把多个流层组合成完整生成模型。}
    堆叠许多可逆层后，简单局部形变可以组合成表达力很强的全局变换。
    似然计算把数据逆向映射回基础空间，而采样则把基础噪声正向映射到数据空间。

    \item[\(\rightarrow\)] \textbf{之后引入自回归分解，作为获得精确似然的第二条路线。}
    自回归模型利用链式法则
    $p(x)=\prod_{i=1}^d p(x_i\mid x_{<i})$
    构造精确似然，无需潜变量、变分下界或 Jacobian 行列式。

    \item[\(\rightarrow\)] \textbf{最后解释条件参数化、训练、采样与掩码。}
    共享模型可以参数化所有条件分布。
    训练时完整数据向量是已知的，因此似然项常可并行计算；采样时，各分量则必须顺序生成。
    掩码用来强制实现使自回归分解可处理的有序信息流。
\end{itemize}

本章的两条路线以及使其保持可处理性的结构原则总结如下图所示。

\begin{center}
\includegraphics[width=0.96\textwidth]{figures/ch07/chapter7_exact_density_at_a_glance.pdf}
\end{center}


\section{归一化流：通过可逆变换实现精确密度}

归一化流通过在简单参考分布与数据分布之间构造显式映射，同时始终保留精确似然计算能力，为生成建模提供了一条直接且完全概率化的路线~\cite{rezende2015variational}。
“归一化流”这一术语指的是这一更广泛的可逆密度模型家族，而不是某一种具体变换；平面流、仿射耦合层以及相关架构，都是设计这类变换的不同方式，其共同目标是同时保持表达力、可逆性与计算可处理性。

归一化流背后的核心数学工具是\emph{变量变换公式}。
该公式精确刻画概率密度在可微、可逆映射下如何变化。
归一化流不引入潜变量或辅助随机过程，而完全依赖确定性变换，并且这些变换对密度的影响可以精确计算。
通过串联许多这样的变换，一个简单基础分布——通常是标准高斯——可以逐渐被重塑成高度结构化的数据分布。

从几何角度看，每个变换都是对空间进行受控形变，通过局部拉伸或压缩体积来改变概率密度，而这种变化由 Jacobian 行列式显式计入。
只要每个变换可逆且其 Jacobian 行列式可处理，所得模型就同时支持精确似然计算和精确采样。
这种表达力与数学透明性的结合，使归一化流区别于其他生成建模范式，也使其成为研究精确密度模型的自然起点。

\subsection{重新审视变量变换公式}

虽然变量变换公式已经在第~5 章从更一般的确定性密度输运角度介绍过，但这里仍值得重新审视，因为它正是归一化流的建模原则。

假设 $X$ 有已知密度 $p_X(x)$，并应用一个可微、可逆函数
$Y=f(X)$。
我们的目标是确定 $Y$ 的密度 $p_Y(y)$。
直观上，$Y$ 的密度需要同时考虑两件事。

第一，它必须考虑\emph{概率质量从哪里来}。
由于 $y=f(x)$，观测到 $y$ 的可能性取决于其对应原像
$x=f^{-1}(y)$
在原分布下有多大概率。

第二，它必须考虑\emph{变换如何拉伸或压缩空间}。
如果一个变换让点彼此离得更远，密度应当降低；如果让点压缩得更紧，密度应当提高。
密度始终是单位长度、面积或体积上的概率，因此局部间距至关重要。

在一维中，这两个思想结合为变量变换公式
\[
\boxed{
p_Y(y)
=
p_X(f^{-1}(y))
\left|
\frac{d f^{-1}(y)}{dy}
\right|.
}
\]
第一项告诉我们哪个 $x$ 生成了 $y$；第二项是缩放因子，用来补偿 $x$ 附近的邻域映射到 $y$ 后被拉伸或压缩的程度。

\paragraph{一个简单几何例子。}

设
$X\sim \mathrm{Uniform}(0,1),
\qquad
Y=f(X)=2X$。

这个变换把所有点扩大两倍，把区间 $[0,1]$ 拉伸成 $[0,2]$。
其逆映射为
\[
x=f^{-1}(y)=\frac{y}{2},
\]
所以
\[
\frac{d f^{-1}(y)}{dy}
=
\frac{1}{2}.
\]
代入公式得到
\[
p_Y(y)
=
p_X\!\left(\frac{y}{2}\right)\frac{1}{2}.
\]
在新区间 $[0,2]$ 内，原密度满足 $p_X(y/2)=1$，因此
\[
p_Y(y)=\frac{1}{2},
\qquad 0\leq y\leq 2.
\]
这与直觉完全一致：同样的总概率现在填充一个宽度变为两倍的区间，所以密度减半。
这个简单逻辑可以直接推广到高维。

\paragraph{高维情形：从长度拉伸到体积拉伸。}

在多维中，普通导数被体积缩放的行列式取代。
对于简单线性映射
\[
f(x)=Ax,
\qquad
A=
\begin{bmatrix}
2 & 0\\
0 & 3
\end{bmatrix},
\]
局部面积会乘以 $|\det A|=6$，因此
\[
p_Y(y)=p_X(A^{-1}y)\frac{1}{|\det A|}.
\]
同样的概率质量占据了更大的局部体积，所以密度按照其倒数因子下降。

\paragraph{Jacobian：观察局部线性变化的透镜。}

为了推广到非线性变换，需要使用\emph{Jacobian 矩阵}。
对于可微函数
\[
f:\mathbb{R}^d\to\mathbb{R}^d,
\]
Jacobian 定义为
\[
[J_f(x)]_{ij}
=
\frac{\partial f_i(x)}{\partial x_j}.
\]
可以把 Jacobian 看作观察局部线性行为的一面透镜。
在点 $x$ 附近，一个很小的位移 $\delta x$ 近似按照
\[
f(x+\delta x)
\approx
f(x)+J_f(x)\delta x
\]
变换。
因此，$J_f(x)$ 告诉我们 $x$ 附近一个无穷小邻域映射到 $y=f(x)$ 后如何发生形变；其行列式则量化对应局部体积膨胀或收缩了多少。

对于一般可微可逆变换 $y=f(x)$，多元变量变换公式为
\[
\boxed{
p_Y(y)
=
p_X(f^{-1}(y))
\left|
\det J_{f^{-1}}(y)
\right|
=
p_X(x)
\frac{1}{|\det J_f(x)|},
\qquad y=f(x).
}
\]

这就是归一化流的数学核心：
\emph{只要每个变换可逆且其 Jacobian 行列式可处理，即使整体映射具有很强表达力，我们仍然能够精确写出变换后变量的密度。}
图~\ref{fig:flow-change-variable-geometry} 展示了这一公式的几何含义：固定的概率质量落在局部被拉伸或压缩的区域中，密度按体积变化因子的倒数进行调整。

\begin{figure}[H]
\centering
\includegraphics[width=0.88\textwidth]{figures/ch07/figure7_1_change_of_variables_geometry.pdf}
\caption{\textbf{变量变换通过局部体积变化修正概率密度。}
同一份概率质量经过可逆变换后会占据不同的局部面积。在线性例子中，行列式给出面积缩放比例；对于非线性流，Jacobian 则逐点给出相应的局部缩放。这就是对数行列式校正项的几何来源。}
\label{fig:flow-change-variable-geometry}
\end{figure}


\begin{tcolorbox}[
  colback=blue!4,
  colframe=blue!50!black,
  title={为什么标准归一化流保持维度不变},
  boxrule=0.6pt,
  arc=2mm,
  left=2mm,
  right=2mm,
  top=1mm,
  bottom=1mm
]
标准归一化流采用形式为
\[
f:\mathbb{R}^d\to\mathbb{R}^d
\]
的可微可逆变换。
相同的输入输出维度是标准变量变换公式所要求的：Jacobian 必须是方阵，$\det J_f(x)$ 才有定义，并能够度量局部体积变化。
如果映射把 $\mathbb{R}^d$ 映到更低维空间，信息会被压塌，不同输入可能映射到同一个输出；如果把 $\mathbb{R}^d$ 映到更高维空间，像会落在低维流形上，在没有额外结构时并不能定义普通的满维密度。
因此，标准归一化流保持维度不变。
若失去这种满维可逆性，就无法使用变量变换公式精确追踪概率质量。
\end{tcolorbox}

\subsection{平面流}

上述例子是有意选择的简单情形，它们展示了变量变换公式如何跟踪线性拉伸与体积缩放下的密度。
然而，线性变换的表达能力有限：它们可以拉伸、旋转或剪切空间，却无法以丰富的非线性方式弯曲或扭曲空间。
为了逼近复杂分布，我们需要能够引入受控非线性，同时仍保持归一化流两项核心要求的变换：可逆性与可处理的 Jacobian 行列式。

平面流（planar flow）是这种权衡的经典例子，最初作为 Rezende 与 Mohamed~\cite{rezende2015variational} 的归一化流框架的一部分提出。
一种基础非线性流层定义为
\[
f(x)
=
x+u\,h(w^\top x+b),
\]
其中 $x\in\mathbb{R}^d$，$u,w\in\mathbb{R}^d$，$b\in\mathbb{R}$，且
\[
h:\mathbb{R}\to\mathbb{R}
\]
是光滑激活函数，例如 $\tanh$ 或 \texttt{sigmoid}。
这里 $x$ 表示当前单个流层的输入；在堆叠流中，它通常是某个中间变量。

项
\[
w^\top x+b
\]
对输入进行仿射投影，把高维向量 $x$ 压缩为一个标量。
单独使用这个投影只会产生线性量。
为了引入曲率和更丰富的行为，我们把该标量送入非线性激活函数 $h(\cdot)$。
$\tanh$ 与 \texttt{sigmoid} 之类的函数会把输入压缩到有界范围，因此 $h$ 的输出可以看成一个软门控值，用来控制变换应该多强地修改输入。

把这个门控标量与向量 $u$ 相乘，会产生一个具有方向的缩放扰动。
形变始终沿 $u$ 的方向发生，但其幅度以非线性方式依赖投影 $w^\top x+b$。
把这一扰动重新加回 $x$，既保证输出仍然位于完整的 $d$ 维空间中，又引入了灵活的非线性修正。

\begin{tcolorbox}[
  colback=blue!4,
  colframe=blue!50!black,
  title={残差结构及其动机},
  boxrule=0.6pt,
  arc=2mm,
  left=2mm,
  right=2mm,
  top=1mm,
  bottom=1mm
]
平面流具有深度学习中残差连接的基本精神：
\[
f(x)=x+\text{学习得到的扰动}.
\]
它保留恒等映射，同时叠加受控形变。

如果只使用非线性项
\[
x\mapsto u\,h(w^\top x+b),
\]
由于 $h(\cdot)$ 输出标量，结果会塌缩到 $u$ 所在的一维方向，从而破坏归一化流所需的满维结构。

把扰动加回 $x$ 后，平面流能够保持整个 $d$ 维空间可达，同时引入有针对性的非线性形变。
\end{tcolorbox}

\paragraph{Jacobian 推导。}
为了计算平面流作用下密度如何变换，我们推导
\[
f(x)=x+u\,h(w^\top x+b)
\]
的 Jacobian。为简化记号，令
\[
a=w^\top x+b,
\qquad
f(x)=x+u\,h(a).
\]
对 $x$ 求导并使用链式法则，得到
\[
\frac{\partial f(x)}{\partial x}
=
I+u\,h'(a)\frac{\partial a}{\partial x}.
\]
由于 $\frac{\partial a}{\partial x}=w^\top$，Jacobian 可写成紧凑形式
\[
J_f(x)
=
I+u\,w^\top h'(a).
\]
等价地，
\[
J_f(x)
=
I+u\psi(x)^\top,
\qquad
\psi(x):=h'(w^\top x+b)w.
\]

这一结构非常便利。Jacobian 是单位矩阵加上一个秩一矩阵扰动。因此，无须计算通常需要 $O(d^3)$ 代价的完整 $d\times d$ 行列式，行列式可以化为单个标量表达式。利用矩阵行列式引理
\[
\det(I+uv^\top)=1+v^\top u,
\]
可得
\[
\det J_f(x)
=
1+\psi(x)^\top u.
\]
代入 $\psi(x)=h'(w^\top x+b)w$，得到
\[
\boxed{
\det J_f(x)
=
1+h'(w^\top x+b)\,w^\top u.
}
\]
因此，平面流对对数行列式的贡献为
\[
\log|\det J_f(x)|
=
\log\left|
1+h'(w^\top x+b)\,w^\top u
\right|.
\]
这一闭式表达式使平面流能够用于基于似然的可处理训练。

\begin{tcolorbox}[
  colback=blue!4,
  colframe=blue!50!black,
  title={理解秩一更新},
  boxrule=0.6pt,
  arc=2mm,
  left=2mm,
  right=2mm,
  top=1mm,
  bottom=1mm
]
矩阵 $u w^\top$ 是 $\mathbb{R}^d$ 中两个向量的外积，而这样的矩阵秩为一。

要理解这一点，可将它作用于任意向量 $v$：
\[
u w^\top v
=
u\,(w^\top v).
\]
标量 $w^\top v$ 会随 $v$ 变化，但输出始终只是 $u$ 的某个标量倍数。因此，该矩阵只有一个独立的输出方向。

这也解释了为什么行列式如此简单：平面流并不试图同时对所有维度实施任意的非线性形变，而是引入一个受控的单方向形变，其体积变化可以解析计算。
\end{tcolorbox}

图~\ref{fig:planar-flow-rank-one} 展示了相应的几何图景：形变幅度随空间位置变化，但残差位移始终平行于单一方向 $u$。

\begin{figure}[H]
\centering
\includegraphics[width=0.92\textwidth]{figures/ch07/figure7_2_planar_flow_rank_one.pdf}
\caption{\textbf{平面流是非线性的，但其方向受到约束。}
位移大小通过标量非线性响应随位置变化，但所有残差位移始终平行于同一个向量 $u$。这正是正文所推导秩一结构的几何结果。}
\label{fig:planar-flow-rank-one}
\end{figure}


\paragraph{可逆性条件。}

除了 Jacobian 行列式必须可处理之外，流层还必须可逆。对于平面流
\[
f(x)=x+u\,h(w^\top x+b),
\]
其逆映射一般没有闭式表达式，因此通常通过参数约束来控制可逆性。

由前面的 Jacobian 推导，
\[
\det J_f(x)
=
1+h'(w^\top x+b)\,w^\top u.
\]
如果该行列式变为零，变换会在局部压缩体积，因而无法局部求逆。因此，一个常见的充分条件是让该行列式对所有 $x$ 都保持为正。

对于 $\tanh$ 这类满足
\[
0\leq h'(a)\leq 1
\]
的激活函数，只需要求
\[
w^\top u>-1.
\]
于是对任意 $a$ 都有
\[
1+h'(a)\,w^\top u>0,
\]
从而平面流可避免局部体积坍缩。实践中，可以通过对 $u$ 进行重参数化来强制满足这一条件。

因此，平面流很好地体现了一项核心设计权衡：通过简单的残差扰动引入非线性，同时使 Jacobian 行列式保持解析可处理，并通过一个标量条件控制可逆性。


\section{实践中的归一化流：耦合层与组合}

上一节将归一化流介绍为建立在可逆变换之上的精确密度模型。平面流给出了一个简单的非线性流层示例，其 Jacobian 行列式仍然易于处理，但也暴露出一个实际限制：其逆映射通常没有闭式形式。

现在，我们从这一教学性构造转向实用流架构背后的核心设计原则之一：\emph{耦合（coupling）}。耦合层每次只变换输入的一部分，并把剩余部分用作条件信号，由此保留精确可逆性。这种结构会产生三角 Jacobian，给出闭式逆映射，并允许高效堆叠多个层。

随后我们将说明如何把这些层组合成一个完整的生成模型。正是这种组合，把简单的局部变换转化为富有表达能力的密度模型，同时保留归一化流的两项定义性操作：精确似然计算，以及从学习到的模型中进行精确采样。

\subsection{仿射耦合层与三角 Jacobian}

仿射耦合层最早用于 NICE、RealNVP 等基于耦合的流模型~\cite{dinh2014nice,dinh2017density}，是构建具有可处理逆映射和对数行列式的高表达能力流模型时最实用的方法之一。

仿射耦合层首先把输入向量划分为两部分。设 $x\in\mathbb{R}^d$ 是某个流层的输入，并写为
\[
x=(x_1,x_2),
\]
其中 $x_1\in\mathbb{R}^{d_1}$、$x_2\in\mathbb{R}^{d_2}$，且 $d_1+d_2=d$。仿射耦合变换定义为
\[
\begin{aligned}
y_1 &= x_1,\\
y_2 &= x_2\odot \exp(s(x_1)) + t(x_1),
\end{aligned}
\]
其中 $s(\cdot)$ 与 $t(\cdot)$ 是函数，通常由神经网络表示，$\odot$ 表示逐元素乘法。

输入的第一部分保持不变。第二部分通过一个逐元素仿射映射进行变换，其缩放量和平移量由第一部分预测。因此，$x_1$ 充当上下文变量，而 $x_2$ 则依据这一上下文进行缩放和平移。

虽然该变换关于 $x_2$ 是仿射的，但函数 $s(x_1)$ 和 $t(x_1)$ 可以具有高度非线性。这使整体变换既有较强表达能力，又保持简单的代数结构。因此，该层有意采用非对称设计：牺牲单层内部所有维度的同时混合，换取精确求逆与高效行列式计算。

\paragraph{Jacobian 行列式。}

为了计算 Jacobian，写成
\[
y=
\begin{bmatrix}
y_1\\
y_2
\end{bmatrix}
=
\begin{bmatrix}
x_1\\
x_2\odot \exp(s(x_1))+t(x_1)
\end{bmatrix}.
\]
其 Jacobian 具有分块形式
\[
J
=
\frac{\partial y}{\partial x}
=
\begin{bmatrix}
\frac{\partial y_1}{\partial x_1}
&
\frac{\partial y_1}{\partial x_2}
\\[4pt]
\frac{\partial y_2}{\partial x_1}
&
\frac{\partial y_2}{\partial x_2}
\end{bmatrix}.
\]

由于 $y_1=x_1$，有
\[
\frac{\partial y_1}{\partial x_1}=I,
\qquad
\frac{\partial y_1}{\partial x_2}=0.
\]
左下角分块 $\frac{\partial y_2}{\partial x_1}$ 可能十分复杂，因为其中包含神经网络 $s(\cdot)$ 与 $t(\cdot)$ 的导数。然而关键在于，这一分块并不会影响行列式。

对于右下角分块，注意 $y_2$ 的每个分量为
\[
y_{2,i}
=
x_{2,i}\exp(s_i(x_1))+t_i(x_1).
\]
在固定 $x_1$ 的情况下，对 $x_{2,j}$ 求导可得
\[
\frac{\partial y_{2,i}}{\partial x_{2,j}}
=
\begin{cases}
\exp(s_i(x_1)), & i=j,\\
0, & i\neq j.
\end{cases}
\]
因此
\[
\frac{\partial y_2}{\partial x_2}
=
\operatorname{diag}\!\big(\exp(s(x_1))\big).
\]

把这些分块组合起来，得到
\[
J
=
\begin{bmatrix}
I & 0\\
* & \operatorname{diag}\!\big(\exp(s(x_1))\big)
\end{bmatrix}.
\]
这是一个分块下三角矩阵。对于任意三角矩阵或分块三角矩阵，其行列式等于对角分块行列式之积。因此
\[
\det J
=
\det(I)\,
\det\!\left(
\operatorname{diag}\!\big(\exp(s(x_1))\big)
\right).
\]
由于 $\det(I)=1$，而对角矩阵的行列式就是各对角元素之积，所以
\[
\det J
=
\prod_i \exp(s_i(x_1)).
\]
取对数得到特别简单的表达式
\[
\boxed{
\log|\det J|
=
\sum_i s_i(x_1).
}
\]

这正是仿射耦合层的主要可处理性来源。稠密且可能很复杂的分块 $\frac{\partial y_2}{\partial x_1}$ 无须参与行列式计算，全部体积变化都由缩放输出 $s(x_1)$ 控制。因此，即使层内部使用了非线性神经网络，对数行列式仍可以关于被变换维数在线性时间内计算。

三角分块结构是这里的关键计算简化：非对角分块可以很复杂，但不会进入行列式。这体现了精确密度建模中的一个反复出现的原则：以精心设计的方式组织非线性依赖，使高维概率计算仍能分解为简单项。

\paragraph{可逆性。}

带来三角 Jacobian 的同一种非对称结构，也同时给出了闭式逆映射。由前向变换
\[
\begin{aligned}
y_1 &= x_1,\\
y_2 &= x_2\odot \exp(s(x_1)) + t(x_1),
\end{aligned}
\]
可以立即恢复
\[
x_1=y_1.
\]
将其代入第二个方程，得到
\[
y_2
=
x_2\odot \exp(s(y_1)) + t(y_1).
\]
由于 $\exp(s(y_1))$ 逐分量严格为正，这个逐元素仿射变换可直接求逆：
\[
x_2
=
\bigl(y_2-t(y_1)\bigr)\odot \exp(-s(y_1)).
\]
因此，逆映射为
\[
\boxed{
\begin{aligned}
x_1 &= y_1,\\
x_2 &= \bigl(y_2-t(y_1)\bigr)\odot \exp(-s(y_1)).
\end{aligned}
}
\]

这里不需要求根或迭代求逆。逆变换与前向变换具有基本相同的计算代价。这使仿射耦合层尤其适合大规模模型，因为采样与似然计算往往需要反复应用许多流变换。图~\ref{fig:affine-coupling-structure} 总结了使这两种性质同时成立的非对称信息流：一部分变量原样通过，同时控制另一部分变量的变换。

\begin{figure}[H]
\centering
\begin{tikzpicture}[
  >=Latex,
  every node/.style={font=\small},
  state/.style={draw,rounded corners=2pt,align=center,inner sep=5pt,text width=2.35cm,minimum height=.82cm},
  note/.style={draw,rounded corners=2pt,align=center,inner sep=5pt,text width=4.45cm,minimum height=.82cm},
  flow/.style={->,thick}
]
\node[font=\normalsize\bfseries] at (6.0,2.75) {仿射耦合：单向依赖带来可处理性};

\node[state] (x1) at (.8,1.20) {$x_1$\\上下文};
\node[state] (x2) at (.8,-.55) {$x_2$\\被变换};
\node[state] (st) at (4.0,.30) {缩放与平移\\由 $x_1$ 决定};
\node[state] (y1) at (7.35,1.20) {$y_1$\\保持不变};
\node[state] (y2) at (7.35,-.55) {$y_2$\\被变换};

\draw[flow] (x1) -- (st);
\draw[flow] (x1) -- (y1);
\draw[flow] (st) -- (y2);
\draw[flow] (x2) -- (y2);

\node[note] (tri) at (3.0,-2.45) {单向依赖\\$\Rightarrow$ 分块三角 Jacobian};
\node[note] (inv) at (8.7,-2.45) {保持不变的分支\\$\Rightarrow$ 可在逆变换中直接恢复};
\end{tikzpicture}
\caption{\textbf{仿射耦合中的非对称信息流。}
一部分变量原样通过，同时决定另一部分变量如何变换。这种单向依赖同时带来了三角 Jacobian 和闭式逆映射。}
\label{fig:affine-coupling-structure}
\end{figure}


\subsection{堆叠流：似然计算与采样}

单个流层只能产生有限的空间形变。平面流沿一个学习到的方向对恒等映射施加扰动，而仿射耦合层每次只变换输入的一部分。因此，归一化流的表达能力来自\emph{组合}：许多简单、可逆的变换相互堆叠，使各自的小幅形变累积为丰富的全局映射。

设
\[
f_1,f_2,\dots,f_K
\]
为一系列可微且可逆的变换。从基础变量
\[
z_0\sim p_Z(z_0)
\]
出发，定义中间变量
\[
z_k=f_k(z_{k-1}),
\qquad k=1,\dots,K,
\]
以及最终输出
\[
x=z_K.
\]
等价地，
\[
x
=
f_K\circ f_{K-1}\circ \cdots \circ f_1(z_0).
\]
由于每一层都可逆，整个组合也可逆。因此我们能够沿两个方向移动：
\[
z_0
\xrightarrow{f_1}
z_1
\xrightarrow{f_2}
\cdots
\xrightarrow{f_K}
z_K=x,
\]
或者
\[
x=z_K
\xrightarrow{f_K^{-1}}
z_{K-1}
\xrightarrow{f_{K-1}^{-1}}
\cdots
\xrightarrow{f_1^{-1}}
z_0.
\]

这种双向结构给出了流模型的两个基本操作：似然计算与采样。

\paragraph{似然计算：从数据空间到基础空间。}

在基于似然的训练中，我们给定一个数据点 $x$。按照上述约定，$x$ 是最终变量 $z_K$。为了计算其密度，我们通过逆变换逐层向后映射：
\[
x=z_K
\xrightarrow{f_K^{-1}}
z_{K-1}
\xrightarrow{}
\cdots
\xrightarrow{f_1^{-1}}
z_0.
\]
如果模型拟合良好，恢复出来的基础空间变量 $z_0$ 应当在简单基础密度 $p_Z$ 下具有较高概率。

逐层应用变量变换公式，得到
\[
p_{Z_k}(z_k)
=
p_{Z_{k-1}}(z_{k-1})
\left|
\det J_{f_k}(z_{k-1})
\right|^{-1}.
\]
从 $k=1$ 到 $K$ 反复使用这一递推，并利用 $z_K=x$，得到
\[
\boxed{
p_X(x)
=
p_Z(z_0)
\prod_{k=1}^K
\left|
\det J_{f_k}(z_{k-1})
\right|^{-1}.
}
\]
取对数可得
\[
\boxed{
\log p_X(x)
=
\log p_Z(z_0)
-
\sum_{k=1}^K
\log
\left|
\det J_{f_k}(z_{k-1})
\right|.
}
\]

这一表达式具有直接的解释。基础密度项 $\log p_Z(z_0)$ 衡量经过逆变换的数据点在简单参考分布下有多合理；对数行列式项则校正整个变换序列所引入的全部体积变化。每一层都会向最终对数似然贡献一个可加的体积校正项。

\begin{tcolorbox}[
  colback=blue!3,
  colframe=blue!50!black,
  title={乘积形式是如何产生的},
  boxrule=0.6pt,
  arc=2mm,
  left=2mm,
  right=2mm,
  top=1mm,
  bottom=1mm
]
对于单个流层 $z_k=f_k(z_{k-1})$，变量变换公式给出
\[
p_{Z_k}(z_k)
=
p_{Z_{k-1}}(z_{k-1})
\left|\det J_{f_k}(z_{k-1})\right|^{-1}.
\]
反复应用这一关系，得到
\[
\begin{aligned}
p_{Z_2}(z_2)
&=
p_{Z_0}(z_0)
\left|\det J_{f_1}(z_0)\right|^{-1}
\left|\det J_{f_2}(z_1)\right|^{-1},\\
&\ \vdots\\
p_{Z_K}(z_K)
&=
p_{Z_0}(z_0)
\prod_{k=1}^K
\left|\det J_{f_k}(z_{k-1})\right|^{-1}.
\end{aligned}
\]
由于 $z_K=x$，这就给出了观测数据点的似然。
\end{tcolorbox}

\paragraph{采样：从基础空间到数据空间。}

采样沿相反方向使用同一组变换。首先从基础分布中抽取一个点，
\[
z_0\sim p_Z(z_0),
\]
然后依次应用学习到的变换：
\[
z_0
\xrightarrow{f_1}
z_1
\xrightarrow{f_2}
\cdots
\xrightarrow{f_K}
z_K=x.
\]
所得 $x$ 就是来自学习到的模型分布的一个样本。

相对于模型本身而言，这一过程是精确的：不需要 Markov 链、变分近似，也不需要数值积分。变换学习完成后，生成过程只是从基础分布采样，再应用一系列确定性的可逆映射。

这也进一步澄清了它与扩散模型的差别。扩散模型通过模拟从噪声到数据的学习到的反向随机过程来生成样本，而归一化流则应用一个学习到的确定性可逆变换。

\paragraph{组合、混合与表达能力。}

堆叠许多流层能够提升模型表达能力，但仅有组合还不够。层的排列还必须使信息能够在不同维度之间混合。如果每一层始终采用相同的划分，一些变量就会反复充当条件变量，却始终没有被直接变换。

为了解决这一问题，实用流模型通常会在非线性流层之间穿插简单的可逆混合操作。这些操作可以是置换
\[
z\mapsto Pz,
\]
也可以是更一般的可逆线性变换
\[
z\mapsto Az,
\]
其中 $P$ 是置换矩阵，$A$ 是具有可处理行列式的可逆矩阵。

这些混合操作的作用，是改变后续层中哪些坐标彼此交互。在一个耦合层中保持不变的变量，可以在下一层被移动到需要变换的子集中。从几何上看，每一次混合操作都改变了下一次非线性形变所处的坐标系。

因此，一个现代流模型通常由以下模式反复构成：
\[
\text{混合坐标}
\quad\longrightarrow\quad
\text{施加可处理的非线性形变}
\quad\longrightarrow\quad
\text{再次混合}.
\]
这种不断交替使局部、结构化的变换可以组合为灵活的全局映射，同时保留精确可逆性和可处理的似然计算。

RealNVP 和 Glow 等架构都建立在这一原则之上。它们在耦合变换形式、缩放参数化和坐标混合细节上有所不同，但底层数学结构相同：组合简单的可逆层，追踪每一层的体积变化，并通过精确最大似然训练整个模型。

\section{作为生成模型的自回归分解}

归一化流通过在简单基础分布与数据分布之间构造可逆映射，实现精确密度建模。其可处理性来自变量变换公式：只要变换可逆且 Jacobian 行列式易于计算，似然就可以被精确求得。

自回归模型则提供了第二条实现精确密度建模的路径。它不把基础随机变量变换为数据，而是直接把联合密度分解为一系列条件密度的乘积。此时，精确似然来自概率论的链式法则。这也改变了结构约束施加的位置：在归一化流中，结构施加在变换本身；在自回归模型中，结构施加在一个数据点各分量之间的依赖关系上。通过选定一个顺序，并规定每个分量只能依赖先前分量，复杂的联合密度就变成了若干较简单条件项的乘积。

因此，一个高维数据点可以被看作按顺序逐个揭示的一列分量。每个分量都在先前分量的条件下建模，由此得到的条件分布共同定义了一个有效联合分布，其似然可以精确计算。由此，自回归模型在不使用潜变量、变分下界或全局可逆变换的情况下，也构成了一套完整的生成建模框架。它付出的代价不同：必须显式规定顺序，并且采样过程变为顺序执行。

\subsection{作为精确密度模型的链式法则}

设
\[
x=(x_1,x_2,\dots,x_d)
\]
是一个 $d$ 维随机向量。概率链式法则可由恒等式
\[
p(a,b)=p(a\mid b)p(b)
\]
反复应用得到。我们可以从完整联合密度出发，先把最后一个变量分离出来：
\[
p(x_1,\dots,x_d)
=
p(x_d\mid x_1,\dots,x_{d-1})p(x_1,\dots,x_{d-1}).
\]
再对剩余联合项递归应用同一个恒等式，得到
\[
\begin{aligned}
p(x_1,\dots,x_d)
&=
p(x_d\mid x_1,\dots,x_{d-1})
p(x_{d-1}\mid x_1,\dots,x_{d-2})
\cdots
p(x_2\mid x_1)p(x_1).
\end{aligned}
\]
重新排列各个因子，即得到标准的自回归分解：
\[
\boxed{
p(x_1,\dots,x_d)
=
\prod_{i=1}^d p(x_i\mid x_{<i}),
}
\]
其中
\[
x_{<i}:=(x_1,\dots,x_{i-1}),
\]
而当 $i=1$ 时，$p(x_1\mid x_{<1})$ 直接理解为 $p(x_1)$。

这一恒等式是精确的。它本身不包含任何近似，也不施加任何建模假设。只要选定变量顺序，任意联合分布都可以用这种方式分解。

真正的建模步骤在于对条件密度进行参数化：
\[
p_\theta(x)
=
\prod_{i=1}^d p_\theta(x_i\mid x_{<i}).
\]
如果每个条件密度都已归一化，那么它们的乘积就定义了一个正确归一化的联合密度。这正是自回归模型按构造即可获得精确似然的原因。

\begin{tcolorbox}[
  colback=blue!4,
  colframe=blue!50!black,
  title={为什么自回归乘积是有效的概率密度},
  boxrule=0.6pt,
  arc=2mm,
  left=2mm,
  right=2mm,
  top=1mm,
  bottom=1mm
]
假设每个条件密度 $p_\theta(x_i\mid x_{<i})$ 都归一化：
\[
\int p_\theta(x_i\mid x_{<i})\,dx_i=1.
\]
那么联合密度
\[
p_\theta(x_1,\dots,x_d)
=
\prod_{i=1}^d p_\theta(x_i\mid x_{<i})
\]
也同样归一化。从 $x_d$ 开始向前积分，首先有
\[
\int p_\theta(x_d\mid x_{<d})\,dx_d=1,
\]
接着有
\[
\int p_\theta(x_{d-1}\mid x_{<d-1})\,dx_{d-1}=1,
\]
依此类推，直到只剩下 $p_\theta(x_1)$，而它的积分同样为一。

因此，自回归乘积不仅是一种方便的分解形式；只要每个条件项本身是有效的，它就定义了一个有效的概率密度。
\end{tcolorbox}

\subsection{条件参数化与生成}

一旦规定如何表示每个条件密度，链式法则分解就成为一个可学习的生成模型。一般写作
\[
p_\theta(x_i\mid x_{<i}),
\qquad i=1,\dots,d,
\]
其中 $\theta$ 表示模型参数。

需要强调的是，这些条件分布通常并不是由 $d$ 个互不相关的模型分别表示。相反，通常使用一个统一的参数化函数——往往是神经网络——在遵守自回归顺序的同时，输出所有条件分布的参数。整个分解过程中复用的是同一组模型参数；从一个条件分布到下一个条件分布，变化的是条件上下文。

对于连续变量，一个常见而简单的选择是高斯条件分布：
\[
p_\theta(x_i\mid x_{<i})
=
\mathcal{N}\!\left(
x_i;
\mu_i(x_{<i}),
\sigma_i^2(x_{<i})
\right),
\]
其中均值和方差都是先前变量的函数。更灵活的条件分布还可以采用高斯混合、离散化 Logistic 混合或其他富有表达能力的一维分布。

对于离散变量，模型可以输出类别分布的概率：
\[
p_\theta(x_i\mid x_{<i})
=
\operatorname{Categorical}\!\left(\pi_i(x_{<i})\right).
\]
无论哪种情形，模型预测的都是完整的条件分布，而不仅是一个点估计。

\begin{tcolorbox}[
  colback=blue!4,
  colframe=blue!50!black,
  title={三变量自回归示例},
  boxrule=0.6pt,
  arc=2mm,
  left=2mm,
  right=2mm,
  top=1mm,
  bottom=1mm
]
对于 $x=(x_1,x_2,x_3)$，链式法则给出
\[
p_\theta(x_1,x_2,x_3)
=
p_\theta(x_1)\,
p_\theta(x_2\mid x_1)\,
p_\theta(x_3\mid x_1,x_2).
\]

若采用高斯条件分布，则每个条件分布都可以由一个均值和一个方差表示。共享的参数生成函数 $f_\theta$ 可以输出这些参数：
\[
(\mu_1,\sigma_1^2)=f_\theta(\varnothing),
\]
\[
(\mu_2,\sigma_2^2)=f_\theta(x_1),
\]
\[
(\mu_3,\sigma_3^2)=f_\theta(x_1,x_2).
\]

这里 $f_\theta(\varnothing)$ 表示第一个条件分布没有任何先前变量可以作为条件，因此它输出 $p_\theta(x_1)$ 的无条件参数。随后，$f_\theta(x_1)$ 使用已知的 $x_1$ 生成 $p_\theta(x_2\mid x_1)$ 的参数，而 $f_\theta(x_1,x_2)$ 则利用前两个变量生成 $p_\theta(x_3\mid x_1,x_2)$ 的参数。

例如，
\[
p_\theta(x_2\mid x_1)
=
\mathcal{N}\!\left(x_2;\mu_2,\sigma_2^2\right),
\qquad
(\mu_2,\sigma_2^2)=f_\theta(x_1),
\]
同理，
\[
p_\theta(x_3\mid x_1,x_2)
=
\mathcal{N}\!\left(x_3;\mu_3,\sigma_3^2\right),
\qquad
(\mu_3,\sigma_3^2)=f_\theta(x_1,x_2).
\]
\end{tcolorbox}

整个过程中复用的是同一组模型参数 $\theta$。从一个条件分布到下一个条件分布，改变的不是模型本身，而是当前步骤可用的条件信息量。这可以消除一个常见误解：自回归分解虽然列出了一系列条件分布，却并不意味着我们为每个因子训练一个彼此无关的独立模型。在现代自回归模型中，这些条件分布通常由一个具有受限信息流的共享架构统一表示。

同一个分解也自然给出了直接的采样过程。给定
\[
p_\theta(x)
=
\prod_{i=1}^d p_\theta(x_i\mid x_{<i}),
\]
可以按顺序生成样本：
\[
x_1\sim p_\theta(x_1),
\]
然后
\[
x_2\sim p_\theta(x_2\mid x_1),
\]
再然后
\[
x_3\sim p_\theta(x_3\mid x_1,x_2),
\]
依此类推，直到
\[
x_d\sim p_\theta(x_d\mid x_1,\dots,x_{d-1}).
\]
所有分量生成后，所得向量
\[
x=(x_1,\dots,x_d)
\]
就是来自该自回归模型的一个样本。

应把这一过程理解为完整的生成机制。模型并不是先采样一个潜变量再对其解码，也不是运行随机去噪过程，而是直接逐分量构造数据点，并使用已经生成的分量作为后续分量的上下文。

\subsection{并行训练、顺序采样与掩码}

自回归模型具有一种重要的计算不对称性：训练期间的似然计算通常可以并行化，而采样在本质上是顺序的。

给定一个观测数据点
\[
x=(x_1,\dots,x_d),
\]
其对数似然为
\[
\log p_\theta(x)
=
\sum_{i=1}^d
\log p_\theta(x_i\mid x_{<i}).
\]
训练时，$x$ 的全部分量都已知。因此，每个条件上下文 $x_{<i}$ 都可以直接从观测数据中读取，无需采样。模型使用真实的先前分量来计算每个条件项。

这就是自回归似然训练精确且稳定的原因。模型不需要先生成 $x_1$，再用生成出的值来计算 $x_2$，如此继续；相反，它直接针对观测数据计算所有条件概率：
\[
\log p_\theta(x_1),
\quad
\log p_\theta(x_2\mid x_1),
\quad
\dots,
\quad
\log p_\theta(x_d\mid x_{<d}).
\]
具体取决于架构，其中许多乃至全部条件项都可以并行计算。

采样则不同。生成新样本时，未来分量尚未产生，因此是未知的。于是必须按顺序进行：
\[
x_1\sim p_\theta(x_1),
\]
\[
x_2\sim p_\theta(x_2\mid x_1),
\]
\[
\cdots
\]
\[
x_d\sim p_\theta(x_d\mid x_{<d}).
\]
每个新采样得到的分量都会成为下一步条件上下文的一部分，因此采样从构造上就是顺序执行的。

为了在神经网络中实现自回归分解，还必须保证第 $i$ 个条件分布的参数只能依赖 $x_{<i}$，不能依赖 $x_i$ 本身或未来变量。这一限制通常通过\emph{掩码（masking）}来强制实现。掩码会阻断模型中的非法连接，确保信息只能从较早变量流向较晚变量。

因此，即使训练时网络一次处理整个向量 $x$，针对 $x_i$ 输出的参数也不能依赖 $x_i,x_{i+1},\dots,x_d$。计算可以并行，但依赖结构仍然保持有序。

抽象地说，掩码施加了一个下三角依赖模式。如果定义依赖矩阵 $D$，其中 $D_{ij}$ 表示 $x_i$ 的条件分布是否允许依赖 $x_j$，那么自回归顺序要求
\[
D_{ij}=0
\qquad
\text{只要}
\qquad
j\geq i.
\]
例如，
\[
D
=
\begin{pmatrix}
0      & 0      & 0      & \cdots & 0 \\
\ast   & 0      & 0      & \cdots & 0 \\
\ast   & \ast   & 0      & \cdots & 0 \\
\vdots & \vdots & \vdots & \ddots & 0 \\
\ast   & \ast   & \ast   & \cdots & 0
\end{pmatrix},
\]
其中 $\ast$ 表示允许存在的依赖关系。

这种三角依赖是模型在实践中保持可处理性的关键原因之一。它避免了循环依赖：$x_i$ 的分布可以依赖更早的分量，但更早分量不会同时反过来依赖 $x_i$。因此，高维联合似然可以分解为条件对数似然之和。

从这个意义上说，自回归模型中的三角依赖，与基于耦合的归一化流中的三角 Jacobian 起着相似作用。两种情况下，依赖关系都沿单一方向组织，使高维概率计算可以分解为简单项。

掌握这一原则之后，许多著名架构都可以在结构层面得到理解。MADE 通过掩码在前馈网络中强制自回归结构~\cite{germain2015made}；PixelRNN 和 PixelCNN 将自回归分解用于图像~\cite{oord2016pixel}；WaveNet 则利用因果卷积生成原始音频~\cite{oord2016wavenet}。

因此，自回归模型提供了第二条实现精确密度建模的重要路径。归一化流依靠可逆变换和可处理的 Jacobian 行列式获得精确似然；自回归模型则依靠有序条件分解获得精确似然。二者都表明：当高维概率模型中的依赖关系经过精心组织时，其计算仍然可以保持可处理。


\section*{总结与参考文献}
\addcontentsline{toc}{section}{总结与参考文献}

本章介绍了精确密度建模这一与前文潜变量模型和扩散模型互补的范式。在这一范式中，似然可以直接求值，不需要变分下界、随机加噪过程或学习得到的得分场。

我们首先研究了归一化流。其核心思想是使用可微、可逆的映射，把简单的基础分布变换为复杂的数据分布。变量变换公式负责追踪局部体积膨胀或压缩如何改变概率密度，由此引出了流层的主要设计要求：可逆性、维度保持以及可处理的 Jacobian 行列式。

随后我们考察了若干代表性流构造。平面流展示了简单的非线性秩一扰动如何在保持 Jacobian 行列式可处理的同时引入形变。仿射耦合层则提供了更实用的机制：每次只变换输入的一部分，并以其余部分作为条件，从而同时获得闭式逆映射和高效的对数行列式计算。将这些层与坐标混合操作共同堆叠，就可以让简单的局部变换形成具有精确似然计算和精确模型采样能力的高表达生成模型。

本章后半部分介绍了自回归模型，将其作为另一条获得精确似然的路径。自回归模型不通过可逆映射变换密度，而是利用概率链式法则把联合密度分解为条件分布的乘积。这样无需潜变量、变分下界或 Jacobian 行列式，也能获得精确似然。代价是必须规定一个顺序，而且采样通常是顺序执行的。我们还看到，掩码如何在神经网络架构中强制所需的有序信息流，使训练时能够高效计算似然项，同时保持自回归分解成立。

贯穿这两类模型的中心观点是：精确似然之所以成为可能，依赖于结构约束。在归一化流中，这些约束体现为可逆性与可处理的 Jacobian 行列式；在自回归模型中，则体现为有序的条件依赖。两种方法都说明，只要对变换或依赖关系进行精心组织，高维概率模型仍可以保持计算可行。

\bigskip

Rezende 和 Mohamed~\cite{rezende2015variational} 将归一化流作为通过一系列可逆变换构造灵活分布的一般框架加以提出。重要的实用流架构包括 Dinh 等人提出的 NICE 和 RealNVP~\cite{dinh2014nice,dinh2017density}，它们发展了具有可处理 Jacobian 的耦合式变换；以及 Kingma 与 Dhariwal 提出的 Glow~\cite{kingma2018glow}，后者通过可逆 $1\times1$ 卷积和可扩展流设计扩展了这些思想。Papamakarios 等人~\cite{papamakarios2021normalizing} 则给出了关于归一化流及其在密度估计中作用的更广泛综述。

自回归神经密度估计同样在现代生成建模中发挥了重要作用。Germain 等人的 MADE~\cite{germain2015made} 展示了如何通过掩码在前馈神经网络中强制自回归结构。van den Oord 等人的 PixelRNN 和 PixelCNN~\cite{oord2016pixel} 发展了自回归图像模型，而 WaveNet~\cite{oord2016wavenet} 则把因果自回归建模用于原始音频。综合来看，这些工作体现了基于精确似然的建模方法持续的重要性——无论这种精确性来自可逆变换，还是来自有序条件分解。

\chapter{超越似然：生成对抗网络与能量模型}

\section*{概述}
\addcontentsline{toc}{section}{概述}
贯穿本书，我们已经研究了以不同数学方式表示复杂数据分布的生成模型。PPCA 与 VAE 引入了潜变量建模与近似推断；扩散模型把这一视角扩展为逐步加噪与去噪过程；基于得分的模型则进一步用学习到的得分场和随机动力学重新解释去噪。归一化流与自回归模型展示了如何借助精心设计的变换或分解，保留精确似然。

最后一章转向一个不同的问题：如果可处理的似然不再是中心对象，会发生什么？在许多重要的生成模型中，数据分布并不是由一个可以逐点计算的、显式归一化的概率密度来表示。相反，学习通过比较、几何关系或未归一化的标量景观来进行。

生成对抗网络（Generative Adversarial Networks, GAN）首先展示了这种转变。生成器学习产生难以与真实数据区分的样本，而判别器提供比较信号。Wasserstein GAN 则进一步改进了这一对抗思想：它不再使用基于分类的比较，而是采用建立在最优传输基础上的分布几何距离。

能量模型（Energy-Based Models, EBM）提供了另一条超越显式似然的路径。它不直接定义归一化密度，而是为每个配置赋予一个标量能量。低能量标记数据空间中合理的区域，高能量则抑制不合理的配置。因此，学习的任务变成塑造一个能量景观，而不是直接最大化某个可处理的似然。

本章旨在补全本书建立起来的概念版图。GAN、Wasserstein GAN 与能量模型共同表明：即使无法得到归一化概率，生成学习仍然可以进行——通过比较样本、度量几何差异，或者学习其梯度能够引导样本在数据空间中移动的标量景观。

\section*{概念图}
\addcontentsline{toc}{section}{概念图}

本章研究不要求可处理似然计算的生成模型。核心主题是：当概率是隐式的、未归一化的，或者被比较关系所取代时，学习如何继续进行。

\begin{itemize}

    \item[\(\rightarrow\)] \textbf{首先，我们把 GAN 介绍为“通过比较来学习”。}  
    生成器把潜在噪声映射为类似数据的样本，而判别器学习区分生成样本与真实数据。生成器则通过试图让这种比较失效来学习。

    \item[\(\rightarrow\)] \textbf{随后，我们从分布层面解释对抗博弈。}  
    对于固定的生成器，最优判别器估计数据分布与生成器分布之间的局部密度比。把这一判别器代回 GAN 目标函数，会得到 Jensen--Shannon 散度。

    \item[\(\rightarrow\)] \textbf{接着，我们讨论 GAN 优化及其局限。}  
    原始极小极大生成器更新在判别器非常确信时可能产生很弱的梯度。非饱和损失可以改善生成器梯度，但模式坍缩以及基于重叠的比较仍然是重要限制。

    \item[\(\rightarrow\)] \textbf{然后，我们从基于重叠的比较转向 Wasserstein 几何。}  
    Wasserstein GAN 用 $1$-Wasserstein 距离取代基于分类的差异度量，衡量把概率质量从一个分布搬运到另一个分布所需的距离。由此引出评论器（critic）、Lipschitz 约束以及具有几何意义的训练。

    \item[\(\rightarrow\)] \textbf{接下来，本章转向把能量模型理解为标量景观。}  
    能量模型为每个配置赋予标量能量 $E_\theta(x)$。低能量标记数据空间中的合理区域，高能量抑制不合理配置。模型通常定义的是未归一化密度，而不是可处理的似然。

    \item[\(\rightarrow\)] \textbf{随后，我们推导 EBM 学习中的对比结构。}  
    能量模型的最大似然会降低数据样本的能量，同时提高来自模型分布的样本能量。这暴露出核心困难：学习需要从当前模型中进行采样。

    \item[\(\rightarrow\)] \textbf{最后，我们连接能量、得分与扩散。}  
    在未归一化能量模型中，低能量对应高概率，而得分等于负能量梯度。因此，可以把扩散模型看成在多个噪声尺度上学习类似能量梯度的方向，从而提供一条稳定的多尺度路径，绕开显式似然。

\end{itemize}

\noindent 本章讨论的三条“超越似然”路线概括如下。
\begin{center}
\begin{tikzpicture}[
  >=Latex,
  every node/.style={font=\small},
  box/.style={draw,rounded corners=2pt,align=center,inner sep=6pt,
             text width=3.15cm,minimum height=1.02cm},
  flow/.style={->,thick}
]
\node[font=\normalsize\bfseries] at (6.0,2.65)
  {超越似然：概览};

\node[box] (gan) at (1.7,0.85) {
  \textbf{GAN 比较}\\[2pt]
  生成器样本\\
  $\leftrightarrow$ 判别器\\
  反馈
};
\node[box] (wgan) at (6.0,0.85) {
  \textbf{WGAN 几何}\\[2pt]
  传输距离\\
  取代\\
  分类
};
\node[box] (ebm) at (10.3,0.85) {
  \textbf{EBM 景观}\\[2pt]
  能量赋予\\
  相对\\
  合理性
};
\draw[flow] (gan) -- (wgan);
\draw[flow] (wgan) -- (ebm);

\node[box] (ratio) at (1.7,-2.15) {
  \textbf{最优}\\[-1pt]
  \textbf{判别器}\\[2pt]
  局部密度比\\
  信息
};
\node[box] (energy) at (6.0,-2.15) {
  \textbf{未归一化密度}\\[2pt]
  $p_\theta(x)\propto e^{-E_\theta(x)}$\\
  全局\\
  配分函数
};
\node[box] (score) at (10.3,-2.15) {
  \textbf{回到得分}\\[2pt]
  $\nabla_x\log p_\theta(x)$\\
  $=-\nabla_xE_\theta(x)$
};
\draw[flow] (gan) -- (ratio);
\draw[flow] (ebm) -- (energy);
\draw[flow] (energy) -- (score);
\end{tikzpicture}
\end{center}

\section{生成对抗网络：通过比较来学习}

到目前为止介绍的许多生成模型都依赖显式概率分布。无论使用精确似然还是精心构造的近似，学习通常都被表述为给观测数据赋予高概率的问题。这提供了原则清晰的基础，但也对模型如何设计与训练施加了强结构约束。

于是自然会产生一个更简单的问题：
\[
\textit{为了生成逼真的数据，我们真的必须计算概率吗？}
\]
如果目标只是生成看起来像来自数据分布的样本，那么也许通过比较进行学习就足够了。我们不再询问一个样本有多大似然，而是询问它能否与真实数据区分。这样，学习就通过匹配样本来进行，而不是通过计算密度。

生成对抗网络（GAN）~\cite{goodfellow2014generative} 正是围绕这一思想构建。它通过生成令人信服的样本这一能力来定义生成模型，而无须规定或计算概率密度。训练由比较驱动：生成样本与真实数据相互对抗，模型不断学习，使二者越来越难以区分。

\subsection{生成器作为隐式采样机制}

GAN 从一个\emph{生成器（generator）}开始，它是参数化函数
\[
x = G_\theta(z),
\qquad z \sim p(z),
\]
其中 $z$ 是从简单固定分布（如标准高斯分布或均匀分布）中抽取的潜变量。生成器 $G_\theta$ 通常由神经网络实现，把这一简单随机源映射到数据空间。

映射
\[
z \longmapsto G_\theta(z)
\]
把简单噪声转化为有结构、类似真实数据的样本。对 $z\sim p(z)$ 采样并送入 $G_\theta$，生成器会在 $x$ 上诱导出一个分布，记为 $p_\theta(x)$。关键在于，这个分布只被\emph{隐式}定义：我们可以从中采样，但通常无法计算给定点处的密度值。

这使 GAN 与前文基于似然的模型有所区别。与 VAE 不同，潜变量 $z$ 不是从数据中推断得到的，也不学习后验分布；与扩散模型不同，$z$ 并非由某个前向加噪过程产生；与归一化流或自回归模型不同，生成器不要求可逆，也无须计算 Jacobian 行列式或显式归一化 $p_\theta(x)$。因此，生成器更像一个确定性模拟器：它规定\emph{如何生成数据}，却不规定\emph{某个给定数据点有多大可能出现}。

训练完成后，从 GAN 采样非常直接。先从潜在噪声分布中抽取 $z\sim p(z)$，再计算 $x=G_\theta(z)$。生成阶段不再需要判别器；判别器只在训练时提供塑造生成器的信号。因此，GAN 能够直接采样，却不能直接计算似然。

\subsection{判别器与对抗博弈}

如果模型只会生成样本，就立即面临一个基本问题：
\[
\textit{我们如何判断这些样本是否足够好？}
\]
在基于似然的模型中，概率回答了这个问题：如果样本在模型下具有高似然，就认为它们是好的。而在 GAN 中，似然不可获得，没有一个可处理的 $p_\theta(x)$ 可供计算或最大化。

为了解决这一问题，GAN 引入第二个学习组件：\emph{判别器（discriminator）}。判别器并不测量概率密度，而是对两种样本来源进行数据驱动的比较：
\[
x\sim p_{\text{data}}(x),
\qquad
x=G_\theta(z),\quad z\sim p(z).
\]
第一种来源是真实数据样本，第二种来源是生成样本。

形式上，判别器是函数
\[
D_\phi:\mathbb{R}^d\to(0,1),
\]
由参数 $\phi$ 决定。对于输入 $x$，$D_\phi(x)$ 被解释为判别器认为 $x$ 来自真实数据分布而非生成器的置信度。因此，判别器被训练为二分类器：
\[
D_\phi(x)\approx 1
\quad\text{对真实数据，}
\qquad
D_\phi(G_\theta(z))\approx 0
\quad\text{对生成数据。}
\]

判别器本身并不是最终目的。它既不定义概率密度，也不是训练结束后我们真正想使用的对象。它的作用，是在真实样本与生成样本之间提供一个不断变化、自适应的差异度量。生成器完全通过这一反馈来学习，不断调整参数，使生成样本越来越难与真实样本区分。

这种相互作用由极小极大目标形式化：
\[
\boxed{
\min_\theta \max_\phi
\;\;
\mathbb{E}_{x \sim p_{\text{data}}}
\big[\log D_\phi(x)\big]
+
\mathbb{E}_{z \sim p(z)}
\big[\log (1 - D_\phi(G_\theta(z)))\big].
}
\]
判别器试图最大化该目标：给真实样本高分、给生成样本低分。生成器则试图最小化它，使 $D_\phi(G_\theta(z))$ 变大，从而让判别器难以把生成样本判为伪造。

\begin{tcolorbox}[
  colback=blue!4,
  colframe=blue!50!black,
  title={如何理解极小极大目标},
  boxrule=0.6pt,
  arc=2mm,
  left=2mm,
  right=2mm,
  top=1mm,
  bottom=1mm
]
表达式
\[
\min_\theta \max_\phi \; \mathcal{L}(\theta,\phi)
\]
应被理解为一个双人博弈。

在生成器参数 $\theta$ 固定时，内层最大化询问：
\[
\textit{面对当前生成器，最优判别器是什么？}
\]
这一步训练 $D_\phi$ 区分真实样本与生成样本。

随后，把判别器视为固定，外层最小化询问：
\[
\textit{生成器应如何改变，才能让真实样本与生成样本更难区分？}
\]
实践中，这两步交替执行：先改进判别器，再改进生成器。
\end{tcolorbox}

尽管这个目标看起来类似分类损失，其作用本质上是生成性的。判别器在两个分布之间定义一个学习到的比较，而生成器不断重塑自己的输出分布，使这种比较逐渐失去区分能力。

\subsection{最优判别器}

为了从分布层面理解对抗训练，考虑一个理想化情形：生成器 $G_\theta$ 固定。对 $z\sim p(z)$ 采样并令 $x=G_\theta(z)$，会诱导出隐式生成器分布 $p_\theta(x)$。在生成器固定后，判别器的任务就退化为在
\[
p_{\text{data}}(x)
\qquad\text{与}\qquad
p_\theta(x)
\]
之间进行二分类。

判别器目标可以写成
\[
\mathbb{E}_{x \sim p_{\text{data}}}\!\big[\log D(x)\big]
+
\mathbb{E}_{x \sim p_\theta}\!\big[\log(1 - D(x))\big].
\]
把期望写成积分，得到
\[
\int
p_{\text{data}}(x)\log D(x)\,dx
+
\int
p_\theta(x)\log(1-D(x))\,dx.
\]
由于目标函数对每个 $x$ 都可以分解，因此最优判别器可以逐点推导。对固定的 $x$，我们关于标量 $D(x)$ 最大化
\[
p_{\text{data}}(x)\log D(x)
+
p_\theta(x)\log(1-D(x)).
\]

对 $D(x)$ 求导，得到
\[
\frac{p_{\text{data}}(x)}{D(x)}
-
\frac{p_\theta(x)}{1-D(x)}.
\]
令导数为零，得到
\[
\frac{p_{\text{data}}(x)}{D(x)}
=
\frac{p_\theta(x)}{1-D(x)}.
\]
解出 $D(x)$，得到最优判别器
\[
\boxed{
D^*(x)
=
\frac{p_{\text{data}}(x)}
{p_{\text{data}}(x)+p_\theta(x)}.
}
\]

这一公式给出一个重要解释：即使两种密度都没有显式提供给模型，最优判别器仍然是一个局部密度比对象。图~\ref{fig:gan-optimal-discriminator-geometry} 直观展示了这一点：数据密度占优处，其输出趋近于一；生成器密度占优处，其输出趋近于零；两种密度相等处，输出恰好为二分之一。

\begin{figure}[H]
\centering
\includegraphics[width=0.98\textwidth]{figures/ch08/figure8_1_optimal_discriminator.pdf}
\caption{\textbf{最优判别器反映局部密度的相对平衡。}
对于固定生成器，$D^*(x)=p_{\mathrm{data}}(x)/(p_{\mathrm{data}}(x)+p_\theta(x))$。在数据密度占优处它趋近于一，在生成器密度占优处趋近于零，而在两种密度相等处恰为 $1/2$。}
\label{fig:gan-optimal-discriminator-geometry}
\end{figure}

\subsection{从最优判别器到 Jensen--Shannon 散度}

得到固定生成器下的最优判别器之后，我们可以进一步问：整个对抗博弈在分布层面究竟代表什么？为此，把判别器目标写成一个\emph{价值函数（value function）}，它通过判别器区分真实样本与生成样本的能力，为一个生成器--判别器组合打分。

定义
\[
V(G_\theta,D)
=
\mathbb{E}_{x \sim p_{\text{data}}}\!\big[\log D(x)\big]
+
\mathbb{E}_{x \sim p_\theta}\!\big[\log(1-D(x))\big].
\]
对于固定的 $G_\theta$，关于 $D$ 最大化 $V(G_\theta,D)$ 就得到针对该生成器的最优判别器。然后在这个最优判别器处计算价值函数，就可以识别生成器实际上在隐式最小化哪一种分布差异。

利用
\[
D^*(x)
=
\frac{p_{\text{data}}(x)}
{p_{\text{data}}(x)+p_\theta(x)},
\]
经过代数化简，得到
\[
V(G_\theta,D^*)
=
-\log 4
+
2\,\mathrm{JS}\!\left(p_{\text{data}}\,\|\,p_\theta\right),
\]
其中 $\mathrm{JS}(\cdot\|\cdot)$ 表示 Jensen--Shannon 散度。

常数项 $-\log 4$ 与生成器无关，因此不影响优化。所以，在判别器达到最优的理想化情形下，GAN 训练等价于最小化数据分布与生成器分布之间的 Jensen--Shannon 散度。

上式引入了概率分布之间一种新的散度：\emph{Jensen--Shannon 散度}。要理解它，可以回忆 VAE 章节在证据下界中出现过的 Kullback--Leibler（KL）散度。KL 散度按如下方式比较两个分布：
\[
\mathrm{KL}(p\,\|\,q)
=
\int p(x)\log\frac{p(x)}{q(x)}\,dx.
\]
它衡量用 $q$ 近似 $p$ 所付出的代价。然而，KL 散度是不对称的：
\[
\mathrm{KL}(p\,\|\,q)
\neq
\mathrm{KL}(q\,\|\,p),
\]
并且对支撑集不匹配十分敏感。如果在某一区域 $p(x)>0$ 而 $q(x)=0$，那么 $\mathrm{KL}(p\,\|\,q)$ 会变为无穷大。

Jensen--Shannon 散度通过让两个分布都与一个共享的中间分布进行比较来修改这一方式。给定两个分布 $p$ 与 $q$，定义它们的混合分布
\[
m(x)
=
\frac{1}{2}\big(p(x)+q(x)\big).
\]
可以把这个混合分布理解为：先以相同概率在 $p$ 与 $q$ 之间选择一个分布，再从被选中的分布中采样。

Jensen--Shannon 散度定义为
\[
\boxed{
\mathrm{JS}(p\,\|\,q)
=
\frac{1}{2}\mathrm{KL}(p\,\|\,m)
+
\frac{1}{2}\mathrm{KL}(q\,\|\,m).
}
\]
因此，它不要求一个分布直接解释另一个分布，而是询问两个分布分别与它们共同的混合分布相距多远。这使其具有对称性：
\[
\mathrm{JS}(p\,\|\,q)
=
\mathrm{JS}(q\,\|\,p),
\]
并且有界：
\[
0
\leq
\mathrm{JS}(p\,\|\,q)
\leq
\log 2.
\]

这正是把最优判别器代回原始 GAN 价值函数后出现的散度。因此，在判别器最优的理想化条件下，生成器并不是在最大化似然；它是在调整自己的隐式分布 $p_\theta$，以减小
\[
\mathrm{JS}\!\left(p_{\text{data}}\,\|\,p_\theta\right).
\]

这揭示了对抗训练在分布层面的含义。GAN 并不构造显式密度，而是执行\emph{通过比较进行的分布匹配}。判别器估计真实分布与生成分布有多容易区分，生成器则通过降低这种可区分性来学习。

全局最优点满足
\[
p_\theta(x)=p_{\text{data}}(x).
\]
此时，处处都有
\[
D^*(x)=\frac{1}{2},
\]
因为真实样本与生成样本已经无法区分。

\section{几何优化与 Wasserstein GAN}

前面的分析表明，在最优判别器下，原始 GAN 目标对应于最小化数据分布与生成器隐式分布之间的 Jensen--Shannon 散度。这给出了清晰的理论解释，同时也暴露出一个局限：当 $p_{\text{data}}$ 与 $p_\theta$ 很少甚至完全不重叠时，JS 散度可能进入饱和状态。此时判别器几乎可以完美地区分真实样本与生成样本，却几乎不能告诉生成器应当如何向数据分布移动。这促成了一个关键转变：从基于分布重叠的比较，转向几何距离。Wasserstein GAN~\cite{arjovsky2017wasserstein} 不再询问两个分布“是否可区分”，而是询问：为了把一个分布变成另一个分布，需要把概率质量移动多远。

\subsection{生成器优化与非饱和损失}

在转向 Wasserstein GAN 之前，我们先讨论原始 GAN 框架中的一种常用实践改进：改变生成器损失，以改善训练时的梯度信号。
原始 GAN 目标写作极小极大问题
\[
\min_\theta \max_\phi
\;
\mathbb{E}_{x \sim p_{\text{data}}}\big[\log D_\phi(x)\big]
+
\mathbb{E}_{z \sim p(z)}\big[\log (1 - D_\phi(G_\theta(z)))\big].
\]
对判别器而言，其解释很直接：真实样本得分越高、生成样本得分越低，奖励越大。但对生成器而言，情况更微妙。

若固定判别器，并关于生成器参数 $\theta$ 最小化同一个价值函数，唯一依赖 $\theta$ 的项是
\[
\mathbb{E}_{z \sim p(z)}
\big[\log (1 - D_\phi(G_\theta(z)))\big].
\]
这一目标要求生成器降低判别器认为生成样本是伪造数据的置信度。

困难常出现在训练早期。此时判别器通常可以非常自信地识别生成样本：
\[
D_\phi(G_\theta(z))\approx 0.
\]
于是
\[
1-D_\phi(G_\theta(z))\approx 1,
\qquad
\log(1-D_\phi(G_\theta(z)))\approx 0.
\]
生成器目标在这一区域因此接近平坦，即使生成样本仍然很差，反向传给生成器的梯度也可能非常小。

因此，GAN 实践中常采用另一个生成器目标，称为\emph{非饱和损失（non-saturating loss）}：
\[
\min_\theta
\;
\mathbb{E}_{z \sim p(z)}
\big[
-\log D_\phi(G_\theta(z))
\big].
\]
这个损失并不是对前一个表达式做代数化简得到的，而是把生成器的目标从降低判别器的“伪造”置信度，改为提高其“真实”置信度。

两种目标不同，但期望达到的平衡点相同：生成样本最终应与真实数据无法区分。它们真正的差异在于梯度行为。当
\[
D_\phi(G_\theta(z))\approx 0
\]
时，非饱和项
\[
-\log D_\phi(G_\theta(z))
\]
很大，因此生成器能获得更强的学习信号。这就是实践中通常更偏好非饱和损失的原因，尽管原始极小极大目标与 Jensen--Shannon 散度之间的理论联系更加直接。

这种优化行为有助于理解 GAN 的优势与弱点。一方面，判别器会惩罚让生成样本显得不真实的特征，因此对抗训练能够鼓励清晰、视觉上更可信的输出，而不是模糊的平均结果。另一方面，同一机制也可能导致\emph{模式坍缩（mode collapse）}。由于生成器的目标是欺骗判别器，它可能只把数据分布中的一部分模式建模得很好。只要这一子集足以迷惑判别器，其他模式就可能被忽略。

基于似然的模型如果不能覆盖观测到的模式，会直接受到惩罚，因为似然要求模型给数据分配概率质量。相比之下，GAN 优化的是可区分性，而不是显式的概率覆盖。因此，非饱和损失虽然改善了生成器梯度，却没有改变对抗比较的底层几何。当两个分布几乎不重叠时，基于分类的差异度量可能只告诉我们它们“可区分”，却没有告诉我们一个分布应该朝哪里移动才能接近另一个分布。这就需要一种更具几何意义的分布差异度量。

\subsection{从重叠关系到 Wasserstein 几何}

为了发展这一思想，可以从几何角度理解概率分布。不要把分布只看成静态密度，而把其概率质量想象成可以在空间中搬运的物质。如果把一个分布变成另一个分布只需要做很小的质量位移，那么两个分布应被认为很接近；如果必须把大量质量搬运很远，它们就应被认为相距很远。

Wasserstein 距离把这种直觉形式化。在 Wasserstein GAN 中，相关的是 $1$-Wasserstein 距离，也称 Earth Mover's distance（推土机距离）。它衡量把概率质量从一个分布搬运到另一个分布的最小代价，其中代价与搬运距离成正比。

一个关键理论结果——Kantorovich--Rubinstein 对偶性——给出了这种距离的实用表达。对于 $\mathbb{R}^d$ 上两个概率分布 $p$ 与 $q$，$1$-Wasserstein 距离可写为
\[
\boxed{
W(p,q)
=
\sup_{\|f\|_{L} \le 1}
\left(
\mathbb{E}_{x \sim p}[f(x)]
-
\mathbb{E}_{x \sim q}[f(x)]
\right).
}
\]
这里，$f$ 是实值评分函数，而约束
\[
\|f\|_{L}\leq 1
\]
表示 $f$ 必须是 $1$-Lipschitz 的：其函数值相对于输入空间中的距离，变化速度不能超过单位速率。

表达式
\[
\mathbb{E}_{x \sim p}[f(x)]
-
\mathbb{E}_{x \sim q}[f(x)]
\]
衡量该评分函数在平均意义下把两个分布区分得有多强。上确界意味着：在满足 Lipschitz 约束的所有评分函数中，选择区分能力最强的那个。

这一约束至关重要。没有它，评分函数可以给一个分布中的样本任意大的值、给另一个分布中的样本任意小的值，从而使比较失去意义。Lipschitz 约束迫使评分差异反映数据空间的几何结构：只有当对应点彼此相距较远时，才允许出现很大的评分差异。

需要注意，这里的“score”不同于基于得分的扩散模型中的得分。在后者中，得分指 $\nabla_x\log p(x)$，即对数密度的梯度；而在 WGAN 中，评论器输出 $f_\phi(x)$ 是用于比较的实值标量，而不是概率。对评论器施加的是 Lipschitz 行为约束，而非概率归一化约束。

\begin{tcolorbox}[
  colback=blue!4,
  colframe=blue!50!black,
  title={Wasserstein 距离的一个计算示例},
  boxrule=0.6pt,
  arc=2mm,
  left=2mm,
  right=2mm,
  top=1mm,
  bottom=1mm
]
考虑一维例子：$p$ 把全部概率质量放在 $x=0$，而 $q$ 把全部概率质量放在 $x=10$。从 $p$ 采样总会得到 $0$，从 $q$ 采样总会得到 $10$，于是
\[
\mathbb{E}_{x\sim p}[f(x)] = f(0),
\qquad
\mathbb{E}_{x\sim q}[f(x)] = f(10),
\qquad
W(p,q)
=
\sup_{\|f\|_L\leq 1}
\big(f(0)-f(10)\big).
\]

约束 $\|f\|_L\leq 1$ 意味着
\[
|f(x_1)-f(x_2)|\leq |x_1-x_2|
\quad
\text{对任意 }x_1,x_2.
\]
因此，每个合法的 $f$ 都满足
\[
f(0)-f(10)\leq |0-10|=10.
\]
函数 $f(x)=-x$ 达到这一上界，因为
\[
f(0)-f(10)=0-(-10)=10.
\]
所以 $W(p,q)=10$，与搬运直觉完全一致：全部概率质量都必须移动距离 $10$。
\end{tcolorbox}

这个例子说明了为什么在分布完全不重叠时，Wasserstein 距离仍然具有信息量。基于重叠的散度能够检测支撑集彼此分离，却未必能反映两个分布相距一个单位还是十个单位。相比之下，Wasserstein 距离会直接度量空间分离程度。图~\ref{fig:gan-js-wasserstein-geometry} 对一族平移高斯分布展示了相同差异：随着重叠消失，JS 散度很快接近饱和，而 Wasserstein 距离仍持续追踪空间上的分离程度。

\begin{figure}[H]
\centering
\begin{minipage}[t]{0.48\textwidth}
  \vspace{0pt}\centering
  \includegraphics[width=\linewidth]{figures/ch08/panels/figure8_2_a.pdf}
\end{minipage}\hfill
\begin{minipage}[t]{0.48\textwidth}
  \vspace{0pt}\centering
  \includegraphics[width=\linewidth]{figures/ch08/panels/figure8_2_b.pdf}
\end{minipage}

\vspace{0.8em}

\begin{minipage}[t]{0.48\textwidth}
  \vspace{0pt}\centering
  \includegraphics[width=\linewidth]{figures/ch08/panels/figure8_2_c.pdf}
\end{minipage}\hfill
\begin{minipage}[t]{0.48\textwidth}
  \vspace{0pt}
  \caption{\textbf{JS 散度会饱和，而 Wasserstein 距离仍保留几何信息。}
  各子图比较方差相同、均值逐渐分离的高斯分布。随着重叠消失，JS 散度逼近上界 $\log 2$，对进一步分离逐渐不敏感。对于这些平移分布，$1$-Wasserstein 距离为 $W_1=\Delta$，因此仍持续度量概率质量需要搬运多远。}
  \label{fig:gan-js-wasserstein-geometry}
\end{minipage}
\end{figure}


\subsection{Wasserstein GAN 目标}

Wasserstein 距离的对偶表征直接导出一个可实践的对抗目标。如果生成器诱导出分布 $p_\theta$，那么数据分布与生成器分布之间的 Wasserstein 距离可以写成
\[
W(p_{\text{data}},p_\theta)
=
\sup_{\|f\|_{L}\leq 1}
\left(
\mathbb{E}_{x\sim p_{\text{data}}}[f(x)]
-
\mathbb{E}_{x\sim p_\theta}[f(x)]
\right).
\]
由于 $p_\theta$ 中的样本通过
\[
x=G_\theta(z),
\qquad z\sim p(z)
\]
生成，因此可以写成
\[
\mathbb{E}_{x\sim p_\theta}[f(x)]
=
\mathbb{E}_{z\sim p(z)}[f(G_\theta(z))].
\]

Wasserstein GAN 用神经网络 $f_\phi$ 表示评分函数 $f$，并优化
\[
\boxed{
\min_\theta \max_{\phi \in \mathcal{L}}
\;
\mathbb{E}_{x \sim p_{\text{data}}}[f_\phi(x)]
-
\mathbb{E}_{z \sim p(z)}[f_\phi(G_\theta(z))],
}
\]
其中 $\mathcal{L}$ 表示所有使 $f_\phi$ 满足所需 Lipschitz 约束的参数集合。

关于 $\phi$ 的最大化试图找到一个评分函数：在尊重数据空间几何结构的前提下，尽可能强地区分真实样本与生成样本。随后，关于 $\theta$ 的最小化更新生成器，使这种分离越来越小。

到这里，术语也随之变化。在原始 GAN 中，负责评价样本的网络称为\emph{判别器}，其输出
\[
D_\phi(x)\in(0,1)
\]
被解释为 $x$ 是真实样本的概率。而在 Wasserstein GAN 中，对应的网络称为\emph{评论器（critic）}。它输出实值分数
\[
f_\phi(x)\in\mathbb{R},
\]
而不是概率。这里没有 sigmoid 输出，不作二分类解释，也不要求分数归一化；真正有意义的是评论器分数之间的差异。

从这个角度看，评论器是 Wasserstein 距离中最优评分函数的神经网络近似。它的目的不是作出一个硬性的真假判断，而是赋予样本分数，使其平均差异反映两个分布之间的几何分离程度。在分数不能跨数据空间任意快速变化的约束下，真实样本的平均得分应高于生成样本。由于 Wasserstein 距离会随着生成器分布的移动而连续变化，即使 $p_\theta$ 与 $p_{\text{data}}$ 没有重叠，生成器仍能获得有用梯度；因此，优化由几何邻近性而非仅由重叠程度来引导。

与原始 GAN 一样，评论器只用于训练。生成器训练完成后，采样仍然只是抽取 $z\sim p(z)$ 并返回 $G_\theta(z)$。

\subsection{施加 Lipschitz 约束}

剩下的实际问题是：训练期间如何对评论器施加 Lipschitz 约束？回忆这一约束的作用：让评论器分数以受控速率变化，使分数差异反映几何距离，而不是任意缩放。普通神经网络并不会自动满足这一性质。

最初的 WGAN 通过\emph{权重裁剪（weight clipping）}施加该约束，把评论器的每个参数限制在固定区间内。这种方法简单，却是一种非常粗糙的 Lipschitz 常数控制方式，可能削弱评论器容量，并导致优化不稳定。

更有效的替代方法是 WGAN-GP~\cite{gulrajani2017improved} 引入的\emph{梯度惩罚（gradient penalty）}。它不直接限制参数，而是限制评论器的输入--输出行为。其核心思想是控制评论器分数对输入微小变化的敏感程度。

在评论器训练目标中加入如下惩罚项：
\[
\lambda
\,
\mathbb{E}_{\hat{x}}
\left[
\big(\|\nabla_{\hat{x}} f_\phi(\hat{x})\|_2-1\big)^2
\right],
\]
其中 $\hat{x}$ 从真实样本与生成样本之间的直线插值点上采样。

梯度
\[
\nabla_{\hat{x}} f_\phi(\hat{x})
\]
衡量评论器分数对输入无穷小变化的敏感程度。因此，其欧氏范数就是评论器在 $\hat{x}$ 处的局部变化率，也就是局部斜率。惩罚 $\|\nabla_{\hat{x}} f_\phi(\hat{x})\|_2$ 偏离一的程度，会鼓励评论器在最相关的比较区域中平滑、受控地改变分数。

插值点 $\hat{x}$ 位于真实样本与生成样本之间，而这里恰好是评论器必须学习如何比较两个分布的区域。沿这些路径施加约束，既能稳定训练，又能保留 Wasserstein 距离的几何解释。系数 $\lambda$ 控制惩罚强度：较大的值会更强地施加约束，较小的值则允许更大灵活性。

通过这种方式施加 Lipschitz 约束，带梯度惩罚的 WGAN 在保留 Wasserstein 距离最优传输含义的同时，获得了更可靠的训练动力学。更广泛地说，它们说明生成学习不必依赖显式似然或归一化密度，也可以由一个学习得到的实值函数来引导；该函数的几何性质会把生成样本推向数据分布。

这自然形成通往能量模型的桥梁。与 WGAN 类似，能量模型也利用实值函数比较、排序或引导样本。但它不通过对抗博弈来表述这一思想，而是通过未归一化的能量景观定义分布。下面转向这一视角。

\section{能量模型：基础与学习}

前面几节讨论了通过比较进行学习的生成模型。在 GAN 中，判别器提供一个学习得到的测试，用来区分真实样本与生成样本；在 Wasserstein GAN 中，这种比较被几何评论器所取代，其函数值反映把概率质量搬运多远才能完成分布变换。两种方法都绕开显式似然计算，转而依赖样本之间的相对评价进行学习。

能量模型~\cite{lecun2006tutorial} 则把这一思想进一步推进。它不从生成器、判别器或显式似然出发，而是从数据空间上的单个标量函数开始：\emph{能量函数（energy function）}。该函数为每个配置 $x$ 赋予一个值
\[
E_\theta(x)\in\mathbb{R}.
\]
较低能量表示该配置与数据结构更相容，较高能量表示该配置不太合理。

能量函数有意不施加太多限制：它无须为正、无须有界，也无须归一化，而且它本身并不定义概率分布。因此，在这一阶段不能把能量直接解释为概率。更合适的理解是：它是一种偏好、代价或相容性度量。类似真实数据的配置应位于低能量区域，而不真实的配置则应位于高能量区域。

这给出了一种几何化的建模视角。能量模型在数据空间上定义一个景观：
\[
\text{低能量}
\quad\Longleftrightarrow\quad
\text{合理配置},
\qquad
\text{高能量}
\quad\Longleftrightarrow\quad
\text{不合理配置}.
\]
学习意味着塑造这一景观，使观测数据处于能量谷底，而远离数据的区域被赋予较高能量。

\subsection{作为建模基本对象的能量}

能量模型中的核心建模对象是标量能量函数
\[
E_\theta:\mathbb{R}^d\to\mathbb{R}.
\]
与归一化流或自回归模型不同，它不需要给出精确密度；与 VAE 不同，它不要求显式潜变量后验；与 GAN 不同，它也不需要在对抗博弈中训练一个独立判别器。模型直接由一个标量景观规定。

只有相对能量差异真正重要。如果给能量函数加上常数 $c$，
\[
\tilde{E}_\theta(x)=E_\theta(x)+c,
\]
所有配置的能量都移动同样的幅度，从低到高的相对排序保持不变。由于模型行为依赖“哪些配置相对于其他配置更受偏好”，这种整体常数平移没有实质影响。

这种相对性十分重要。能量函数并不是试图为每个点赋予绝对概率值，而是在编码偏好：哪些配置应当被鼓励，哪些应当被抑制。从这个意义上说，能量建模提出的是一个不同于显式密度建模的问题：
\[
\textit{模型应当把低能量区域放在数据空间的什么位置？}
\]
答案由数据学习得到。

\subsection{从能量到未归一化分布}

尽管能量函数本身不是概率分布，但可以通过 Boltzmann 形式把它与概率联系起来：
\[
p_\theta(x)
=
\frac{1}{Z_\theta}
\exp\big(-E_\theta(x)\big),
\]
其中
\[
Z_\theta
=
\int
\exp\big(-E_\theta(x)\big)\,dx
\]
称为归一化常数或\emph{配分函数（partition function）}。

负指数把低能量转化为较大的正权重，而配分函数执行一次全局归一化。图~\ref{fig:ebm-energy-to-density} 展示了从标量能量到未归一化权重，再到归一化密度的映射；它也清楚说明了为什么给所有能量加同一个常数不会改变最终密度。

\begin{figure}[H]
\centering
\begin{minipage}[t]{0.48\textwidth}
  \vspace{0pt}\centering
  \includegraphics[width=\linewidth]{figures/ch08/panels/figure8_3_a.pdf}
\end{minipage}\hfill
\begin{minipage}[t]{0.48\textwidth}
  \vspace{0pt}\centering
  \includegraphics[width=\linewidth]{figures/ch08/panels/figure8_3_b.pdf}
\end{minipage}

\vspace{0.8em}

\begin{minipage}[t]{0.48\textwidth}
  \vspace{0pt}\centering
  \includegraphics[width=\linewidth]{figures/ch08/panels/figure8_3_c.pdf}
\end{minipage}\hfill
\begin{minipage}[t]{0.48\textwidth}
  \vspace{0pt}
  \caption{\textbf{从能量到归一化密度。}
  低能量区域获得较大的未归一化权重 $\exp(-E_\theta(x))$；除以 $Z_\theta$ 后即可变成概率密度而不改变其形状。给所有能量加同一个常数，会按相同比例缩放权重与 $Z_\theta$，因此归一化后的密度保持不变。}
  \label{fig:ebm-energy-to-density}
\end{minipage}
\end{figure}


困难在于，高维空间中的 $Z_\theta$ 通常不可处理。计算它要求对所有可能配置进行积分：
\[
Z_\theta
=
\int
\exp\big(-E_\theta(x)\big)\,dx.
\]
对于图像、音频、文本表示或其他高维数据，这个积分一般不可能精确求得。

这正是能量建模的核心权衡。由于避免了精确密度模型所需的结构限制，能量函数可以非常灵活；但这种灵活性以归一化困难为代价。模型通常只定义\emph{未归一化}密度
\[
p_\theta(x)
\propto
\exp\big(-E_\theta(x)\big),
\]
而不是可处理的归一化似然。

这使能量模型不同于上一章讨论的精确密度模型。归一化流与自回归模型经过精心设计，使 $p_\theta(x)$ 可以被精确计算；能量模型则把重点放在能量景观本身。即使归一化常数未知，相对能量差异和能量梯度仍然具有实际用途。

\subsection{学习能量景观}

能量模型中的学习，就是塑造能量函数，使其反映数据结构。假设我们观察到样本
\[
x\sim p_{\text{data}}(x).
\]
理想的能量函数应给这些观测样本较低能量：
\[
E_\theta(x)
\quad\text{较小，当}\quad
x\sim p_{\text{data}}.
\]
但仅满足这一条件还不够。如果我们只是把所有位置的能量都降低，模型就无法区分类似数据的配置与不合理配置。因此，学习必须是对比性的：既要降低数据样本的能量，也要提高其他位置的能量。

非正式地说，理想行为是
\[
E_\theta(x)
\quad\text{在数据上较低},
\qquad
E_\theta(x)
\quad\text{远离数据时较高}.
\]
正是这种对比结构赋予能量学习以意义。模型不仅要奖励观测数据，还必须学习哪些备选配置应该被抑制。

理解这一点的一个有用工具是能量梯度
\[
\nabla_x E_\theta(x).
\]
它描述能量在数据空间局部如何变化。沿
\[
-\nabla_x E_\theta(x)
\]
方向移动会降低能量，从而把一个配置引向模型认为更合理的区域。因此，能量景观不仅定义标量偏好值，还定义局部的改进方向。

这把能量模型与前文 GAN 和基于得分模型中的几何视角联系起来。在 Wasserstein GAN 中，评论器赋予实值分数，分数差异用于引导生成器；在能量模型中，能量函数赋予标量值，其梯度定义配置应当如何在数据空间中移动。两者相关但并不相同：WGAN 评论器为了区分两个分布而训练，而能量函数的目标是定义一个合理性的景观。

\subsection{最大似然与对比梯度}

尽管能量模型通常作为未归一化模型来讨论，Boltzmann 形式仍然给出了一个原则明确的学习目标。如果
\[
p_\theta(x)
=
\frac{1}{Z_\theta}\exp(-E_\theta(x)),
\]
那么最大似然试图最大化
\[
\mathbb{E}_{x\sim p_{\text{data}}}
\big[\log p_\theta(x)\big].
\]
代入能量形式的密度，得到
\[
\log p_\theta(x)
=
-E_\theta(x)-\log Z_\theta,
\]
因此
\[
\mathbb{E}_{x\sim p_{\text{data}}}
\big[\log p_\theta(x)\big]
=
-\mathbb{E}_{x\sim p_{\text{data}}}
\big[E_\theta(x)\big]
-
\log Z_\theta.
\]

第一项的作用很直观：最大化它会降低数据样本的能量。第二项则阻止“到处降低能量”这一平凡解，因为给过大的空间区域分配低能量会使配分函数增大。因此，最大似然平衡两股力量：
\[
\text{降低数据处的能量}
\qquad
\text{与}
\qquad
\text{控制其他区域的能量}.
\]

关于 $\theta$ 对目标函数求导，这一点会更加清楚。数据项的梯度为
\[
-\nabla_\theta
\mathbb{E}_{x\sim p_{\text{data}}}
\big[E_\theta(x)\big]
=
-
\mathbb{E}_{x\sim p_{\text{data}}}
\big[\nabla_\theta E_\theta(x)\big].
\]
对配分函数项，使用
\[
Z_\theta
=
\int
\exp(-E_\theta(x))\,dx.
\]
求导得到
\[
\begin{aligned}
\nabla_\theta \log Z_\theta
&=
\frac{1}{Z_\theta}
\nabla_\theta
\int
\exp(-E_\theta(x))\,dx \\[4pt]
&=
\frac{1}{Z_\theta}
\int
\exp(-E_\theta(x))
\big(-\nabla_\theta E_\theta(x)\big)\,dx \\[4pt]
&=
-
\mathbb{E}_{x\sim p_\theta}
\big[\nabla_\theta E_\theta(x)\big].
\end{aligned}
\]
由于对数似然中包含 $-\log Z_\theta$，其梯度贡献为
\[
-\nabla_\theta \log Z_\theta
=
\mathbb{E}_{x\sim p_\theta}
\big[\nabla_\theta E_\theta(x)\big].
\]

合并两项，最大似然梯度为
\[
\boxed{
\nabla_\theta
\mathbb{E}_{x\sim p_{\text{data}}}
\big[\log p_\theta(x)\big]
=
-
\mathbb{E}_{x\sim p_{\text{data}}}
\big[\nabla_\theta E_\theta(x)\big]
+
\mathbb{E}_{x\sim p_\theta}
\big[\nabla_\theta E_\theta(x)\big].
}
\]

这一表达式是能量学习的核心。它包含两个部分：
\[
\text{数据项}
\quad
-
\mathbb{E}_{p_{\text{data}}}
\big[\nabla_\theta E_\theta(x)\big],
\qquad
\text{模型项}
\quad
+
\mathbb{E}_{p_\theta}
\big[\nabla_\theta E_\theta(x)\big].
\]

数据项降低观测样本的能量，模型项提高从当前模型分布抽取的样本能量。这通常被描述为\emph{正相位（positive phase）}与\emph{负相位（negative phase）}之间的对比：正相位使数据更受偏好，而负相位则对模型当前分配了过多概率质量的配置施加反向压力。

这一对比梯度暴露出能量模型最核心的实际困难。对数据分布的期望可以直接用训练集近似；而对模型分布的期望
\[
\mathbb{E}_{x\sim p_\theta}
\big[\nabla_\theta E_\theta(x)\big]
\]
要困难得多，因为从 $p_\theta(x)\propto \exp(-E_\theta(x))$ 本身采样就是一个非平凡问题。

这使学习与采样耦合在一起。为了改进能量函数，我们需要当前模型的样本；但要生成好的模型样本，又需要一个已经被良好塑造的能量景观。这种循环依赖，是能量模型虽然强大却难以训练的主要原因之一。

下一节将更明确地展开这一联系。由于能量模型本身不提供直接采样器，生成通常被表述为在能量景观上的移动。这自然导向能量梯度动力学、Langevin 采样，以及能量建模与前文基于得分扩散方法之间的联系。

\section{能量梯度、得分场与扩散}

上一节表明，能量模型通过未归一化的能量景观间接定义概率分布：
\[
p_\theta(x)
\propto
\exp(-E_\theta(x)).
\]
这带来很强的建模灵活性，却也产生一个实际困难：模型本身不提供直接采样机制。为了生成样本，我们必须在能量景观中移动。

这正是能量梯度成为核心对象的地方。标量能量 $E_\theta(x)$ 告诉我们一个配置有多合理，而其梯度
\[
\nabla_x E_\theta(x)
\]
则告诉我们能量在数据空间局部如何变化。沿
\[
-\nabla_x E_\theta(x)
\]
方向移动会降低能量，因此把配置推向模型认为更加合理的区域。从这个意义上说，负能量梯度就是一种局部驱动力。

本节把这一思想与本书前面建立的基于得分方法连接起来。核心结论简单但十分重要：对于形如 $p(x)\propto \exp(-E(x))$ 的分布，其得分恰好就是负能量梯度：
\[
\nabla_x\log p(x)
=
-\nabla_x E(x).
\]
因此，能量模型与基于得分模型关系非常紧密。二者真正的差异，与其说是所使用的局部方向不同，不如说在于该方向如何参数化、如何学习以及如何用于生成。

\subsection{把采样看作能量引导的动力学}

能量模型给合理配置较低能量，但并不会直接告诉我们如何从对应分布中抽取相互独立的样本。与归一化流不同，它不存在从基础分布到数据空间的可逆映射；与自回归模型不同，它也没有可以从左到右顺序采样的条件分布乘积。因此，采样通常被表述为一个迭代动力学过程。

基本思想是先初始化一个配置 $x_0$，然后反复更新，使其向低能量区域移动。纯确定性的更新可以写为
\[
x_{t+1}
=
x_t
-
\eta\nabla_x E_\theta(x_t),
\]
其中 $\eta>0$ 是步长。这很像梯度下降，但此时“下降”发生在数据空间中，而不是参数空间中；被移动的是配置 $x_t$ 本身，使其走向更低能量。

然而，仅有确定性下降还不足以构成采样。它可能陷入能量景观的某个局部极小点，而无法探索多个合理区域。为了得到采样过程，需要加入随机性：
\[
\boxed{
x_{t+1}
=
x_t
-
\eta \nabla_x E_\theta(x_t)
+
\sqrt{2\eta}\,\xi_t,
\qquad
\xi_t\sim\mathcal{N}(0,I).
}
\]
这就是与能量函数相对应的离散 Langevin 型更新。

确定性项
\[
-\eta\nabla_x E_\theta(x_t)
\]
推动当前配置走向低能量，而噪声项
\[
\sqrt{2\eta}\,\xi_t
\]
注入随机性，防止动力学简单坍缩到单一极小点，并使其能够探索更广泛的能量景观。在适当条件与合适的连续时间解释下，这种“梯度加噪声”的动力学可以以能量模型分布为平稳分布。

这正是前面基于得分章节研究过的 Langevin 思想。改变的只是漂移场的写法：在那里，漂移使用得分 $\nabla_x\log p(x)$ 表示；这里则使用负能量梯度 $-\nabla_xE_\theta(x)$ 表示。下一小节会明确给出二者的联系。

\subsection{能量梯度与得分}

回忆密度 $p(x)$ 的得分定义为
\[
s(x)
=
\nabla_x\log p(x).
\]
对于能量模型，
\[
p_\theta(x)
=
\frac{1}{Z_\theta}\exp(-E_\theta(x)).
\]
取对数得到
\[
\log p_\theta(x)
=
-E_\theta(x)-\log Z_\theta.
\]
现在对 $x$ 求导。由于配分函数 $Z_\theta$ 依赖 $\theta$，但不依赖 $x$，所以
\[
\nabla_x\log Z_\theta=0.
\]
因此
\[
\boxed{
\nabla_x\log p_\theta(x)
=
-\nabla_xE_\theta(x).
}
\]

这一恒等式是能量建模与基于得分建模之间的桥梁。能量分布的得分场恰好等于能量函数的负梯度。因此，朝高概率区域移动的局部方向，与朝低能量区域移动的方向完全相同。

区别在于对象如何表示。在经典能量模型中，我们参数化一个标量函数 $E_\theta(x)$，再通过求导得到驱动向量场：
\[
x
\quad\longmapsto\quad
-\nabla_xE_\theta(x).
\]
而在基于得分的模型中，通常直接参数化向量场：
\[
s_\theta(x)
\approx
\nabla_x\log p(x).
\]
因此，能量模型从标量景观出发并导出向量场，而得分模型往往直接学习向量场本身。

这一差别很有意义。一个标量能量天然会产生保守梯度场：向量场来源于同一个底层标量函数的梯度。直接学习的得分网络作为向量场更加灵活，但除非额外施加结构，它未必被显式表示为某个标量能量的梯度。不过，就生成建模而言，二者在概念上的联系依然很强：它们都在学习把样本引向更高数据密度区域的局部方向。

\subsection{把得分匹配理解为能量梯度学习}

上一节表明，能量模型的最大似然会产生一个包含模型分布样本的对比梯度。这造成实际困难，因为从模型本身采样就很难。得分匹配（score matching）~\cite{hyvarinen2005estimation} 提供了另一条路径：不去匹配归一化概率，而是匹配局部得分场。

在基于得分的章节中，我们已经利用 Fisher 散度和分部积分详细推导了得分匹配。这里仅回顾理解能量视角所需的结果。目标是学习模型得分
\[
s_\theta(x)
\approx
\nabla_x\log p_{\text{data}}(x).
\]
对于能量模型，
\[
s_\theta(x)
=
\nabla_x\log p_\theta(x)
=
-\nabla_xE_\theta(x).
\]
因此，可以把得分匹配解释为学习一个能量景观，使其负梯度与数据得分一致。

忽略与 $\theta$ 无关的常数项，经典得分匹配目标可以写为
\[
J_{\mathrm{SM}}(\theta)
=
\mathbb{E}_{x\sim p_{\text{data}}}
\left[
\frac{1}{2}\|s_\theta(x)\|_2^2
+
\nabla_x\cdot s_\theta(x)
\right].
\]
代入
\[
s_\theta(x)=-\nabla_xE_\theta(x)
\]
得到能量形式的目标：
\[
\boxed{
J_{\mathrm{SM}}(\theta)
=
\mathbb{E}_{x\sim p_{\text{data}}}
\left[
\frac{1}{2}
\|\nabla_xE_\theta(x)\|_2^2
-
\nabla_x^2E_\theta(x)
\right],
}
\]
其中
\[
\nabla_x^2E_\theta(x)
=
\sum_{j=1}^d
\frac{\partial^2E_\theta(x)}{\partial x_j^2}
\]
是能量函数的 Laplacian。

这里最重要的不是具体推导——前文已经完成了这些细节——而是“被匹配的对象”发生了变化。最大似然比较归一化概率分布，并会产生一个模型期望项；得分匹配比较对数密度的局部梯度，并绕开配分函数，因为
\[
\nabla_x\log Z_\theta=0.
\]
用能量的语言说，它通过导数学习数据样本附近能量景观的形状，而无须精确计算归一化密度。

这就是为什么从能量视角看，得分匹配如此自然。它不询问某个数据点的概率值，而是询问：概率应当沿哪个局部方向增大？等价地，能量应当沿哪个方向减小？

\subsection{把扩散理解为多尺度能量梯度学习}

现在可以清楚表述能量模型与基于得分扩散之间的关系。能量模型定义标量景观 $E_\theta(x)$。如果该景观对应未归一化密度
\[
p_\theta(x)\propto \exp(-E_\theta(x)),
\]
那么其得分为
\[
\nabla_x\log p_\theta(x)
=
-\nabla_xE_\theta(x).
\]
因此，Langevin 采样中使用的局部方向，与基于得分模型试图直接学习的对象完全相同：指向更高概率、也就是更低能量的方向。

从这一视角出发，可以把基于得分的扩散模型~\cite{song2019generative,song2021scorebased} 看作在一系列含噪分布上学习随时间变化的“类能量梯度”方向。
真正不同的并不是局部几何本身，而是这种几何如何表示和学习。能量模型参数化标量能量，再通过求导得到向量场；基于得分模型通常直接参数化向量场；扩散模型则更进一步，学习时间依赖的得分场
\[
s_\theta(x,t)
\approx
\nabla_x\log p_t(x),
\]
其中 $p_t(x)$ 是数据经过加噪直到时间 $t$ 后得到的分布。

从能量视角看，每个含噪分布 $p_t$ 都可以与一个隐式能量景观联系起来：
\[
E_t(x)
=
-\log p_t(x)+C_t,
\]
其中 $C_t$ 是任意常数。由于只有梯度起作用，该常数没有影响，于是
\[
\nabla_x\log p_t(x)
=
-\nabla_xE_t(x).
\]
因此，可以把扩散模型解释为学习
\[
s_\theta(x,t)
\approx
-\nabla_xE_t(x),
\]
即一整个由噪声水平索引的能量景观族的负梯度。

这个视角之所以有用，是因为前向加噪过程创造了一套多尺度分布层级。在高噪声水平下，$p_t$ 平滑并接近简单分布，因此对应的能量景观宽缓而规则；在低噪声水平下，$p_t$ 更接近数据分布，因此能量景观包含更精细的结构。扩散训练就是在学习这条序列的局部几何：
\[
\text{粗尺度能量景观}
\quad\longrightarrow\quad
\text{精细的数据级能量景观}.
\]

从这个意义上说，扩散模型可以被看作一种实用且稳定的多尺度能量梯度学习形式。它不是一次性学习一个结构极其尖锐的数据分布能量景观，再直接从中采样；而是把问题分散到许多经过平滑的分布上。高噪声层级提供全局引导，低噪声层级提供细致的局部修正。

这也解释了为什么扩散模型能够绕开传统能量学习中的一些困难。经典 EBM 最大似然需要从当前模型分布采样，以估计负相位；扩散模型则通过已知的前向扰动过程构造含噪训练对，并跨不同噪声层级学习去噪方向。所得方法依然与能量梯度有深刻联系，但训练过程稳定得多。

因此，能量模型为这一部分生成建模提供了更广阔的“景观”视角：它通过标量相容性函数描述分布；而基于得分的扩散，则可以看作现代生成建模中一种尤其有效的方法，用来沿从噪声到数据的多尺度路径，学习并利用相应的能量梯度方向。

\section*{总结与参考文献}
\addcontentsline{toc}{section}{总结与参考文献}

本章超越基于似然的生成建模，介绍了通过比较、几何以及未归一化景观进行学习的模型。GAN 展示了生成器如何通过对抗比较而非似然最大化进行训练。最优判别器把原始 GAN 目标与 Jensen--Shannon 散度联系起来，而实践中通常采用非饱和损失，为生成器提供更强的梯度。

Wasserstein GAN 进一步发展了这一对抗视角，用分布之间的几何距离取代基于重叠的分类。通过对偶形式，可以利用实值的 Lipschitz 评论器估计 $1$-Wasserstein 距离。这个评论器不输出概率；其分数差异反映概率质量如何在数据空间中搬运。

能量模型则提供了更一般的景观视角。EBM 为每个配置赋予标量能量 $E_\theta(x)$，低能量对应合理数据，高能量对应不合理配置。通过 $p_\theta(x)\propto \exp(-E_\theta(x))$，能量定义了未归一化密度。因此，学习与采样围绕能量差异和能量梯度展开，而不是围绕可处理的归一化似然展开。

最后，我们利用恒等式
\[
\nabla_x\log p_\theta(x)
=
-\nabla_xE_\theta(x)
\]
把 EBM 与基于得分的扩散联系起来。这表明能量模型与得分模型共享相同的局部几何：二者都使用把样本引向更高密度、较低能量区域的方向。扩散模型把这一思想扩展到一系列含噪分布，从粗糙的高噪声层级一直学习到精细的数据级结构。

\bigskip

生成对抗网络由 Goodfellow 等人~\cite{goodfellow2014generative} 提出，他们把生成表述为生成器与判别器之间的极小极大博弈，并把最优判别器与 Jensen--Shannon 散度联系起来。Arjovsky 等人~\cite{arjovsky2017wasserstein} 提出了 Wasserstein GAN，以基于最优传输的距离取代基于分类的比较；Gulrajani 等人~\cite{gulrajani2017improved} 随后发展了梯度惩罚形式。

能量建模在机器学习和统计物理中有很长的发展历史。LeCun 等人~\cite{lecun2006tutorial} 从“塑造能量景观”的角度给出了广泛教程。经典例子包括 Hopfield 网络~\cite{hopfield1982neural} 和 Boltzmann 机~\cite{hinton1985learning}。Hyv\"arinen~\cite{hyvarinen2005estimation} 提出得分匹配，用于在不计算配分函数的情况下估计未归一化统计模型；Vincent~\cite{vincent2011connection} 则把得分匹配与去噪联系起来。前面的基于得分扩散章节已经详细讨论了 Song 与 Ermon~\cite{song2019generative} 以及 Song 等人~\cite{song2021scorebased} 的工作；这里再次回顾它们，是因为这些工作把能量梯度思想与多尺度得分学习及扩散模型连接了起来。

\chapter*{结语}
\addcontentsline{toc}{chapter}{结语}

本书有意选择了一条有所取舍的路径，来梳理生成建模的数学基础。我们从线性表示开始：把矩阵理解为变换，把 PCA 理解为保持方差的投影，并把自编码器理解为基于重构的压缩。随后，概率主成分分析把这一确定性图景转化为潜变量生成模型，并由此引入先验、似然、边缘似然和后验这些核心建模对象。

变分自编码器把这一潜变量视角推广到了非线性情形，此时精确的边缘似然和后验已不再可得。这自然引出了证据下界、变分推断、近似后验以及重参数化技巧。扩散模型则进一步把生成过程从单一潜变量转变为一整条带噪潜在状态轨迹，使生成成为逐步加噪过程的一个学习得到的逆过程。

连续时间和基于得分的章节从第二种视角考察了同一个广义问题。我们不再只把生成描述为离散链，而是研究运动、密度演化、布朗噪声、Liouville 方程、Fokker--Planck 方程、Langevin 动力学、得分匹配以及逆时间扩散。这说明，生成建模也可以被理解为学习概率空间中的运动方向。

后续章节考察了两条互补路径。归一化流和自回归模型展示了如何通过结构设计保留精确似然：要么使用具有易于计算的 Jacobian 的可逆变换，要么采用有序的条件分解。GAN、Wasserstein GAN 和能量模型则说明，生成学习也可以超越易于计算的似然，转而依赖比较、几何差异或标量能量景观。

综合来看，这些章节把生成式人工智能呈现为一组反复出现的数学思想，而不是一堆彼此无关的架构。这些思想包括：表示、变换、边缘化、近似、加噪与去噪、密度演化、得分场、似然构造、分布比较以及能量塑形。它们也仍在不断涌现的新型生成建模框架中反复出现。

例如，流匹配（flow matching）~\cite{lipman2023flow} 和整流流（rectified flows）~\cite{liu2022rectified} 可以看作连续时间输运视角的进一步发展：它们学习一个向量场，将概率质量从简单分布输运到数据分布。一致性模型（consistency models）~\cite{song2023consistency} 则可看作对扩散模型采样效率问题的另一种回应，其目标是在减少采样步数的同时，尽可能保留迭代去噪所带来的生成质量。

因此，本书的目的，是建立一套紧凑的数学基础，使读者能够更加从容地理解这些后续发展。一旦掌握潜变量、变分界、扩散路径、密度演化、得分场、精确似然模型、对抗式比较以及能量景观，许多现代扩展方法就更容易被放置在同一张概念导图中理解。

\appendix
\chapter{高斯代数与配方法}
\label{app:completing-square}

第 2 章通过直接的线性高斯论证推导了 PPCA 的边缘分布 \(p(\mathbf{x})\)。对于该章的主线而言，这种方法最为快捷、透明。不过，还有另一种推导方式值得考察，因为同样的代数结构会在后续其他高斯潜变量模型中反复出现，即\emph{配方法}（completing the square）。

本附录集中整理若干依赖显式高斯代数的推导，尤其着重把正文中隐式使用的配方恒等式说明清楚。在 PPCA 情形下，这一方法的价值在于：它能够在同一套推导中同时得到两个关键的高斯量：
\begin{itemize}
    \item 后验分布 \(p(\mathbf{z}\mid \mathbf{x})\)；
    \item 以及对 \(\mathbf{z}\) 积分消去之后得到的边缘分布 \(p(\mathbf{x})\)。
\end{itemize}

\paragraph{配方法。}
高斯密度是二次型的指数函数。因此，只要某个变量的对数密度可以写成该变量的二次函数，我们就可以尝试把它改写成标准高斯形式。配方法正是完成这种改写的代数工具：它把二次表达式重新组织为
\[
-\frac12 (\text{变量} - \text{均值})^\top
(\text{精度矩阵})
(\text{变量} - \text{均值})
+\text{常数},
\]
于是便可以直接读出高斯分布的均值和协方差。

为了便于后续复用，有必要显式给出这一技巧的矩阵形式。若 \(A\) 为对称正定矩阵，则
\[
-\frac12 \mathbf{x}^\top A \mathbf{x} + \mathbf{b}^\top \mathbf{x}
=
-\frac12 (\mathbf{x}-A^{-1}\mathbf{b})^\top A (\mathbf{x}-A^{-1}\mathbf{b})
+\frac12 \mathbf{b}^\top A^{-1}\mathbf{b}.
\]
这是因为将右侧展开可得
\[
\begin{aligned}
&-\frac12 (\mathbf{x}-A^{-1}\mathbf{b})^\top A (\mathbf{x}-A^{-1}\mathbf{b})
+\frac12 \mathbf{b}^\top A^{-1}\mathbf{b} \\[4pt]
&\qquad=
-\frac12 \mathbf{x}^\top A \mathbf{x}
+\mathbf{b}^\top \mathbf{x}
-\frac12 \mathbf{b}^\top A^{-1}\mathbf{b}
+\frac12 \mathbf{b}^\top A^{-1}\mathbf{b} \\[4pt]
&\qquad=
-\frac12 \mathbf{x}^\top A \mathbf{x}+\mathbf{b}^\top \mathbf{x}.
\end{aligned}
\]
因此，这两个表达式完全等价；第二个只是把第一个改写成以高斯分布中心为中心的形式。向量 \(A^{-1}\mathbf{b}\) 是该二次型的中心，而 \(A\) 是精度矩阵，其逆矩阵给出协方差。

\paragraph{步骤 1：从联合密度开始。}
回顾 PPCA 模型：
\[
p(\mathbf{z})=\mathcal{N}(\mathbf{z}\mid \mathbf{0},\mathbf{I}),
\qquad
p(\mathbf{x}\mid \mathbf{z})
=
\mathcal{N}(\mathbf{x}\mid \mathbf{Wz}+\boldsymbol{\mu},\alpha^2\mathbf{I}).
\]
因此
\[
p(\mathbf{x},\mathbf{z})
=
p(\mathbf{x}\mid \mathbf{z})\,p(\mathbf{z}).
\]

由于高斯密度是二次型的指数函数，因此处理联合对数密度更加方便：
\[
\log p(\mathbf{x},\mathbf{z})
=
-\frac{1}{2\alpha^2}\|\mathbf{x}-\boldsymbol{\mu}-\mathbf{Wz}\|^2
-\frac12 \|\mathbf{z}\|^2
+\text{常数}.
\]

\paragraph{步骤 2：展开关于 \(\mathbf{z}\) 的二次项。}
展开第一个范数得
\[
\|\mathbf{x}-\boldsymbol{\mu}-\mathbf{Wz}\|^2
=
(\mathbf{x}-\boldsymbol{\mu})^\top(\mathbf{x}-\boldsymbol{\mu})
-2\,\mathbf{z}^\top \mathbf{W}^\top(\mathbf{x}-\boldsymbol{\mu})
+\mathbf{z}^\top \mathbf{W}^\top\mathbf{W}\mathbf{z}.
\]
代回可得
\[
\log p(\mathbf{x},\mathbf{z})
=
-\frac{1}{2\alpha^2}(\mathbf{x}-\boldsymbol{\mu})^\top(\mathbf{x}-\boldsymbol{\mu})
+\frac{1}{\alpha^2}\mathbf{z}^\top\mathbf{W}^\top(\mathbf{x}-\boldsymbol{\mu})
-\frac12 \mathbf{z}^\top
\Bigl(\frac{1}{\alpha^2}\mathbf{W}^\top\mathbf{W}+\mathbf{I}\Bigr)\mathbf{z}
+\text{常数}.
\]

现在把 \(\mathbf{x}\) 视为固定。此时联合对数密度是 \(\mathbf{z}\) 的二次函数，因此可以通过配方法把它改写为关于 \(\mathbf{z}\) 的高斯形式。

\paragraph{步骤 3：配方。}
定义
\[
\mathbf{M}
=
\mathbf{W}^\top\mathbf{W}+\alpha^2\mathbf{I},
\qquad
\mathbf{m}
=
\mathbf{M}^{-1}\mathbf{W}^\top(\mathbf{x}-\boldsymbol{\mu}).
\]
由于
\[
\frac{1}{\alpha^2}\mathbf{W}^\top\mathbf{W}+\mathbf{I}
=
\frac{1}{\alpha^2}\mathbf{M},
\]
联合对数密度中依赖于 \(\mathbf{z}\) 的部分变为
\[
-\frac{1}{2\alpha^2}\mathbf{z}^\top\mathbf{M}\mathbf{z}
+
\frac{1}{\alpha^2}\mathbf{z}^\top\mathbf{W}^\top(\mathbf{x}-\boldsymbol{\mu}).
\]
此时它与上面的通用恒等式完全同形，其中
\[
A=\frac{1}{\alpha^2}\mathbf{M},
\qquad
\mathbf{b}=\frac{1}{\alpha^2}\mathbf{W}^\top(\mathbf{x}-\boldsymbol{\mu}).
\]
应用该恒等式得到
\[
-\frac{1}{2\alpha^2}\mathbf{z}^\top\mathbf{M}\mathbf{z}
+
\frac{1}{\alpha^2}\mathbf{z}^\top\mathbf{W}^\top(\mathbf{x}-\boldsymbol{\mu})
=
-\frac{1}{2\alpha^2}(\mathbf{z}-\mathbf{m})^\top\mathbf{M}(\mathbf{z}-\mathbf{m})
+\frac{1}{2\alpha^2}\mathbf{m}^\top\mathbf{M}\mathbf{m}.
\]
将其代回可得
\[
\log p(\mathbf{x},\mathbf{z})
=
-\frac{1}{2\alpha^2}(\mathbf{z}-\mathbf{m})^\top\mathbf{M}(\mathbf{z}-\mathbf{m})
-\frac{1}{2\alpha^2}(\mathbf{x}-\boldsymbol{\mu})^\top(\mathbf{x}-\boldsymbol{\mu})
+\frac{1}{2\alpha^2}\mathbf{m}^\top\mathbf{M}\mathbf{m}
+\text{常数}.
\]

到这里，关于 \(\mathbf{z}\) 的高斯结构已经显式呈现出来。

\paragraph{为什么这能够确定一个高斯后验。}
一旦 \(\mathbf{x}\) 固定，所有与 \(\mathbf{z}\) 无关的项都可以并入归一化常数。因此，真正需要关注的唯一 \(\mathbf{z}\) 相关项是
\[
-\frac{1}{2\alpha^2}(\mathbf{z}-\mathbf{m})^\top\mathbf{M}(\mathbf{z}-\mathbf{m}),
\]
它恰好就是关于 \(\mathbf{z}\) 的高斯密度的对数核。由于 \(\mathbf{M}\) 为正定矩阵，该二次函数关于 \(\mathbf{z}\) 是凹的，因此它确实定义了一个适当的高斯分布。

\paragraph{步骤 4：读出后验 \(p(\mathbf{z}\mid \mathbf{x})\)。}
因此
\[
p(\mathbf{z}\mid \mathbf{x})
=
\mathcal{N}\!\bigl(
\mathbf{z}\mid
\mathbf{m},
\alpha^2\mathbf{M}^{-1}
\bigr),
\]
即
\[
p(\mathbf{z}\mid \mathbf{x})
=
\mathcal{N}\!\bigl(
\mathbf{z}\mid
\mathbf{M}^{-1}\mathbf{W}^\top(\mathbf{x}-\boldsymbol{\mu}),
\alpha^2\mathbf{M}^{-1}
\bigr),
\qquad
\mathbf{M}=\mathbf{W}^\top\mathbf{W}+\alpha^2\mathbf{I}.
\]

所以，配方法直接给出了后验分布：其均值是潜变量估计 \(\mathbf{m}\)，协方差则为 \(\alpha^2\mathbf{M}^{-1}\)。

\paragraph{步骤 5：对 \(\mathbf{z}\) 积分消去以得到 \(p(\mathbf{x})\)。}
边缘分布为
\[
p(\mathbf{x})
=
\int p(\mathbf{x},\mathbf{z})\,d\mathbf{z}.
\]
从配方后的形式可见，依赖于 \(\mathbf{z}\) 的部分是一个以 \(\mathbf{m}\) 为中心的关于 \(\mathbf{z}\) 的高斯密度。因此，对 \(\mathbf{z}\) 积分会消去这一部分，只留下一个归一化常数。剩余部分是 \((\mathbf{x}-\boldsymbol{\mu})\) 的二次函数，所以边缘分布也必然是高斯分布。

继续完成余下的代数运算可得
\[
p(\mathbf{x})
=
\mathcal{N}\!\bigl(
\mathbf{x}\mid
\boldsymbol{\mu},
\mathbf{W}\mathbf{W}^\top+\alpha^2\mathbf{I}
\bigr).
\]

\bigskip
因此，只需完成一次配方，我们便同时得到
\[
p(\mathbf{z}\mid \mathbf{x})
=
\mathcal{N}\!\bigl(
\mathbf{z}\mid
\mathbf{M}^{-1}\mathbf{W}^\top(\mathbf{x}-\boldsymbol{\mu}),
\alpha^2\mathbf{M}^{-1}
\bigr)
\]
以及
\[
p(\mathbf{x})
=
\mathcal{N}\!\bigl(
\mathbf{x}\mid
\boldsymbol{\mu},
\mathbf{W}\mathbf{W}^\top+\alpha^2\mathbf{I}
\bigr).
\]

对于 PPCA 的主要概念主线而言，这一推导并非必需；但它仍然很有用，因为此后每当需要显式处理高斯潜变量模型时，同样的二次项配方技巧都会再次出现。

\chapter{扩散模型逆后验推导}
\label{app:diffusion-posterior}

在第 4.4 节中，我们指出真实逆后验
\[
q(\mathbf{x}_{t-1}\mid \mathbf{x}_t,\mathbf{x}_0)
\]
是高斯分布，并给出了其均值与方差的闭式表达式。这里给出的推导建立在与此前 PPCA 中相同的基本代数思想之上：一旦所有相关项都是高斯的，就可以把关于目标变量的二次项收集起来，再通过配方法得到结果。

\subsection*{步骤 1：从贝叶斯公式开始}

根据前向过程，有
\[
q(\mathbf{x}_{t-1}\mid \mathbf{x}_t,\mathbf{x}_0)
=
\frac{
q(\mathbf{x}_t\mid \mathbf{x}_{t-1},\mathbf{x}_0)\,
q(\mathbf{x}_{t-1}\mid \mathbf{x}_0)
}{
q(\mathbf{x}_t\mid \mathbf{x}_0)
}.
\]
由于分母不依赖于 \(\mathbf{x}_{t-1}\)，当把该式视为关于 \(\mathbf{x}_{t-1}\) 的函数时，分母只起归一化常数的作用。因此，在比例意义下，
\[
q(\mathbf{x}_{t-1}\mid \mathbf{x}_t,\mathbf{x}_0)
\propto
q(\mathbf{x}_t\mid \mathbf{x}_{t-1},\mathbf{x}_0)\,
q(\mathbf{x}_{t-1}\mid \mathbf{x}_0).
\]

由于前向过程满足 Markov 性，
\[
q(\mathbf{x}_t\mid \mathbf{x}_{t-1},\mathbf{x}_0)
=
q(\mathbf{x}_t\mid \mathbf{x}_{t-1}),
\]
因此两个高斯因子分别为
\[
q(\mathbf{x}_t\mid \mathbf{x}_{t-1})
=
\mathcal{N}\!\bigl(
\mathbf{x}_t\mid
\sqrt{\alpha_t}\,\mathbf{x}_{t-1},
(1-\alpha_t)\mathbf{I}
\bigr),
\]
以及
\[
q(\mathbf{x}_{t-1}\mid \mathbf{x}_0)
=
\mathcal{N}\!\bigl(
\mathbf{x}_{t-1}\mid
\sqrt{\bar{\alpha}_{t-1}}\,\mathbf{x}_0,
(1-\bar{\alpha}_{t-1})\mathbf{I}
\bigr).
\]

\subsection*{步骤 2：把乘积写成指数形式}

忽略与 \(\mathbf{x}_{t-1}\) 无关的归一化常数，有
\[
\begin{aligned}
q(\mathbf{x}_{t-1}\mid \mathbf{x}_t,\mathbf{x}_0)
&\propto
\exp\!\left(
-\frac{1}{2(1-\alpha_t)}
\left\|
\mathbf{x}_t-\sqrt{\alpha_t}\,\mathbf{x}_{t-1}
\right\|^2
\right)
\\[4pt]
&\qquad\times
\exp\!\left(
-\frac{1}{2(1-\bar{\alpha}_{t-1})}
\left\|
\mathbf{x}_{t-1}-\sqrt{\bar{\alpha}_{t-1}}\,\mathbf{x}_0
\right\|^2
\right).
\end{aligned}
\]
合并两个指数项得到
\[
\begin{aligned}
q(\mathbf{x}_{t-1}\mid \mathbf{x}_t,\mathbf{x}_0)
\propto
\exp\!\Bigg(
&-\frac{1}{2(1-\alpha_t)}
\left\|
\mathbf{x}_t-\sqrt{\alpha_t}\,\mathbf{x}_{t-1}
\right\|^2
\\
&-\frac{1}{2(1-\bar{\alpha}_{t-1})}
\left\|
\mathbf{x}_{t-1}-\sqrt{\bar{\alpha}_{t-1}}\,\mathbf{x}_0
\right\|^2
\Bigg).
\end{aligned}
\]

\subsection*{步骤 3：展开关于 \texorpdfstring{$\mathbf{x}_{t-1}$}{x\_\{t-1\}} 的二次项}

现在展开两个平方范数，只保留包含 \(\mathbf{x}_{t-1}\) 的项。

对于第一项，
\[
\begin{aligned}
\left\|
\mathbf{x}_t-\sqrt{\alpha_t}\,\mathbf{x}_{t-1}
\right\|^2
&=
\mathbf{x}_t^\top\mathbf{x}_t
-2\sqrt{\alpha_t}\,\mathbf{x}_t^\top \mathbf{x}_{t-1}
+\alpha_t\,\mathbf{x}_{t-1}^\top \mathbf{x}_{t-1}.
\end{aligned}
\]
对于第二项，
\[
\begin{aligned}
\left\|
\mathbf{x}_{t-1}-\sqrt{\bar{\alpha}_{t-1}}\,\mathbf{x}_0
\right\|^2
&=
\mathbf{x}_{t-1}^\top \mathbf{x}_{t-1}
-2\sqrt{\bar{\alpha}_{t-1}}\,\mathbf{x}_{t-1}^\top \mathbf{x}_0
+\bar{\alpha}_{t-1}\,\mathbf{x}_0^\top\mathbf{x}_0.
\end{aligned}
\]

将这些结果代入指数部分，并去掉与 \(\mathbf{x}_{t-1}\) 无关的项，得到
\[
\begin{aligned}
q(\mathbf{x}_{t-1}\mid \mathbf{x}_t,\mathbf{x}_0)
\propto
\exp\!\Bigg(
&-\frac{1}{2}
\left[
\frac{\alpha_t}{1-\alpha_t}
+
\frac{1}{1-\bar{\alpha}_{t-1}}
\right]
\mathbf{x}_{t-1}^\top \mathbf{x}_{t-1}
\\
&+
\left[
\frac{\sqrt{\alpha_t}}{1-\alpha_t}\,\mathbf{x}_t^\top
+
\frac{\sqrt{\bar{\alpha}_{t-1}}}{1-\bar{\alpha}_{t-1}}\,\mathbf{x}_0^\top
\right]
\mathbf{x}_{t-1}
\Bigg).
\end{aligned}
\]

此时它已经具有标准的高斯二次型形式
\[
-\frac12 A\,\mathbf{x}_{t-1}^\top \mathbf{x}_{t-1}
+
\mathbf{b}^\top \mathbf{x}_{t-1},
\]
其中
\[
A
=
\frac{\alpha_t}{1-\alpha_t}
+
\frac{1}{1-\bar{\alpha}_{t-1}},
\qquad
\mathbf{b}
=
\frac{\sqrt{\alpha_t}}{1-\alpha_t}\,\mathbf{x}_t
+
\frac{\sqrt{\bar{\alpha}_{t-1}}}{1-\bar{\alpha}_{t-1}}\,\mathbf{x}_0.
\]

\subsection*{步骤 4：配方}

利用标准恒等式
\[
-\frac12 A\,\mathbf{x}^\top\mathbf{x}
+
\mathbf{b}^\top\mathbf{x}
=
-\frac12 A
\left\|
\mathbf{x}-A^{-1}\mathbf{b}
\right\|^2
+\text{常数},
\]
可以看出，该后验是高斯分布，其协方差系数为
\[
\tilde{\beta}_t = A^{-1},
\]
均值为
\[
\tilde{\boldsymbol{\mu}}_t = A^{-1}\mathbf{b}.
\]

因此先化简 \(A\)：
\[
\begin{aligned}
A
&=
\frac{\alpha_t}{1-\alpha_t}
+
\frac{1}{1-\bar{\alpha}_{t-1}}
\\[4pt]
&=
\frac{
\alpha_t(1-\bar{\alpha}_{t-1}) + (1-\alpha_t)
}{
(1-\alpha_t)(1-\bar{\alpha}_{t-1})
}.
\end{aligned}
\]
由于 \(1-\alpha_t=\beta_t\)，且 \(\bar{\alpha}_t=\alpha_t\bar{\alpha}_{t-1}\)，分子变为
\[
\alpha_t - \alpha_t\bar{\alpha}_{t-1} + 1 - \alpha_t
=
1-\bar{\alpha}_t.
\]
因此
\[
A
=
\frac{
1-\bar{\alpha}_t
}{
\beta_t(1-\bar{\alpha}_{t-1})
},
\]
从而
\[
\boxed{
\tilde{\beta}_t
=
A^{-1}
=
\frac{1-\bar{\alpha}_{t-1}}{1-\bar{\alpha}_t}\,\beta_t.
}
\]

接下来计算均值：
\[
\tilde{\boldsymbol{\mu}}_t
=
A^{-1}\mathbf{b}
=
\tilde{\beta}_t
\left(
\frac{\sqrt{\alpha_t}}{1-\alpha_t}\,\mathbf{x}_t
+
\frac{\sqrt{\bar{\alpha}_{t-1}}}{1-\bar{\alpha}_{t-1}}\,\mathbf{x}_0
\right).
\]
代入
\(
\tilde{\beta}_t
=
\frac{1-\bar{\alpha}_{t-1}}{1-\bar{\alpha}_t}\beta_t
\)
以及 \(1-\alpha_t=\beta_t\)，可得
\[
\begin{aligned}
\tilde{\boldsymbol{\mu}}_t
&=
\frac{1-\bar{\alpha}_{t-1}}{1-\bar{\alpha}_t}\beta_t
\left(
\frac{\sqrt{\alpha_t}}{\beta_t}\,\mathbf{x}_t
+
\frac{\sqrt{\bar{\alpha}_{t-1}}}{1-\bar{\alpha}_{t-1}}\,\mathbf{x}_0
\right)
\\[6pt]
&=
\frac{\sqrt{\alpha_t}(1-\bar{\alpha}_{t-1})}{1-\bar{\alpha}_t}\,\mathbf{x}_t
+
\frac{\sqrt{\bar{\alpha}_{t-1}}\beta_t}{1-\bar{\alpha}_t}\,\mathbf{x}_0.
\end{aligned}
\]
因此
\[
\boxed{
\tilde{\boldsymbol{\mu}}_t(\mathbf{x}_t,\mathbf{x}_0)
=
\frac{
\sqrt{\bar{\alpha}_{t-1}}\,\beta_t
}{
1-\bar{\alpha}_t
}\,\mathbf{x}_0
+
\frac{
\sqrt{\alpha_t}(1-\bar{\alpha}_{t-1})
}{
1-\bar{\alpha}_t
}\,\mathbf{x}_t.
}
\]

\subsection*{步骤 5：最终结果}

综合以上各式，
\[
\boxed{
q(\mathbf{x}_{t-1}\mid \mathbf{x}_t,\mathbf{x}_0)
=
\mathcal{N}\!\Bigl(
\mathbf{x}_{t-1}\mid
\tilde{\boldsymbol{\mu}}_t(\mathbf{x}_t,\mathbf{x}_0),
\tilde{\beta}_t\mathbf{I}
\Bigr),
}
\]
其中
\[
\tilde{\boldsymbol{\mu}}_t(\mathbf{x}_t,\mathbf{x}_0)
=
\frac{
\sqrt{\bar{\alpha}_{t-1}}\,\beta_t
}{
1-\bar{\alpha}_t
}\,\mathbf{x}_0
+
\frac{
\sqrt{\alpha_t}(1-\bar{\alpha}_{t-1})
}{
1-\bar{\alpha}_t
}\,\mathbf{x}_t,
\]
并且
\[
\tilde{\beta}_t
=
\frac{
1-\bar{\alpha}_{t-1}
}{
1-\bar{\alpha}_t
}\,\beta_t.
\]

这就是 ELBO 中学习得到的去噪转移 \(p_\theta(\mathbf{x}_{t-1}\mid \mathbf{x}_t)\) 所需要拟合的精确局部逆向目标。

\backmatter
\printbibliography[heading=bibintoc,title={参考文献}]
\end{document}
